\documentclass[11pt,oneside]{book}
\usepackage[utf8]{inputenc}
\usepackage[T2A]{fontenc}
\usepackage[russian]{babel}
\usepackage{graphicx}
\graphicspath{{target/img/},{target/img/tall/},{target/img/wide/}}

\usepackage{wrapfig}
\usepackage{indentfirst}
\usepackage{url}
\usepackage{ifpdf}

\renewcommand\sectionmark[1]{}
\newcommand{\partornament}{cross}
\newcommand{\ornament}{uzor_begin_10}
\newcommand{\ornamentending}{uzor_end_3}

% calculate header offset from text body
\setlength\headheight{14pt}
\newlength\headoff
\addtolength\headoff{\headheight}
\addtolength\headoff{\headsep}
\addtolength\headoff{14pt}

\title{Молитвослов на всякую потребу}
\author{www.molitvoslov.com}
\date{2011-10-13}
\setcounter{secnumdepth}{-3}
%\setcounter{tocdepth}{-1}

\newcommand{\mypart}[1]{
    \part{#1}
}

\newcommand{\mychapter}[1]{
    \renewcommand{\thechapter}{}
    \chapter{#1}
}

\newcommand{\mychapterz}[2]{
    \renewcommand{\thechapter}{}
    \chapter[#1]{#2}
}

\newcommand\mysubtitle[1]{
    \subsubsection{#1}
}

\newcommand\mysubsection[1]{
    \subsection{#1}
}

\newcommand\myparagraph[1]{
    \paragraph{#1}
}

\newenvironment{mymulticols}[1]{}{}

% the idea of identity environment is to provide a mechanism
% to assign some id to the element it surrounds. Its main use is
% with htlatex configs
\newcommand\myid{}
\newenvironment{identity}{}{}

\newenvironment{Parallel}{}{}

\newenvironment{centericon}{}{}

\newcommand{\myfig}[2][0.43]{
    \begin{wrapfigure}{l}{0pt}
    \includegraphics[width=#1\textwidth]{#2}
    \end{wrapfigure}
}

\newcommand{\myfigh}[3][0.43]{
    \begin{wrapfigure}[#3]{l}{0pt}
    \includegraphics[width=#1\textwidth]{#2}
    \end{wrapfigure}
}

\newcommand{\myfigr}[2][0.43]{
    \begin{wrapfigure}{r}{0pt}
    \includegraphics[width=#1\textwidth]{#2}
    \end{wrapfigure}
}

\newcommand{\myfigrh}[3][0.43]{
    \begin{wrapfigure}[#3]{r}{0pt}
    \includegraphics[width=#1\textwidth]{#2}
    \end{wrapfigure}
}

\newcommand{\myfigure}[2][0.43]{\myfig[#1]{#2}}

\newcommand{\mychapterending}[1][2.5]{}

\newcommand{\bukvaending}{\mychapterending}

\newcommand\longpage[1][1]{}
\newcommand\shortpage[1][1]{}

% save the value of current parindent to be used
% later in minipage environments
\newlength\myparindent
\setlength\myparindent{\parindent}
\newcommand\noparindent{\setlength\parindent{0pt}}
\newcommand\restoreparindent{\setlength\parindent{\myparindent}}

\long\def\symbolfootnote[#1]#2{\begingroup%
\def\thefootnote{\fnsymbol{footnote}}\footnote[#1]{#2}\endgroup}

\newcommand{\firstletter}[1]{#1}%{\myheadingcolor #1}}

\newcommand\pripev[2][Припев:]{{\small\myemph{#1} #2}}
\newcommand\irmos[1]{\pripev[Ирм\'{о}с:]{#1}}
\newcommand\Bogorodichen[1]{\myemph{Богородичен:} #1}
\newcommand\slavan{Слава Отцу и Сыну и Святому Духу.}
\newcommand\slava{{\small\slavan}}
\newcommand\inynen{И ныне и присно и во веки веков. Аминь.}
\newcommand\inyne{{\small\inynen}}
\newcommand\slavainynen{Слава Отцу и Сыну и Святому Духу. И ныне и присно и во веки веков. Аминь.}
\newcommand\slavainyne{{\small\slavainynen}}

\newcommand{\pripevc}[1]{{\small \centerline{#1} \nopagebreak}}
\newcommand{\pripevmskipc}[1]{\medskip\pripevc{#1}}
\newcommand{\pripevpomiluj}{\pripevmskipc{\pripev{\firstletter{П}омилуй мя, Боже, помилуй мя.}}}
\newcommand{\slavac}{\pripevmskipc{\slavan}}
\newcommand{\inynec}{\pripevmskipc{\inynen}}

\newcommand{\TsariuNebesnyj}{%
Царю Небесный, Утешителю, Душе истины, Иже везде сый и вся исполняяй, Сокровище благих и жизни Подателю, прииди и вселися в ны, и очисти ны от всякия скверны, и спаси, Блаже, души наша.}

\newcommand{\TrisviatoePoOtcheNash}{%
Святый Боже, Святый Крепкий, Святый Безсмертный, помилуй нас. \myemph{ (Tрижды)}

Слава Отцу и Сыну и Святому Духу, и ныне и присно и во веки веков. Аминь.

Пресвятая Троице, помилуй нас; Господи, очисти грехи наша; Владыко, прости беззакония наша; Святый, посети и исцели немощи наша, имене Твоего ради.

Господи, помилуй. \myemph{ (Трижды)}

Слава Отцу и Сыну и Святому Духу, и ныне и присно и во веки веков. Аминь.

Отче наш, Иже еси на небесех! Да святится имя Твое, да приидет Царствие Твое, да будет воля Твоя, яко на небеси и на земли. Хлеб наш насущный даждь нам днесь; и остави нам долги наша, якоже и мы оставляем должником нашим; и не введи нас во искушение, но избави нас от лукаваго.
}

\newcommand{\priiditepoklonimsia}{%
Приидите, поклонимся Цареви нашему Богу. \myemph{(Поклон)}

Приидите, поклонимся и припадем Христу, Цареви нашему Богу. \myemph{(Поклон)}

Приидите, поклонимся и припадем Самому Христу, Цареви и Богу нашему. \myemph{(Поклон)}}


\newcommand{\PsalmFifty}{%
Помилуй мя, Боже, по велицей милости Твоей, и по множеству щедрот Твоих очисти беззаконие мое. Наипаче омый мя от беззакония моего, и от греха моего очисти мя; яко беззаконие мое аз знаю, и грех мой предо мною есть выну. Тебе Единому согреших и лукавое пред Тобою сотворих, яко да оправдишися во словесех Твоих, и победиши внегда судити Ти. Се бо, в беззакониих зачат есмь, и во гресех роди мя мати моя. Се бо, истину возлюбил еси; безвестная и тайная премудрости Твоея явил ми еси. Окропиши мя иссопом, и очищуся; омыеши мя, и паче снега убелюся. Слуху моему даси радость и веселие; возрадуются кости смиренныя. Отврати лице Твое от грех моих и вся беззакония моя очисти. Сердце чисто созижди во мне, Боже, и дух прав обнови во утробе моей. Не отвержи мене от лица Твоего и Духа Твоего Святаго не отыми от мене. Воздаждь ми радость спасения Твоего и Духом владычним утверди мя. Научу беззаконыя путем Твоим, и нечестивии к Тебе обратятся. Избави мя от кровей, Боже, Боже спасения моего; возрадуется язык мой правде Твоей. Господи, устне мои отверзеши, и уста моя возвестят хвалу Твою. Яко аще бы восхотел еси жертвы, дал бых убо: всесожжения не благоволиши. Жертва Богу дух сокрушен; сердце сокрушенно и смиренно Бог не уничижит. Ублажи, Господи, благоволением Твоим Сиона, и да созиждутся стены Иерусалимския. Тогда благоволиши жертву правды, возношение и всесожегаемая; тогда возложат на oлтарь Твой тельцы.\par}

\newcommand{\Chestneyshuyu}{%
Достойно есть яко воистинну блажити Тя, Богородицу, Присноблаженную и Пренепорочную и Матерь Бога нашего. Честнейшую Херувим и славнейшую без сравнения Серафим, без истления Бога Слова рождшую, сущую Богородицу Тя величаем.}

\newcommand{\MolitvamiSviatyhOtecNashih}{%
Молитвами святых отец наших, Господи Иисусе Христе, Боже наш, помилуй нас. Аминь.}

\newcommand{\tolkopoblagosloveniyu}{%
{\centering\myemph{\normalfont Читаются только по благословению духовника}\par}}

\newcommand{\SymbolOfFaith}{%
  Верую во единаго Бога Отца, Вседержителя, Творца небу и земли, видимым же всем и невидимым.
  И во единаго Господа Иисуса Христа, Сына Божия, Единороднаго, Иже от Отца рожденнаго прежде всех век; Света от Света, Бога истинна от Бога истинна, рожденна, несотворенна, единосущна Отцу, Имже вся быша.
  Нас ради человек и нашего ради спасения сшедшаго с небес и воплотившагося от Духа Свята и Марии Девы и вочеловечшася.
  Распятаго же за ны при Понтийстем Пилате, и страдавша, и погребенна.
  И воскресшаго в третий день по Писанием.
  И возшедшаго на небеса, и седяща одесную Отца.
  И паки грядущаго со славою судити живым и мертвым, Егоже Царствию не будет конца.
  И в Духа Святаго, Господа, Животворящаго, Иже от Отца исходящего, Иже со Отцем и Сыном спокланяема и сславима, глаголавшаго пророки.
  Во едину Святую, Соборную и Апостольскую Церковь.
  Исповедую едино крещение во оставление грехов.
  Чаю воскресения мертвых, и жизни будущаго века. Аминь.}

\newcommand{\TroparPomilujNas}{Помилуй нас, Господи, помилуй нас; всякаго бо ответа недоумеюще, сию Ти молитву яко Владыце грешнии приносим: помилуй нас.

\slavan

 Господи, помилуй нас, на Тя бо уповахом; не прогневайся на ны зело, ниже помяни беззаконий наших, но призри и ныне яко благоутробен, и избави ны от враг наших; Ты бо еси Бог наш, и мы людие Твои, вси дела руку Твоею, и имя Твое призываем.

\inynen

 Милосердия двери отверзи нам, благословенная Богородице, надеющиися на Тя да не погибнем, но да избавимся Тобою от бед: Ты бо еси спасение рода христианскаго.}



\begin{document}

\renewcommand{\ornament}{uzor_psaltyr}

\vspace*{3cm}

\renewcommand\myid{secondpagewarning}

\begin{identity}

\subsection*{Дорогие братья и сестры!}

Прежде чем приступать к молитвам,
пожалуйста, внимательно изучите
наставления и правила, изложенные в
разделе «О Молитве», в особенности
«Поучение святителя Игнатия
Брянчанинова о молитвенном правиле».

\end{identity}


\tableofcontents[part,chapter,section,bukva]

\renewcommand{\ornament}{uzor_begin_10}



\mypart{МОЛИТВЫ}\label{_content_molitvi}
%http://www.molitvoslov.com/content/molitvi

\mychapter{Молитва Господня. Отче наш}\begin{mymulticols}

%http://www.molitvoslov.com/node/37

\myfigure{Spasnatrone}

Отче наш, Иже еси на небесех! Да святится имя Твое, да приидет Царствие Твое, да будет воля Твоя, яко на небеси и на земли. Хлеб наш насущный даждь нам днесь; и остави нам долги наша, якоже и мы оставляем должником нашим; и не введи нас во искушение, но избави нас от лукаваго.

Отче наш, сущий на небесах! да святится имя Твое; да приидет Царствие Твое; да будет воля Твоя и на земле, как на небе; хлеб наш насущный дай нам на сей день; и прости нам долги наши, как и мы прощаем должникам нашим; и не введи нас в искушение, но избавь нас от лукавого. Ибо Твое есть Царство и сила и слава во веки. Аминь. (Матф. 6:9--13)


\end{mymulticols}

\mychapterending


\mychapter{Иисусова молитва}
%http://www.molitvoslov.com/text596.htm

\myfigure[0.5]{456}

\begin{center}
Господи Иисусе Христе, Сыне Божий, помилуй мя, грешнаго.
\end{center}

\mychapterending


\mychapter{Благодарение за всякое благодеяние Божие}\begin{mymulticols}
%http://www.molitvoslov.com/text889.htm

\myfigure{777}

\mysubsubsection{Тропарь, глас 4-й}

Благодарни суще недостойнии раби Твои, Господи, о Твоих великих благодеяниих на нас бывших, славяще Тя хвалим, благословим, благодарим, поем и величаем Твое благоутробие, и рабски любовию вопием Ти: Благодетелю Спасе наш, слава Тебе.

\mysubsubsection{Кондак, глас 3-й}

Твоих благодеяний и даров туне, яко раби непотребнии, сподобльшеся, Владыко, к Тебе усердно притекающе, благодарение по силе приносим, и Тебе яко Благодетеля и Творца славяще, вопием: слава Тебе, Боже Всещедрый.

\slavainynen

\Bogorodichen{Богородице, христианом Помощнице, Твое предстательство стяжавше раби Твои, благодарно Тебе вопием: радуйся, Пречистая Богородице Дево, и от всех нас бед Твоими молитвами всегда избави, Едина вскоре предстательствующая.}

\end{mymulticols}

\mychapterending


\mychapter{Молитвы утренние}\begin{mymulticols}
%http://www.molitvoslov.com/text893.htm

\myfigure{795}

\footnote{Напечатанное {\itshape курсивом} или цветом (\myemph{пояснения} и {\bfseries\myheadingcolor названия молитв}) не читается во время молитвы.}\myemph{  Востав от сна, прежде всякого другого дела, стань благоговейно, представляя себя пред Всевидящим Богом, и, совершая крестное знамение, произнеси:}

Во имя Отца, и Сына, и Святаго Духа, Аминь.

\medskip\myemph{ Затем немного подожди, пока все чувства твои не придут в тишину и мысли твои не оставят все земное, и тогда произноси следующие молитвы, без поспешности и со вниманием сердечным:}

\mysubsubsection{Молитва мытаря \myemph{(Евангелие от Луки, глава 18, стих 13)}}

Боже, милостив буди мне грешному. \myemph{ (Поклон)}

\mysubsubsection{Молитва предначинательная}

Господи Иисусе Христе, Сыне Божий, молитв ради Пречистыя Твоея Матере и всех святых, помилуй нас. Аминь.

Слава Тебе, Боже наш, слава Тебе.

\mysubsubsection{Молитва Святому Духу}

Царю Небесный, Утешителю, Душе истины, Иже везде сый и вся исполняяй, Сокровище благих и жизни Подателю, прииди и вселися в ны, и очисти ны от всякия скверны, и спаси, Блаже, души наша\footnote{От Пасхи до Вознесения вместо этой молитвы читается тропарь: «Христос воскресе из мертвых, смертию смерть поправ, и сущим во гробех живот даровав». (Трижды) От Вознесения до Троицы начинаем молитвы со «Святый Боже…», опуская все предшествующие.

Это замечание относится и к молитвам на сон грядущим.
}.

\mysubsubsection{Трисвятое}

Святый Боже, Святый Крепкий, Святый Безсмертный, помилуй нас. \myemph{ (Читается трижды, с крестным знамением и поясным поклоном.) }

Слава Отцу и Сыну и Святому Духу, и ныне и присно и во веки веков. Аминь.

\mysubsubsection{Молитва ко Пресвятой Троице}

Пресвятая Троице, помилуй нас; Господи, очисти грехи наша; Владыко, прости беззакония наша; Святый, посети и исцели немощи наша, имене Твоего ради.

Господи, помилуй. \myemph{ (Трижды)}.

Слава Отцу и Сыну и Святому Духу, и ныне и присно и во веки веков. Аминь\footnote{Когда написано «Слава», «И ныне», надо читать полностью: «Слава Отцу и Сыну и Святому Духу», «И ныне и присно и во веки веков. Аминь»}.

\mysubsubsection{Молитва Господня}

Отче наш, Иже еси на небесех! Да святится имя Твое, да приидет Царствие Твое, да будет воля Твоя, яко на небеси и на земли. Хлеб наш насущный даждь нам днесь; и остави нам долги наша, якоже и мы оставляем должником нашим; и не введи нас во искушение, но избави нас от лукаваго.

\mysubsubsection{Тропари Троичные}

Воставше от сна, припадаем Ти, Блаже, и ангельскую песнь вопием Ти, Сильне: Свят, Свят, Свят еси, Боже, Богородицею помилуй нас.

\slavan

 От одра и сна воздвигл мя еси, Господи, ум мой просвети и сердце, и устне мои отверзи, во еже пети Тя, Святая Троице: Свят, Свят, Свят еси, Боже, Богородицею помилуй нас.

\inynen

 Внезапно Судия приидет, и коегождо деяния обнажатся, но страхом зовем\footnote{В церковнославянском языке нет звука ё, а поэтому надо читать «зовем», а не «зовём», «твое», а не «твоё», «мое», а не «моё» и т.~д.} в полунощи: Свят, Свят, Свят еси, Боже, Богородицею помилуй нас.

Господи, помилуй. \myemph{ (12 раз)}

\mysubsubsection{Молитва ко Пресвятой Троице}

От сна востав, благодарю Тя, Святая Троице, яко многия ради Твоея благости и долготерпения не прогневался еси на мя, лениваго и грешнаго, ниже погубил мя еси со беззаконьми моими; но человеколюбствовал еси обычно и в нечаянии лежащаго воздвигл мя еси, во еже утреневати и славословити державу Твою. И ныне просвети мои очи мысленныя, отверзи моя уста поучатися словесем Твоим, и разумети заповеди Твоя, и творити волю Твою, и пети Тя во исповедании сердечнем, и воспевати всесвятое имя Твое, Отца и Сына и Святаго Духа, ныне и присно и во веки веков. Аминь.

Приидите, поклонимся Цареви нашему Богу. \myemph{ (Поклон)}

Приидите, поклонимся и припадем Христу, Цареви нашему Богу. \myemph{ (Поклон)}

Приидите, поклонимся и припадем Самому Христу, Цареви и Богу нашему. \myemph{ (Поклон)}

\mysubsubsection{Псалом 50}

Помилуй мя, Боже, по велицей милости Твоей, и по множеству щедрот Твоих очисти беззаконие мое. Наипаче омый мя от беззакония моего, и от греха моего очисти мя; яко беззаконие мое аз знаю, и грех мой предо мною есть выну. Тебе Единому согреших и лукавое пред Тобою сотворих, яко да оправдишися во словесех Твоих, и победиши внегда судити Ти. Се бо, в беззакониих зачат есмь, и во гресех роди мя мати моя. Се бо, истину возлюбил еси; безвестная и тайная премудрости Твоея явил ми еси. Окропиши мя иссопом, и очищуся; омыеши мя, и паче снега убелюся. Слуху моему даси радость и веселие; возрадуются кости смиренныя. Отврати лице Твое от грех моих и вся беззакония моя очисти. Сердце чисто созижди во мне, Боже, и дух прав обнови во утробе моей. Не отвержи мене от лица Твоего и Духа Твоего Святаго не отыми от мене. Воздаждь ми радость спасения Твоего и Духом владычним утверди мя. Научу беззаконыя путем Твоим, и нечестивии к Тебе обратятся. Избави мя от кровей, Боже, Боже спасения моего; возрадуется язык мой правде Твоей. Господи, устне мои отверзеши, и уста моя возвестят хвалу Твою. Яко аще бы восхотел еси жертвы, дал бых убо: всесожжения не благоволиши. Жертва Богу дух сокрушен; сердце сокрушенно и смиренно Бог не уничижит. Ублажи, Господи, благоволением Твоим Сиона, и да созиждутся стены Иерусалимския. Тогда благоволиши жертву правды, возношение и всесожегаемая; тогда возложат на oлтарь Твой тельцы.

\mysubsubsection{Символ веры}

  Верую во единаго Бога Отца, Вседержителя, Творца небу и земли, видимым же всем и невидимым.
  И во единаго Господа Иисуса Христа, Сына Божия, Единороднаго, Иже от Отца рожденнаго прежде всех век; Света от Света, Бога истинна от Бога истинна, рожденна, несотворенна, единосущна Отцу, Имже вся быша.
  Нас ради человек и нашего ради спасения сшедшаго с небес и воплотившагося от Духа Свята и Марии Девы и вочеловечшася.
  Распятаго же за ны при Понтийстем Пилате, и страдавша, и погребенна.
  И воскресшаго в третий день по Писанием.
  И возшедшаго на небеса, и седяща одесную Отца.
  И паки грядущаго со славою судити живым и мертвым, Егоже Царствию не будет конца.
  И в Духа Святаго, Господа, Животворящаго, Иже от Отца исходящего, Иже со Отцем и Сыном спокланяема и сславима, глаголавшаго пророки.
  Во едину Святую, Соборную и Апостольскую Церковь.
  Исповедую едино крещение во оставление грехов.
  Чаю воскресения мертвых, и жизни будущаго века. Аминь.

\mysubsubsection{Молитва первая, святого Макария Великого}

Боже, очисти мя грешнаго, яко николиже сотворих благое пред Тобою; но избави мя от лукаваго, и да будет во мне воля Твоя, да неосужденно отверзу уста моя недостойная и восхвалю имя Твое святое, Отца и Сына и Святаго Духа, ныне и присно и во веки веков Аминь.

\mysubsubsection{ Молитва вторая, того же святого}

От сна востав, полунощную песнь приношу Ти, Спасе, и припадая вопию Ти: не даждь ми уснути во греховней смерти, но ущедри мя, распныйся волею, и лежащаго мя в лености ускорив возстави, и спаси мя в предстоянии и молитве, и по сне нощнем возсияй ми день безгрешен, Христе Боже, и спаси мя.

\mysubsubsection{Молитва третья, того же святого}

К Тебе, Владыко Человеколюбче, от сна востав прибегаю, и на дела Твоя подвизаюся милосердием Твоим, и молюся Тебе: помози ми на всякое время, во всякой вещи, и избави мя от всякия мирския злыя вещи и диавольскаго поспешения, и спаси мя, и введи в Царство Твое вечное. Ты бо еси мой Сотворитель и всякому благу Промысленник и Податель, о Тебе же все упование мое, и Тебе славу возсылаю, ныне и присно и во веки веков. Аминь.

\mysubsubsection{Молитва четвертая, того же святого}

Господи, Иже многою Твоею благостию и великими щедротами Твоими дал еси мне, рабу Твоему, мимошедшее время нощи сея без напасти прейти от всякаго зла противна; Ты Сам, Владыко, всяческих Творче, сподоби мя истинным Твоим светом и просвещенным сердцем творити волю Твою, ныне и присно и во веки веков. Аминь.

\mysubsubsection{Молитва пятая, святого Василия Великого}

Господи Вседержителю, Боже сил и всякия плоти, в вышних живый и на смиренныя призираяй, сердца же и утробы испытуяй и сокровенная человеков яве предведый, Безначальный и Присносущный Свете, у Него же несть пременение, или преложения осенение; Сам, Безсмертный Царю, приими моления наша, яже в настоящее время, на множество Твоих щедрот дерзающе, от скверных к Тебе устен творим, и остави нам прегрешения наша, яже делом, и словом, и мыслию, ведением, или неведением согрешенная нами; и очисти ны от всякия скверны плоти и духа. И даруй нам бодренным сердцем и трезвенною мыслию всю настоящаго жития нощь прейти, ожидающим пришествия светлаго и явленнаго дне Единороднаго Твоего Сына, Господа и Бога и Спаса нашего Иисуса Христа, в оньже со славою Судия всех приидет, комуждо отдати по делом его; да не падше и обленившеся, но бодрствующе и воздвижени в делание обрящемся готови, в радость и Божественный чертог славы Его совнидем, идеже празднующих глас непрестанный, и неизреченная сладость зрящих Твоего лица доброту неизреченную. Ты бо еси истинный Свет, просвещаяй и освящаяй всяческая, и Тя поет вся тварь во веки веков. Аминь.

\mysubsubsection{Молитва шестая, того же святого}

Тя благословим, вышний Боже и Господи милости, творящаго присно с нами великая же и неизследованная, славная же и ужасная, ихже несть числа, подавшаго нам сон во упокоение немощи нашея, и ослабление трудов многотрудныя плоти. Благодарим Тя, якo не погубил еси нас со беззаконьми нашими, но человеколюбствовал еси обычно, и в нечаянии лежащия ны воздвигл еси, во еже славословити державу Твою. Темже молим безмерную Твою благость, просвети наша мысли, очеса, и ум наш от тяжкаго сна лености возстави: отверзи наша уста, и исполни я Твоего хваления, яко да возможем непоколеблемо пети же и исповедатися Тебе, во всех, и от всех славимому Богу, Безначальному Отцу, со Единородным Твоим Сыном, и Всесвятым и Благим и Животворящим Твоим Духом, ныне и присно и во веки веков. Аминь.

\mysubsubsection{Молитва седьмая, ко Пресвятой Богородице}

Воспеваю благодать Твою, Владычице, молю Тя, ум мой облагодати. Ступати право мя настави, путем Христовых заповедей. Бдети к песни укрепи, уныния сон отгоняющи. Связана пленицами грехопадений, мольбами Твоими разреши, Богоневесто. В нощи мя и во дни сохраняй, борющих враг избавляющи мя. Жизнодателя Бога рождшая, умерщвлена мя страстьми оживи. Яже Свет невечерний рождшая, душу мою ослепшую просвети. О дивная Владычня палато, дом Духа Божественна мене сотвори. Врача рождшая, уврачуй души моея многолетныя страсти. Волнующася житейскою бурею, ко стези мя покаяния направи. Избави мя огня вечнующаго, и червия же злаго, и тартара. Да мя не явиши бесом радование, иже многим грехом повинника. Нова сотвори мя, обетшавшаго нечувственными, Пренепорочная, согрешении. Странна муки всякия покажи мя, и всех Владыку умоли. Небесная ми улучити веселия, со всеми святыми, сподоби. Пресвятая Дево, услыши глас непотребнаго раба Твоего. Струю давай мне слезам, Пречистая, души моея скверну очищающи. Стенания от сердца приношу Ти непрестанно, усердствуй, Владычице. Молебную службу мою приими, и Богу благоутробному принеси. Превышшая Ангел, мирскаго мя превышша слития сотвори. Светоносная Сене небесная, духовную благодать во мне направи. Руце воздею и устне к похвалению, осквернены скверною, Всенепорочная. Душетленных мя пакостей избави, Христа прилежно умоляющи; Емуже честь и поклонение подобает, ныне и присно и во веки веков. Аминь.

\mysubsubsection{Молитва восьмая, ко Господу нашему Иисусу Христу}

Многомилостиве и Всемилостиве Боже мой, Господи Иисусе Христе, многия ради любве сшел и воплотился еси, яко да спасеши всех. И паки, Спасе, спаси мя по благодати, молю Тя; аще бо от дел спасеши мя, несть се благодать, и дар, но долг паче. Ей, многий в щедротах и неизреченный в милости! Веруяй бо в Мя, рекл еси, о Христе мой, жив будет и не узрит смерти во веки. Аще убо вера, яже в Тя, спасает отчаянныя, се верую, спаси мя, яко Бог мой еси Ты и Создатель. Вера же вместо дел да вменится мне, Боже мой, не обрящеши бо дел отнюд оправдающих мя. Но та вера моя да довлеет вместо всех, та да отвещает, та да оправдит мя, та да покажет мя причастника славы Твоея вечныя. Да не убо похитит мя сатана, и похвалится, Слове, еже отторгнути мя от Твоей руки и ограды; но или хощу, спаси мя, или не хощу, Христе Спасе мой, предвари скоро, скоро погибох: Ты бо еси Бог мой от чрева матере моея. Сподоби мя, Господи, ныне возлюбити Тя, якоже возлюбих иногда той самый грех; и паки поработати Тебе без лености тощно, якоже поработах прежде сатане льстивому. Наипаче же поработаю Тебе, Господу и Богу моему Иисусу Христу, во вся дни живота моего, ныне и присно и во веки веков. Аминь.

\mysubsubsection{Молитва девятая, к Ангелу хранителю}

Святый Ангеле, предстояй окаянной моей души и страстной моей жизни, не остави мене грешнаго, ниже отступи от мене за невоздержание мое. Не даждь места лукавому демону обладати мною, насильством смертнаго сего телесе; укрепи бедствующую и худую мою руку и настави мя на путь спасения. Ей, святый Ангеле Божий, хранителю и покровителю окаянныя моея души и тела, вся мне прости, еликими тя оскорбих во вся дни живота моего, и аще что согреших в прешедшую нощь сию, покрый мя в настоящий день, и сохрани мя от всякаго искушения противнаго, да ни в коем гресе прогневаю Бога, и молися за мя ко Господу, да утвердит мя в страсе Своем, и достойна покажет мя раба Своея благости. Аминь.

\mysubsubsection{Молитва десятая, ко Пресвятой Богородице}

Пресвятая Владычице моя Богородице, святыми Твоими и всесильными мольбами отжени от мене, смиреннаго и окаяннаго раба Твоего, уныние, забвение, неразумие, нерадение, и вся скверная, лукавая и хульная помышления от окаяннаго моего сердца и от помраченнаго ума моего; и погаси пламень страстей моих, яко нищ есмь и окаянен. И избави мя от многих и лютых воспоминаний и предприятий, и от всех действ злых свободи мя. Яко благословена еси от всех родов, и славится пречестное имя Твое во веки веков. Аминь.

\mysubsubsection{Молитвенное призывание святого, имя которого носишь}

Моли Бога о мне, святый угодниче Божий \myemph{ (имя)}, яко аз усердно к тебе прибегаю, скорому помощнику и молитвеннику о душе моей.

\mysubsubsection{Песнь Пресвятой Богородице}

Богородице Дево, радуйся, Благодатная Марие, Господь с Тобою; благословена Ты в женах и благословен плод чрева Твоего, яко Спаса родила еси душ наших.

\mysubsubsection{Тропарь Кресту и молитва за отечество}

Спаси, Господи, люди Твоя, и благослови достояние Твое, победы православным христианом на сопротивныя даруя, и Твое сохраняя Крестом Твоим жительство.

\mysubsubsection{Молитва о живых}

Спаси, Господи, и помилуй отца моего духовнаго \myemph{ (имя)}, родителей моих \myemph{ (имена)}, сродников \myemph{ (имена)}, начальников, наставников, благодетелей \myemph{ (имена их)} и всех православных христиан.

\mysubsubsection{Молитва о усопших}

Упокой, Господи, души усопших раб Твоих: родителей моих, сродников, благодетелей (имена их), и всех православных христиан, и прости им вся согрешения вольная и невольная, и даруй им Царствие Небесное.

\myemph{ Если можешь, читай вместо кратких молитв о живых и усопших этот помянник:}

\mysubsubsection{О живых}

Помяни, Господи Иисусе Христе, Боже наш, милости и щедроты Твоя от века сущия, ихже ради и вочеловечился еси, и распятие и смерть, спасения ради право в Тя верующих, претерпети изволил еси; и воскрес из мертвых, вознеслся еси на небеса и седиши одесную Бога Отца, и призираеши на смиренныя мольбы всем сердцем призывающих Тя: приклони ухо Твое, и услыши смиренное моление мене, непотребнаго раба Твоего, в воню благоухания духовнаго, Тебе за вся люди Твоя приносящаго. И в первых помяни Церковь Твою Святую, Соборную и Апостольскую, юже снабдел еси честною Твоею Кровию, и утверди, и укрепи, и разшири, умножи, умири, и непреобориму адовыми враты во веки сохрани; раздирания Церквей утиши, шатания языческая угаси, и ересей востания скоро разори и искорени, и в ничтоже силою Святаго Твоего Духа обрати. \myemph{ (Поклон)}

Спаси, Господи, и помилуй Богом хранимую страну нашу, власти и воинство ея, огради миром державу их, и покори под нозе Православных всякаго врага и супостата, и глаголи мирная и благая в сердцах их о Церкви Твоей Святей, и о всех людех Твоих: да тихое и безмолвное житие поживем во правоверии, и во всяком благочестии и чистоте. \myemph{ (Поклон)}

Спаси, Господи, и помилуй Великаго Господина и Отца нашего Святейшего Патриарха Кирилла, преосвященныя митрополиты, архиепископы и епископы православныя, иереи же и диаконы, и весь причет церковный, яже поставил еси пасти словесное Твое стадо, и молитвами их помилуй и спаси мя грешнаго. \myemph{ (Поклон)}

Спаси, Господи, и помилуй отца моего духовнаго \myemph{ (имя его)}, и святыми его молитвами прости моя согрешения. \myemph{ (Поклон)}

Спаси, Господи, и помилуй родители моя \myemph{ (имена их)}, братию и сестры, и сродники моя по плоти, и вся ближния рода моего, и други, и даруй им мирная Твоя и премирная благая. \myemph{ (Поклон)}

Спаси, Господи, и помилуй по множеству щедрот Твоих вся священноиноки, иноки же и инокини, и вся в девстве же и благоговении и постничестве живущия в монастырех, в пустынях, в пещерах, горах, столпех, затворех, разселинах каменных, островех же морских, и на всяком месте владычествия Твоего правоверно живущия, и благочестно служащия Ти, и молящияся Тебе: облегчи им тяготу, и утеши их скорбь, и к подвигу о Тебе силу и крепость им подаждь, и молитвами их даруй ми оставление грехов. \myemph{ (Поклон)}

Спаси, Господи, и помилуй старцы и юныя, нищия и сироты и вдовицы, и сущия в болезни и в печалех, бедах же и скорбех, обстояниих и пленениих, темницах же и заточениих, изряднее же в гонениих, Тебе ради и веры православныя, от язык безбожных, от отступник и от еретиков, сущия рабы Твоя, и помяни я, посети, укрепи, утеши, и вскоре силою Твоею ослабу, свободу и избаву им подаждь. \myemph{ (Поклон) }

Спаси, Господи, и помилуй благотворящия нам, милующия и питающия нас, давшия нам милостыни, и заповедавшия нам недостойным молитися о них, и упокоевающия нас, и сотвори милость Твою с ними, даруя им вся, яже ко спасению прошения, и вечных благ восприятие. \myemph{ (Поклон)}

Спаси, Господи, и помилуй посланныя в службу, путешествующия, отцы и братию нашу, и вся православныя христианы. \myemph{ (Поклон)}

Спаси, Господи, и помилуй ихже аз безумием моим соблазних, и от пути спасительнаго отвратих, к делом злым и неподобным приведох; Божественным Твоим Промыслом к пути спасения паки возврати. \myemph{ (Поклон) }

Спаси, Господи, и помилуй ненавидящия и обидящия мя, и творящия ми напасти, и не остави их погибнути мене ради, грешнаго. \myemph{ (Поклон)}

Отступившия от православныя веры и погибельными ересьми ослепленныя, светом Твоего познания просвети и Святей Твоей Апостольстей Соборней Церкви причти. \myemph{ (Поклон) }

\mysubsubsection{О усопших}

Помяни, Господи, от жития сего отшедшия правоверныя цари и царицы, благоверныя князи и княгини, святейшия патриархи, преосвященныя митрополиты, архиепископы и епископы православныя, во иерейстем же и в причте церковнем, и монашестем чине Тебе послужившия, и в вечных Твоих селениих со святыми упокой. \myemph{ (Поклон.)}

Помяни, Господи, души усопших рабов Твоих, родителей моих \myemph{ (имена их)}, и всех сродников по плоти; и прости их вся согрешения вольная и невольная, даруя им Царствие и причастие вечных Твоих благих и Твоея безконечныя и блаженныя жизни наслаждение. \myemph{ (Поклон) }

Помяни, Господи, и вся в надежди воскресения и жизни вечныя усопшия, отцы и братию нашу, и сестры, и зде лежащия и повсюду, православныя христианы, и со святыми Твоими, идеже присещает свет лица Твоего, всели, и нас помилуй, яко Благ и Человеколюбец. Аминь. \myemph{ (Поклон) }

Подаждь, Господи, оставление грехов всем прежде отшедшим в вере и надежди воскресения, отцем, братиям и сестрам нашим и сотвори им вечную память. \myemph{ (Трижды)}

\mysubsubsection{Окончание молитв}

Достойно есть яко воистину блажити Тя Богородицу, Присноблаженную и Пренепорочную и Матерь Бога нашего. Честнейшую Херувим и славнейшую без сравнения Серафим, без истления Бога Слова рождшую, сущую Богородицу Тя величаем\footnote{От Пасхи до Вознесения вместо этой молитвы читается припев и ирмос 9-й песни пасхального канона:

\begin{quote}

«Ангел вопияше Благодатней: Чистая Дево, радуйся! И паки реку: радуйся! Твой Сын воскресе тридневен от гроба и мертвыя воздвигнувый; людие, веселитеся!

Светися, светися, новый Иерусалиме, слава бо Господня на тебе возсия. Ликуй ныне и веселися, Сионе. Ты же, Чистая, красуйся, Богородице, о востании Рождества Твоего».

\end{quote}

Это замечание относится и к вечерним молитвам. }.

\slavainynen

 Господи, помилуй. \myemph{ (Трижды)}

Господи, Иисусе Христе, Сыне Божий, молитв ради Пречистыя Твоея Матере, преподобных и богоносных отец наших и всех святых помилуй нас. Аминь.

\end{mymulticols}

\mychapterending


\mychapter{Молитвы на сон грядущим}\begin{mymulticols}
%http://www.molitvoslov.com/text2.htm

\myfigure{1_1}

Во имя Отца, и Сына, и Святаго Духа. Аминь.

Господи Иисусе Христе, Сыне Божий, молитв ради Пречистыя Твоея Матере, преподобных и богоносных отец наших и всех святых, помилуй нас. Аминь.

Слава Тебе, Боже наш, слава Тебе.

Царю Небесный, Утешителю, Душе истины, Иже везде сый и вся исполняяй, Сокровище благих и жизни Подателю, прииди и вселися в ны, и очисти ны от всякия скверны, и спаси, Блаже, души наша.

Святый Боже, Святый Крепкий, Святый Безсмертный, помилуй нас.\myemph{  (Трижды)}

Слава Отцу и Сыну и Святому Духу, и ныне и присно и во веки веков. Аминь.

Пресвятая Троице, помилуй нас; Господи, очисти грехи наша; Владыко, прости беззакония наша; Святый, посети и исцели немощи наша, имене Твоего ради.

Господи, помилуй. \myemph{ (Трижды)}

Слава Отцу и Сыну и Святому Духу, и ныне и присно и во веки веков. Аминь.

Отче наш, Иже еси на небесех! Да святится имя Твое, да приидет Царствие Твое, да будет воля Твоя, яко на небеси и на земли. Хлеб наш насущный даждь нам днесь; и остави нам долги наша, якоже и мы оставляем должником нашим; и не введи нас во искушение, но избави нас от лукаваго.

\mysubsubsection{Тропари}

Помилуй нас, Господи, помилуй нас; всякаго бо ответа недоумеюще, сию Ти молитву яко Владыце грешнии приносим: помилуй нас.

\slavan

 Господи, помилуй нас, на Тя бо уповахом; не прогневайся на ны зело, ниже помяни беззаконий наших, но призри и ныне яко благоутробен, и избави ны от враг наших; Ты бо еси Бог наш, и мы людие Твои, вси дела руку Твоею, и имя Твое призываем.

\inynen

 Милосердия двери отверзи нам, благословенная Богородице, надеющиися на Тя да не погибнем, но да избавимся Тобою от бед: Ты бо еси спасение рода христианскаго.

Господи, помилуй. \myemph{ (12 раз)}

\mysubsubsection{Молитва 1-я, святого Макария Великого, к Богу Отцу}

Боже вечный и Царю всякаго создания, сподобивый мя даже в час сей доспети, прости ми грехи, яже сотворих в сей день делом, словом и помышлением, и очисти, Господи, смиренную мою душу от всякия скверны плоти и духа. И даждь ми, Господи, в нощи сей сон прейти в мире, да востав от смиреннаго ми ложа, благоугожду пресвятому имени Твоему, во вся дни живота моего, и поперу борющия мя враги плотския и безплотныя. И избави мя, Господи, от помышлений суетных, оскверняющих мя, и похотей лукавых. Яко Твое есть царство, и сила и слава, Отца и Сына и Святаго Духа, ныне и присно и во веки веков. Аминь.

\mysubsubsection{Молитва 2-я, святого Антиоха, ко Господу нашему Иисусу Христу}

Вседержителю, Слово Отчее, Сам совершен сый, Иисусе Христе, многаго ради милосердия Твоего никогдаже отлучайся мене, раба Твоего, но всегда во мне почивай. Иисусе, добрый Пастырю Твоих овец, не предаждь мене крамоле змиине, и желанию сатанину не остави мене, яко семя тли во мне есть. Ты убо, Господи Боже покланяемый, Царю Святый, Иисусе Христе, спяща мя сохрани немерцающим светом, Духом Твоим Святым, Имже освятил еси Твоя ученики. Даждь, Господи, и мне, недостойному рабу Твоему, спасение Твое на ложи моем: просвети ум мой светом разума святаго Евангелия Твоего, душу любовию Креста Твоего, сердце чистотою словесе Твоего, тело мое Твоею страстию безстрастною, мысль мою Твоим смирением сохрани, и воздвигни мя во время подобно на Твое славословие. Яко препрославлен еси со Безначальным Твоим Отцем и с Пресвятым Духом во веки. Аминь.

\mysubsubsection{Молитва 3-я, ко Пресвятому Духу}

Господи, Царю Небесный, Утешителю, Душе истины, умилосердися и помилуй мя грешнаго раба Твоего, и отпусти ми недостойному, и прости вся, елика Ти согреших днесь яко человек, паче же и не яко человек, но и горее скота, вольныя моя грехи и невольныя, ведомыя и неведомыя: яже от юности и науки злы, и яже суть от нагльства и уныния. Аще именем Твоим кляхся, или похулих е в помышлении моем; или кого укорих; или оклеветах кого гневом моим, или опечалих, или о чем прогневахся; или солгах, или безгодно спах, или нищ прииде ко мне, и презрех его; или брата моего опечалих, или свадих, или кого осудих; или развеличахся, или разгордехся, или разгневахся; или стоящу ми на молитве, ум мой о лукавствии мира сего подвижеся, или развращение помыслих; или объядохся, или опихся, или без ума смеяхся; или лукавое помыслих, или доброту чуждую видев, и тою уязвлен бых сердцем; или неподобная глаголах, или греху брата моего посмеяхся, моя же суть безчисленная согрешения; или о молитве не радих, или ино что содеях лукавое, не помню, та бо вся и больша сих содеях. Помилуй мя, Творче мой Владыко, унылаго и недостойнаго раба Твоего, и остави ми, и отпусти, и прости мя, яко Благ и Человеколюбец, да с миром лягу, усну и почию, блудный, грешный и окаянный аз, и поклонюся, и воспою, и прославлю пречестное имя Твое, со Отцем, и Единородным Его Сыном, ныне и присно и во веки. Аминь.

\mysubsubsection{Молитва 4-я, святого Макария Великого}

Что Ти принесу, или что Ти воздам, великодаровитый Безсмертный Царю, щедре и человеколюбче Господи, яко ленящася мене на Твое угождение, и ничтоже благо сотворша, привел еси на конец мимошедшаго дне сего, обращение и спасение души моей строя? Милостив ми буди грешному и обнаженному всякаго дела блага, возстави падшую мою душу, осквернившуюся в безмерных согрешениих, и отыми от мене весь помысл лукавый видимаго сего жития. Прости моя согрешения, едине Безгрешне, яже Ти согреших в сей день, ведением и неведением, словом, и делом, и помышлением, и всеми моими чувствы. Ты Сам, покрывая, сохрани мя от всякаго сопротивнаго обстояния Божественною Твоею властию, и неизреченным человеколюбием, и силою. Очисти, Боже, очисти множество грехов моих. Благоволи, Господи, избавити мя от сети лукаваго, и спаси страстную мою душу, и осени мя светом лица Твоего, егда приидеши во славе, и неосужденна ныне сном уснути сотвори, и без мечтания, и несмущен помысл раба Твоего соблюди, и всю сатанину детель отжени от мене, и просвети ми разумныя очи сердечныя, да не усну в смерть. И посли ми Ангела мирна, хранителя и наставника души и телу моему, да избавит мя от враг моих; да востав со одра моего, принесу Ти благодарственныя мольбы. Ей, Господи, услыши мя грешнаго и убогаго раба Твоего, изволением и совестию; даруй ми воставшу словесем Твоим поучитися, и уныние бесовское далече от мене отгнано быти сотвори Твоими Ангелы; да благословлю имя Твое святое, и прославлю, и славлю Пречистую Богородицу Марию, Юже дал еси нам грешным заступление, и приими Сию молящуюся за ны; вем бо, яко подражает Твое человеколюбие, и молящися не престает. Тоя заступлением, и Честнаго Креста знамением, и всех святых Твоих ради, убогую душу мою соблюди, Иисусе Христе Боже наш, яко Свят еси, и препрославлен во веки. Аминь.

\mysubsubsection{Молитва 5-я}

Господи Боже наш, еже согреших во дни сем словом, делом и помышлением, яко Благ и Человеколюбец прости ми. Мирен сон и безмятежен даруй ми. Ангела Твоего хранителя посли, покрывающа и соблюдающа мя от всякаго зла, яко Ты еси хранитель душам и телесем нашим, и Тебе славу возсылаем, Отцу и Сыну и Святому Духу, ныне и присно и во веки веков. Аминь.

\mysubsubsection{Молитва 6-я}

Господи Боже наш, в Негоже веровахом, и Егоже имя паче всякаго имене призываем, даждь нам, ко сну отходящим, ослабу души и телу, и соблюди нас от всякаго мечтания, и темныя сласти кроме; устави стремление страстей, угаси разжжения востания телеснаго. Даждь нам целомудренно пожити делы и словесы; да добродетельное жительство восприемлюще, обетованных не отпадем благих Твоих, яко благословен еси во веки. Аминь.

\mysubsubsection{Молитва 7-я, святого Иоанна Златоуста (24 молитвы, по числу часов дня и ночи)}

Господи, не лиши мене небесных Твоих благ.

Господи, избави мя вечных мук.

Господи, умом ли или помышлением, словом или делом согреших, прости мя.

Господи, избави мя всякаго неведения и забвения, и малодушия, и окамененнаго нечувствия.

Господи, избави мя от всякаго искушения.

Господи, просвети мое сердце, еже помрачи лукавое похотение.

Господи, аз яко человек согреших, Ты же яко Бог щедр, помилуй мя, видя немощь души моея.

Господи, посли благодать Твою в помощь мне, да прославлю имя Твое святое.

Господи Иисусе Христе, напиши мя раба Твоего в книзе животней и даруй ми конец благий.

Господи, Боже мой, аще и ничтоже благо сотворих пред Тобою, но даждь ми по благодати Твоей положити начало благое.

Господи, окропи в сердце моем росу благодати Твоея.

Господи небесе и земли, помяни мя грешнаго раба Твоего, студнаго и нечистаго, во Царствии Твоем. Аминь.

Господи, в покаянии приими мя.

Господи, не остави мене.

Господи, не введи мене в напасть.

Господи, даждь ми мысль благу.

Господи, даждь ми слезы и память смертную, и умиление.

Господи, даждь ми помысл исповедания грехов моих.

Господи, даждь ми смирение, целомудрие и послушание.

Господи, даждь ми терпение, великодушие и кротость.

Господи, всели в мя корень благих, страх Твой в сердце мое.

Господи, сподоби мя любити Тя от всея души моея и помышления и творити во всем волю Твою.

Господи, покрый мя от человек некоторых, и бесов, и страстей, и от всякия иныя неподобныя вещи.

Господи, веси, яко твориши, якоже Ты волиши, да будет воля Твоя и во мне грешнем, яко благословен еси во веки. Аминь.

\mysubsubsection{Молитва 8-я, ко Господу нашему Иисусу Христу}

Господи Иисусе Христе, Сыне Божий, ради честнейшия Матере Твоея, и безплотных Твоих Ангел, Пророка же и Предтечи и Крестителя Твоего, богоглаголивых же апостол, светлых и добропобедных мученик, преподобных и богоносных отец, и всех святых молитвами, избави мя настоящаго обстояния бесовскаго. Ей, Господи мой и Творче, не хотяй смерти грешнаго, но якоже обратитися и живу быти ему, даждь и мне обращение окаянному и недостойному; изми мя от уст пагубнаго змия, зияющаго пожрети мя и свести во ад жива. Ей, Господи мой, утешение мое, Иже мене ради окаяннаго в тленную плоть оболкийся, исторгни мя от окаянства, и утешение подаждь души моей окаянней. Всади в сердце мое творити Твоя повеления, и оставити лукавая деяния, и получити блаженства Твоя: на Тя бо, Господи, уповах, спаси мя.

\mysubsubsection{Молитва 9-я, ко Пресвятой Богородице, Петра Студийского}

К Тебе Пречистей Божией Матери аз окаянный припадая молюся: веси, Царице, яко безпрестани согрешаю и прогневляю Сына Твоего и Бога моего, и многажды аще каюся, лож пред Богом обретаюся, и каюся трепеща: неужели Господь поразит мя, и по часе паки таяжде творю; ведущи сия, Владычице моя Госпоже Богородице, молю, да помилуеши, да укрепиши, и благая творити да подаси ми. Веси бо, Владычице моя Богородице, яко отнюд имам в ненависти злая моя дела, и всею мыслию люблю закон Бога моего; но не вем, Госпоже Пречистая, откуду яже ненавижду, та и люблю, а благая преступаю. Не попущай, Пречистая, воли моей совершатися, не угодна бо есть, но да будет воля Сына Твоего и Бога моего: да мя спасет, и вразумит, и подаст благодать Святаго Духа, да бых аз отселе престал сквернодейства, и прочее пожил бых в повелении Сына Твоего, Емуже подобает всякая слава, честь и держава, со Безначальным Его Отцем, и Пресвятым и Благим и Животворящим Его Духом, ныне и присно, и во веки веков. Aминь.

\mysubsubsection{Молитва 10-я, ко Пресвятой Богородице}

Благаго Царя благая Мати, Пречистая и Благословенная Богородице Марие, милость Сына Твоего и Бога нашего излей на страстную мою душу и Твоими молитвами настави мя на деяния благая, да прочее время живота моего без порока прейду и Тобою рай да обрящу, Богородице Дево, едина Чистая и Благословенная.

\mysubsubsection{Молитва 11-я, ко святому Ангелу хранителю}

Ангеле Христов, хранителю мой святый и покровителю души и тела моего, вся ми прости, елика согреших во днешний день, и от всякаго лукавствия противнаго ми врага избави мя, да ни в коемже гресе прогневаю Бога моего; но моли за мя грешнаго и недостойнаго раба, яко да достойна мя покажеши благости и милости Всесвятыя Троицы и Матере Господа моего Иисуса Христа и всех святых. Аминь.

\mysubsubsection{Кондак Богородице}

Взбранной Воеводе победительная, яко избавльшеся от злых, благодарственная восписуем Ти раби Твои, Богородице, но яко имущая державу непобедимую, от всяких нас бед свободи, да зовем Ти; радуйся, Невесто Неневестная.

Преславная Приснодево, Мати Христа Бога, принеси нашу молитву Сыну Твоему и Богу нашему, да спасет Тобою души наша.

Все упование мое на Тя возлагаю, Мати Божия, сохрани мя под кровом Твоим.

Богородице Дево, не презри мене, грешнаго, требующа Твоея помощи и Твоего заступления, на Тя бо упова душа моя, и помилуй мя.

\mysubsubsection{Молитва святого Иоанникия}

Упование мое Отец, прибежище мое Сын, покров мой Дух Святый: Троице Святая, слава Тебе.

Достойно есть яко воистинну блажити Тя, Богородицу, Присноблаженную и Пренепорочную и Матерь Бога нашего. Честнейшую Херувим и славнейшую без сравнения Серафим, без истления Бога Слова рождшую, сущую Богородицу Тя величаем.

Слава Отцу и Сыну и Святому Духу, и ныне и присно и во веки веков. Аминь.

Господи, помилуй. \myemph{ (Трижды)}

Господи Иисусе Христе, Сыне Божий, молитв ради Пречистыя Твоея Матере, преподобных и богоносных отец наших и всех святых, помилуй нас. Аминь.

\mysubsubsection{Молитва святого Иоанна Дамаскина}

Владыко Человеколюбче, неужели мне одр сей гроб будет, или еще окаянную мою душу просветиши днем? Се ми гроб предлежит, се ми смерть предстоит. Суда Твоего, Господи, боюся и муки безконечныя, злое же творя не престаю: Тебе Господа Бога моего всегда прогневляю, и Пречистую Твою Матерь, и вся Небесныя силы, и святаго Ангела хранителя моего. Вем убо, Господи, яко недостоин есмь человеколюбия Твоего, но достоин есмь всякаго осуждения и муки. Но, Господи, или хощу, или не хощу, спаси мя. Аще бо праведника спасеши, ничтоже велие; и аще чистаго помилуеши, ничтоже дивно: достойни бо суть милости Твоея. Но на мне грешнем удиви милость Твою: о сем яви человеколюбие Твое, да не одолеет моя злоба Твоей неизглаголанней благости и милосердию: и якоже хощеши, устрой о мне вещь.

Просвети очи мои, Христе Боже, да не когда усну в смерть, да не когда речет враг мой: укрепихся на него.

\slavan

 Заступник души моея буди, Боже, яко посреде хожду сетей многих; избави мя от них и спаси мя, Блаже, яко Человеколюбец.

\inynen

 Преславную Божию Матерь, и святых Ангел Святейшую, немолчно воспоим сердцем и усты, Богородицу сию исповедающе, яко воистинну рождшую нам Бога воплощенна, и молящуюся непрестанно о душах наших.

\mysubsubsection{Знаменуй себя крестом и говори молитву Честному Кресту:}

Да воскреснет Бог, и расточатся врази Его, и да бежат от лица Его ненавидящии Его. Яко исчезает дым, да исчезнут; яко тает воск от лица огня, тако да погибнут беси от лица любящих Бога и знаменующихся крестным знамением, и в веселии глаголющих: радуйся, Пречестный и Животворящий Кресте Господень, прогоняяй бесы силою на тебе пропятаго Господа нашего Иисуса Христа, во ад сшедшаго и поправшего силу диаволю, и даровавшаго нам тебе Крест Свой Честный на прогнание всякаго супостата. О, Пречестный и Животворящий Кресте Господень! Помогай ми со Святою Госпожею Девою Богородицею и со всеми святыми во веки. Аминь.

\myemph{ Или кратко:}

Огради мя, Господи, силою Честнаго и Животворящаго Твоего Креста, и сохрани мя от всякаго зла.

\mysubsubsection{Молитва}

Ослаби, остави, прости, Боже, прегрешения наша, вольная и невольная, яже в слове и в деле, яже в ведении и не в ведении, яже во дни и в нощи, яже во уме и в помышлении: вся нам прости, яко Благ и Человеколюбец.

\mysubsubsection{Молитва}

Ненавидящих и обидящих нас прости, Господи Человеколюбче. Благотворящим благосотвори. Братиям и сродником нашим даруй яже ко спасению прошения и жизнь вечную. В немощех сущия посети и исцеление даруй. Иже на мори управи. Путешествующим спутешествуй. Православным христианом споб\'{о}рствуй. Служащим и милующим нас грехов оставление даруй. Заповедавших нам недостойным молитися о них помилуй по велицей Твоей милости. Помяни, Господи, прежде усопших отец и братий наших и упокой их, идеже присещает свет лица Твоего. Помяни, Господи, братий наших плененных и избави я от всякаго обстояния. Помяни, Господи, плодоносящих и доброделающих во святых Твоих церквах, и даждь им яже ко спасению прошения и жизнь вечную. Помяни, Господи, и нас, смиренных и грешных и недостойных раб Твоих, и просвети наш ум светом разума Твоего, и настави нас на стезю заповедей Твоих, молитвами Пречистыя Владычицы нашея Богородицы и Приснодевы Марии и всех Твоих святых: яко благословен еси во веки веков. Аминь.

\mysubsubsection{Исповедание грехов повседневное}

Исповедаю Тебе Господу Богу моему и Творцу, во Святей Троице Единому, славимому и покланяемому, Отцу и Сыну и Святому Духу, вся моя грехи, яже содеях во вся дни живота моего, и на всякий час, и в настоящее время, и в прешедшия дни и нощи, делом, словом, помышлением, объядением, пиянством, тайноядением, празднословием, унынием, леностию, прекословием, непослушанием, оклеветанием, осуждением, небрежением, самолюбием, многостяжанием, хищением, неправдоглаголанием, скверноприбытчеством, мшелоимством, ревнованием, завистию, гневом, памятозлобием, ненавистию, лихоимством и всеми моими чувствы: зрением, слухом, обонянием, вкусом, осязанием и прочими моими грехи, душевными вкупе и телесными, имиже Тебе Бога моего и Творца прогневах, и ближняго моего онеправдовах: о сих жалея, винна себе Тебе Богу моему представляю, и имею волю каятися: точию, Господи Боже мой, помози ми, со слезами смиренно молю Тя: прешедшая же согрешения моя милосердием Твоим прости ми, и разреши от всех сих, яже изглаголах пред Тобою, яко Благ и Человеколюбец.

\mysubsubsection{\myemph{ Когда отходишь ко сну, произноси:}}

В руце Твои, Господи Иисусе Христе, Боже мой, предаю дух мой: Ты же мя благослови, Ты мя помилуй и живот вечный даруй ми. Аминь.

\end{mymulticols}

\mychapterending


\mychapter{Канон покаянный ко Господу нашему Иисусу Христу}\begin{mymulticols}
%http://www.molitvoslov.com/text3.htm 

\myfigure{748}

\mysubsubsection{Глас 6-й, Песнь 1}

\irmos{Яко по суху пешешествовав Израиль, по бездне стопами, гонителя фараона видя потопляема, Богу победную песнь поим, вопияше.}

\pripev{Помилуй мя, Боже, помилуй мя.}

Ныне приступих аз грешный и обремененный к Тебе, Владыце и Богу моему; не смею же взирати на небо, токмо молюся, глаголя: даждь ми, Господи, ум, да плачуся дел моих горько.

\pripev{Помилуй мя, Боже, помилуй мя.}

О, горе мне грешному! Паче всех человек окаянен есмь, покаяния несть во мне; даждь ми, Господи, слезы, да плачуся дел моих горько.

\slava

Безумне, окаянне человече, в лености время губиши; помысли житие твое, и обратися ко Господу Богу, и плачися о делех твоих горько.

\inyne

Мати Божия Пречистая, воззри на мя грешного, и от сети диаволи избави мя, и на путь покаяния настави мя, да плачуся дел моих горько.

\mysubsubsection{Песнь 3}

\irmos{Несть свят, якоже Ты, Господи Боже мой, вознесый рог верных Твоих, Блаже, и утвердивый нас на камени исповедания Твоего.}

\pripev{Помилуй мя, Боже, помилуй мя.}

Внегда поставлени будут престоли на судищи страшнем, тогда всех человек дела обличатся; горе тамо будет грешным, в муку отсылаемым; и то ведущи, душе моя, покайся от злых дел твоих.

\pripev{Помилуй мя, Боже, помилуй мя.}

Праведницы возрадуются, а грешнии восплачутся, тогда никтоже возможет помощи нам, но дела наша осудят нас, темже прежде конца покайся от злых дел твоих.

\slava

Увы мне великогрешному, иже делы и мысльми осквернився, ни капли слез имею от жестосердия; ныне возникни от земли, душе моя, и покайся от злых дел твоих.

\inyne

Се, взывает, Госпоже, Сын Твой, и поучает нас на доброе, аз же грешный добра всегда бегаю; но Ты, Милостивая, помилуй мя, да покаюся от злых моих дел.

\mysubsubsection{Седален, глас 6-й}

Помышляю день страшный и плачуся деяний моих лукавых: како отвещаю Безсмертному Царю, или коим дерзновением воззрю на Судию, блудный аз? Благоутробный Отче, Сыне Единородный и Душе Святый, помилуй мя.

Слава Отцу и Сыну и Святому Духу. И ныне и присно и во веки веков. Аминь.

\Bogorodichen{Связан многими ныне пленицами грехов и содержим лютыми страстьми и бедами, к Тебе прибегаю, моему спасению, и вопию: помози ми, Дево, Мати Божия.}

\mysubsubsection{Песнь 4}

\irmos{Христос моя сила, Бог и Господь, честная Церковь боголепно поет, взывающи от смысла чиста, о Господе празднующи.}

\pripev{Помилуй мя, Боже, помилуй мя.}

Широк путь зде и угодный сласти творити, но горько будет в последний день, егда душа от тела разлучатися будет: блюдися от сих, человече, Царствия ради Божия.

\pripev{Помилуй мя, Боже, помилуй мя.}

Почто убогаго обидиши, мзду наемничу удержуеши, брата твоего не любиши, блуд и гордость гониши? Остави убо сия, душе моя, и покайся Царствия ради Божия.

\slava

О, безумный человече, доколе углебаеши, яко пчела, собирающи богатство твое? Вскоре бо погибнет, яко прах и пепел: но более взыщи Царствия Божия.

\inyne

Госпоже Богородице, помилуй мя грешного, и в добродетели укрепи, и соблюди мя, да наглая смерть не похитит мя неготоваго, и доведи мя, Дево, Царствия Божия.

\mysubsubsection{Песнь 5}

\irmos{Божиим светом Твоим, Блаже, утренюющих Ти души любовию озари, молюся, Тя ведети, Слове Божий, истиннаго Бога, от мрака греховнаго взывающа.}

\pripev{Помилуй мя, Боже, помилуй мя.}

Воспомяни, окаянный человече, како лжам, клеветам, разбою, немощем, лютым зверем, грехов ради порабощен еси; душе моя грешная, того ли восхотела еси?

\pripev{Помилуй мя, Боже, помилуй мя.}

Трепещут ми уди, всеми бо сотворих вину: очима взираяй, ушима слышай, языком злая глаголяй, всего себе геенне предаяй; душе моя грешная, сего ли восхотела еси?

\slava

Блудника и разбойника кающася приял еси, Спасе, аз же един леностию греховною отягчихся и злым делом поработихся, душе моя грешная, сего ли восхотела еси?

\inyne

Дивная и скорая помощнице всем человеком, Мати Божия, помози мне недостойному, душа бо моя грешная того восхоте.

\mysubsubsection{Песнь 6}

\irmos{Житейское море, воздвизаемое зря напастей бурею, к тихому пристанищу Твоему притек, вопию Ти: возведи от тли живот мой, Многомилостиве.}

\pripev{Помилуй мя, Боже, помилуй мя.}

Житие на земли блудно пожих и душу во тьму предах, ныне убо молю Тя, Милостивый Владыко: свободи мя от работы сея вражия, и даждь ми разум творити волю Твою.

\pripev{Помилуй мя, Боже, помилуй мя.}

Кто творит таковая, якоже аз? Якоже бо свиния лежит в калу, тако и аз греху служу. Но Ты, Господи, исторгни мя от гнуса сего и даждь ми сердце творити заповеди Твоя.

\slava

Воспряни, окаянный человече, к Богу, воспомянув своя согрешения, припадая ко Творцу, слезя и стеня; Той же, яко милосерд, даст ти ум знати волю Свою.

\inyne

Богородице Дево, от видимаго и невидимаго зла сохрани мя, Пречистая, и приими молитвы моя, и донеси я Сыну Твоему, да даст ми ум творити волю Его.

\mysubsubsection{Кондак}

Душе моя, почто грехами богатееши, почто волю диаволю твориши, в чесом надежду полагаеши? Престани от сих и обратися к Богу с плачем, зовущи: милосерде Господи, помилуй мя грешнаго.

\mysubsubsection{Икос}

Помысли, душе моя, горький час смерти и страшный суд Творца твоего и Бога: Ангели бо грознии поймут тя, душе, и в вечный огнь введут: убо прежде смерти покайся, вопиющи: Господи, помилуй мя грешнаго.

\mysubsubsection{Песнь 7}

\irmos{Росодательну убо пещь содела Ангел преподобным отроком, халдеи же опаляющее веление Божие, мучителя увеща вопити: благословен еси, Боже отец наших.}

\pripev{Помилуй мя, Боже, помилуй мя.}

Не надейся, душе моя, на тленное богатство и на неправедное собрание, вся бо сия не веси кому оставиши, но возопий: помилуй мя, Христе Боже, недостойнаго.

\pripev{Помилуй мя, Боже, помилуй мя.}

Не уповай, душе моя, на телесное здравие и на скоромимоходящую красоту, видиши бо, яко сильнии и младии умирают; но возопий: помилуй мя, Христе Боже, недостойнаго.

\slava

Воспомяни, душе моя, вечное житие, Царство Небесное, уготованное святым, и тьму кромешную и гнев Божий злым, и возопий: помилуй мя, Христе Боже, недостойнаго.

\inyne

Припади, душе моя, к Божией Матери и помолися Той, есть бо скорая помощница кающимся, умолит Сына Христа Бога, и помилует мя недостойнаго.

\mysubsubsection{Песнь 8}

\irmos{Из пламене преподобным росу источил еси и праведнаго жертву водою попалил еси: вся бо твориши, Христе, токмо еже хотети. Тя превозносим во вся веки.}

\pripev{Помилуй мя, Боже, помилуй мя.}

Како не имам плакатися, егда помышляю смерть, видех бо во гробе лежаща брата моего, безславна и безобразна? Что убо чаю, и на что надеюся? Токмо даждь ми, Господи, прежде конца покаяние. \myemph{ (Дважды)}

\slava

Верую, яко приидеши судити живых и мертвых, и вси во своем чину станут, старии и младии, владыки и князи, девы и священницы; где обрящуся аз? Сего ради вопию: даждь ми, Господи, прежде конца покаяние.

\inyne

Пречистая Богородице, приими недостойную молитву мою и сохрани мя от наглыя смерти, и даруй ми прежде конца покаяние.

\mysubsubsection{Песнь 9}

\irmos{Бога человеком невозможно видети, на Негоже не смеют чини Ангельстии взирати; Тобою же, Всечистая, явися человеком Слово Воплощенно, Егоже величающе, с небесными вои Тя ублажаем.}

\pripev{Помилуй мя, Боже, помилуй мя.}

Ныне к вам прибегаю, Ангели, Архангели и вся небесныя силы, у Престола Божия стоящии, молитеся ко Творцу своему, да избавит душу мою от муки вечныя.

\pripev{Помилуй мя, Боже, помилуй мя.}

Ныне плачуся к вам, святии патриарси, царие и пророцы, апостоли и святителие и вси избраннии Христовы: помозите ми на суде, да спасет душу мою от силы вражия.

\slava

Ныне к вам воздежу руце, святии мученицы, пустынницы, девственницы, праведницы и вси святии, молящиися ко Господу за весь мир, да помилует мя в час смерти моея.

\inyne

Мати Божия, помози ми, на Тя сильне надеющемуся, умоли Сына Своего, да поставит мя недостойнаго одесную Себе, егда сядет судяй живых и мертвых, аминь.

\mysubsubsection{Молитва}

Господи Иисусе Христе, Сыне Божий, помилуй мя грешнаго.

Владыко Христе Боже, Иже страстьми Своими страсти моя исцеливый и язвами Своими язвы моя уврачевавый, даруй мне, много Тебе прегрешившему, слезы умиления; сраствори моему телу от обоняния Животворящаго Тела Твоего, и наслади душу мою Твоею Честною Кровию от горести, еюже мя сопротивник напои; возвыси мой ум к Тебе, долу поникший, и возведи от пропасти погибели: яко не имам покаяния, не имам умиления, не имам слезы утешительныя, возводящия чада ко своему наследию. Омрачихся умом в житейских страстех, не могу воззрети к Тебе в болезни, не могу согретися слезами, яже к Тебе любве. Но, Владыко Господи Иисусе Христе, сокровище благих, даруй мне покаяние всецелое и сердце люботрудное во взыскание Твое, даруй мне благодать Твою и обнови во мне зраки Твоего образа. Оставих Тя, не остави мене; изыди на взыскание мое, возведи к пажити Твоей и сопричти мя овцам избраннаго Твоего стада, воспитай мя с ними от злака Божественных Твоих Таинств, молитвами Пречистыя Твоея Матере и всех святых Твоих. Аминь.

\end{mymulticols}

\mychapterending


\mychapter{Канон молебный ко Пресвятой Богородице}\begin{mymulticols}
%http://www.molitvoslov.com/text4.htm 

\myfigure{3}

\mysubsubsection{Поемый во всякой скорби душевной и обстоянии.}

\mysubsubsection{Tворение Феостирикта монаха.}

\mysubsubsection{Тропaрь Богородице, глас 4-й}

К Богородице прилежно ныне притецем, грешнии и смиреннии, и припадем, в покаянии зовуще из глубины души: Владычице, помози, на ны милосердовавши, потщися, погибaем от множества прегрешений, не отврати Твоя рабы тщи, Тя бо и едину надежду имамы. \myemph{ (Дважды)}

Слава Отцу и Сыну и Святому Духу. И ныне и присно и во веки веков. Аминь.

Не умолчим никогда, Богородице, силы Твоя глаголати, недостойнии: aще бо Ты не бы предстояла молящи, кто бы нас избaвил от толиких бед, кто же бы сохранил до ныне свободны? Не отступим, Владычице, от Тебе: Твоя бо рабы спасaеши присно от всяких лютых.

\mysubsubsection{Псалом 50}

Помилуй мя, Боже, по велицей милости Твоей, и по множеству щедрот Твоих очисти беззаконие мое. Наипaче омый мя от беззакония моего, и от греха моего очисти мя; яко беззаконие мое аз знaю, и грех мой предо мною есть выну. Тебе единому согреших и лукaвое пред Тобою сотворих; яко да оправдишися во словесех Твоих, и победиши внегдa судити Ти. Се бо, в беззакониих зачaт есмь, и во гресех роди мя мaти моя. Се бо, истину возлюбил еси; безвестная и тайная премудрости Твоея явил ми еси. Окропиши мя иссопом, и очищуся; омыеши мя, и пaче снега убелюся. Слуху моему дaси рaдость и веселие; возрaдуются кости смиренныя. Отврати лице Твое от грех моих и вся беззакония моя очисти. Сердце чисто созижди во мне, Боже, и дух прав обнови во утробе моей. Не отвержи мене от лица Твоего и Духа Твоего Святaго не отыми от мене. Воздaждь ми рaдость спасения Твоего и Духом Владычним утверди мя. Научу беззаконныя путем Твоим, и нечестивии к Тебе обратятся. Избaви мя от кровей, Боже, Боже спасения моего; возрaдуется язык мой правде Твоей. Господи, устне мои отверзеши, и устa моя возвестят хвалу Твою. Яко aще бы восхотел еси жертвы, дал бых убо: всесожжения не благоволиши. Жертва Богу дух сокрушен; сердце сокрушенно и смиренно Бог не уничижит. Ублажи, Господи, благоволением Твоим Сиона, и да созиждутся стены Иерусалимския. Тогда благоволиши жертву правды, возношение и всесожигaемая; тогда возложaт на олтaрь Твой тельцы.

\mysubsubsection{Канон ко Пресвятой Богородице, глас 8-й}

\mysubsubsection{Песнь 1}

\irmos{Воду прошед яко сушу, и египетскаго зла избежaв, изрaильтянин вопияше: избaвителю и Богу нашему поим.}

\pripev{Пресвятая Богородице, спаси нас.}

Многими содержимь напaстьми, к Тебе прибегаю, спасения иский: о, Мaти Слова и Дево, от тяжких и лютых мя спаси.

\pripev{Пресвятая Богородице, спаси нас.}

Страстей мя смущaют прилози, многаго уныния исполнити мою душу; умири, Отроковице, тишиною Сына и Бога Твоего, Всенепорочная.

\slava

Спaса рождшую Тя и Бога, молю, Дево, избaвитися ми лютых: к Тебе бо ныне прибегaя, простирaю и душу и помышление.

\inyne

Недугующа телом и душею, посещения Божественнаго и промышления от Тебе сподоби, едина Богомaти, яко благая, Благaго же Родительница.

\mysubsubsection{Песнь 3}

\irmos{Небеснаго круга Верхотворче, Господи, и Церкве Зиждителю, Ты мене утверди в любви Твоей, желaний крaю, верных утверждение, едине Человеколюбче.}

\pripev{Пресвятая Богородице, спаси нас.}

Предстaтельство и покров жизни моея полагaю Тя, Богородительнице Дево: Ты мя окорми ко пристaнищу Твоему, благих виновна; верных утверждение, едина Всепетая.

\pripev{Пресвятая Богородице, спаси нас.}

Молю, Дево, душевное смущение и печали моея бурю разорити: Ты бо, Богоневестная, начальника тишины Христа родилa еси, едина Пречистая.

\slava

Благодетеля рождши добрых виновнаго, благодеяния богатство всем источи, вся бо можеши, яко сильнаго в крепости Христа рождши, Богоблаженная.

\inyne

Лютыми недуги и болезненными страстьми истязaему, Дево, Ты ми помози: исцелений бо неоскудное Тя знаю сокровище, Пренепорочная, неиждивaемое.

Спаси от бед рабы Твоя, Богородице, яко вси по Бозе к Тебе прибегaем, яко нерушимей стене и предстaтельству.

Призри благосердием, всепетая Богородице, на мое лютое телесе озлобление, и исцели души моея болезнь.

\mysubsubsection{Тропарь, глас 2-й}

Моление теплое и стенa необоримая, милости источниче, мирови прибежище, прилежно вопием Ти: Богородице Владычице, предвари, и от бед избaви нас, едина вскоре предстaтельствующая.

\mysubsubsection{Песнь 4}

\irmos{Услышах, Господи, смотрения Твоего тaинство, разумех дела Твоя и прослaвих Твое Божество.}

\pripev{Пресвятая Богородице, спаси нас.}

Страстей моих смущение, кормчию рождшая Господа, и бурю утиши моих прегрешений, Богоневестная.

\pripev{Пресвятая Богородице, спаси нас.}

Милосердия Твоего бездну призывaющу подaждь ми, яже Благосердаго рождшая и Спaса всех поющих Тя.

\pripev{Пресвятая Богородице, спаси нас.}

Наслаждaющеся, Пречистая, Твоих даровaний, благодaрственное воспевaем пение, ведуще Тя Богомaтерь.

\slava

На одре болезни моея и немощи низлежaщу ми, яко Благолюбива, помози, Богородице, едина Приснодево.

\inyne

Надежду и утверждение и спасения стену недвижиму имуще Тя, Всепетая, неудобства всякаго избавляемся.

\mysubsubsection{Песнь 5}

\irmos{Просвети нас повелении Твоими, Господи, и мышцею Твоею высокою Твой мир подaждь нам, Человеколюбче.}

\pripev{Пресвятая Богородице, спаси нас.}

Исполни, Чистая, веселия сердце мое, Твою нетленную дающи радость, веселия рождшая виновнаго.

\pripev{Пресвятая Богородице, спаси нас.}

Избaви нас от бед, Богородице чистая, вечное рождши избавление, и мир, всяк ум преимущий.

\slava

Разреши мглу прегрешений моих, Богоневесто, просвещением Твоея светлости, Свет рождшая Божественный и превечный.

\inyne

Исцели, Чистая, души моея неможение, посещения Твоего сподобльшая, и здрaвие молитвами Твоими подaждь ми.

\mysubsubsection{Песнь 6}

\irmos{Молитву пролию ко Господу, и Тому возвещу печали моя, яко зол душа моя исполнися, и живот мой аду приближися, и молюся яко Иона: от тли, Боже, возведи мя.}

\pripev{Пресвятая Богородице, спаси нас.}

Смерти и тли яко спасл есть, Сам Ся издaв смерти, тлением и смертию мое естество, ято бывшее, Дево, моли Господа и Сына Твоего, врагов злодействия мя избaвити.

\pripev{Пресвятая Богородице, спаси нас.}

Предстaтельницу Тя живота вем и хранительницу тверду, Дево, и напaстей решaщу молвы, и налоги бесов отгоняющу; и молюся всегда, от тли страстей моих избaвити мя.

\slava

Яко стену прибежища стяжaхом, и душ всесовершенное спасение, и прострaнство в скорбех, Отроковице, и просвещением Твоим присно рaдуемся: о, Владычице, и ныне нас от страстей и бед спаси.

\inyne

На одре ныне немощствуяй лежу, и несть исцеления плоти моей: но, Бога и Спaса миру и Избaвителя недугов рождшая, Тебе молюся, Благой: от тли недуг возстaви мя.

\mysubsubsection{Кондaк, глас 6-й}

Предстaтельство христиан непостыдное, ходaтайство ко Творцу непреложное, не презри грешных молений глaсы, но предвари, яко Благaя, на помощь нас, верно зовущих Ти; ускори на молитву, и потщися на умоление, предстaтельствующи присно, Богородице, чтущих Тя.

\mysubsubsection{Другой кондaк, глас тот же}

 Не имамы иныя помощи, не имамы иныя надежды, разве Тебе, Пречистая Дево. Ты нам помози, на Тебе надеемся, и Тобою хвaлимся, Твои бо есмы рабы, да не постыдимся.

\mysubsubsection{Стихира, глас тот же}

Не ввери мя человеческому предстaтельству, Пресвятая Владычице, но приими моление раба Твоего: скорбь бо обдержит мя, терпети не могу демонскаго стреляния, покрова не имам, ниже где прибегну, окаянный, всегда побеждaемь, и утешения не имам, разве Тебе, Владычице мира, уповaние и предстaтельство верных, не презри моление мое, полезно сотвори.

\mysubsubsection{Песнь 7}

\irmos{От Иудеи дошедше отроцы, в Вавилоне иногдa, верою Троическою плaмень пещный попрaша, поюще: отцев Боже, благословен еси.}

\pripev{Пресвятая Богородице, спаси нас.}

Наше спасение якоже восхотел еси, Спaсе, устроити, во утробу Девыя вселился еси, Юже миру предстaтельницу показал еси: отец наших Боже, благословен еси.

\pripev{Пресвятая Богородице, спаси нас.}

Волителя милости, Егоже родилa еси, Мaти чистая, умоли избaвитися от прегрешений и душевных скверн верою зовущим: отец наших Боже, благословен еси.

\slava

Сокровище спасения и Источник нетления, Тя рождшую, и столп утверждения, и дверь покаяния, зовущим показал еси: отец наших Боже, благословен еси.

\inyne

Телесныя слабости и душевныя недуги, Богородительнице, любовию приступaющих к крову Твоему, Дево, исцелити сподоби, Спaса Христа нам рождшая.

\mysubsubsection{Песнь 8}

\irmos{Царя Небеснаго, Егоже поют вои aнгельстии, хвалите и превозносите во вся веки.}

\pripev{Пресвятая Богородице, спаси нас.}

Помощи яже от Тебе требующия не презри, Дево, поющия и превозносящия Тя во веки.

\pripev{Пресвятая Богородице, спаси нас.}

Неможение души моея исцеляеши и телесныя болезни, Дево, да Тя прослaвлю, Чистая, во веки.

\slava

Исцелений богатство изливaеши верно поющим Тя, Дево, и превозносящим неизреченное Твое рождество.

\inyne

Напaстей Ты прилоги отгоняеши и страстей находы, Дево: темже Тя поем во вся веки.

\mysubsubsection{Песнь 9}

\irmos{Воистинну Богородицу Тя исповедуем, спасеннии Тобою, Дево чистая, с безплотными лики Тя величaюще.}

\pripev{Пресвятая Богородице, спаси нас.}

Тока слез моих не отвратися, Яже от всякаго лица всяку слезу отъемшаго, Дево, Христа рождшая.

\pripev{Пресвятая Богородице, спаси нас.}

Радости мое сердце исполни, Дево, Яже радости приемшая исполнение, греховную печаль потребляющи.

\pripev{Пресвятая Богородице, спаси нас.}

Пристaнище и предстaтельство к Тебе прибегaющих буди, Дево, и стена нерушимая, прибежище же и покров и веселие.

\slava

Света Твоего зарями просвети, Дево, мрак неведения отгоняющи, благоверно Богородицу Тя исповедающих.

\inyne

На месте озлобления немощи смирившагося, Дево, исцели, из нездрaвия во здрaвие претворяющи.

\mysubsubsection{Стихиры, глас 2-й}

Высшую небес и чистшую светлостей солнечных, избaвльшую нас от клятвы, Владычицу мира песньми почтим.

От многих моих грехов немощствует тело, немощствует и душа моя; к Тебе прибегaю, Благодaтней, надеждо ненадежных, Ты ми помози.

Владычице и Мaти Избaвителя, приими моление недостойных раб Твоих, да ходaтайствуеши к Рождшемуся от Тебе; о, Владычице мира, буди Ходaтаица!

Поем прилежно Тебе песнь ныне, всепетой Богородице, рaдостно: со Предтечею и всеми святыми моли, Богородице, еже ущедрити ны.

Вся aнгелов воинства, Предтече Господень, апостолов двоенадесятице, святии вси с Богородицею, сотворите молитву, во еже спастися нам.

\mysubsubsection{Молитвы ко Пресвятой Богородице}

Пресвятая Богородице, спаси мя.

Царице моя преблагaя, надеждо моя Богородице, приятелище сирых и странных предстaтельнице, скорбящих рaдосте, обидимых покровительнице! Зриши мою беду, зриши мою скорбь, помози ми яко немощну, окорми мя яко стрaнна. Обиду мою веси, разреши ту, яко волиши: яко не имам иныя помощи разве Тебе, ни иныя предстaтельницы, ни благия утешительницы, токмо Тебе, о Богомaти, яко да сохраниши мя и покрыеши во веки веков. Аминь.

К кому возопию, Владычице? К кому прибегну в горести моей, aще не к Тебе, Царице Небесная? Кто плач мой и воздыхaние мое приимет, aще не Ты, Пренепорочная, надеждо христиан и прибежище нам, грешным? Кто пaче Тебе в напaстех защитит? Услыши убо стенaние мое, и приклони ухо Твое ко мне, Владычице Мaти Бога моего, и не презри мене, требующаго Твоея помощи, и не отрини мене, грешнаго. Вразуми и научи мя, Царице Небесная; не отступи от мене, раба Твоего, Владычице, за роптaние мое, но буди мне Мaти и заступница. Вручaю себе милостивому покрову Твоему: приведи мя, грешнаго, к тихой и безмятежной жизни, да плaчуся о гресех моих. К кому бо прибегну повинный аз, aще не к Тебе, уповaнию и прибежищу грешных, надеждою на неизреченную милость Твою и щедроты Твоя окриляемь? О, Владычице Царице Небесная! Ты мне уповaние и прибежище, покров и заступление и помощь. Царице моя преблагaя и скорая заступнице! Покрый Твоим ходaтайством моя прегрешения, защити мене от враг видимых и невидимых; умягчи сердца злых человек, возстающих на мя. О, Мaти Господа моего Творцa! Ты еси корень девства и неувядaемый цвет чистоты. О, Богородительнице! Ты подaждь ми помощь немощствующему плотскими страстьми и болезнующему сердцем, едино бо Твое и с Тобою Твоего Сына и Бога нашего имам заступление; и Твоим пречудным заступлением да избaвлюся от всякия беды и напaсти, о пренепорочная и преслaвная Божия Мaти Марие. Темже со уповaнием глаголю и вопию: радуйся, благодaтная, радуйся, обрaдованная; радуйся, преблагословенная, Господь с Тобою.

\end{mymulticols}

\mychapterending


\mychapter{Канон Ангелу Хранителю}\begin{mymulticols}
%http://www.molitvoslov.com/text5.htm

\myfigure{8}

\mysubsubsection{Тропарь, глас 6-й}

Ангеле Божий, хранителю мой святый, живот мой соблюди во страсе Христа Бога, ум мой утверди во истиннем пути, и к любви горней уязви душу мою, да тобою направляемь, получу от Христа Бога велию милость.

Слава Отцу и Сыну и Святому Духу. И ныне и присно и во веки веков. Аминь.

\mysubsubsection{Богородичен}

Святая Владычице, Христа Бога нашего Мати, яко всех Творца недоуменно рождшая, моли благость Его всегда, со хранителем моим ангелом, спасти душу мою, страстьми одержимую, и оставление грехов даровати ми.

\mysubsubsection{Канон, глас 8-й}

\mysubsubsection{Песнь 1}

\irmos{Поим Господеви, проведшему люди Своя сквозе Чермное море, яко един славно прославися.}

\pripev[Иисусу:]{Господи Иисусе Христе Боже мой, помилуй мя.}

Песнь воспети и восхвалити, Спасе, Твоего раба достойно сподоби, безплотному Aнгелу, наставнику и хранителю моему.

\pripev{Святый Aнгеле Божий, хранителю мой, моли Бога о мне.}

Един аз в неразумии и в лености ныне лежу, наставниче мой и хранителю, не остави мене, погибающа.

\slava

Ум мой твоею молитвою направи, творити ми Божия повеления, да получу от Бога отдание грехов, и ненавидети ми злых настави мя, молюся ти.

\inyne

Молися, Девице, о мне, рабе Твоем, ко Благодателю, со хранителем моим Aнгелом, и настави мя творити заповеди Сына Твоего и Творца моего.

\mysubsubsection{Песнь 3}

\irmos{Ты еси утверждение притекающих к Тебе, Господи, Ты еси свет омраченных, и поет Тя дух мой.}

\pripev{Святый Aнгеле Божий, хранителю мой, моли Бога о мне.}

Все помышление мое и душу мою к тебе возложих, хранителю мой; ты от всякия мя напасти вражия избави.

\pripev{Святый Aнгеле Божий, хранителю мой, моли Бога о мне.}

Враг попирает мя, и озлобляет, и поучает всегда творити своя хотения; но ты, наставниче мой, не остави мене погибающа.

\slava

Пети песнь со благодарением и усердием Творцу и Богу даждь ми, и тебе, благому Aнгелу хранителю моему: избавителю мой, изми мя от враг озлобляющих мя.

\inyne

Исцели, Пречистая, моя многонедужныя струпы, яже в души, прожени враги, иже присно борются со мною.

\mysubsubsection{Седален, глас 2-й}

От любве душевныя вопию ти, хранителю моея души, всесвятый мой Aнгеле: покрый мя и соблюди от лукаваго ловления всегда, и к жизни настави небесней, вразумляя и просвещая и укрепляя мя.

Слава Отцу и Сыну и Святому Духу. И ныне и присно и во веки веков. Аминь.

\mysubsubsection{Богородичен:}

Богородице безневестная Пречистая, Яже без семени рождши всех Владыку, Того со Aнгелом хранителем моим моли, избавити ми ся всякаго недоумения, и дати умиление и свет души моей и согрешением очищение, Яже едина вскоре заступающи.

\mysubsubsection{Песнь 4}

\irmos{Услышах, Господи, смотрения Твоего таинство, разумех дела Твоя, и прославих Твое Божество.}

\pripev{Святый Aнгеле Божий, хранителю мой, моли Бога о мне.}

Моли Человеколюбца Бога ты, хранителю мой, и не остави мене, но присно в мире житие мое соблюди и подаждь ми спасение необоримое.

\pripev{Святый Aнгеле Божий, хранителю мой, моли Бога о мне.}

Яко заступника и хранителя животу моему прием тя от Бога, Aнгеле, молю тя, святый, от всяких мя бед свободи.

\slava

Мою скверность твоею святынею очисти, хранителю мой, и от части шуия да отлучен буду молитвами твоими и причастник славы явлюся.

\inyne

Недоумение предлежит ми от обышедших мя зол, Пречистая, но избави мя от них скоро: к Тебе бо единей прибегох.

\mysubsubsection{Песнь 5}

\irmos{Утренююще вопием Ти: Господи, спаси ны; Ты бо еси Бог наш, разве Тебе иного не вемы.}

\pripev{Святый Aнгеле Божий, хранителю мой, моли Бога о мне.}

Яко имея дерзновение к Богу, хранителю мой святый, Сего умоли от оскорбляющих мя зол избавити.

\pripev{Святый Aнгеле Божий, хранителю мой, моли Бога о мне.}

Свете светлый, светло просвети душу мою, наставниче мой и хранителю, от Бога данный ми Aнгеле.

\slava

Спяща мя зле тяготою греховною, яко бдяща сохрани, Aнгеле Божий, и возстави мя на славословие молением твоим.

\inyne

Марие, Госпоже Богородице безневестная, надеждо верных, вражия возношения низложи, поющия же Тя возвесели.

\mysubsubsection{Песнь 6}

\irmos{Ризу ми подаждь светлу, одеяйся светом яко ризою, многомилостиве Христе Боже наш.}

\pripev{Святый Aнгеле Божий, хранителю мой, моли Бога о мне.}

Всяких мя напастей свободи, и от печалей спаси, молюся ти, святый Aнгеле, данный ми от Бога, хранителю мой добрый.

\pripev{Святый Aнгеле Божий, хранителю мой, моли Бога о мне.}

Освети ум мой, блаже, и просвети мя, молюся ти, святый Aнгеле, и мыслити ми полезная всегда настави мя.

\slava

Устави сердце мое от настоящаго мятежа, и бдети укрепи мя во благих, хранителю мой, и настави мя чудно к тишине животней.

\inyne

Слово Божие в Тя вселися, Богородице, и человеком Тя показа небесную лествицу; Тобою бо к нам Вышний сошел есть.

\mysubsubsection{Кондак, глас 4-й}

Явися мне милосерд, святый Aнгеле Господень, хранителю мой, и не отлучайся от мене, сквернаго, но просвети мя светом неприкосновенным и сотвори мя достойна Царствия Небеснаго.

\mysubsubsection{Икос}

Уничиженную душу мою многими соблазны, ты, святый предстателю, неизреченныя славы небесныя сподоби, и певец с лики безплотных сил Божиих, помилуй мя и сохрани, и помыслы добрыми душу мою просвети, да твоею славою, Aнгеле мой, обогащуся, и низложи зломыслящия мне враги, и сотвори мя достойна Царствия Небеснаго.

\mysubsubsection{Песнь 7}

\irmos{От Иудеи дошедше отроцы, в Вавилоне иногда, верою Троическою пламень пещный попраша, поюще: отцев Боже, благословен еси.}

\pripev{Святый Aнгеле Божий, хранителю мой, моли Бога о мне.}

Милостив буди ми, и умоли Бога, Господень Aнгеле, имею бо тя заступника во всем животе моем, наставника же и хранителя, от Бога дарованнаго ми во веки.

\pripev{Святый Aнгеле Божий, хранителю мой, моли Бога о мне.}

Не остави в путь шествующия души моея окаянныя убити разбойником, святый Aнгеле, яже ти от Бога предана бысть непорочне; но настави ю на путь покаяния.

\slava

Всю посрамлену душу мою привожду от лукавых ми помысл и дел: но предвари, наставниче мой, и исцеление ми подаждь благих помысл, уклоняти ми ся всегда на правыя стези.

\inyne

Премудрости исполни всех и крепости Божественныя, Ипостасная Премудросте Вышняго, Богородицы ради, верою вопиющих: отец наших Боже, благословен еси.

\mysubsubsection{Песнь 8}

\irmos{Царя Небеснаго, Егоже поют вои ангельстии, хвалите и превозносите во вся веки.}

\pripev{Святый Aнгеле Божий, хранителю мой, моли Бога о мне.}

От Бога посланный, утверди живот мой, раба твоего, преблагий Aнгеле, и не остави мене во веки.

\pripev{Святый Aнгеле Божий, хранителю мой, моли Бога о мне.}

Ангела тя суща блага, души моея наставника и хранителя, преблаженне, воспеваю во веки.

\slava

Буди ми покров и забрало в день испытания всех человек, воньже огнем искушаются дела благая же и злая.

\inyne

Буди ми помощница и тишина, Богородице Приснодево, рабу Твоему, и не остави мене лишена быти Твоего владычества.

\mysubsubsection{Песнь 9}

\irmos{Воистинну Богородицу Тя исповедуем, спасеннии Тобою, Дево чистая, с безплотными лики Тя величающе.}

\pripev[Иисусу:]{Господи Иисусе Христе Боже мой, помилуй мя.}

Помилуй мя, едине Спасе мой, яко милостив еси и милосерд, и праведных ликов сотвори мя причастника.

\pripev{Святый Aнгеле Божий, хранителю мой, моли Бога о мне.}

Мыслити ми присно и творити, Господень Aнгеле, благая и полезная даруй, яко сильна яви в немощи и непорочна.

\slava

Яко имея дерзновение к Царю Небесному, Того моли, с прочими безплотными, помиловати мя, окаяннаго.

\inyne

Много дерзновение имущи, Дево, к Воплощшемуся из Тебе, преложи мя от уз и разрешение ми подаждь и спасение, молитвами Твоими.

\mysubsubsection{Молитва к Aнгелу Xранителю}

Святый Aнгеле Божий, хранителю мой, моли Бога о мне.

Ангеле Христов святый, к тебе припадая молюся, хранителю мой святый, приданный мне на соблюдение души и телу моему грешному от святаго крещения, аз же своею леностию и своим злым обычаем прогневах твою пречистую светлость и отгнах тя от себе всеми студными делы: лжами, клеветами, завистию, осуждением, презорством, непокорством, братоненавидением, и злопомнением, сребролюбием, прелюбодеянием, яростию, скупостию, объядением без сытости и опивством, многоглаголанием, злыми помыслы и лукавыми, гордым обычаем и блудным возбешением, имый самохотение на всякое плотское вожделение. О, злое мое произволение, егоже и скоти безсловеснии не творят! Да како возможеши воззрети на мя, или приступити ко мне, аки псу смердящему? Которыма очима, ангеле Христов, воззриши на мя, оплетшася зле во гнусных делех? Да како уже возмогу отпущения просити горьким и злым моим и лукавым деянием, в няже впадаю по вся дни и нощи и на всяк час? Но молюся ти припадая, хранителю мой святый, умилосердися на мя грешнаго и недостойнаго раба твоего \myemph{ (имя)}, буди ми помощник и заступник на злаго моего сопротивника, святыми твоими молитвами, и Царствия Божия причастника мя сотвори со всеми святыми, всегда, и ныне и присно и во веки веков. Аминь.

\end{mymulticols}

\mychapterending


\mychapter{Последование ко Святому Причащению}\begin{mymulticols}
%http://www.molitvoslov.com/text207.htm

\myfigure{132}

Молитвами святых отец наших, Господи Иисусе Христе Боже наш, помилуй нас. Аминь.

Царю Небесный, Утешителю, Душе истины, Иже везде сый и вся исполняяй, Сокровище благих и жизни Подателю, прииди и вселися в ны, и очисти ны от всякия скверны, и спаси, Блаже, души наша.

Святый Боже, Святый Крепкий, Святый Безсмертный, помилуй нас. \myemph{ (Tрижды)}

Слава Отцу и Сыну и Святому Духу, и ныне и присно и во веки веков. Аминь.

Пресвятая Троице, помилуй нас; Господи, очисти грехи наша; Владыко, прости беззакония наша; Святый, посети и исцели немощи наша, имене Твоего ради.

Господи, помилуй. \myemph{ (Трижды)}

Слава Отцу и Сыну и Святому Духу, и ныне и присно и во веки веков. Аминь.

Отче наш, Иже еси на небесех! Да святится имя Твое, да приидет Царствие Твое, да будет воля Твоя, яко на небеси и на земли. Хлеб наш насущный даждь нам днесь; и остави нам долги наша, якоже и мы оставляем должником нашим; и не введи нас во искушение, но избави нас от лукаваго.

Господи, помилуй.\myemph{  (12раз)}

Приидите, поклонимся Цареви нашему Богу. \myemph{ (Поклон)}

Приидите, поклонимся и припадем Христу, Цареви нашему Богу. \myemph{ (Поклон)}

Приидите, поклонимся и припадем Самому Христу, Цареви и Богу нашему.\myemph{ (Поклон)}

\mysubsubsection{Псалом 22}

Господь пасет мя, и ничтоже мя лишит. На месте злачне, тамо всели мя, на воде покойне воспита мя. Душу мою обрати, настави мя на стези правды, имене ради Своего. Аще бо и пойду посреде сени смертныя, не убоюся зла, яко Ты со мною еси, жезл Твой и палица Твоя, та мя утешиста. Уготовал еси предо мною трапезу сопротив стужающим мне, умастил еси елеом главу мою, и чаша Твоя упоявающи мя, яко державна. И милость Твоя поженет мя вся дни живота моего, и еже вселити ми ся в дом Господень, в долготу дний.

\mysubsubsection{Псалом 23}

Господня земля, и исполнение ея, вселенная, и вси живущии на ней. Той на морях основал ю есть, и на реках уготовал ю есть. Кто взыдет на гору Господню? Или кто станет на месте святем Его? Неповинен рукама и чист сердцем, иже не прият всуе душу свою, и не клятся лестию искреннему своему. Сей приимет благословение от Господа, и милостыню от Бога, Спаса своего. Сей род ищущих Господа, ищущих лице Бога Иаковля. Возмите врата князи ваша, и возмитеся врата вечная; и внидет Царь Славы. Кто есть сей Царь Славы? Господь крепок и силен, Господь силен в брани. Возмите врата князи ваша, и возмитеся врата вечная, и внидет Царь Славы. Кто есть сей Царь Славы? Господь сил, Той есть Царь Славы.

\mysubsubsection{Псалом 115}

Веровах, темже возглаголах, аз же смирихся зело. Аз же рех во изступлении моем: всяк человек ложь. Что воздам Господеви о всех, яже воздаде ми? Чашу спасения прииму, и имя Господне призову, молитвы моя Господеви воздам пред всеми людьми Его. Честна пред Господем смерть преподобных Его. О, Господи, аз раб Твой, аз раб Твой и сын рабыни Твоея; растерзал еси узы моя. Тебе пожру жертву хвалы, и во имя Господне призову. Молитвы моя Господеви воздам пред всеми людьми Его, во дворех дому Господня, посреди тебе, Иерусалиме.

Слава Отцу и Сыну и Святому Духу, и ныне и присно и во веки веков. Аминь.

Аллилуия. \myemph{ (Трижды с тремя поклонами)}

\mysubsubsection{Тропари, глас 8-й}

Беззакония моя презри, Господи, от Девы рождейся, и сердце мое очисти, храм то творя пречистому Твоему Телу и Крови, ниже отрини мене от Твоего лица, без числа имеяй велию милость.

\slava

Во причастие святынь Твоих како дерзну [вниду], недостойный? Аше бо дерзну к Тебе приступити с достойными, хитон мя обличает, яко несть вечерний, и осуждение исходатайствую многогрешной души моей. Очисти, Господи, скверну души моея, и спаси мя, яко Человеколюбец.

\inyne

Многая множества моих, Богородице, прегрешений, к Тебе прибегох, Чистая, спасения требуя: посети немощствующую мою душу, и моли Сына Твоего и Бога нашего, дати ми оставление, яже содеях лютых, Едина благословенная.

\mysubsubsection{[Во Святую же Четыредесятницу:}

Егда славнии ученицы на умовении вечери просвещахуся, тогда Иуда злочестивый сребролюбием недуговав омрачашеся, и беззаконным судиям Тебе праведнаго Судию предает. Виждь, имений рачителю, сих ради удавление употребивша: бежи несытыя души, Учителю таковая дерзнувшия. Иже о всех благий Господи, слава Тебе. ]

\mysubsubsection{Псалом 50}

Помилуй мя, Боже, по велицей милости Твоей, и по множеству щедрот Твоих очисти беззаконие мое. Наипаче омый мя от беззакония моего, и от греха моего очисти мя; яко беззаконие мое аз знаю, и грех мой предо мною есть выну. Тебе единому согреших и лукавое пред Тобою сотворих; яко да оправдишися во словесех Твоих, и победиши внегда судити Ти. Се бо, в беззакониих зачат есмь, и во гресех роди мя мати моя. Се бо, истину возлюбил еси; безвестная и тайная премудрости Твоея явил ми еси. Окропиши мя иссопом, и очищуся; омыеши мя, и паче снега убелюся. Слуху моему даси радость и веселие; возрадуются кости смиренныя. Отврати лице Твое от грех моих и вся беззакония моя очисти. Сердце чисто созижди во мне, Боже, и дух прав обнови во утробе моей. Не отвержи мене от лица Твоего и Духа Твоего Святаго не отыми от мене. Воздаждь ми радость спасения Твоего и Духом Владычним утверди мя. Научу беззаконныя путем Твоим, и нечестивии к Тебе обратятся. Избави мя от кровей, Боже, Боже спасения моего; возрадуется язык мой правде Твоей. Господи, устне мои отверзеши, и уста моя возвестят хвалу Твою. Яко аще бы восхотел еси жертвы, дал бых убо: всесожжения не благоволиши. Жертва Богу дух сокрушен; сердце сокрушенно и смиренно Бог не уничижит. Ублажи, Господи, благоволением Твоим Сиона, и да созиждутся стены Иерусалимския. Тогда благоволиши жертву правды, возношение и всесожегаемая; тогда возложат на олтарь Твой тельцы.

\mysubsubsection{Канон, глас 2-й}

\mysubsubsection{Песнь 1}

\irmos{Грядите людие, поим песнь Христу Богу, раздельшему море, и наставльшему люди, яже изведе из работы египетския, яко прославися.}

\pripev{Сердце чисто созижди во мне, Боже, и дух прав обнови во утробе моей.}

Хлеб живота вечнующаго да будет ми Тело Твое Святое, благоутробне Господи, и Честная Кровь, и недуг многообразных исцеление.

\pripev{Не отвержи мене от лица Твоего, и Духа Твоего Святаго не отыми от мене.}

Осквернен делы безместными окаянный, Твоего Пречистаго Тела и Божественныя Крове недостоин есмь, Христе, причащения, егоже мя сподоби.

\pripev{Пресвятая Богородице, спаси нас.}

\Bogorodichen{Земле благая, благословенная Богоневесто, клас прозябшая неоранный и спасительный миру, сподоби мя сей ядуща спастися.}

\mysubsubsection{Песнь 3}

\irmos{На камени мя веры утвердив, разширил еси уста моя на враги моя. Возвесели бо ся дух мой, внегда пети: несть свят, якоже Бог наш, и несть праведен паче Тебе, Господи.}

\pripev{Сердце чисто созижди во мне, Боже, и дух прав обнови во утробе моей.}

Слезныя ми подаждь, Христе, капли, скверну сердца моего очищающия: яко да благою совестию очищен, верою прихожду и страхом, Владыко, ко причащению Божественных Даров Твоих.

\pripev{Не отвержи мене от лица Твоего, и Духа Твоего Святаго не отыми от мене.}

Во оставление да будет ми прегрешений Пречистое Тело Твое, и Божественная Кровь, Духа же Святаго общение, и в жизнь вечную, Человеколюбче, и страстей и скорбей отчуждение.

\pripev{Пресвятая Богородице, спаси нас.}

\Bogorodichen{Хлеба животнаго Tрапеза Пресвятая, свыше милости ради сшедшаго, и мирови новый живот дающаго, и мене ныне сподоби недостойнаго, со страхом вкусити сего, и живу быти.}

\mysubsubsection{Песнь 4}

\irmos{Пришел еси от Девы, не ходатай, ни Ангел, но Сам, Господи, воплощься, и спася еси всего мя человека. Тем зову Ти: слава силе Твоей, Господи.}

\pripev{Сердце чисто созижди во мне, Боже, и дух прав обнови во утробе моей.}

Восхотел еси, нас ради воплощься, Многомилостиве, заклан быти яко овча, грех ради человеческих: темже молю Тя, и моя очисти согрешения.

\pripev{Не отвержи мене от лица Твоего, и Духа Твоего Святаго не отыми от мене.}

Исцели души моея язвы, Господи, и всего освяти: и сподоби, Владыко, яко да причащуся тайныя Твоея Божественныя вечери, окаянный.

\pripev{Пресвятая Богородице, спаси нас.}

\Bogorodichen{Умилостиви и мне сущаго от утробы Твоея, Владычице, и соблюди мя нескверна раба Твоего и непорочна, яко да прием умнаго бисера, освящуся.}

\mysubsubsection{Песнь 5}

\irmos{Света Подателю и веков Творче, Господи, во свете Твоих повелений настави нас; разве бо Тебе иного бога не знаем.}

\pripev{Сердце чисто созижди во мне, Боже, и дух прав обнови во утробе моей.}

Якоже предрекл еси, Христе, да будет убо худому рабу Твоему, и во мне пребуди, якоже обещался еси: се бо Тело Твое ям Божественное, и пию Кровь Твою.

\pripev{Не отвержи мене от лица Твоего, и Духа Твоего Святаго не отыми от мене.}

Слове Божий и Боже, угль Тела Твоего да будет мне помраченному в просвещение, и очищение оскверненной души моей Кровь Твоя.

\pripev{Пресвятая Богородице, спаси нас.}

\Bogorodichen{Марие, Мати Божия, благоухания честное селение, Твоими молитвами сосуд мя избранный соделай, яко да освящений причащуся Сына Твоего.}

\mysubsubsection{Песнь 6}

\irmos{В бездне греховней валяяся, неизследную милосердия Твоего призываю бездну: от тли, Боже, мя возведи.}

\pripev{Сердце чисто созижди во мне, Боже, и дух прав обнови во утробе моей.}

Ум, душу и сердце освяти, Спасе, и тело мое, и сподоби неосужденно, Владыко, к страшным Тайнам приступити.

\pripev{Не отвержи мене от лица Твоего, и Духа Твоего Святаго не отыми от мене.}

Да бых устранился от страстей, и Твоея благодати имел бы приложение, живота же утверждение, причащением Святых, Христе, Таин Твоих.

\pripev{Пресвятая Богородице, спаси нас.}

\Bogorodichen{Божие, Боже, Слово Святое, всего мя освяти, ныне приходящаго к Божественным Твоим Тайнам, Святыя Матере Твоея мольбами.}

\mysubsubsection{Кондак, глас 2-й}

Хлеб, Христе, взяти не презри мя, Тело Твое, и Божественную Твою ныне Кровь, пречистых, Владыко, и страшных Твоих Таин причаститися окаяннаго, да не будет ми в суд, да будет же ми в живот вечный и безсмертный.

\mysubsubsection{Песнь 7}

\irmos{Телу златому премудрыя дети не послужиша, и в пламень сами поидоша, и боги их обругаша, среди пламене возопиша, и ороси я Ангел: услышася уже уст ваших молитва.}

\pripev{Сердце чисто созижди во мне, Боже, и дух прав обнови во утробе моей.}

Источник благих, причащение, Христе, безсмертных Твоих ныне Таинств да будет ми свет, и живот, и безстрастие, и к преспеянию же и умножению добродетели Божественнейшия ходатайственно, едине Блаже, яко да славлю Тя.

\pripev{Не отвержи мене от лица Твоего, и Духа Твоего Святаго не отыми от мене.}

Да избавлюся от страстей, и врагов, и нужды, и всякия скорби, трепетом и любовию со благоговением, Человеколюбче, приступаяй ныне к Твоим безсмертным и Божественным Тайнам, и пети Тебе сподоби: благословен еси, Господи, Боже отец наших.

\pripev{Пресвятая Богородице, спаси нас.}

\Bogorodichen{Спаса Христа рождшая паче ума, Богоблагодатная, молю Тя ныне, раб Твой, Чистую нечистый: хотящаго мя ныне к пречистым Тайнам приступити, очисти всего от скверны плоти и духа.}

\mysubsubsection{Песнь 8}

\irmos{В пещь огненную ко отроком еврейским снизшедшаго, и пламень в росу преложшаго Бога, пойте дела яко Господа, и превозносите во вся веки.}

\pripev{Сердце чисто созижди во мне, Боже, и дух прав обнови во утробе моей.}

Небесных, и страшных, и святых Твоих, Христе, ныне Таин, и Божественныя Твоея и тайныя вечери общника быти и мене сподоби отчаяннаго, Боже, Спасе мой.

\pripev{Не отвержи мене от лица Твоего, и Духа Твоего Святаго не отыми от мене.}

Под Твое прибег благоутробие, Блаже, со страхом зову Ти: во мне пребуди, Спасе, и аз, якоже рекл еси, в Тебе; се бо дерзая на милость Твою, ям Тело Твое, и пию Кровь Твою.

\pripev{Пресвятая Троице, Боже наш, слава Тебе.}

\myemph{ Троичен:} Трепещу, приемля огнь, да не опалюся яко воск и яко трава; oле страшнаго таинства! oле благоутробия Божия! Како Божественнаго Тела и Крове брение причащаюся, и нетленен сотворяюся?

\mysubsubsection{Песнь 9}

\irmos{Безначальна Родителя Сын, Бог и Господь, воплощся от Девы нам явися, омраченная просветити, собрати расточенная: тем всепетую Богородицу величаем.}

\pripev{Сердце чисто созижди во мне, Боже, и дух прав обнови во утробе моей.}

Христос есть, вкусите и видите: Господь нас ради, по нам бо древле бывый, единою Себе принес, яко приношение Отцу Своему, присно закалается, освящаяй причащающияся.

\pripev{Не отвержи мене от лица Твоего, и Духа Твоего Святаго не отыми от мене.}

Душею и телом да освящуся, Владыко, да просвещуся, да спасуся, да буду дом Твой причащением священных Таин, живущаго Тя имея в себе со Отцем и Духом, Благодетелю Многомилостиве.

\pripev{Воздаждь ми радость спасения Твоего и Духом Владычним утверди мя.}

Якоже огнь да будет ми, и яко свет, Тело Твое и Кровь, Спасе мой, пречестная, опаляя греховное вещество, сжигая же страстей терние, и всего мя просвещая, покланятися Божеству Твоему.

\pripev{Пресвятая Богородице, спаси нас.}

\Bogorodichen{Бог воплотися от чистых кровей Твоих; темже всякий род поет Тя, Владычице, умная же множества славят, яко Тобою яве узреша всеми Владычествующаго, осуществовавшагося человечеством.}

\mysubsubsection{Далее}

Достойно есть яко воистину блажити Тя Богородицу, Присноблаженную и Пренепорочную и Матерь Бога нашего. Честнейшую Херувим и славнейшую без сравнения Серафим, без истления Бога Слова рождшую, сущую Богородицу Тя величаем.

Святый Боже, Святый Крепкий, Святый Безсмертный, помилуй нас. \myemph{ (Tрижды)}

Слава Отцу и Сыну и Святому Духу, и ныне и присно и во веки веков. Аминь.

Пресвятая Троице, помилуй нас; Господи, очисти грехи наша; Владыко, прости беззакония наша; Святый, посети и исцели немощи наша, имене Твоего ради.

Господи, помилуй. \myemph{ (Трижды)}

Слава Отцу и Сыну и Святому Духу, и ныне и присно и во веки веков. Аминь.

Отче наш, Иже еси на небесех! Да святится имя Твое, да приидет Царствие Твое, да будет воля Твоя, яко на небеси и на земли. Хлеб наш насущный даждь нам днесь; и остави нам долги наша, якоже и мы оставляем должником нашим; и не введи нас во искушение, но избави нас от лукаваго.

\myemph{ Если неделя, тропарь воскресный по гласу. Если же нет, настоящие тропари, глас 6-й:}

Помилуй нас, Господи, помилуй нас; всякаго бо ответа недоумеюще, сию Ти молитву, яко Владыце, грешнии приносим: помилуй нас.

\slava

Господи, помилуй нас, на Тя бо уповахом; не прогневайся на ны зело, ниже помяни беззаконий наших, но призри и ныне яко благоутробен, и избави ны от враг наших. Ты бо еси Бог наш, и мы людие Твои, вси дела руку Твоею, и имя Твое призываем.

\inyne

Милосердия двери отверзи нам, благословенная Богородице, надеющиися на Тя да не погибнем, но да избавимся Тобою от бед: Ты бо еси спасение рода христианскаго.

Господи, помилуй. \myemph{ (40 раз) И поклоны, сколько хочешь.}

\myemph{ И стихи:}

Хотя ясти, человече, Тело Владычне,

Страхом приступи, да не опалишися: огнь бо есть.

Божественную же пия Кровь ко общению,

Первее примирися тя опечалившим.

Таже дерзая, таинственное брашно яждь.

Прежде причастия страшныя жертвы,

Животворящаго Тела Владычня,

Сим помолися образом со трепетом:

\mysubsubsection{Молитва 1-я, Василия Великого}

Владыко Господи Иисусе Христе, Боже наш, Источниче жизни и безсмертия, всея твари видимыя и невидимыя Содетелю, безначальнаго Отца соприсносущный Сыне и собезначальный, премногия ради благости в последния дни в плоть оболкийся, и распныйся, и погребыйся за ны неблагодарныя и злонравныя, и Твоею Кровию обновивый растлевшее грехом естество наше, Сам, Безсмертный Царю, приими и мое грешнаго покаяние, и приклони ухо Твое мне, и услыши глаголы моя. Согреших бо, Господи, согреших на небо и пред Тобою, и несмь достоин воззрети на высоту славы Твоея: прогневах бо Твою благость, Твоя заповеди преступив, и не послушав Твоих повелений. Но Ты, Господи, незлобив сый, долготерпелив же и многомилостив, не предал еси мя погибнути со беззаконьми моими, моего всячески ожидая обращения. Ты бо рекл еси, Человеколюбче, пророком Твоим: яко хотением не хощу смерти грешника, но еже обратится и живу быти ему. Не хощеши бо, Владыко, создания Твоею руку погубити, ниже благоволиши о погибели человечестей, но хощеши всем спастися, и в разум истины приити. Темже и аз, аще и недостоин есмь небесе и земли, и сея привременныя жизни, всего себе повинув греху, и сластем поработив, и Твой осквернив образ; но творение и создание Твое быв, не отчаяваю своего спасения окаянный, на Твое же безмерное благоутробие дерзая, прихожду. Приими убо и мене, Человеколюбче Господи, якоже блудницу, яко разбойника, яко мытаря и яко блуднаго, и возми мое тяжкое бремя грехов, грех вземляй мира, и немощи человеческия исцеляяй, труждающияся и обремененныя к Себе призываяй и упокоеваяй, не пришедый призвати праведныя, но грешныя на покаяние. И очисти мя от всякия скверны плоти и духа, и научи мя совершати святыню во страсе Твоем: яко да чистым сведением совести моея, святынь Твоих часть приемля, соединюся святому Телу Твоему и Крови, и имею Тебе во мне живуща и пребывающа, со Отцем, и Святым Твоим Духом. Ей, Господи Иисусе Христе, Боже мой, и да не в суд ми будет причастие пречистых и животворящих Таин Твоих, ниже да немощен буду душею же и телом, от еже недостойне тем причащатися, но даждь ми, даже до конечнаго моего издыхания, неосужденно восприимати часть святынь Твоих, в Духа Святаго общение, в напутие живота вечнаго, и во благоприятен ответ на Страшнем судищи Твоем: яко да и аз со всеми избранными Твоими общник буду нетленных Твоих благ, яже уготовал еси любящим Тя, Господи, в нихже препрославлен еси во веки. Аминь.

\mysubsubsection{Молитва 2-я , святого Иоанна Златоустого}

Господи Боже мой, вем, яко несмь достоин, ниже доволен, да под кров внидеши храма души моея, занеже весь пуст и пался есть, и не имаши во мне места достойна еже главу подклонити: но якоже с высоты нас ради смирил еси Себе, смирися и ныне смирению моему; и якоже восприял еси в вертепе и в яслех безсловесных возлещи, сице восприими и в яслех безсловесныя моея души, и во оскверненное мое тело внити. И якоже не неудостоил еси внити, и свечеряти со грешники в дому Симона прокаженнаго, тако изволи внити и в дом смиренныя моея души, прокаженныя и грешныя; и якоже не отринул еси подобную мне блудницу и грешную, пришедшую и прикоснувшуюся Тебе, сице умилосердися и о мне грешнем, приходящем и прикасающем Ти ся; и якоже не возгнушался еси скверных ея уст и нечистых, целующих Тя, ниже моих возгнушайся сквернших оныя уст и нечистших, ниже мерзких моих и нечистых устен, и сквернаго и нечистейшаго моего языка. Но да будет ми угль пресвятаго Твоего Тела, и честныя Твоея Крове, во освящение и просвещение и здравие смиренней моей души и телу, во облегчение тяжестей многих моих согрешений, в соблюдение от всякаго диавольскаго действа, во отгнание и возбранение злаго моего и лукаваго обычая, во умерщвление страстей, в снабдение заповедей Твоих, в приложение Божественныя Твоея благодати, и Твоего Царствия присвоение. Не бо яко презираяй прихожду к Тебе, Христе Боже, но яко дерзая на неизреченную Твою благость, и да не на мнозе удаляяйся общения Твоего, от мысленнаго волка звероуловлен буду. Темже молюся Тебе: яко един сый Свят, Владыко, освяти мою душу и тело, ум и сердце, чревеса и утробы, и всего мя обнови, и вкорени страх Твой во удесех моих, и освящение Твое неотъемлемо от мене сотвори; и буди ми помощник и заступник, окормляя в мире живот мой, сподобляя мя и одесную Тебе предстояния со святыми Твоими, молитвами и моленьми Пречистыя Твоея Матере, невещественных Твоих служителей и пречистых сил, и всех святых, от века Тебе благоугодивших. Аминь.

\mysubsubsection{Молитва 3-я, Симеона Метафраста}

Едине чистый и нетленный Господи, за неизреченную милость человеколюбия наше все восприемый смешение, от чистых и девственных кровей паче естества рождшия Тя, Духа Божественнаго нашествием, и благоволением Отца присносущнаго, Христе Иисусе, премудросте Божия, и мире, и сило; Твоим восприятием животворящая и спасительная страдания восприемый, крест, гвоздия, копие, смерть, умертви моя душетленныя страсти телесныя. Погребением Твоим адова пленивый царствия, погреби моя благими помыслы лукавая советования, и лукавствия духи разори. Тридневным Твоим и живоносным воскресением падшаго праотца возставивый, возстави мя грехом поползшагося, образы мне покаяния предлагая. Преславным Твоим вознесением плотское обоживый восприятие, и сие десным Отца седением почтый, сподоби мя причастием святых Твоих Таин десную часть спасаемых получити. Снитием Утешителя Твоего Духа сосуды честны священныя Твоя ученики соделавый, приятелище и мене покажи Того пришествия. Хотяй паки прийти судити вселенней правдою, благоволи и мне усрести Тя на облацех, Судию и Создателя моего, со всеми святыми Твоими: да безконечно славословлю и воспеваю Тя, со безначальным Твоим Отцем, и Пресвятым и Благим и Животворящим Твоим Духом, ныне и присно, и во веки веков. Аминь.

\mysubsubsection{Молитва 4-я, святого Иоанна Дамаскина}

Владыко Господи Иисусе Христе, Боже наш, едине имеяй власть человеком оставляти грехи, яко благ и Человеколюбец презри моя вся в ведении и не в ведении прегрешения, и сподоби мя неосужденно причаститися Божественных, и преславных, и пречистых, и животворящих Твоих Таин, не в тяжесть, ни в муку, ни в приложение грехов, но во очищение, и освящение, и обручение будущаго Живота и царствия, в стену и помощь, и в возражение сопротивных, во истребление многих моих согрешений. Ты бо еси Бог милости, и щедрот, и человеколюбия, и Тебе славу возсылаем, со Отцем, и Святым Духом, ныне и присно, и во веки веков. Аминь.

\mysubsubsection{Молитва 5-я, святого Василия Великого}

Вем, Господи, яко недостойне причащаюся пречистаго Твоего Тела и честныя Твоея Крове, и повинен есмь, и суд себе ям и пию, не разсуждая Тела и Крове Тебе Христа и Бога моего, но на щедроты Твоя дерзая прихожду к Тебе рекшему: ядый Мою плоть, и пияй Мою кровь, во Мне пребывает, и Аз в нем. Умилосердися убо, Господи, и не обличи мя грешнаго, но сотвори со мною по милости Твоей; и да будут ми святая сия во исцеление, и очищение, и просвещение, и сохранение, и спасение, и во освящение души и тела; во отгнание всякаго мечтания, и лукаваго деяния, и действа диавольскаго, мысленнe во удесех моих действуемаго, в дерзновение и любовь, яже к Тебе; во исправление жития и утверждение, в возращение добродетели и совершенства; во исполнение заповедей, в Духа Святаго общение, в напутие живота вечнаго, во ответ благоприятен на Страшнем судищи Твоем: не в суд или во осуждение.

\mysubsubsection{Молитва 6-я,святого Симеона Нового Богослова}

От скверных устен, от мерзкаго сердца, от нечистаго языка, от души осквернены, приими моление, Христе мой, и не презри моих ни словес, ниже образов, ниже безстудия. Даждь ми дерзновенно глаголати, яже хощу, Христе мой, паче же и научи мя, что ми подобает творити и глаголати. Согреших паче блудницы, яже уведе, где обитаеши, миро купивши, прииде дерзостне помазати Твои нозе, Бога моего, Владыки и Христа моего. Якоже ону не отринул еси пришедшую от сердца, ниже мене возгнушайся, Слове: Твои же ми подаждь нозе, и держати и целовати, и струями слезными, яко многоценным миром, сия дерзостно помазати. Омый мя слезами моими, очисти мя ими, Слове. Остави и прегрешения моя, и прощение ми подаждь. Веси зол множество, веси и струпы моя, и язвы зриши моя, но и веру веси, и произволение зриши, и воздыхание слышиши. Не таится Тебе, Боже мой, Творче мой, Избавителю мой, ниже капля слезная, ниже капли часть некая. Несоделанное мое видесте очи Твои, в книзе же Твоей и еще несодеянная написана Тебе суть. Виждь смирение мое, виждь труд мой елик, и грехи вся остави ми, Боже всяческих: да чистым сердцем, притрепетною мыслию, и душею сокрушенною, нескверных Твоих причащуся и пресвятых Таин, имиже оживляется и обожается всяк ядый же и пияй чистым сердцем; Ты бо рекл еси, Владыко мой: всяк ядый Мою Плоть, и пияй Мою Кровь, во Мне убо сей пребывает, в немже и Аз есмь. Истинно слово всяко Владыки и Бога моего: божественных бо причащаяйся и боготворящих благодатей, не убо есмь един, но с Тобою, Христе мой, Светом трисолнечным, просвещающим мир. Да убо не един пребуду кроме Тебе Живодавца, дыхания моего, живота моего, радования моего, спасения миру. Сего ради к Тебе приступих, якоже зриши, со слезами, и душею сокрушенною избавления моих прегрешений прошу прияти ми, и Твоих живодательных и непорочных Таинств причаститися неосужденно, да пребудеши, якоже рекл еси, со мною треокаянным: да не кроме обрет мя Твоея благодати, прелестник восхитит мя льстивне, и прельстив отведет боготворящих Твоих словес. Сего ради к Тебе припадаю, и тепле вопию Ти: якоже блуднаго приял еси, и блудницу пришедшую, тако приими мя блуднаго и сквернаго, Щедре. Душею сокрушенною, ныне бо к Тебе приходя, вем, Спасе, яко иный, якоже аз, не прегреши Тебе, ниже содея деяния, яже аз содеях. Но сие паки вем, яко не величество прегрешений, ни грехов множество превосходит Бога моего многое долготерпение, и человеколюбие крайнее; но милостию сострастия тепле кающияся, и чистиши, и светлиши, и света твориши причастники, общники Божества Твоего соделоваяй независтно, и странное и Ангелом, и человеческим мыслем, беседуеши им многажды, якоже другом Твоим истинным. Сия дерзостна творят мя, сия вперяют мя, Христе мой. И дерзая Твоим богатым к нам благодеянием, радуяся вкупе и трепеща, огневи причащаюся трава сый, и странно чудо, орошаем неопально, якоже убо купина древле неопальне горящи. Ныне благодарною мыслию, благодарным же сердцем, благодарными удесы моими, души и тела моего, покланяюся и величаю, и славословлю Тя, Боже мой, яко благословенна суща, ныне же и во веки.

\mysubsubsection{Молитва 7-я, святого Иоанна Златоустого}

Боже, ослаби, остави, прости ми согрешения моя, елика Ти согреших, аще словом, аще делом, аще помышлением, волею или неволею, разумом или неразумием, вся ми прости яко благ и Человеколюбец, и молитвами Пречистыя Твоея Матере, умных Твоих служителей и святых сил, и всех святых от века Тебе благоугодивших, неосужденно благоволи прияти ми святое и пречистое Твое Тело и честную Кровь, во исцеление души же и тела, и во очищение лукавых моих помышлений. Яко Твое есть царство и сила и слава, со Отцем и Святым Духом, ныне и присно, и во веки веков. Аминь.

\mysubsubsection{Его же, 8-я}

Несмь доволен, Владыко Господи, да внидеши под кров души моея; но понеже хощеши Ты, яко Человеколюбец, жити во мне, дерзая приступаю; повелеваеши, да отверзу двери, яже Ты един создал еси, и внидеши со человеколюбием якоже еси, внидеши и просвещаеши помраченный мой помысл. Верую, яко сие сотвориши: не бо блудницу, со слезами пришедшую к Тебе, отгнал еси; ниже мытаря отвергл еси покаявшася; ниже разбойника, познавша Царство Твое, отгнал еси; ниже гонителя покаявшася оставил еси, еже бе: но от покаяния Тебе пришедшия вся, в лице Твоих другов вчинил еси, Един сый благословенный всегда, ныне и в безконечныя веки. Аминь.

\mysubsubsection{Его же, 9-я}

Господи Иисусе Христе Боже мой, ослаби, остави, очисти и прости ми грешному, и непотребному, и недостойному рабу Твоему, прегрешения, и согрешения, и грехопадения моя, елика Ти от юности моея, даже до настоящего дне и часа согреших: аще в разуме и в неразумии, аще в словесех или делех, или помышлениих и мыслех, и начинаниих, и всех моих чувствах. И молитвами безсеменно рождшия Тя, Пречистая и Приснодевы Марии, Матере Твоея, единыя непостыдныя надежды и предстательства и спасения моего, сподоби мя неосужденно причаститися пречистых, безсмертных, животворящих и страшных Твоих Таинств, во оставление грехов и в жизнь вечную: во освящение, и просвещение, крепость, исцеление, и здравие души же и тела, и в потребление и всесовершенное погубление лукавых моих помыслов, и помышлений, и предприятий, и нощных мечтаний, темных и лукавых духов; яко Твое есть царство, и сила, и слава, и честь, и поклонение, со Отцем и Святым Твоим Духом, ныне и присно, и во веки веков. Аминь.

\mysubsubsection{Молитва 10-я, святого Иоанна Дамаскина}

Пред дверьми храма Твоего предстою, и лютых помышлений не отступаю; но Ты, Христе Боже, мытаря оправдивый, и хананею помиловавый, и разбойнику рая двери отверзый, отверзи ми утробы человеколюбия Твоего и приими мя приходяща и прикасающася Тебе, яко блудницу, и кровоточивую: ова убо края ризы Твоея коснувшися, удобь исцеление прият, ова же пречистеи Твои нозе удержавши, разрешение грехов понесе. Аз же, окаянный, все Твое Тело дерзая восприяти, да не опален буду; но приими мя, якоже оныя, и просвети моя душевныя чувства, попаляя моя греховныя вины, молитвами безсеменно Рождшия Тя, и Небесных сил; яко благословен еси во веки веков. Аминь.

\mysubsubsection{Молитва святого Иоанна Златоустого}

Верую, Господи, и исповедую, яко Ты еси воистинну Христос, Сын Бога живаго, пришедый в мир грешныя спасти, от нихже первый есмь аз. Еще верую, яко сие есть самое пречистое Тело Твое, и сия самая есть честная Кровь Твоя. Молюся убо Тебе: помилуй мя, и прости ми прегрешения моя, вольная и невольная, яже словом, яже делом, яже ведением и неведением, и сподоби мя неосужденно причаститися пречистых Твоих Таинств, во оставление грехов, и в жизнь вечную. Аминь.

\mysubsubsection{Приходя же причаститься, произноси мысленно эти стихи Метафраста:}

Се приступаю к Божественному Причащению.

Содетелю, да не опалиши мя приобщением:

Огнь бо еси, недостойныя попаляяй.

Но убо очисти мя от всякия скверны.

\mysubsubsection{Затем:}

Вечери Твоея тайныя днесь, Сыне Божий, причастника мя приими; не бо врагом Твоим тайну повем, ни лобзания Ти дам, яко Иуда, но яко разбойник исповедаю Тя: помяни мя, Господи, во Царствии Твоем.

\mysubsubsection{И стихи:}

Боготворящую Кровь ужаснися, человече, зря:

Огнь бо есть, недостойныя попаляяй.

Божественное Тело и обожает мя и питает:

Обожает дух, ум же питает странно.

\mysubsubsection{Потом тропари:}

Усладил мя еси любовию, Христе, и изменил мя еси Божественным Твоим рачением; но попали огнем невещественным грехи моя, и насытитися еже в Тебе наслаждения сподоби: да ликуя возвеличаю, Блаже, два пришествия Твоя.

Во светлостех Святых Твоих како вниду недостойный? Аще бо дерзну совнити в чертог, одежда мя обличает, яко несть брачна, и связан извержен буду от Ангелов. Очисти, Господи, скверну души моея, и спаси мя, яко Человеколюбец.

\mysubsubsection{Также молитву:}

Владыко Человеколюбче, Господи Иисусе Христе Боже мой, да не в суд ми будут Святая сия, за еже недостойну ми быти: но во очищение и освящение души же и тела, и во обручение будущия жизни и царствия. Мне же, еже прилеплятися Богу, благо есть, полагати во Господе упование спасения моего.

\mysubsubsection{И еще:}

Вечери Твоея тайныя днесь, Сыне Божий, причастника мя приими; не бо врагом Твоим тайну повем, ни лобзания Ти дам, яко Иуда, но яко разбойник исповедаю Тя: помяни мя, Господи, во Царствии Твоем.

\end{mymulticols}


\medskip\myemph{ Желающий причаститься должен достойно приготовится к этому святому таинству. Приготовление это (в церковной практике оно называется говением) продолжается несколько дней и касается как телесной, так и духовной жизни человека. Телу предписывается воздержание, т.~е. телесная чистота (воздержание от супружеских отношений) и ограничение в пищи (пост). В дни поста исключается пища животного происхождения "--- мясо, молоко, яйца и, про строгом посте, рыба. Хлеб, овощи, фрукты употребляются в умеренном количестве. Ум не должен рассеиваться по мелочам житейским и развлекаться.

В дни говения надлежит посещать богослужения в храме, если позволят обстоятельства, и более прилежно выполнять домашнее молитвенное правило: кто читает обычно не все утренние и вечерние молитвы, пусть читает все полностью, кто не читает каноны, пусть в эти дни читает хотя бы по одному канону. Накануне причащения надо быть на вечернем богослужении и прочитать дома, кроме обычных молитв на сон грядущим, канон покаянный, канон Богородице и Ангелу хранителю. Каноны читают или один за другим полностью, или соединяя таким образом: читается ирмос первой песни покаянного канона («Яко по суху петешествовав Израиль, по бездне стопами, гонителя фараона видя потопляема, Богу победную песнь поим, вопияше») и тропари, затем тропари первой песни канона Богородице («Многими содержимь напaстьми, к Тебе прибегаю, спасения иский: о, Мaти Слова и Дево, от тяжких и лютых мя спаси»), опуская ирмос «Воду прошед…», и тропари канона Ангелу хранителю, тоже без ирмоса («Поим Господеви, проведшему люди Своя сквозе Чермное море, яко един славно прославися»). Так же читают и следующие песни. Тропари перед каноном Богородице и Ангелу хранителю, а также стихиры после канона Богородице в таком случае опускаются.

Читается также канон ко причащению и, кто пожелает,"--- акафист Иисусу Сладчайшему. После полуночи уже не едят и не пьют, ибо принято приступать к Таинству Причащения натощак. Утром прочитываются утренние молитвы и все последование ко Святому Причащению, кроме канона, прочитанного накануне.

Перед причащением необходима исповедь "--- вечером ли, или утром, перед литургией.
}


\mychapterending


\mychapter{Акафист Иисусу Сладчайшему}\begin{mymulticols}
%http://www.molitvoslov.com/text11.htm

\myfigure{B-2757_01-IV}

\mysubsubsection{Кондак 1}

Возбранный  Воеводо и Господи, ада победителю, яко избавлься от вечныя смерти, похвальная восписую Ти, создание и раб Твой; но, яко имеяй милосердие неизреченное, от всяких мя бед свободи, зовуща: Иисусе, Сыне Божий, помилуй мя.

\mysubsubsection{Икос 1}


Ангелов
Творче и Господи cил, отверзи ми недоуменный ум и язык на похвалу пречистаго Твоего имене, якоже глухому и гугнивому древле слух и язык отверзл еси, и, глаголаше зовый таковая: Иисусе пречудный, aнгелов удивление; Иисусе пресильный, прародителей избавление. Иисусе пресладкий, патриархов величание; Иисусе преславный, царей укрепление. Иисусе прелюбимый, пророков исполнение; Иисусе предивный, мучеников крепосте. Иисусе претихий, монахов радосте; Иисусе премилостивый,
пресвитеров сладосте. Иисусе премилосердый, постников воздержание; Иисусе пресладостный, преподобных радование. Иисусе пречестный, девственных целомудрие; Иисусе предвечный, грешников спасение. Иисусе, Сыне Божий, помилуй мя.

\mysubsubsection{Кондак 2}

Видя
вдовицу зельне плачущу, Господи, якоже бо тогда умилосердився, сына ея на погребение несома воскресил еси; сице и о мне умилосердися, Человеколюбче, и грехми умерщвленную мою душу воскреси, зовущую: Аллилуиа.

\mysubsubsection{Икос 2}

Разум
неуразуменный разумети Филипп ища, Господи, покажи нам Отца, глаголаше; Ты же к нему: толикое время сый со Мною, не познал ли еси, яко Отец во Мне, и Аз во Отце есмь? Темже, Неизследованне, со страхом зову Ти: Иисусе, Боже предвечный; Иисусе, Царю пресильный. Иисусе, Владыко долготерпеливый; Иисусе, Спасе премилостивый. Иисусе, хранителю мой преблагий; Иисусе, очисти грехи моя. Иисусе, отыми беззакония моя; Иисусе, отпусти неправды моя. Иисусе, надеждо моя, не остави мене; Иисусе, помощниче мой, не отрини мене. Иисусе, Создателю мой, не забуди
мене; Иисусе, Пастырю мой, не погуби мене. Иисусе, Сыне Божий, помилуй мя.

\mysubsubsection{Кондак 3}

Силою
свыше апостолы облекий, Иисусе, во Иерусалиме седящия, облецы и мене, обнаженнаго от всякаго благотворения, теплотою Духа Святаго Твоего и даждь ми с любвью пети Тебе: Аллилуиа.

\mysubsubsection{Икос 3}

Имеяй
богатство милосердия, мытари и грешники, и неверныя призвал еси, Иисусе; не презри и мене ныне, подобнаго им, но, яко многоценное миро, приими песнь сию: Иисусе, сило непобедимая; Иисусе, милосте
безконечная. Иисусе, красото пресветлая; Иисусе, любы неизреченная. Иисусе, Сыне Бога Живаго; Иисусе, помилуй мя грешнаго. Иисусе, услыши мя в беззакониих зачатаго; Иисусе, очисти мя во гресех рожденнаго. Иисусе, научи мя непотребнаго; Иисусе, освети мя темнаго. Иисусе, очисти мя сквернаго; Иисусе, возведи мя блуднаго. Иисусе, Сыне Божий, помилуй мя.

\mysubsubsection{Кондак 4}

Бурю
внутрь имеяй помышлений сумнительных, Петр утопаше; узрев же во плоти Тя суща, Иисусе, и по водам ходяща, позна Тя Бога истиннаго и, руку спасения получив, рече: Аллилуиа.

\mysubsubsection{Икос 4}

Слыша
слепый мимоходяща Тя, Господи, путем вопияше: Иисусе, Сыне Давидов, помилуй мя! И, призвав, отверзл еси очи его. Просвети убо милостию Твоею очи мысленныя сердца и мене, вопиюща Ти и глаголюща: Иисусе, вышних Создателю; Иисусе, нижних Искупителю. Иисусе, преисподних потребителю; Иисусе, всея твари украсителю. Иисусе, души моея утешителю; Иисусе, ума моего просветителю. Иисусе, сердца моего веселие; Иисусе, тела моего здравие. Иисусе, Спасе мой, спаси мя; Иисусе, свете мой, просвети мя. Иисусе, муки всякия избави мя; Иисусе, спаси мя, недостойнаго. Иисусе, Сыне Божий, помилуй мя.

\mysubsubsection{Кондак 5}

Боготочною
Кровию якоже искупил еси нас древле от законныя клятвы, Иисусе, сице изми нас от сети, еюже змий запят ны страстьми плотскими, и блудным наваждением, и злым унынием, вопиющия Ти: Аллилуиа.

\mysubsubsection{Икос 5}

Видевше
отроцы еврейстии во образе человечестем Создавшаго рукою человека, и Владыку разумевше Его, потщашася ветвьми угодити Ему, осанна вопиюще. Мы же песнь приносим Ти, глаголюще: Иисусе, Боже истинный; Иисусе, Сыне Давидов. Иисусе, Царю преславный; Иисусе, Агнче непорочный. Иисусе, Пастырю предивный; Иисусе, хранителю во младости моей. Иисусе, кормителю во юности моей; Иисусе, похвало в старости моей. Иисусе, надежде в смерти моей; Иисусе, животе по смерти моей. Иисусе, утешение мое на суде Твоем; Иисусе, желание мое, не посрами мене тогда. Иисусе, Сыне Божий, помилуй мя.

\mysubsubsection{Кондак 6}

Проповедник
богоносных вещание и глаголы исполняя, Иисусе, на земли явлься и с человеки Невместимый пожил еси, и болезни наша подъял еси, отнюдуже ранами Твоими мы исцелевше, пети навыкохом: Аллилуиа.

\mysubsubsection{Икос 6}

Возсия
вселенней просвещение истины Твоея, и отгнася лесть бесовская: идоли бо, Спасе наш, не терпяще Твоея крепости, падоша. Мы же, спасение получивше, вопием Ти: Иисусе, истино, лесть отгонящая; Иисусе, свете, превышший всех светлостей. Иисусе, Царю, премогаяй всех крепости; Иисусе, Боже, пребываяй в милости. Иисусе, Хлебе Животный, насыти мя алчущаго; Иисусе, источниче разума, напой мя жаждущаго. Иисусе, одеждо веселия, одей мя тленнаго; Иисусе, покрове радости, покрый мя недостойнаго. Иисусе, подателю просящим, даждь мне плач за грехи моя; Иисусе, обретение ищущим, обрящи душу мою. Иисусе, отверзителю толкущим, отверзи сердце мое окаянное; Иисусе, Искупителю грешных, очисти беззакония моя. Иисусе, Сыне Божий, помилуй мя.

\mysubsubsection{Кондак 7}

Хотя
сокровенную тайну от века открыти, яко овча на заколение веден был еси, Иисусе, и яко агнец прямо стригущаго его безгласен, и яко Бог из мертвых воскресл еси, и со славою на небеса вознеслся еси, и нас совоздвигл еси, зовущих: Аллилуиа.

\mysubsubsection{Икос 7}

Дивную
показа тварь, явлейся Творец нам: без семене от Девы воплотися, из гроба, печати не рушив, воскресе, и ко апостолом, дверем затворенным, с плотию вниде. Темже чудящеся, воспоим: Иисусе, Слове необыменный; Иисусе, Слове несоглядаемый. Иисусе, сило непостижимая; Иисусе,
мудросте недомыслимая. Иисусе, Божество неописанное; Иисусе, господство неисчетное. Иисусе, царство непобедимое; Иисусе, владычество безконечное. Иисусе, крепосте высочайшая; Иисусе, власте вечная. Иисусе, Творче мой, ущедри мя; Иисусе, Спасе мой, спаси мя. Иисусе, Сыне Божий, помилуй мя.

\mysubsubsection{Кондак 8}

Странно
Бога вочеловечшася видяще, устранимся суетнаго мира и ум на Божественная возложим. Сего бо ради Бог на землю сниде, да нас на небеса возведет, вопиющих Ему: Аллилуиа.

\mysubsubsection{Икос 8}

Весь
бе в нижних, и вышних никакоже отступи Неисчетный, егда волею нас ради пострада, и смертию Своею нашу смерть умертви, и воскресением живот дарова поющим: Иисусе, сладосте сердечная; Иисусе, крепосте телесная. Иисусе, светлосте душевная; Иисусе, быстрото умная. Иисусе, радосте совестная; Иисусе, надеждо известная. Иисусе, памяте предвечная; Иисусе, похвало высокая. Иисусе, славо моя превознесенная; Иисусе, желание мое, не отрини мене. Иисусе, Пастырю мой, взыщи мене; Иисусе, Спасе мой, спаси мене. Иисусе, Сыне Божий, помилуй мя.

\mysubsubsection{Кондак 9}

Все
естество ангельское безпрестани славит пресвятое имя Твое, Иисусе, на небеси: Свят, Свят, Свят, вопиюще; мы же, грешнии на земли бренными устнами вопием: Аллилуиа.

\mysubsubsection{Икос 9}

Ветия
многовещанны, якоже рыбы безгласныя видим о Тебе, Иисусе, Спасе наш: недоумеют бо глаголати, како Бог непреложний и человек совершенный пребываеши? Мы же таинству дивящеся, вопием верно: Иисусе, Боже предвечный; Иисусе, Царю царствующих. Иисусе, Владыко владеющих; Иисусе, Судие живых и мертвых. Иисусе, надеждо ненадежных; Иисусе, утешение плачущих. Иисусе, славо нищих; Иисусе, не осуди мя по делом моим. Иисусе, очисти мя по милости Твоей; Иисусе, отжени от мене уныние. Иисусе, просвети моя мысли сердечныя; Иисусе, даждь ми память смертную. Иисусе, Сыне Божий, помилуй мя.

\mysubsubsection{Кондак 10}

Спасти
хотя мир, Восточе востоком, к темному западу "--- естеству нашему пришед, смирился еси до смерти; темже превознесеся имя Твое паче всякаго имене, и от всех колен небесных и земных слышиши: Аллилуиа.

\mysubsubsection{Икос 10}

Царю
Превечный, Утешителю, Христе истинный, очисти ны от всякия скверны, якоже очистил еси десять прокаженных, и исцели ны, якоже исцелил еси сребролюбивую душу Закхеа мытаря, да вопием Ти, во умилении зовуще: Иисусе, сокровище нетленное; Иисусе, богатство неистощимое. Иисусе, пище крепкая; Иисусе, питие неисчерпаемое. Иисусе, нищих одеяние; Иисусе, вдов заступление. Иисусе, сирых защитниче; Иисусе, труждающихся помоще. Иисусе, странных наставниче; Иисусе, плавающих кормчий. Иисусе, бурных отишие; Иисусе Боже, воздвигни мя падшаго. Иисусе, Сыне Божий, помилуй мя.

\mysubsubsection{Кондак 11}

Пение
всеумиленное приношу Ти недостойный, вопию Ти яко хананеа: Иисусе, помилуй мя; не дщерь бо, но плоть имам страстьми люте бесящуюся и яростию палимую, и исцеление даждь вопиющу Ти: Аллилуиа.

\mysubsubsection{Икос 11}

Светоподательна
светильника сущим во тьме неразумия, прежде гоняй Тя Павел, богоразумнаго гласа силу внуши и душевную быстроту уясни; сице и мене темныя зеницы душевныя просвети, зовуща: Иисусе, Царю мой прекрепкий; Иисусе, Боже мой пресильный. Иисусе, Господи мой пребезсмертный; Иисусе, Создателю мой преславный. Иисусе, Наставниче мой предобрый; Иисусе, Пастырю мой прещедрый. Иисусе, Владыко мой премилостивый; Иисусе, Спасе мой премилосердый. Иисусе, просвети моя чувствия, потемненныя страстьми; Иисусе, исцели мое тело, острупленное грехми. Иисусе, очисти мой ум от помыслов суетных; Иисусе, сохрани сердце мое от похотей лукавых. Иисусе, Сыне Божий, помилуй мя.

\mysubsubsection{Кондак 12}

Благодать
подаждь ми, всех долгов решителю, Иисусе, и приими мя кающася, якоже приял еси Петра, отвергшагося Тебе, и призови мя унывающаго, якоже древле Павла, гоняща Тя, и услыши мя, вопиюща Ти: Аллилуиа.

\mysubsubsection{Икос 12}

Поюще
Твое вочеловечение, восхваляем Тя вси, и веруем со Фомою, яко Господь и Бог еси, седяй со Отцем и хотяй судити живым и мертвым. Тогда убо сподоби мя деснаго стояния, вопиющаго: Иисусе, Царю предвечный, помилуй мя; Иисусе, цвете благовонный, облагоухай мя. Иисусе, теплото любимая, огрей мя; Иисусе, храме предвечный, покрый мя. Иисусе, одеждо светлая, украси мя; Иисусе, бисере честный, осияй мя. Иисусе, каменю драгий, просвети мя; Иисусе, солнце правды, освети мя. Иисусе, свете святый, облистай мя; Иисусе, болезни душевныя и телесныя избави мя. Иисусе, из руки сопротивныя изми мя; Иисусе, огня неугасимаго и прочих вечных мук свободи мя. Иисусе, Сыне Божий, помилуй мя.

\mysubsubsection{Кондак 13}

О,
пресладкий и всещедрый Иисусе! Приими ныне малое моление сие наше, якоже приял еси вдовицы две лепте, и сохрани достояние Твое от враг видимых и невидимых, от нашествия иноплеменних, от недуга и глада, от всякия скорби и смертоносныя раны, и грядущия изми муки всех, вопиющих Ти: Аллилуиа, aллилуиа, aллилуиа.

\mysubsubsection{(Kондак читается трижды)}

\mysubsubsection{Икос 1}

Ангелов
Творче и Господи cил, отверзи ми недоуменный ум и язык на похвалу пречистаго Твоего имене, якоже глухому и гугнивому древле слух и язык отверзл еси, и, глаголаше зовый таковая: Иисусе пречудный, aнгелов удивление; Иисусе пресильный, прародителей избавление. Иисусе пресладкий, патриархов величание; Иисусе преславный, царей укрепление. Иисусе прелюбимый, пророков исполнение; Иисусе предивный, мучеников крепосте. Иисусе претихий, монахов радосте; Иисусе премилостивый, пресвитеров сладосте. Иисусе премилосердый, постников воздержание; Иисусе пресладостный, преподобных радование. Иисусе пречестный, девственных целомудрие; Иисусе предвечный, грешников спасение. Иисусе, Сыне Божий, помилуй мя.

\mysubsubsection{Кондак 1}

Возбранный
Воеводо и Господи, ада победителю, яко избавлься от вечныя смерти, похвальная восписую Ти, создание и раб Твой; но, яко имеяй милосердие неизреченное, от всяких мя бед свободи, зовуща: Иисусе, Сыне Божий, помилуй мя.

\mysubsubsection{Молитва}

Владыко
Господи Иисусе Христе Боже мой, иже неизреченнаго ради Твоего человеколюбия на конец веков во плоть оболкийся от Приснодевы Марии, славлю о мне Твое спасительное промышление, раб Твой, Владыко;
песнословлю Тя, яко Тебе ради Отца познах; благословлю Тя, Егоже ради и Дух Святый в мир прииде; покланяюся Твоей по плоти Пречистой Матери, таковей страшней тайне послужившей; восхваляю Твоя Ангельская ликостояния, яко воспеватели и служители Твоего величествия; ублажаю Предтечу Иоанна, Тебе крестившаго, Господи; почитаю и провозвестившия Тя пророки, прославляю апостолы Твоя святыя; торжествую же и мученики, священники же Твоя славлю; покланяюся преподобным Твоим, и вся Твоя праведники пестунствую. Таковаго и толикаго многаго и неизреченнаго лика Божественнаго в молитву привожду Тебе, всещедрому Богу, раб Твой, и сего ради прошу моим согрешением прощения, еже даруй ми всех Твоих ради святых, изряднее же святых Твоих щедрот, яко благословен еси во веки. Аминь.


\end{mymulticols}

\mychapterending[1.5]


\mychapter{Акафист Пресвятой Богородице}\begin{mymulticols}
%http://www.molitvoslov.com/text13.htm

\myfigure{5_0}

\mysubsubsection{Кондак 1}

Взбранной Воеводе победительная, яко избавльшеся от злых, благодарственная восписуем Ти раби Твои, Богородице; но яко имущая державу непобедимую, от всяких нас бед свободи, да зовем Ти: радуйся, Невесто Неневестная.

\mysubsubsection{Икос 1}

Ангел предстатель с небесе послан бысть рещи Богородице: радуйся, и со безплотным гласом воплощаема Тя зря, Господи, ужасашеся и стояше, зовый к Ней таковая: Радуйся, Еюже рaдocть возсияет; радуйся, Еюже клятва изчезнет. Радуйся, падшаго Адама воззвание; радуйся, слез Евиных избавление. Радуйся, высото неудобовосходимая человеческими помыслы; радуйся, глубино неудобозримая и ангельскима очима. Радуйся, яко еси Царево седалище; радуйся, яко носиши Носящаго вся. Радуйся, Звездо, являющая Солнце; радуйся, утробо Божественнаго воплощения. Радуйся, Еюже обновляется тварь; радуйся, Еюже покланяемся Творцу. Радуйся, Невесто Неневестная.

\mysubsubsection{Кондак 2}

Видящи Святая Себе в чистоте, глаголет Гавриилу дерзостно: преславное твоего гласа неудобоприятельно души Моей является: безсеменнаго бо зачатия рождество како глаголеши, зовый: Аллилуиа.

\mysubsubsection{Икос 2}

Разум недоразумеваемый разумети Дева ищущи, возопи к служащему: из боку чисту, Сыну како есть родитися мощно, рцы Ми? К Нейже он рече со страхом, обаче зовый сице: Радуйся, совета неизреченнаго Таиннице; радуйся, молчания просящих веро. Радуйся, чудес Христовых начало; радуйся, велений Его главизно. Радуйся, лествице небесная, Еюже сниде Бог; радуйся, мосте, преводяй сущих от земли на небо. Радуйся, Ангелов многословущее чудо; радуйся, бесов многоплачевное поражение. Радуйся, Свет неизреченно родившая; радуйся, еже како, ни единаго же научившая. Радуйся, премудрых превосходящая разум; радуйся, верных озаряющая смыслы. Радуйся, Невесто Неневестная.

\mysubsubsection{Кондак 3}

Сила Вышняго осени тогда к зачатию Браконеискусную, и благоплодная Тоя ложесна, яко село показа сладкое, всем хотящим жати спасение, всегда пети сице: Аллилуиа.

\mysubsubsection{Икос 3}

Имущи Богоприятную Дева утробу, востече ко Елисавети: младенец же оноя абие познав Сея целование, радовашеся, и играньми яко песньми вопияше к Богородице: Радуйся, отрасли неувядаемыя розго; радуйся, Плода безсмертнаго стяжание. Радуйся, Делателя делающая Человеколюбца; радуйся, Садителя жизни нашея рождшая. Радуйся, ниво, растящая гобзование щедрот; радуйся, трапезо, носящая обилие очищения. Радуйся, яко рай пищный процветаеши; радуйся, яко пристанище душам готовиши. Радуйся, приятное молитвы кадило; радуйся, всего мира очищение. Радуйся, Божие к смертным благоволение; радуйся, смертных к Богу дерзновение. Радуйся, Невесто Неневестная.

\mysubsubsection{Кондак 4}

Бурю внутрь имея помышлений сумнительных, целомудренный Иосиф смятеся, к Тебе зря небрачней, и бракоокрадованную помышляя, Непорочная; уведев же Твое зачатие от Духа Свята, рече: Аллилуиа.

\mysubsubsection{Икос 4}

Слышаша пастырие Ангелов поющих плотское Христово пришествие, и текше яко к Пастырю видят Сего яко агнца непорочна, во чреве Мариине упасшася, Юже поюще реша: Радуйся, Агнца и Пастыря Мати; радуйся, дворе словесных овец. Радуйся, невидимых врагов мучение; радуйся, райских дверей отверзение. Радуйся, яко небесная срадуются земным; радуйся, яко земная сликовствуют небесным. Радуйся, апостолов немолчная уста; радуйся, страстотерпцев непобедимая дерзосте. Радуйся, твердое веры утверждение; радуйся, светлое благодати познание. Радуйся, Еюже обнажися ад; радуйся, Еюже облекохомся славою. Радуйся, Невесто Неневестная.

\mysubsubsection{Кондак 5}

Боготечную звезду узревше волсви, тоя последоваша зари, и яко светильник держаще ю, тою испытаху крепкаго Царя, и достигше Непостижимаго, возрадовашася, Ему вопиюще: Аллилуиа.

\mysubsubsection{Икос 5}

Видеша отроцы халдейстии на руку Девичу Создавшаго руками человеки, и Владыку разумевающе Его, аще и рабий прият зрак, потщашася дарми послужити Ему, и возопити Благословенней: Радуйся, Звезды незаходимыя Мати; радуйся, заре таинственнаго дне. Радуйся, прелести пещь угасившая; радуйся, Троицы таинники просвещающая. Радуйся, мучителя безчеловечнаго изметающая от начальства; радуйся, Господа Человеколюбца показавшая Христа. Радуйся, варварскаго избавляющая служения; радуйся, тимения изымающая дел. Радуйся, огня поклонение угасившая; радуйся, пламене страстей изменяющая. Радуйся, верных наставнице целомудрия; радуйся, всех родов веселие. Радуйся, Невесто Неневестная.

\mysubsubsection{Кондак 6}

Проповедницы богоноснии, бывше волсви, возвратишася в Вавилон, скончавше Твое пророчество, и проповедавше Тя Христа всем, оставиша Ирода яко буесловна, не ведуща пети: Аллилуиа.

\mysubsubsection{Икос 6}

Возсиявый во Египте просвещение истины, отгнал еси лжи тьму: идоли бо его, Спасе, не терпяще Твоея крепости, падоша, сих же избавльшиися вопияху к Богородице: Радуйся, исправление человеков; радуйся, низпадение бесов. Радуйся, прелести державу поправшая; радуйся, идольскую лесть обличившая. Радуйся, море, потопившее фараона мысленнаго; радуйся, каменю, напоивший жаждущия жизни. Радуйся, огненный столпе, наставляяй сущия во тьме; радуйся, покрове миру, ширший облака. Радуйся, пище, манны приемнице; радуйся, сладости святыя служительнице. Радуйся, земле обетования; радуйся, из неяже течет мед и млеко. Радуйся, Невесто Неневестная.

\mysubsubsection{Кондак 7}

Хотящу Симеону от нынешняго века преставитися прелестнаго, вдался еси яко младенец тому, но познался еси ему и Бог совершенный. Темже удивися Твоей неизреченней премудрости, зовый: Аллилуиа.

\mysubsubsection{Икос 7}

Новую показа тварь, явлься Зиждитель нам от Него бывшим, из безсеменныя прозяб утробы, и сохранив Ю, якоже бе, нетленну, да чудо видяще, воспоем Ю, вопиюще: Радуйся, цвете нетления; радуйся, венче воздержания. Радуйся, воскресения образ облистающая; радуйся, ангельское житие являющая. Радуйся, древо светлоплодовитое, от негоже питаются вернии; радуйся, древо благосеннолиственное, имже покрываются мнози. Радуйся, во чреве носящая Избавителя плененным; радуйся, рождшая Наставника заблудшим. Радуйся, Судии праведнаго умоление; радуйся, многих согрешений прощение. Радуйся, одеждо нагих дерзновения; радуйся, любы, всякое желание побеждающая. Радуйся, Невесто Неневестная.

\mysubsubsection{Кондак 8}

Странное рождество видевше, устранимся мира, ум на небеса преложше: сего бо ради высокий Бог на земли явися смиренный человек, хотяй привлещи к высоте Тому вопиющия: Аллилуиа.

\mysubsubsection{Икос 8}

Весь бе в нижних и вышних никакоже отступи неописанное Слово: снизхождение бо Божественное, не прехождение же местное бысть, и рождество от Девы Богоприятныя, слышащия сия: Радуйся, Бога невместимаго вместилище; радуйся, честнаго таинства двери. Радуйся, неверных сумнительное слышание; радуйся, верных известная похвало. Радуйся, колеснице пресвятая Сущаго на Херувимех; радуйся, селение преславное Сущаго на Серафимех. Радуйся, противная в тожде собравшая; радуйся, девство и рождество сочетавшая. Радуйся, Еюже разрешися преступление; радуйся, Еюже отверзеся рай. Радуйся, ключу Царствия Христова; радуйся, надеждо благ вечных. Радуйся, Невесто Неневестная.

\mysubsubsection{Кондак 9}

Всякое естество ангельское удивися великому Твоего вочеловечения делу; неприступнаго бо яко Бога, зряще всем приступнаго Человека, нам убо спребывающа, слышаща же от всех: Аллилуиа.

\mysubsubsection{Икос 9}

Ветия многовещанныя, яко рыбы безгласныя видим о Тебе, Богородице, недоумевают бо глаголати, еже како и Дева пребываеши, и родити возмогла еси. Мы же, таинству дивящеся, верно вопием: Радуйся, премудрости Божия приятелище, радуйся, промышления Его сокровище. Радуйся, любомудрыя немудрыя являющая; радуйся, хитрословесныя безсловесныя обличающая. Радуйся, яко обуяша лютии взыскателе; радуйся, яко увядоша баснотворцы. Радуйся, афинейская плетения растерзающая; радуйся, рыбарския мрежи исполняющая. Радуйся, из глубины неведения извлачающая; радуйся, многи в разуме просвещающая. Радуйся, кораблю хотящих спастися; радуйся, пристанище житейских плаваний. Радуйся, Невесто Неневестная.

\mysubsubsection{Кондак 10}

Спасти хотя мир, Иже всех Украситель, к сему самообетован прииде, и Пастырь сый, яко Бог, нас ради явися по нам человек: подобным бо подобное призвав, яко Бог слышит: Аллилуиа.

\mysubsubsection{Икос 10}

Стена еси девам, Богородице Дево, и всем к Тебе прибегающим: ибо небесе и земли Творец устрои Тя, Пречистая, вселься во утробе Твоей, и вся приглашати Тебе научи: Радуйся, столпе девства; радуйся, дверь спасения. Радуйся, начальнице мысленнаго наздания; радуйся, подательнице Божественныя благости. Радуйся, Ты бо обновила еси зачатыя студно; радуйся, Ты бо наказала еси окраденныя умом. Радуйся, тлителя смыслов упраждняющая; радуйся, Сеятеля чистоты рождшая. Радуйся, чертоже безсеменнаго уневещения; радуйся, верных Господеви сочетавшая. Радуйся, добрая младопитательнице девам; радуйся, невестокрасительнице душ святых. Радуйся, Невесто Неневестная.

\mysubsubsection{Кондак 11}

Пение всякое побеждается, спростретися тщащееся ко множеству многих щедрот Твоих: равночисленныя бо песка песни аще приносим Ти, Царю Святый, ничтоже совершаем достойно, яже дал еси нам, Тебе вопиющим: Аллилуиа.

\mysubsubsection{Икос 11}

Светоприемную свещу, сущим во тьме явльшуюся, зрим Святую Деву, невещественный бо вжигающи огнь, наставляет к разуму Божественному вся, зарею ум просвещающая, званием же почитаемая, сими: Радуйся, луче умнаго Солнца; радуйся, светило незаходимаго Света. Радуйся, молние, души просвещающая; радуйся, яко гром враги устрашающая. Радуйся, яко многосветлое возсияваеши просвещение; радуйся, яко многотекущую источаеши реку. Радуйся, купели живописующая образ; радуйся, греховную отъемлющая скверну. Радуйся, бане, омывающая совесть; радуйся, чаше, черплющая рaдocть. Радуйся, обоняние Христова благоухания; радуйся, животе тайнаго веселия. Радуйся, Невесто Неневестная.

\mysubsubsection{Кондак 12}

Благодать дати восхотев, долгов древних, всех долгов Решитель человеком, прииде Собою ко отшедшим Того благодати, и раздрав рукописание, слышит от всех сице: Аллилуиа.

\mysubsubsection{Икос 12}

Поюще Твое Рождество, хвалим Тя вси, яко одушевленный храм, Богородице: во Твоей бо вселився утробе содержай вся рукою Господь, освяти, прослави и научи вопити Тебе всех: Радуйся, селение Бога и Слова; радуйся, святая святых большая. Радуйся, ковчеже, позлащенный Духом; радуйся, сокровище живота неистощимое. Радуйся, честный венче людей благочестивых; радуйся, честная похвало иереев благоговейных. Радуйся, церкве непоколебимый столпе; радуйся, Царствия нерушимая стено. Радуйся, Еюже воздвижутся победы; радуйся, Еюже низпадают врази. Радуйся, тела моего врачевание; радуйся, души моея спасение. Радуйся, Невесто Неневестная.

\mysubsubsection{Кондак 13}

О, Всепетая Мати, рождшая всех святых Святейшее Слово! Нынешнее приемши приношение, от всякия избави напасти всех, и будущия изми муки, о Тебе вопиющих: Аллилуиа, aллилуиа, aллилуиа.

\myemph{ (Kондак читается трижды)}

\mysubsubsection{Икос 1}

Ангел предстатель с небесе послан бысть рещи Богородице: радуйся, и со безплотным гласом воплощаема Тя зря, Господи, ужасашеся и стояше, зовый к Ней таковая: Радуйся, Еюже рaдocть возсияет; радуйся, Еюже клятва изчезнет. Радуйся, падшаго Адама воззвание; радуйся, слез Евиных избавление. Радуйся, высото неудобовосходимая человеческими помыслы; радуйся, глубино неудобозримая и ангельскима очима. Радуйся, яко еси Царево седалище; радуйся, яко носиши Носящаго вся. Радуйся, Звездо, являющая Солнце; радуйся, утробо Божественнаго воплощения. Радуйся, Еюже обновляется тварь; радуйся, Еюже покланяемся Творцу. Радуйся, Невесто Неневестная.

\mysubsubsection{Кондак 1}

Взбранной Воеводе победительная, яко избавльшеся от злых, благодарственная восписуем Ти раби Твои, Богородице; но яко имущая державу непобедимую, от всяких нас бед свободи, да зовем Ти: радуйся, Невесто Неневестная.

\mysubsubsection{Молитвы}

О, Пресвятая Госпоже Владычице Богородице, вышши еси всех Ангел и Архангел, и всея твари честнейши, помощнице еси обидимых, ненадеющихся надеяние, убогих заступнице, печальных утешение, алчущих кормительнице, нагих одеяние, больных исцеление, грешных спасение, христиан всех поможение и заступление. О, Всемилостивая Госпоже, Дево Богородице Владычице, милостию Твоею спаси и помилуй святейшия патриархи православныя, преосвященныя митрополиты, архиепископы и епископы и весь священнический и иноческий чин, и вся православныя христианы ризою Твоею честною защити; и умоли, Госпоже, из Тебе без семене воплотившагося Христа Бога нашего, да препояшет нас силою Своею свыше, на невидимыя и видимыя враги наша. О, Всемилостивая Госпоже Владычице Богородице! Воздвигни нас из глубины греховныя и избави нас от глада, губительства, от труса и потопа, от огня и меча, от нахождения иноплеменных и междоусобныя брани, и от напрасныя смерти, и от нападения вражия, и от тлетворных ветр, и от смертоносныя язвы, и от всякаго зла. Подаждь, Госпоже, мир и здравие рабом Твоим, всем православным христианом, и просвети им ум, и очи сердечнии, еже ко спасению; и сподоби ны, грешныя рабы Твоя, Царствия Сына Твоего, Христа Бога нашего; яко держава Его благословена и препрославлена, со Безначальным Его Отцем, и с Пресвятым, и Благим, и Животворящим Его Духом, ныне и присно, и во веки веков. Аминь.

О, Пресвятая Дево Мати Господа, Царице Небесе и земли! Вонми многоболезненному воздыханию души нашея, призри с высоты святыя Твоея на нас, с верою и любовию покланяющихся пречистому образу Твоему. Се бо грехми погружаемии и скорбьми обуреваемии, взирая на Твой образ, яко живей Ти сущей с нами, приносим смиренныя моления наша. Не имамы бо ни иныя помощи, ни инаго предстательства, ни утешения, токмо Тебе, о, Мати всех скорбящих и обремененных. Помози нам немощным, утоли скорбь нашу, настави на путь правый нас, заблуждающих, уврачуй и спаси безнадежных, даруй нам прочее время живота нашего в мире и тишине проводити, подаждь христианскую кончину, и на страшнем суде Сына Твоего явися нам милосердая Заступница, да всегда поем, величаем и славим Тя, яко благую Заступницу рода христианскаго, со всеми угодившими Богу. Аминь.


\end{mymulticols}

\mychapterending


\mychapter{Благодарственные молитвы по Святом Причащении}\begin{mymulticols}
%http://www.molitvoslov.com/text206.htm

Слава Тебе, Боже. Слава Тебе, Боже. Слава Тебе, Боже.

\mysubsubsection{Благодарственная молитва, 1-я}

Благодарю тя, Господи, Боже мой, яко не отринул мя еси грешнаго, но общника мя быти святынь Твоих сподобил еси. Благодарю Тя, яко мене недостойнаго причаститися Пречистых Твоих и Небесных Даров сподобил еси. Но Владыко Человеколюбче, нас ради умерый же и воскресый, и даровавый нам страшная сия и животворящая Таинства во благодеяние и освящение душ и телес наших, даждь быти сим и мне во исцеление души же и тела, во отгнание всякаго сопротивнаго, в просвещение очию сердца моего, в мир душевных моих сил, в веру непостыдну, в любовь нелицемерну, во исполнение премудрости, в соблюдение заповедей Твоих, в приложение Божественныя Твоея благодати и Твоего Царствия присвоение; да во святыни Твоей теми сохраняемь, Твою благодать поминаю всегда, и не к тому себе живу, но Тебе, нашему Владыце и Благодетелю; и тако сего жития исшед о надежди живота вечнаго, в присносущный достигну покой, идеже празднующих глас непрестанный, и безконечная сладость, зрящих Твоего лица доброту неизреченную. Ты бо еси истинное желание, и неизреченное веселие любящих Тя, Христе Боже наш, и Тя поет вся тварь во веки. Аминь.

\mysubsubsection{Молитва 2-я, святого Василия Великого}

Владыко Христе Боже, Царю веков, и Содетелю всех, Благодарю Тя о всех, яже ми подал благих, и о причащении пречистых и животворящих Твоих Таинств. Молю убо Тя, Блаже и Человеколюбче: сохрани мя под кровом Твоим, и в сени крилу Твоею; и даруй ми чистою совестию, даже до последняго моего издыхания, достойно причащатися святынь Твоих, во оставление грехов, и в жизнь вечную. Ты бо еси хлеб животный, источник святыни, Податель благих, и Тебе славу возсылаем, со Отцем и Святым Духом, ныне и присно, и во веки веков. Аминь.

\mysubsubsection{Молитва 3-я , Симеона Метафраста}

Давый пищу мне плоть Твою волею, огнь сый и опаляяй недостойныя, да не опалиши мене, Содетелю мой; паче же пройди во уды моя, во вся составы, во утробу, в сердце. Попали терние всех моих прегрешений. Душу очисти, освяти помышления. Составы утверди с костьми вкупе. Чувств просвети простую пятерицу. Всего мя спригвозди страху Твоему. Присно покрый, соблюди же, и сохрани мя от всякаго дела и слова душетленнаго. Очисти и омый, и украси мя; удобри, вразуми, и просвети мя. Покажи мя Твое селение единаго Духа, и не ктому селение греха. Да яко Твоего дому, входом причащения, яко огня мене бежит всяк злодей, всяка страсть. Молитвенники Тебе приношу вся святыя, чиноначалия же безплотных, Предтечу Твоего, премудрыя Апостолы, к сим же Твою нескверную чистую Матерь, ихже мольбы Благоутробне приими, Христе мой, и сыном света соделай Твоего служителя. Ты бо еси освящение и единый наших, Блаже, душ и светлость; и Тебе лепоподобно, яко Богу и Владыце, славу вси возсылаем на всяк день.

\mysubsubsection{Молитва 4-я}

Тело Твое Святое, Господи, Иисусе Христе, Боже наш, да будет ми в живот вечный, и Кровь Твоя Честная во оставление грехов: буди же ми благодарение сие в радость, здравие и веселие; в страшное же и второе пришествие Твое сподоби мя грешнаго стати одесную славы Твоея, молитвами Пречистыя Твоея Матере, и всех святых.

\mysubsubsection{Молитва 5-я, ко Пресвятой Богородице}

Пресвятая Владычице Богородице, свете помраченныя моея души, надеждо, покрове, прибежище, утешение, радование мое, благодарю Тя, яко сподобила мя еси недостойнаго, причастника быти Пречистаго Тела и Честныя Крове Сына Твоего. Но Рождшая истинный Свет, просвети моя умныя очи сердца; Яже Источник безсмертия рождшая, оживотвори мя умерщвленнаго грехом; Яже милостиваго Бога любоблагоутробная Мати, помилуй мя, и даждь ми умиление и сокрушение в сердце моем, и смирение в мыслех моих, и воззвание в пленениих помышлений моих; и сподоби мя до последняго издыхания, неосужденно приимати пречистых Таин освящение, во исцеление души же и тела. И подаждь ми слезы покаяния и исповедания, во еже пети и славити Тя во вся дни живота моего, яко благословенна и препрoславленна еси во веки. Аминь.

Ныне отпущаеши раба Твоего, Владыко, по глаголу Твоему, с миром: яко видеста очи мои спасение Твое, еже еси уготовал пред лицем всех людей, свет во откровение языков и славу людей Твоих, Израиля.

Святый Боже, Святый Крепкий, Святый Безсмертный, помилуй нас. \myemph{ (Tрижды)}

Слава Отцу и Сыну и Святому Духу, и ныне и присно и во веки веков. Аминь.

Пресвятая Троице, помилуй нас; Господи, очисти грехи наша; Владыко, прости беззакония наша; Святый, посети и исцели немощи наша, имене Твоего ради.

Господи, помилуй. \myemph{ (Tрижды)}

Слава Отцу и Сыну и Святому Духу, и ныне и присно и во веки веков. Аминь.

Отче наш, Иже еси на небесех! Да святится имя Твое, да приидет Царствие Твое, да будет воля Твоя, яко на небеси и на земли. Хлеб наш насущный даждь нам днесь; и остави нам долги наша, якоже и мы оставляем должником нашим; и не введи нас во искушение, но избави нас от лукаваго.

\mysubsubsection{Тропарь св. Иоанну Златоустому, глас 8-й:}

Уст твоих, якоже светлость огня, возсиявши благодать, вселенную просвети: не сребролюбия мирови сокровища сниска, высоту нам смиренномудрия показа, но твоими словесы наказуя, отче Иоанне Златоусте, моли Слова Христа Бога спастися душам нашим.

\mysubsubsection{Кондак, глас 6-й:}

\slavan

От небес приял еси Божественную благодать, и твоими устнами вся учиши покланятися в Троице единому Богу, Иоанне Златоусте всеблаженне, преподобне, достойно хвалим тя: еси бо наставник, яко божественная являя.

\inynen

\Bogorodichen{Предстательство христиан непостыдное, ходатайство ко Творцу непреложное, не презри грешных молений гласы, но предвари, яко Благая, на помощь нас, верно зовущих Ти: ускори на молитву, и потщися на умоление, предстательствующи присно, Богородице, чтущих Тя.}

\medskip\myemph{ Если совершалась литургия святого Василия Великого, читай}

\mysubsubsection{Тропарь Василию Великому, глас 1-й:}

Во всю землю изыде вещание твое, яко приемшую слово твое, имже боголепно научил еси, естество сущих уяснил еси, человеческия обычаи украсил еси, царское священие, отче преподобне, моли Христа Бога, спастися душам нашим.

\mysubsubsection{Кондак, глас 4-й:}

\slavan

Явился еси основание непоколебимое церкве, подая всем некрадомое господство человеком, запечатлея твоими веленьми, небоявленне Василие преподобне.

\inynen

\Bogorodichen{Предстательство христиан непостыдное, ходатайство ко Творцу непреложное, не презри грешных молений гласы, но предвари, яко Благая, на помощь нас, верно зовущих Ти: ускори на молитву, и потщися на умоление, предстательствующи присно, Богородице, чтущих Тя.}

\medskip\myemph{ Если совершалась Литургия Преждеосвященных Даров, читай}

\mysubsubsection{Тропарь святому Григорию Двоеслову, глас 4-й:}

Иже от Бога свыше божественную благодать восприем, славне Григорие, и Того силою укрепляемь, евангельски шествовати изволил еси, отонудуже у Христа возмездие трудов приял еси всеблаженне: Егоже моли, да спасет души наша.

\mysubsubsection{Кондак, глас 3-й:}

\slavan

Подобоначальник показался еси Начальника пастырем Христа, иноков чреды, отче Григорие, ко ограде небесней наставляя, и оттуду научил еси стадо Христово заповедем Его: ныне же с ними радуешися, и ликуеши в небесных кровех.

\inynen

\Bogorodichen{Предстательство христиан непостыдное, ходатайство ко Творцу непреложное, не презри грешных молений гласы, но предвари, яко Благая, на помощь нас, верно зовущих Ти: ускори на молитву, и потщися на умоление, предстательствующи присно, Богородице, чтущих Тя.}

Господи, помилуй. \myemph{ (12 раз)}

Слава Отцу и Сыну и Святому Духу, и ныне и присно и во веки веков. Аминь.

Честнейшую херувим и славнейшую без сравнения Серафим, без истления Бога Слова рождшую, сущую Богородицу Тя величаем.

\end{mymulticols}

\myemph{ После Причащения да пребывает каждый в чистоте, воздержании и немногословии, чтобы достойно сохранить в себе принятого Христа.}

\mychapterending





\mypart{ВЕЛИКИЙ КАНОН (ТВОРЕНИЕ СВЯТОГО АНДРЕЯ КРИТСКОГО)}\label{_content_velikiy-kanon}
%http://www.molitvoslov.com/content/velikiy-kanon

 

%\mychapter{Весь Канон целиком на церковнославянском языке (pdf)}
%http://www.molitvoslov.com/content/Ves-Kanon-tselikom-na-tserkovnoslayavyanskom-yazyke-pdf 
 
\myfig{img/300px-Andreas_cretensis_0.jpg}

Первая седмица Великого Поста с древних времен называется «зарей воздержания» и «неделей чистой». В эту неделю Церковь убеждает своих чад выйти из того греховного состояния, в которое невоздержанием наших прародителей ниспал весь род человеческий, утратив райское блаженство, и которое каждый из нас умножает сам своими грехами, "--- выйти путем веры, молитвы, смирения и Богоугодного поста. Се время покаяния, говорит Церковь, се день спасительный, пощения вход: душе, бодрствуй, и страстей входа затвори, ко Господу взирающе  (из первой песни трипеснца на утрени понедельника первой седмицы Великого Поста).


Подобно ветхозаветной церкви, которая особенно святила первый и последний день некоторых великих праздников, православные христиане, приготовленные и воодушевленные матерними внушениями своей Церкви, издревле, согласно ее уставу, с особенным усердием и строгостью проводят первую и последнюю седмицу Великого Поста.


На первой седмице совершаются особенно продолжительные богослужения и подвиг телесного воздержания значительно более строгий, чем в последующие дни Святой Четыредесятницы. В первые четыре дня Великого Поста великое повечерие совершается с чтением Великого Покаянного Канона преп. Андрея Критского, который как бы задает «тон», определяет всю последующую тональность, «мелодию» Великого Поста. На первой седмице Поста Канон делится на четыре части. Дивное творение св. Андрея Критского полностью предлагается нашему вниманию и в четверг (точнее в среду вечером) пятой седмицы Св. Четыредесятницы для того, чтобы мы, видя приближающееся окончание Поста, не разленились к духовным подвигам, не сделались небрежными, не забылись и не перестали во всем строго следить за собой.


Каждый стих Великого Канона сопровождается псаломским припевом Помилуй мя, Боже, помилуй мя! К канону присоединяется несколько тропарей в честь самого автора "--- св. Андрея и преп. Марии Египетской. Еще при жизни св. Андрея Иерусалимская Церковь ввела у себя в употребление Великий Канон. Отправляясь в 680 году на Шестой Вселенский Собор в Константинополь, св. Андрей принес туда и сделал известным свое великое творение и житие преп. Марии Египетской, написанное его соотечественником и учителем патриархом Иерусалимским Софронием. Житие египетской подвижницы читается совместно с Великим Каноном на утрени в среду пятой седмицы Великого Поста.


Из всех молитвословий Великого Поста, больше всех других поражает душу Великий Канон. Великий Канон "--- это чудо церковной гимнографии, это тексты удивительной силы и поэтической красоты. Канон составлен в 7-м веке св. Андреем, Архиепископом Критским, составившим также много других канонов, которые Церковь использует в течение всего богослужебного года. Церковь именовала этот канон великим, не столько из-за его объема (в нем 250 тропарей или стихов), сколько по его внутреннему достоинству и силе.


Великий канон представляет собой беседу кающегося с собственной душой. Вот как он начинается:


Откуда начну плакати окаянного моего жития деяний? Кое ли положу начало, Христе, нынешнему рыданию? Но яко благоутробен даждь ми прегрешений оставление "--- с чего же начать мне каяться, ведь это так трудно.


Затем следует чудный тропарь:


Гряди, окаянная душе, с плотию твоею. Зиждителю всех исповеждься и останися прочее преждняго безсловесия, и принеси Богу в покаянии слезы.


Удивительные слова, тут и христианская антропология, и аскетика: плоть тоже должна участвовать в покаянии, как неотъемлемая часть человеческого естества.


Своего апогея эта беседа с душой, постоянные уговоры ее, призывы покаяться, достигают в кондаке, который поется после 6-й песни Канона:


Душе моя, душе моя, востани, что спиши? Конец приближается, и имаши смутитися; воспряни убо, да пощадит тя Христос Бог, везде сый, и вся исполняяй.


Эти слова произносит, обращаясь к себе, великий светильник Церкви, тот, к кому самому применимо было бы выражение, употребленное им относительно преп. Марии Египетской, которая действительно была «ангел во плоти». И вот он так к себе обратился, упрекая себя за то, что душа его спит. Если он видел себя таким, то какими должны видеть себя мы? Погруженные уже не только в беспробудный духовный сон, но в какое-то омертвение…


Когда мы внимаем словам кондака из канона святителя Андрея Критского, нужно спросить себя: что мне делать? Если бы человек как должно исполнял Божий закон, его жизнь была бы заполнена совсем другим содержанием. Вот для того Церковь и предлагает нам этот глубокий, проникновенный великопостный покаянный канон, чтобы мы заглянули поглубже в свою душу и посмотрели бы, что там. А душа-то спит… В этом и горе наше и беда наша.


В замечательной молитве преп. Ефрема Сирина, которую мы повторяем в продолжение всего Великого Поста, говорится: Господи Царю, даруй ми зрети моя прегрешения! "--- Я их не вижу, душа моя заснула, задремала и я этих грехов, как должно, даже и не вижу. Как же я буду каяться в них! И вот потому нужно в дни Великого Поста побольше сосредотачиваться на себе, оценивая свою жизнь и ее содержание евангельской мерой, а не какой-нибудь другой.


К основным особенностям Великого канона относится очень широкое использование образов и сюжетов из Священного Писания, как Ветхого, так и Нового Завета. Жаль, что мы плохо знакомы со Святой Библией. Многим из нас имена людей, упомянутых в Великом Каноне, ничего не говорят, потому что мы плохо знаем Библию.


А между тем Библия "--- это не только история израильского народа, но и грандиозная летопись души человеческой "--- души, которая падала и вставала перед лицом Бога, которая согрешала и каялась. Если мы посмотрим на жизнь людей, о которых говорится в Библии, то увидим, что каждый из них представлен не столько как исторический персонаж, не столько как личность, совершившая те или иные дела, сколько как человек, предстоящий перед лицом Живого Бога. Исторические и иные заслуги человека отходят на второй план, остается то, что всего важнее: сохранил человек верность Богу или нет. Если мы будем читать Библию и Великий Канон под таким углом зрения, то увидим, что многое из того, что говорится о древних праведниках и грешниках, является не чем иным, как летописью нашей души, наших падений и восстаний, наших грехов и покаяния.


Один церковный писатель по этому поводу очень кстати замечает: «Если в наши дни столь многие находят его (Великий Канон) скучным и не относящимся к нашей жизни, это происходит оттого, что вера их не питается из источника Священного Писания, которое для Отцов Церкви было именно источником их веры. Мы должны вновь научиться воспринимать мир таким, каким он открывается нам в Библии, научиться жить в этом библейском мире; и нет лучшего способа научиться этому, как именно через церковное богослужение, которое не только передает нам библейское учение, но и открывает нам библейский образ жизни» (Протопр. Александр Шмеман, «Великий Пост», стр. 97).


Итак, в Великом Каноне перед нами проходит в лицах и событиях вся ветхозаветная и новозаветная история. Автор указывает на грехопадение прародителей и растление первобытного мира, на добродетели Ноя и нераскаянность и ожесточение жителей Содома и Гоморры, воскрешает перед нами память благочестивых патриархов и доблестных мужей: Моисея, Иисуса Навина, Гедеона и Иефая, представляет нашему взору благочестие царя Давида, его падение и умиленное покаяние, указывает на нечестие Ахава и Иезавели и на великие образцы покаяния "--- неневитян, Манассию, блудницу и благоразумного разбойника, и в особенности на Марию Египетскую, неоднократно останавливает читателя у Креста и Гроба Господня "--- везде поучая покаянию, смирению, молитве, самоотвержению. На этих примерах постоянно происходит увещание души "--- вспомни этого праведника, он так угодил Богу, вспомни и этого праведника, он так угодил "--- ты ничего подобного не сделала.


Об одних перснонажах Библии говорится в положительном смысле, о других в отрицательном, кому-то нужно подражать, а кому-то нет.


Колесничник Илия, колесницею добродетелей вшед, яко на Небеса, ношашеся превыше иногда от земных. Сего убо, душе моя, восход помышляй "--- помышляй, душа моя, о восхождении ветхозаветных праведников.


Гиезиев подражала еси, окаянная, разум скверный всегда, душе, егоже сребролюбие отложи поне на старость, бегай геенскаго огня, отступивши злых твоих "--- хотя бы в старости отринь сребролюбие Гиезии, душа, и оставив свои злодеяния, избегай геенского огня.


Как видите, тексты довольно трудные, поэтому к восприятию Великого канона необходимо заранее готовиться.


В заключительной песни первого дня после всех воспоминаний следуют тропари удивительной силы:


Закон изнеможе, празднует Евангелие, писание же все в тебе небрежено бысть, пророцы изнемогоша, и все праведное слово: струпи твои, о душе, умножишася, не сущу врачу исцеляющему тя "--- нечего вспоминать из Ветхого Завета, все бесполезно. Буду приводить тебе примеры из Нового Завета, может быть, тогда ты покаешься:


Новаго привожду ти писания указания, вводящая тя, душе, ко умилению: праведным убо поревнуй, грешных же отвращайся, и умилостиви Христа молитвами же и пощеньми и чистотою и говением.


Наконец, духовный писатель, представив все ветхозаветное, восходит к Жизнодавцу, Спасителю душ наших, восклицая, как разбойник: Помяни мя!, взывая, как мытарь: Боже милостив буди мне грешному!, подражая в неотступности Хананеянке и слепцам на распутии: Помилуй мя, сыне Давидов!, источая слезы, вместо мира, на главу и ноги Христа, подобно блуднице, и горько плача над собою, как Марфа и Мария над Лазарем.


Далее в Каноне подчеркивается, что самые страшные грешники покаялись и придут в Царство Небесное прежде нас: Христос вочеловечися, призвав к покаянию разбойники и блудницы: душе, покайся, дверь отверзеся Царствия уже, и предвосхищают о фарисее и мытарии и прелюбодеи кающиися.


Когда же, в некоем духовном ужасе, следуя издали за чудесами Спасителя и умиляясь над каждым подвигом Его земной жизни, автор Канона доходит до страшного заколения Христа, "--- силы сердца его оскудевает и, вместе со всей тварью, он умолкает на трепещущей Голгофе, в последний раз воскликнув: Судие мой и ведче мой, хотяй паки приити со ангелы, судити миру всему, милостивным твоим оком тогда, видев мя пощади и ущедри мя, Иисусе, паче всякого естества человеча согрешивша.


Великий канон, всеми средствами подвигая нас к покаянию, в последних тропарях как бы открывает нам свою «методику»: как я с тобой беседовал, душа, и праведников ветхозаветных тебе напоминал, и новозаветные образы тебе в пример приводил, и все напрасно: ихже не поревновала еси, душе, ни деянием, ни житию: но горе тебе, внегда будеши судитися "--- горе тебе, когда предстанешь на суд!


Вслушиваясь в слова Великого Канона, всматриваясь в историю жизни людей, бежавших от Бога, но настигнутых Им, людей, которые оказывались в безднах, но которых Бог выводил оттуда, подумаем о том, как Бог выводит каждого из нас из бездны греха и отчаяния для того, чтобы мы принесли Ему плоды покаяния.


Не нужно думать, что покаяние заключается в том, чтобы копаться в собственных грехах, заниматься самобичеванием, стараться выявить в себе как можно больше злого и темного. Истинное покаяние "--- это когда мы обращаемся от тьмы к свету, от греха к праведности; когда мы понимаем, что прежняя наша жизнь была недостойна высокого призвания, когда перед лицом Бога мы сознаем, как ничтожны мы сами, и сознаем, что единственная наша надежда "--- Сам Бог. Истинное покаяние "--- это когда перед лицом Бога, по слову Апостола Петра «призвавшего нас из тьмы в чудный Свой свет» (1 Пет. 2: 9), мы понимаем, что жизнь дана нам для того, чтобы стать детьми Божиими, чтобы приобщиться к Божественному свету. Истинное покаяние "--- то, которое выражается не столько в словах, но и в делах, в готовности прийти на помощь людям, в открытости по отношению к ближним, а не в обращенности на себя. Истинное покаяние "--- это когда мы понимаем, что, хотя и не в наших силах стать настоящими христианами, Бог в силах сделать нас таковыми. Как говорится в Великом Каноне, идеже бо хощет Бог, побеждается естества чин. Другими словами там, где Бог хочет, происходят сверхъестественные события: Савл становится Павлом, Иона изводится из чрева кита, Моисей проходит через море по суше, умерший Лазарь воскресает, Мария Египетская из блудницы превращается в великую праведницу. Ибо, по словам Спасителя, «человекам это невозможно, Богу же все возможно» (Мф. 19: 26).


©Протоиерей Виктор Потапов


www.stjohndc.org


февраль 2001 г.
\mychapterending

\mychapter{В понедельник первой седмицы Великого поста}
%http://www.molitvoslov.com/text571.htm 
 
\medskip\bfseries Песнь 1\normalfont{}\nopagebreak

\itshape Ирмос\normalfont{}: Помощник и Покровитель бысть мне во спасение, Сей мой Бог, и прославлю Его, Бог Отца моего, и вознесу Его: славно бо прославися. 

Помилуй мя, Боже, помилуй мя. 

Откуду начну плакати окаяннаго моего жития деяний? кое ли положу начало, Христе, нынешнему рыданию? но, яко благоутробен, даждь ми прегрешений оставление. 

Гряди, окаянная душе, с плотию твоею, Зиждителю всех исповеждься и останися прочее преждняго безсловесия, и принеси Богу в покаянии слезы. 

Первозданнаго Адама преступлению поревновав, познах себе обнажена от Бога и присносущнаго Царствия и сладости, грех ради моих. 

Увы мне, окаянная душе, что уподобилася еси первей Еве? видела бо еси зле, и уязвилася еси горце, и коснулася еси древа, и вкусила еси дерзостно безсловесныя снеди. 

Вместо Евы чувственныя мысленная ми бысть Ева, во плоти страстный помысл, показуяй сладкая и вкушаяй присно горькаго напоения. 

Достойно из Едема изгнан бысть, яко не сохранив едину Твою, Спасе, заповедь Адам: аз же что постражду, отметая всегда животная Твоя словеса? 

\itshape Слава\normalfont{}: Пресущная Троице, во Единице покланяемая, возьми бремя от мене тяжкое греховное и, яко благоутробна, даждь ми слезы умиления. 

\itshape И ныне\normalfont{}: Богородице, Надежде и Предстательство Тебе поющих, возьми бремя от мене тяжкое греховное, и, яко Владычица Чистая, кающася приими мя. 

\medskip\bfseries Песнь 2\normalfont{}\nopagebreak

\itshape Ирмос\normalfont{}: Вонми, Небо, и возглаголю, и воспою Христа, от Девы плотию пришедшаго. 

Вонми, Небо, и возглаголю, земле, внушай глас, кающийся к Богу и воспевающий Его. 

Вонми ми, Боже, Спасе мой, милостивым Твоим оком и приими мое теплое исповедание. 

Согреших паче всех человек, един согреших Тебе; но ущедри, яко Бог, Спасе, творение Твое. 

Вообразив моих страстей безобразие, любосластными стремленьми погубих ума красоту. 

Буря мя злых обдержит, благоутробне Господи; но яко Петру и мне руку простри. 

Оскверних плоти моея ризу и окалях, еже по образу, Спасе, и по подобию. 

Омрачих душевную красоту страстей сластьми и всячески весь ум персть сотворих. 

Раздрах ныне одежду мою первую, юже ми изтка Зиждитель из начала, и оттуду лежу наг. 

Облекохся в раздранную ризу, юже изтка ми змий советом, и стыждуся. 

Слезы блудницы, Щедре, и аз предлагаю, очисти мя, Спасе, благоутробием Твоим. 

Воззрех на садовную красоту и прельстихся умом: и оттуду лежу наг и срамляюся. 

Делаша на хребте моем вси начальницы страстей, продолжающе на мя беззаконие их. 

\itshape Слава\normalfont{}: Единаго Тя в Триех Лицех, Бога всех пою, Отца и Сына и Духа Святаго. 

\itshape И ныне\normalfont{}: Пречистая Богородице Дево, Едина Всепетая, моли прилежно, во еже спастися нам. 

\medskip\bfseries Песнь 3\normalfont{}\nopagebreak

\itshape Ирмос\normalfont{}: На недвижимом, Христе, камени заповедей Твоих утверди мое помышление. 

Огнь от Господа иногда Господь одождив, землю содомскую прежде попали. 

На горе спасайся, душе, якоже Лот оный, и в Сигор угонзай. 

Бегай запаления, о душе, бегай содомскаго горения, бегай тления Божественнаго пламене. 

Согреших Тебе един аз, согреших паче всех, Христе Спасе, да не презриши мене. 

Ты еси Пастырь добрый, взыщи мене, агнца, и заблуждшаго да не презриши мене. 

Ты еси сладкий Иисусе, Ты еси Создателю мой, в Тебе, Спасе, оправдаюся. 

Исповедаюся Тебе, Спасе, согреших, согреших Ти; но ослаби, остави ми, яко благоутробен. 

\itshape Слава\normalfont{}: О Троице Единице Боже, спаси нас от прелести, и искушений, и обстояний. 

\itshape И ныне\normalfont{}: Радуйся, Богоприятная утробо, радуйся, престоле Господень, радуйся, Мати Жизни нашея. 

\medskip\bfseries Песнь 4\normalfont{}\nopagebreak

\itshape Ирмос\normalfont{}: Услыша пророк пришествие Твое, Господи, и убояся, яко хощеши от Девы родитися и человеком явитися, и глаголаше: услышах слух Твой и убояхся, слава силе Твоей, Господи. 

Дел Твоих да не презриши, создания Твоего да не оставиши, Правосуде. Аще и един согреших, яко человек, паче всякаго человека, Человеколюбче; но имаши, яко Господь всех, власть оставляти грехи. 

Приближается, душе, конец, приближается, и нерадиши, ни готовишися, время сокращается: востани, близ при дверех Судия есть. Яко соние, яко цвет, время жития течет: что всуе мятемся? 

Воспряни, о душе моя, деяния твоя, яже соделала еси, помышляй, и сия пред лице твое принеси, и капли испусти слез твоих; рцы со дерзновением деяния и помышления Христу и оправдайся. 

Не бысть в житии греха, ни деяния, ни злобы, еяже аз, Спасе, не согреших, умом, и словом, и произволением, и предложением, и мыслию, и деянием согрешив, яко ин никтоже когда. 

Отсюду и осужден бых, отсюду препрен бых аз, окаянный, от своея совести, еяже ничтоже в мире нужнейше: Судие, Избавителю мой и Ведче, пощади, и избави, и спаси мя, раба Твоего. 

Лествица, юже виде древле великий в патриарсех, указание есть, душе моя, деятельнаго восхождения, разумнаго возшествия: аще хощеши убо деянием, и разумом, и зрением пожити, обновися. 

Зной дневный претерпе лишения ради патриарх и мраз нощный понесе, на всяк день снабдения творя, пасый, труждаяйся, работаяй, да две жене сочетает. 

Жены ми две разумей, деяние же и разум в зрении, Лию убо деяние, яко многочадную, Рахиль же разум, яко многотрудную; ибо кроме трудов ни деяние, ни зрение, душе, исправится. 

\itshape Слава\normalfont{}: Нераздельное Существом, Неслитное Лицы богословлю Тя, Троическое Едино Божество, яко Единоцарственное и Сопрестольное, вопию Ти песнь великую, в вышних трегубо песнословимую. 

\itshape И ныне\normalfont{}: И раждаеши, и девствуеши, и пребываеши обоюду естеством Дева, Рождейся обновляет законы естества, утроба же раждает нераждающая. Бог идеже хощет, побеждается естества чин: творит бо, елика хощет. 

\medskip\bfseries Песнь 5\normalfont{}\nopagebreak

\itshape Ирмос\normalfont{}: От нощи утренююща, Человеколюбче, просвети, молюся, и настави и мене на повеления Твоя, и научи мя, Спасе, творити волю Твою. 

В нощи житие мое преидох присно, тьма бо бысть, и глубока мне мгла, нощь греха, но яко дне сына, Спасе, покажи мя. 

Рувима подражая, окаянный аз, содеях беззаконный и законопреступный совет на Бога Вышняго, осквернив ложе мое, яко отчее он. 

Исповедаюся Тебе, Христе Царю: согреших, согреших, яко прежде Иосифа братия продавшии, чистоты плод и целомудрия. 

От сродников праведная душа связася, продася в работу сладкий, во образ Господень: ты же вся, душе, продалася еси злыми твоими. 

Иосифа праведнаго и целомудреннаго ума подражай, окаянная и неискусная душе, и не оскверняйся безсловесными стремленьми, присно беззаконнующи. 

Аще и в рове поживе иногда Иосиф, Владыко Господи, но во образ погребения и востания Твоего: аз же что Тебе когда сицевое принесу? 

\itshape Слава\normalfont{}: Тя, Троице, славим Единаго Бога: Свят, Свят, Свят еси, Отче, Сыне и Душе, Простое Существо, Единице присно покланяемая. 

\itshape И ныне\normalfont{}: Из Тебе облечеся в мое смешение, нетленная, безмужная Мати Дево, Бог, создавый веки, и соедини Себе человеческое естество. 

\medskip\bfseries Песнь 6\normalfont{}\nopagebreak

\itshape Ирмос\normalfont{}: Возопих всем сердцем моим к щедрому Богу, и услыша мя от ада преисподняго, и возведе от тли живот мой. 

Слезы, Спасе, очию моею и из глубины воздыхания чисте приношу, вопиющу сердцу: Боже, согреших Ти, очисти мя. 

Уклонилася еси, душе, от Господа твоего, якоже Дафан и Авирон, но пощади, воззови из ада преисподняго, да не пропасть земная тебе покрыет. 

Яко юница, душе, разсвирепевшая, уподобилася еси Ефрему, яко серна от тенет сохрани житие, вперивши деянием ум и зрением. 

Рука нас Моисеова да уверит, душе, како может Бог прокаженное житие убелити и очистити, и не отчайся сама себе, аще и прокаженна еси. 

\itshape Слава\normalfont{}: Троица есмь Проста, Нераздельна, раздельна Личне, и Единица есмь естеством соединена, Отец глаголет, и Сын, и Божественный Дух. 

\itshape И ныне\normalfont{}: Утроба Твоя Бога нам роди, воображена по нам: Егоже, яко Создателя всех, моли, Богородице, да молитвами Твоими оправдимся. 

Господи, помилуй. (Трижды.) 

\itshape Слава, и ныне\normalfont{}: 

\itshape Кондак, глас 6: 

\normalfont{}Душе моя, душе моя, востани, что спиши? конец приближается, и имаши смутитися: воспряни убо, да пощадит тя Христос Бог, везде сый и вся исполняяй. 

\medskip\bfseries Песнь 7\normalfont{}\nopagebreak

\itshape Ирмос\normalfont{}: Согрешихом, беззаконновахом, неправдовахом пред Тобою, ниже соблюдохом, ниже сотворихом, якоже заповедал еси нам; но не предаждь нас до конца, отцев Боже. 

Согреших, беззаконновах и отвергох заповедь Твою, яко во гресех произведохся, и приложих язвам струпы себе; но Сам мя помилуй, яко благоутробен, отцев Боже. 

Тайная сердца моего исповедах Тебе, Судии моему, виждь мое смирение, виждь и скорбь мою, и вонми суду моему ныне, и Сам мя помилуй, яко благоутробен, отцев Боже. 

Саул иногда, яко погуби отца своего, душе, ослята, внезапу царство обрете к прослутию; но блюди, не забывай себе, скотския похоти твоя произволивши паче Царства Христова. 

Давид иногда Богоотец, аще и согреши сугубо, душе моя, стрелою убо устрелен быв прелюбодейства, копием же пленен быв убийства томлением; но ты сама тяжчайшими делы недугуеши, самохотными стремленьми. 

Совокупи убо Давид иногда беззаконию беззаконие, убийству же любодейство растворив, покаяние сугубое показа абие; но сама ты, лукавнейшая душе, соделала еси, не покаявшися Богу. 

Давид иногда вообрази, списав яко на иконе песнь, еюже деяние обличает, еже содея, зовый: помилуй мя, Тебе бо единому согреших всех Богу, Сам очисти мя. 

\itshape Слава\normalfont{}: Троице Простая, Нераздельная, Единосущная и Естество Едино, Светове и Свет, и Свята Три, и Едино Свято поется Бог Троица; но воспой, прослави Живот и Животы, душе, всех Бога. 

\itshape И ныне\normalfont{}: Поем Тя, благословим Тя, покланяемся Ти, Богородительнице, яко Нераздельныя Троицы породила еси Единаго Христа Бога, и Сама отверзла еси нам, сущим на земли, Небесная. 

\medskip\bfseries Песнь 8\normalfont{}\nopagebreak

\itshape Ирмос\normalfont{}: Егоже воинства Небесная славят, и трепещут херувими и серафими, всяко дыхание и тварь, пойте, благословите и превозносите во вся веки. 

Согрешивша, Спасе, помилуй, воздвигни мой ум ко обращению, приими мя кающагося, ущедри вопиюща: согреших Ти, спаси, беззаконновах, помилуй мя. 

Колесничник Илия колесницею добродетелей вшед, яко на небеса, ношашеся превыше иногда от земных: сего убо, душе моя, восход помышляй. 

Елиссей иногда прием милоть Илиину, прият сугубую благодать от Бога; ты же, о душе моя, сея не причастилася еси благодати за невоздержание. 

Иорданова струя первее милотию Илииною Елиссеем ста сюду и сюду; ты же, о душе моя, сея не причастилася еси благодати за невоздержание. 

Соманитида иногда праведнаго учреди, о душе, нравом благим; ты же не ввела еси в дом ни странна, ни путника. Темже чертога изринешися вон, рыдающи. 

Гиезиев подражала еси, окаянная, разум скверный всегда, душе, егоже сребролюбие отложи поне на старость; бегай геенскаго огня, отступивши злых твоих. 

\itshape Слава\normalfont{}: Безначальне Отче, Сыне Собезначальне, Утешителю Благий, Душе Правый, Слова Божия Родителю, Отца Безначальна Слове, Душе Живый и Зиждяй, Троице Единице, помилуй мя. 

\itshape И ныне\normalfont{}: Яко от оброщения червленицы, Пречистая, умная багряница Еммануилева внутрь во чреве Твоем плоть исткася. Темже Богородицу воистинну Тя почитаем. 

\medskip\bfseries Песнь 9\normalfont{}\nopagebreak

\itshape Ирмос\normalfont{}: Безсеменнаго зачатия Рождество несказанное, Матере безмужныя нетленен Плод, Божие бо Рождение обновляет естества. Темже Тя вси роди, яко Богоневестную Матерь, православно величаем. 

Ум острупися, тело оболезнися, недугует дух, слово изнеможе, житие умертвися, конец при дверех. Темже, моя окаянная душе, что сотвориши, егда приидет Судия испытати твоя? 

Моисеово приведох ти, душе, миробытие, и от того все заветное Писание, поведающее тебе праведныя и неправедныя: от нихже вторыя, о душе, подражала еси, а не первыя, в Бога согрешивши. 

Закон изнеможе, празднует Евангелие, Писание же все в тебе небрежено бысть, пророцы изнемогоша и все праведное слово; струпи твои, о душе, умножишася, не сущу врачу, исцеляющему тя. 

Новаго привожду ти Писания указания, вводящая тя, душе, ко умилению: праведным убо поревнуй, грешных же отвращайся и умилостиви Христа молитвами же, и пощеньми, и чистотою, и говением. 

Христос вочеловечися, призвав к покаянию разбойники и блудницы; душе, покайся, дверь отверзеся Царствия уже, и предвосхищают е фарисее и мытари и прелюбодеи кающиися. 

Христос вочеловечися, плоти приобщився ми, и вся елика суть естества хотением исполни греха кроме, подобие тебе, о душе, и образ предпоказуя Своего снисхождения. 

Христос волхвы спасе, пастыри созва, младенец множества показа мученики, старцы прослави и старыя вдовицы, ихже не поревновала еси, душе, ни деянием, ни житию, но горе тебе, внегда будеши судитися. 

Постився Господь дний четыредесять в пустыни, последи взалка, показуя человеческое; душе, да не разленишися, аще тебе приложится враг, молитвою же и постом от ног твоих да отразится. 

\itshape Слава\normalfont{}: Отца прославим, Сына превознесем, Божественному Духу верно поклонимся, Троице Нераздельней, Единице по существу, яко Свету и Светом, и Животу и Животом, животворящему и просвещающему концы. 

\itshape И ныне\normalfont{}: Град Твой сохраняй, Богородительнице Пречистая, в Тебе бо сей верно царствуяй, в Тебе и утверждается, и Тобою побеждаяй, побеждает всякое искушение, и пленяет ратники, и проходит послушание. 

\bfseries 

Преподобне отче Андрее, моли Бога о нас.

\normalfont{}

Андрее честный и отче треблаженнейший, пастырю Критский, не престай моляся о воспевающих тя: да избавимся вси гнева и скорби, и тления, и прегрешений безмерных, чтущии твою память верно. 

Таже оба лика вкупе поют\bfseries  \normalfont{}\itshape Ирмос: 

\normalfont{}Безсеменнаго зачатия Рождество несказанное, Матере безмужныя нетленен Плод, Божие бо Рождение обновляет естества. Темже Тя вси роди, яко Богоневестную Матерь, православно величаем. \mychapterending

\mychapter{Во вторник первой седмицы Великого Поста}
%http://www.molitvoslov.com/text572.htm 
 
\medskip\bfseries Песнь 1\normalfont{}\nopagebreak

\itshape Ирмос\normalfont{}: Помощник и Покровитель бысть мне во спасение, Сей мой Бог, и прославлю Его, Бог отца моего, и вознесу Его: славно бо прославися. 

Каиново прешед убийство, произволением бых убийца совести душевней, оживив плоть и воевав на ню лукавыми моими деяньми. 

Авелеве, Иисусе, не уподобихся правде, дара Тебе приятна не принесох когда, ни деяния божественна, ни жертвы чистыя, ни жития непорочнаго. 

Яко Каин и мы, душе окаянная, всех Содетелю деяния скверная, и жертву порочную, и непотребное житие принесохом вкупе: темже и осудихомся. 

Брение Здатель живосоздав, вложил еси мне плоть и кости, и дыхание, и жизнь; но, о Творче мой, Избавителю мой и Судие, кающася приими мя. 

Извещаю Ти, Спасе, грехи, яже содеях, и души и тела моего язвы, яже внутрь убийственнии помыслы разбойнически на мя возложиша. 

Аще и согреших, Спасе, но вем, яко Человеколюбец еси, наказуеши милостивно и милосердствуеши тепле: слезяща зриши и притекаеши, яко отец, призывая блуднаго. 

\itshape Слава\normalfont{}: Пресущная Троице, во Единице покланяемая, возьми бремя от мене тяжкое греховное и, яко благоутробна, даждь ми слезы умиления. 

\itshape И ныне\normalfont{}: Богородице, Надежде и Предстательство Тебе поющих, возьми бремя от мене тяжкое греховное и, яко Владычица Чистая, кающася приими мя. 

\medskip\bfseries Песнь 2\normalfont{}\nopagebreak

\itshape Ирмос\normalfont{}: Вонми, Небо, и возглаголю, и воспою Христа, от Девы плотию пришедшаго. 

Сшиваше кожныя ризы грех мне, обнаживый мя первыя боготканныя одежды. 

Обложен есмь одеянием студа, якоже листвием смоковным, во обличение моих самовластных страстей. 

Одеяхся в срамную ризу и окровавленную студно течением страстнаго и любосластнаго живота. 

Впадох в страстную пагубу и в вещественную тлю, и оттоле до ныне враг мне досаждает. 

Любовещное и любоименное житие невоздержанием, Спасе, предпочет ныне, тяжким бременем обложен есмь. 

Украсих плотский образ скверных помышлений различным обложением и осуждаюся. 

Внешним прилежно благоукрашением единем попекохся, внутреннюю презрев Богообразную скинию. 

Погребох перваго образа доброту, Спасе, страстьми, юже, яко иногда драхму, взыскав, обрящи. 

Согреших, якоже блудница, вопию Ти: един согреших Тебе, яко миро, приими, Спасе, и моя слезы. 

Очисти, якоже мытарь, вопию Ти, Спасе, очисти мя: никтоже бо сущих из Адама, якоже аз, согреших Тебе. 

\itshape Слава\normalfont{}: Единаго Тя в Триех Лицех, Бога всех пою, Отца и Сына и Духа Святаго. 

\itshape И ныне\normalfont{}: Пречистая Богородице Дево, Едина Всепетая, моли прилежно, во еже спастися нам. 

\medskip\bfseries Песнь 3\normalfont{}\nopagebreak

\itshape Ирмос\normalfont{}: Утверди, Господи, на камени заповедей Твоих подвигшееся сердце мое, яко Един Свят еси и Господь. 

Источник живота стяжах Тебе, смерти Низложителя, и вопию Ти от сердца моего прежде конца: согреших, очисти и спаси мя. 

Согреших, Господи, согреших Тебе, очисти мя: несть бо иже кто согреши в человецех, егоже не превзыдох прегрешеньми. 

При Нои, Спасе, блудствовавшия подражах, онех наследствовах осуждение в потопе погружения. 

Хама онаго, душе, отцеубийца подражавши, срама не покрыла еси искренняго, вспять зря возвратившися. 

Запаления, якоже Лот, бегай, душе моя, греха: бегай Содомы и Гоморры, бегай пламене всякаго безсловеснаго желания. 

Помилуй, Господи, помилуй мя, вопию Ти, егда приидеши со ангелы Твоими воздати всем по достоянию деяний. 

\itshape Слава\normalfont{}: Троице Простая, Несозданная, Безначальное Естество, в Троице певаемая Ипостасей, спаси ны, верою покланяющияся державе Твоей. 

\itshape И ныне\normalfont{}: От Отца безлетна Сына в лето, Богородительнице, неискусомужно родила еси, странное чудо, пребывши Дева доящи. 

\medskip\bfseries Песнь 4\normalfont{}\nopagebreak

\itshape Ирмос\normalfont{}: Услыша пророк пришествие Твое, Господи, и убояся, яко хощеши от Девы родитися и человеком явитися, и глаголаше: услышах слух Твой и убояхся, слава силе Твоей, Господи. 

Бди, о душе моя, изрядствуй, якоже древле великий в патриарсех, да стяжеши деяние с разумом, да будеши ум, зряй Бога, и достигнеши незаходящий мрак в видении, и будеши великий купец. 

Дванадесять патриархов великий в патриарсех детотворив, тайно утверди тебе лествицу деятельнаго, душе моя, восхождения: дети, яко основания, степени, яко восхождения, премудренно подложив. 

Исава возненавиденнаго подражала еси, душе, отдала еси прелестнику твоему первыя доброты первенство и отеческия молитвы отпала еси, и дважды поползнулася еси, окаянная, деянием и разумом: темже ныне покайся. 

Едом Исав наречеся, крайняго ради женонеистовнаго смешения: невоздержанием бо присно разжигаем и сластьми оскверняем, Едом именовася, еже глаголется разжжение души любогреховныя. 

Иова на гноищи слышавши, о душе моя, оправдавшагося, того мужеству не поревновала еси, твердаго не имела еси предложения во всех, яже веси, и имиже искусилася еси, но явилася еси нетерпелива. 

Иже первее на престоле, наг ныне на гноище гноен, многий в чадех и славный, безчаден и бездомок напрасно: палату убо гноище и бисерие струпы вменяше. 

\itshape Слава\normalfont{}: Нераздельное Существом, Неслитное Лицы богословлю Тя, Троическое Едино Божество, яко Единоцарственное и Сопрестольное, вопию Ти песнь великую, в вышних трегубо песнословимую. 

\itshape И ныне\normalfont{}: И раждаеши, и девствуеши, и пребываеши обоюду естеством Дева, Рождейся обновляет законы естества, утроба же раждает нераждающая. Бог идеже хощет, побеждается естества чин: творит бо, елика хощет. 

\medskip\bfseries Песнь 5\normalfont{}\nopagebreak

\itshape Ирмос\normalfont{}: От нощи утренююща, Человеколюбче, просвети, молюся, и настави и мене на повеления Твоя, и научи мя, Спасе, творити волю Твою. 

Моисеов слышала еси ковчежец, душе, водами, волнами носим речными, яко в чертозе древле бегающий дела, горькаго совета фараонитска. 

Аще бабы слышала еси, убивающия иногда безвозрастное мужеское, душе окаянная, целомудрия деяние, ныне, яко великий Моисей, сси премудрость. 

Яко Моисей великий египтянина, ума, уязвивши, окаянная, не убила еси, душе; и како вселишися, глаголи, в пустыню страстей покаянием? 

В пустыню вселися великий Моисей; гряди убо, подражай того житие, да и в купине Богоявления, душе, в видении будеши. 

Моисеов жезл воображай, душе, ударяющий море и огустевающий глубину, во образ Креста Божественнаго: имже можеши и ты великая совершити. 

Аарон приношаше огнь Богу непорочный, нелестный; но Офни и Финеес, яко ты, душе, приношаху чуждее Богу, оскверненное житие. 

\itshape Слава\normalfont{}: Тя, Троице, славим Единаго Бога: Свят, Свят, Свят еси, Отче, Сыне и Душе, Простое Существо, Единице присно покланяемая. 

\itshape И ныне\normalfont{}: Из Тебе облечеся в мое смешение, нетленная, безмужная Мати Дево, Бог, создавый веки, и соедини Себе человеческое естество. 

\medskip\bfseries Песнь 6\normalfont{}\nopagebreak

\itshape Ирмос\normalfont{}: Возопих всем сердцем моим к щедрому Богу, и услыша мя от ада преисподняго, и возведе от тли живот мой. 

Волны, Спасе, прегрешений моих, яко в мори Чермнем возвращающеся, покрыша мя внезапу, яко египтяны иногда и тристаты. 

Неразумное, душе, произволение имела еси, яко прежде Израиль: Божественныя бо манны предсудила еси безсловесно любосластное страстей объядение. 

Кладенцы, душе, предпочла еси хананейских мыслей паче жилы камене, из негоже премудрости река, яко чаша, проливает токи богословия. 

Свиная мяса и котлы и египетскую пищу, паче Небесныя, предсудила еси, душе моя, якоже древле неразумнии людие в пустыни. 

Яко удари Моисей, раб Твой, жезлом камень, образно животворивая ребра Твоя прообразоваше, из нихже вси питие жизни, Спасе, почерпаем. 

Испытай, душе, и смотряй, якоже Иисус Навин, обетования землю, какова есть, и вселися в ню благозаконием. 

\itshape Слава\normalfont{}: Троица есмь Проста, Нераздельна, раздельна Личне и Единица есмь естеством соединена, Отец глаголет, и Сын, и Божественный Дух. 

\itshape И ныне\normalfont{}: Утроба Твоя Бога нам роди, воображена по нам: Егоже, яко Создателя всех, моли, Богородице, да молитвами Твоими оправдимся. 

Господи, помилуй. (Трижды.) 

\itshape Слава, и ныне\normalfont{}: 

\medskip\bfseries Кондак, глас 6: \normalfont{}\nopagebreak

Душе моя, душе моя, востани, что спиши? конец приближается, и имаши смутитися: воспряни убо, да пощадит тя Христос Бог, везде сый и вся исполняяй. 

\medskip\bfseries Песнь7\normalfont{}\nopagebreak

\itshape Ирмос\normalfont{}: Согрешихом, беззаконновахом, неправдовахом пред Тобою, ниже соблюдохом, ниже сотворихом, якоже заповедал еси нам; но не предаждь нас до конца, отцев Боже. 

Кивот яко ношашеся на колеснице, Зан оный, егда превращшуся тельцу, точию коснуся, Божиим искусися гневом; но того дерзновения убежавши, душе, почитай Божественная честне. 

Слышала еси Авессалома, како на естество воста, познала еси того скверная деяния, имиже оскверни ложе Давида отца; но ты подражала еси того страстная и любосластная стремления. 

Покорила еси неработное твое достоинство телу твоему, иного бо Ахитофела обретше врага, душе, снизшла еси сего советом; но сия разсыпа Сам Христос, да ты всяко спасешися. 

Соломон чудный и благодати премудрости исполненный, сей лукавое иногда пред Богом сотворив, отступи от Него; емуже ты проклятым твоим житием, душе, уподобилася еси. 

Сластьми влеком страстей своих, оскверняшеся, увы мне, рачитель премудрости, рачитель блудных жен и странен от Бога: егоже ты подражала еси умом, о душе, сладострастьми скверными. 

Ровоаму поревновала еси, не послушавшему совета отча, купно же и злейшему рабу Иеровоаму, прежнему отступнику, душе, но бегай подражания и зови Богу: согреших, ущедри мя. 

\itshape Слава\normalfont{}: Троице Простая, Нераздельная, Единосущная и Естество Едино, Светове и Свет, и Свята Три, и Едино Свято поется Бог Троица; но воспой, прослави Живот и Животы, душе, всех Бога. 

\itshape И ныне\normalfont{}: Поем Тя, благословим Тя, покланяемся Ти, Богородительнице, яко Неразлучныя Троицы породила еси Единаго Христа Бога и Сама отверзла еси нам, сущим на земли, Небесная. 

\medskip\bfseries Песнь 8\normalfont{}\nopagebreak

\itshape Ирмос\normalfont{}: Егоже воинства Небесная славят, и трепещут херувими и серафими, всяко дыхание и тварь, пойте, благословите и превозносите во вся веки. 

Ты Озии, душе, поревновавши, сего прокажение в себе стяжала еси сугубо: безместная бо мыслиши, беззаконная же дееши; остави, яже имаши, и притецы к покаянию. 

Ниневитяны, душе, слышала еси кающияся Богу, вретищем и пепелом, сих не подражала еси, но явилася еси злейшая всех, прежде закона и по законе прегрешивших. 

В рове блата слышала еси Иеремию, душе, града Сионя рыданьми вопиюща и слез ищуща: подражай сего плачевное житие и спасешися. 

Иона в Фарсис побеже, проразумев обращение ниневитянов, разуме бо, яко пророк, Божие благоутробие: темже ревноваше пророчеству не солгатися. 

Даниила в рове слышала еси, како загради уста, о душе, зверей; уведела еси, како отроцы, иже о Азарии, погасиша верою пещи пламень горящий. 

Ветхаго Завета вся приведох ти, душе, к подобию; подражай праведных боголюбивая деяния, избегни же паки лукавых грехов. 

\itshape Слава\normalfont{}: Безначальне Отче, Сыне Собезначальне, Утешителю Благий, Душе Правый, Слова Божия Родителю, Отца Безначальна Слове, Душе Живый и Зиждяй, Троице Единице, помилуй мя. 

\itshape И ныне\normalfont{}: Яко от оброщения червленицы, Пречистая, умная багряница Еммануилева внутрь во чреве Твоем плоть исткася. Темже Богородицу воистинну Тя почитаем. 

\medskip\bfseries Песнь 9\normalfont{}\nopagebreak

\itshape Ирмос\normalfont{}: Безсеменнаго зачатия Рождество несказанное, Матере безмужныя нетленен Плод, Божие бо Рождение обновляет естества. Темже Тя вси роди, яко Богоневестную Матерь, православно величаем. 

Христос искушашеся, диавол искушаше, показуя камение, да хлеби будут, на гору возведе видети вся царствия мира во мгновении; убойся, о душе, ловления, трезвися, молися на всякий час Богу. 

Горлица пустыннолюбная, глас вопиющаго возгласи, Христов светильник, проповедуяй покаяние, Ирод беззаконнова со Иродиадою. Зри, душе моя, да не увязнеши в беззаконныя сети, но облобызай покаяние. 

В пустыню вселися благодати Предтеча, и Иудея вся и Самария, слышавше, течаху и исповедаху грехи своя, крещающеся усердно: ихже ты не подражала еси, душе. 

Брак убо честный и ложе нескверно, обоя бо Христос прежде благослови, плотию ядый и в Кане же на браце воду в вино совершая, и показуя первое чудо, да ты изменишися, о душе. 

Разслабленнаго стягну Христос, одр вземша, и юношу умерша воздвиже, вдовиче рождение, и сотнича отрока, и самаряныне явися, в дусе службу тебе, душе, предживописа. 

Кровоточивую исцели прикосновением края ризна Господь, прокаженныя очисти, слепыя и хромыя просветив, исправи, глухия же, и немыя, и ничащия низу исцели словом: да ты спасешися, окаянная душе. 

\itshape Слава\normalfont{}: Отца прославим, Сына превознесем, Божественному Духу верно поклонимся, Троице Нераздельней, Единице по существу, яко Свету и Светом, и Животу и Животом, Животворящему и Просвещающему концы. 

\itshape И ныне\normalfont{}: Град Твой сохраняй, Богородительнице Пречистая, в Тебе бо сей верно царствуяй, в Тебе и утверждается, и Тобою побеждаяй, побеждает всякое искушение, и пленяет ратники, и проходит послушание. 

Преподобне отче Андрее, моли Бога о нас. 

Андрее честный и отче треблаженнейший, пастырю Критский, не престай моляся о воспевающих тя, да избавимся вси гнева, и скорби, и тления, и прегрешений безмерных, чтущии твою память верно. 

Таже оба лика вкупе поют \itshape Ирмос\normalfont{}: 

Безсеменнаго зачатия Рождество несказанное, Матере безмужныя нетленен Плод, Божие бо Рождение обновляет естества. Темже Тя вси роди, яко Богоневестную Матерь, православно величаем. \mychapterending

\mychapter{В среду первой седмицы Великого Поста}
%http://www.molitvoslov.com/text573.htm 
 
\medskip\bfseries Песнь 1\normalfont{}\nopagebreak

\itshape Ирмос\normalfont{}: Помощник и Покровитель бысть мне во спасение, Сей мой Бог, и прославлю Его, Бог отца моего, и вознесу Его: славно бо прославися. 

От юности, Христе, заповеди Твоя преступих, всестрастно небрегий, унынием преидох житие. Темже зову Ти, Спасе: поне на конец спаси мя. 

Повержена мя, Спасе, пред враты Твоими, поне на старость не отрини мене во ад тща, но прежде конца, яко Человеколюбец, даждь ми прегрешений оставление. 

Богатство мое, Спасе, изнурив в блуде, пуст есмь плодов благочестивых, алчен же зову: Отче щедрот, предварив, Ты мя ущедри. 

В разбойники впадый аз есмь помышленьми моими, весь от них уязвихся ныне и исполнихся ран, но, Сам ми представ, Христе Спасе, исцели. 

Священник, мя предвидев, мимо иде, и левит, видя в лютых нага, презре, но из Марии возсиявый Иисусе, Ты, представ, ущедри мя. 

Преподобная мати Марие, моли Бога о нас. 

Ты ми даждь светозарную благодать от Божественнаго свыше промышления избежати страстей омрачения и пети усердно Твоего, Марие, жития красная исправления. 

\itshape Слава\normalfont{}: Пресущная Троице, во Единице покланяемая, возьми бремя от мене тяжкое греховное и, яко благоутробна, даждь ми слезы умиления. 

\itshape И ныне\normalfont{}: Богородице, Надежде и Предстательство Тебе поющих, возьми бремя от мене тяжкое греховное и, яко Владычица Чистая, кающася приими мя. 

\medskip\bfseries Песнь 2\normalfont{}\nopagebreak

\itshape Ирмос\normalfont{}: Вонми, Небо, и возглаголю, и воспою Христа, от Девы плотию пришедшаго. 

Поползохся, яко Давид, блудно и осквернихся, но омый и мене, Спасе, слезами. 

Ни слез, ниже покаяния имам, ниже умиления. Сам ми сия, Спасе, яко Бог, даруй. 

Погубих первозданную доброту и благолепие мое и ныне лежу наг и стыждуся. 

Дверь Твою не затвори мне тогда, Господи, Господи, но отверзи ми сию, кающемуся Тебе. 

Внуши воздыхания души моея и очию моею приими капли, Спасе, и спаси мя. 

Человеколюбче, хотяй всем спастися, Ты воззови мя и приими, яко благ, кающагося. 

Пресвятая Богородице, спаси нас. 

Пречистая Богородице Дево, Едина Всепетая, моли прилежно, во еже спастися нам. 

Иный. \itshape Ирмос\normalfont{}: Видите, видите, яко Аз есмь Бог, манну одождивый и воду из камене источивый древле в пустыни людем Моим, десницею единою и крепостию Моею. 

Видите, видите, яко Аз есмь Бог, внушай, душе моя, Господа вопиюща, и удалися прежняго греха, и бойся, яко неумытнаго и яко Судии и Бога. 

Кому уподобилася еси, многогрешная душе? токмо первому Каину и Ламеху оному, каменовавшая тело злодействы и убившая ум безсловесными стремленьми. 

Вся прежде закона претекши, о душе, Сифу не уподобилася еси, ни Еноса подражала еси, ни Еноха преложением, ни Ноя, но явилася еси убога праведных жизни. 

Едина отверзла еси хляби гнева Бога Твоего, душе моя, и потопила еси всю, якоже землю, плоть, и деяния, и житие, и пребыла еси вне спасительнаго ковчега. 

Преподобная мати Марие, моли Бога о нас. 

Всем усердием и любовию притекла еси Христу, первый греха путь отвращши, и в пустынях непроходимых питающися, и Того чисте совершающи Божественныя заповеди. 

\itshape Слава\normalfont{}: Безначальная, Несозданная Троице, Нераздельная Единице, кающася мя приими, согрешивша спаси, Твое есмь создание, не презри, но пощади и избави огненнаго мя осуждения. 

\itshape И ныне\normalfont{}: Пречистая Владычице, Богородительнице, Надеждо к Тебе притекающих и пристанище сущих в бури, Милостиваго и Создателя и Сына Твоего умилостиви и мне молитвами Твоими.

\medskip\bfseries Песнь 3\normalfont{}\nopagebreak

\itshape Ирмос\normalfont{}: Утверди, Господи, на камени заповедей Твоих подвигшееся сердце мое, яко Един Свят еси и Господь. 

Благословения Симова не наследовала еси, душе окаянная, ни пространное одержание, якоже Иафеф, имела еси на земли оставления. 

От земли Харран изыди от греха, душе моя, гряди в землю, точащую присноживотное нетление, еже Авраам наследствова. 

Авраама слышала еси, душе моя, древле оставльша землю отечества и бывша пришельца, сего произволению подражай. 

У дуба Мамврийскаго учредив патриарх ангелы, наследствова по старости обетования ловитву. 

Исаака, окаянная душе моя, разумевши новую жертву, тайно всесожженную Господеви, подражай его произволению. 

Исмаила слышала еси, трезвися, душе моя, изгнана, яко рабынино отрождение, виждь, да не како подобно что постраждеши, ласкосердствующи. 

Преподобная мати Марие, моли Бога о нас. 

Содержим есмь бурею и треволнением согрешений, но сама мя, мати, ныне спаси и к пристанищу Божественнаго покаяния возведи. 

Преподобная мати Марие, моли Бога о нас. 

Рабское моление \itshape И ныне\normalfont{}, преподобная, принесши ко благоутробней молитвами твоими Богородице, отверзи ми Божественныя входы. 

\itshape Слава\normalfont{}: Троице Простая, Несозданная, Безначальное Естество, в Троице певаемая Ипостасей, спаси ны, верою покланяющияся державе Твоей. 

\itshape И ныне\normalfont{}: От Отца безлетна Сына в лето, Богородительнице, неискусомужно родила еси, странное чудо, пребывши Дева доящи. 

\medskip\bfseries Песнь 4\normalfont{}\nopagebreak

\itshape Ирмос\normalfont{}: Услыша пророк пришествие Твое, Господи, и убояся, яко хощеши от Девы родитися и человеком явитися, и глаголаше: услышах слух Твой и убояхся, слава силе Твоей, Господи. 

Тело осквернися, дух окаляся, весь острупихся, но яко врач, Христе, обоя покаянием моим уврачуй, омый, очисти, покажи, Спасе мой, паче снега чистейша. 

Тело Твое и кровь, распинаемый о всех, положил еси, Слове: тело убо, да мя обновиши, кровь, да омыеши мя. Дух же предал еси, да мя приведеши, Христе, Твоему Родителю. 

Соделал еси спасение посреде земли, Щедре, да спасемся. Волею на древе распялся еси, Едем затворенный отверзеся, горняя и дольняя тварь, языцы вси, спасени, покланяются Тебе. 

Да будет ми купель кровь из ребр Твоих, вкупе и питие, источившее воду оставления, да обоюду очищаюся, помазуяся и пия, яко помазание и питие, Слове, животочная Твоя словеса. 

Чашу Церковь стяжа, ребра Твоя живоносная, из нихже сугубыя нам источи токи оставления и разума во образ древняго и новаго, двоих вкупе заветов, Спасе наш. 

Наг есмь чертога, наг есмь и брака, купно и вечери; светильник угасе, яко безъелейный, чертог заключися мне спящу, вечеря снедеся, аз же по руку и ногу связан, вон низвержен есмь. 

Слава: Нераздельное Существом, Неслитное Лицы богословлю Тя, Троическое Едино Божество, яко Единоцарственное и Сопрестольное, вопию Ти песнь великую, в вышних трегубо песнословимую. 

\itshape И ныне\normalfont{}: И раждаеши, и девствуеши, и пребываеши обоюду естеством Дева, Рождейся обновляет законы естества, утроба же раждает нераждающая. Бог идеже хощет, побеждается естества чин: творит бо, елика хощет. 

\medskip\bfseries Песнь 5\normalfont{}\nopagebreak

\itshape Ирмос\normalfont{}: От нощи утренююща, Человеколюбче, просвети, молюся, и настави и мене на повеления Твоя, и научи мя, Спасе, творити волю Твою. 

Яко тяжкий нравом, фараону горькому бых, Владыко, Ианни и Иамври, душею и телом, и погружен умом, но помози ми. 

Калом смесихся, окаянный, умом, омый мя, Владыко, банею моих слез, молю Тя, плоти моея одежду убелив, яко снег. 

Аще испытаю моя дела, Спасе, всякаго человека превозшедша грехами себе зрю, яко разумом мудрствуяй, согреших не неведением. 

Пощади, пощади, Господи, создание Твое, согреших, ослаби ми, яко естеством чистый Сам сый Един, и ин разве Тебе никтоже есть кроме скверны. 

Мене ради Бог сый, вообразился еси в мя, показал еси чудеса, исцелив прокаженныя и разслабленнаго стягнув, кровоточивыя ток уставил еси, Спасе, прикосновением риз. 

Преподобная мати Марие, моли Бога о нас. 

Струи Иорданския прешедши, обрела еси покой безболезненный, плоти сласти избежавши, еяже и нас изми твоими молитвами, преподобная. 

Слава: Тя, Троице, славим Единаго Бога: Свят, Свят, Свят еси, Отче, Сыне и Душе, Простое Существо, Единице присно покланяемая. 

\itshape И ныне\normalfont{}: Из Тебе облечеся в мое смешение, нетленная, безмужная Мати Дево, Бог, создавый веки, и соедини Себе человеческое естество.

\medskip\bfseries Песнь 6\normalfont{}\nopagebreak

\itshape Ирмос\normalfont{}: Возопих всем сердцем моим к щедрому Богу, и услыша мя от ада преисподняго, и возведе от тли живот мой. 

Востани и побори, яко Иисус Амалика, плотския страсти, и гаваониты, лестныя помыслы, присно побеждающи. 

Преиди, времене текущее естество, яко прежде ковчег, и земли оныя буди во одержании обетования, душе, Бог повелевает. 

Яко спасл еси Петра, возопивша, спаси, предварив мя, Спасе, от зверя избави, простер Твою руку, и возведи из глубины греховныя. 

Пристанище Тя вем утишное, Владыко, Владыко Христе, но от незаходимых глубин греха и отчаяния мя, предварив, избави. 

\itshape Слава\normalfont{}: Троица есмь Проста, Нераздельна, раздельна Личне, и Единица есмь естеством соединена, Отец глаголет, и Сын, и Божественный Дух. 

\itshape И ныне\normalfont{}: Утроба Твоя Бога нам роди, воображена по нам: Егоже, яко Создателя всех, моли, Богородице, да молитвами Твоими оправдимся. 

Господи, помилуй. (Трижды.) 

\itshape Слава, и ныне\normalfont{}: 

\itshape Кондак, глас 6:

\normalfont{}

Душе моя, душе моя, востани, что спиши? конец приближается, и имаши смутитися: воспряни убо, да пощадит тя Христос Бог, везде сый и вся исполняяй. 

\medskip\bfseries Песнь 7\normalfont{}\nopagebreak

\itshape Ирмос\normalfont{}: Согрешихом, беззаконновахом, неправдовахом пред Тобою, ниже соблюдохом, ниже сотворихом, якоже заповедал еси нам; но не предаждь нас до конца, отцев Боже. 

Манассиева собрала еси согрешения изволением, поставльши яко мерзости страсти и умноживши, душе, негодования, но того покаянию ревнующи тепле, стяжи умиление. 

Ахаавовым поревновала еси сквернам, душе моя, увы мне, была еси плотских скверн пребывалище и сосуд срамлен страстей, но из глубины твоея воздохни и глаголи Богу грехи твоя. 

Заключися тебе небо, душе, и глад Божий постиже тя, егда Илии Фесвитянина, якоже Ахаав, не покорися словесем иногда, но Сараффии уподобився, напитай пророчу душу. 

Попали Илия иногда дващи пятьдесят Иезавелиных, егда студныя пророки погуби, во обличение Ахаавово, но бегай подражания двою, душе, и укрепляйся. 

\itshape Слава\normalfont{}: Троице Простая, Нераздельная, Единосущная, и Естество Едино, Светове и Свет, и Свята Три, и Едино Свято поется Бог Троица; но воспой, прослави Живот и Животы, душе, всех Бога. 

\itshape И ныне\normalfont{}: Поем Тя, благословим Тя, покланяемся Ти, Богородительнице, яко Нераздельныя Троицы породила еси Единаго Христа Бога и Сама отверзла еси нам, сущим на земли, Небесная. 

\medskip\bfseries Песнь 8\normalfont{}\nopagebreak

\itshape Ирмос\normalfont{}: Егоже воинства Небесная славят, и трепещут херувими и серафими, всяко дыхание и тварь, пойте, благословите и превозносите во вся веки. 

Правосуде Спасе, помилуй и избави мя огня и прещения, еже имам на суде праведно претерпети; ослаби ми прежде конца добродетелию и покаянием. 

Яко разбойник, вопию Ти: помяни мя; яко Петр, плачу горце: ослаби ми, Спасе; зову, яко мытарь, слезю, яко блудница; приими мое рыдание, якоже иногда хананеино. 

Гноение, Спасе, исцели смиренныя моея души, Едине Врачу, пластырь мне наложи, и елей, и вино, дела покаяния, умиление со слезами. 

Хананею и аз подражая, помилуй мя, вопию, Сыне Давидов; касаюся края ризы, яко кровоточивая, плачу, яко Марфа и Мария над Лазарем. 

\itshape Слава\normalfont{}: Безначальне Отче, Сыне Собезначальне, Утешителю Благий, Душе Правый, Слова Божия Родителю, Отца Безначальна Слове, Душе Живый и Зиждяй, Троице Единице, помилуй мя. 

\itshape И ныне\normalfont{}: Яко от оброщения червленицы, Пречистая, умная багряница Еммануилева внутрь во чреве Твоем плоть исткася. Темже Богородицу воистинну Тя почитаем. 

\medskip\bfseries Песнь 9\normalfont{}\nopagebreak

\itshape Ирмос\normalfont{}: Безсеменнаго зачатия Рождество несказанное, Матере безмужныя нетленен Плод, Божие бо Рождение обновляет естества. Темже Тя вси роди, яко Богоневестную Матерь, православно величаем. 

Недуги исцеляя, нищим благовествоваше Христос Слово, вредныя уврачева, с мытари ядяше, со грешники беседоваше, Иаировы дщере душу предумершую возврати осязанием руки. 

Мытарь спасашеся, и блудница целомудрствоваше, и фарисей, хваляся, осуждашеся. Ов убо: очисти мя; ова же: помилуй мя; сей же величашеся вопия: Боже, благодарю Тя, и прочия безумныя глаголы. 

Закхей мытарь бе, но обаче спасашеся, и фарисей Симон соблажняшеся, и блудница приимаше оставительная разрешения от Имущаго крепость оставляти грехи, юже, душе, потщися подражати. 

Блуднице, о окаянная душе моя, не поревновала еси, яже приимши мира алавастр, со слезами мазаше нозе Спасове, отре же власы, древних согрешений рукописание Раздирающаго ея. 

Грады, имже даде Христос благовестие, душе моя, уведала еси, како прокляти быша. Убойся указания, да не будеши якоже оны, ихже содомляном Владыка уподобив, даже до ада осуди. 

Да не горшая, о душе моя, явишися отчаянием, хананеи веру слышавшая, еяже дщи словом Божиим исцелися; Сыне Давидов, спаси и мене, воззови из глубины сердца, якоже она Христу. 

Слава: Отца прославим, Сына превознесем, Божественному Духу верно поклонимся, Троице Нераздельней, Единице по Существу, яко Свету и Светом, и Животу и Животом, Животворящему и Просвещающему концы. 

\itshape И ныне\normalfont{}: Град Твой сохраняй, Богородительнице Пречистая, в Тебе бо сей верно царствуяй, в Тебе и утверждается, и Тобою побеждаяй, побеждает всякое искушение, и пленяет ратники, и проходит послушание. 

\bfseries 

Преподобне отче Андрее, моли Бога о нас.\normalfont{}

Андрее честный и отче треблаженнейший, пастырю Критский, не престай моляся о воспевающих тя, да избавимся вси гнева, и скорби, и тления, и прегрешений безмерных, чтущии твою память верно. 

\itshape 

Таже оба лика вкупе поют Ирмос\normalfont{}:

Безсеменнаго зачатия Рождество несказанное, Матере безмужныя нетленен Плод, Божие бо Рождение обновляет естества. Темже Тя вси роди, яко Богоневестную Матерь, православно величаем.

 

\mychapterending

\mychapter{В четверг первой седмицы Великого Поста}
%http://www.molitvoslov.com/text574.htm 
 


\medskip\bfseries Песнь 1\normalfont{}\nopagebreak

\itshape Ирмос\normalfont{}: Помощник и Покровитель бысть мне во спасение, Сей мой Бог, и прославлю Его, Бог Отца моего, и вознесу Его: славно бо прославися. 

Агнче Божий, вземляй грехи всех, возьми бремя от мене тяжкое греховное, и, яко благоутробен, даждь ми слезы умиления. 

Тебе припадаю, Иисусе, согреших Ти, очисти мя, возьми бремя от мене тяжкое греховное и, яко благоутробен, даждь ми слезы умиления. 

Не вниди со мною в суд, нося моя деяния, словеса изыскуя и исправляя стремления. Но в щедротах Твоих презирая моя лютая, спаси мя, Всесильне. 

Покаяния время, прихожду Ти, Создателю моему: возьми бремя от мене тяжкое греховное и, яко благоутробен, даждь ми слезы умиления. 

Богатство душевное иждив грехом, пуст есмь добродетелей благочестивых, гладствуя же зову: милости подателю Господи, спаси мя. 

\bfseries 

Преподобная Мати Марие, моли Бога о нас.

\normalfont{}

Приклоньшися Христовым Божественным законом, к сему приступила еси, сладостей неудержимая стремления оставивши, и всякую добродетель всеблагоговейно, яко едину, исправила еси. 

\itshape Слава\normalfont{}: Пресущественная Троице, во Единице покланяемая, возьми бремя от мене тяжкое греховное и, яко благоутробна, даждь ми слезы умиления. 

\itshape И ныне\normalfont{}: Богородице, Надежде и Предстательство Тебе поющих, возьми бремя от мене тяжкое греховное и, яко Владычица Чистая, кающася приими мя. 


\medskip\bfseries Песнь 2\normalfont{}\nopagebreak

\itshape Ирмос\normalfont{}: Видите, видите, яко Аз есмь Бог, манну одождивый и воду из камене источивый древле в пустыни людем Моим, десницею единою и крепостию Моею. 

Мужа убих, глаголет, в язву мне и юношу в струп, Ламех, рыдая, вопияше; ты же не трепещеши, о душе моя, окалявши плоть и ум осквернивши. 

Столп умудрила еси создати, о душе, и утверждение водрузити твоими похотьми, аще не бы Зиждитель удержал советы твоя и низвергл на землю ухищрения твоя. 

О како поревновах Ламеху, первому убийце, душу, яко мужа, ум, яко юношу, яко брата же моего, тело убив, яко Каин убийца, любосластными стремленьми. 

Одожди Господь от Господа огнь иногда на беззаконие гневающее, сожег содомляны; ты же огнь вжегла еси геенский, в немже имаши, о душе, сожещися. 

Уязвихся, уранихся, се стрелы вражия, уязвившия мою душу и тело; се струпи, гноения, омрачения вопиют, раны самовольных моих страстей. 

\bfseries 

Преподобная Мати Марие, моли Бога о нас.

\normalfont{}

Простерла еси руце твои к щедрому Богу, Марие, в бездне зол погружаемая; и якоже Петру человеколюбно руку Божественную простре твое обращение всячески Иский. 

\itshape Слава\normalfont{}: Безначальная, Несозданная Троице, Нераздельная Единице, кающася мя приими, согрешивша спаси, Твое есмь создание, не презри, но пощади и избави мя огненнаго осуждения. 

\itshape И ныне\normalfont{}: Пречистая Владычице, Богородительнице, Надеждо к Тебе притекающих и пристанище сущих в бури, Милостиваго и Создателя и Сына Твоего умилостиви и мне молитвами Твоими.


\medskip\bfseries Песнь 3\normalfont{}\nopagebreak

\itshape Ирмос\normalfont{}: Утверди, Господи, на камени заповедей Твоих подвигшееся сердце мое, яко Един Свят еси и Господь. 

Агаре древле, душе, египтяныне уподобилася еси, поработившися произволением и рождши новаго Исмаила, презорство. 

Иаковлю лествицу разумела еси, душе моя, являемую от земли к Небесем: почто не имела еси восхода тверда, благочестия. 

Священника Божия и царя уединена, Христово подобие в мире жития, в человецех подражай. 

Обратися, постени, душе окаянная, прежде даже не приимет конец жития торжество, прежде даже дверь не заключит чертога Господь. 

Не буди столп сланый, душе, возвратившися вспять, образ да устрашит тя содомский, горе в Сигор спасайся. 

Моления, Владыко, Тебе поющих не отвержи, но ущедри, Человеколюбче, и подаждь верою просящим оставление. 

\itshape Слава\normalfont{}: Троица Простая, Несозданная, Безначальное Естество, в Троице певаемая Ипостасей, спаси ны, верою покланяющияся державе Твоей. 

\itshape И ныне\normalfont{}: От Отца безлетна Сына в лето, Богородительнице, неискусомужно родила еси, странное чудо, пребывши Дева доящи.


\medskip\bfseries Песнь 4\normalfont{}\nopagebreak

\itshape Ирмос\normalfont{}: Услыша пророк пришествие Твое, Господи, и убояся, яко хощеши от Девы родитися и человеком явитися, и глаголаше: услышах слух Твой и убояхся, слава силе Твоей, Господи. 

Время живота моего мало и исполнено болезней и лукавства, но в покаянии мя приими и в разум призови, да не буду стяжание ни брашно чуждему, Спасе, Сам мя ущедри. 

Царским достоинством, венцем и багряницею одеян, многоименный человек и праведный, богатством кипя и стады, внезапу богатства, славы царства, обнищав, лишися. 

Аще праведен бяше он и непорочен паче всех, и не убеже ловления льстиваго и сети; ты же, грехолюбива сущи, окаянная душе, что сотвориши, аще чесому от недоведомых случится наити тебе? 

Высокоглаголив ныне есмь, жесток же и сердцем, вотще и всуе, да не с фарисеем осудиши мя. Паче же мытарево смирение подаждь ми, Едине Щедре, Правосуде, и сему мя сочисли. 

Согреших, досадив сосуду плоти моея, вем, Щедре, но в покаянии мя приими и в разум призови, да не буду стяжание ни брашно чуждему, Спасе, Сам мя ущедри. 

Самоистукан бых страстьми, душу мою вредя, Щедре, но в покаянии мя приими и в разум призови, да не буду стяжание ни брашно чуждему, Спасе, Сам мя ущедри. 

Не послушах гласа Твоего, преслушах Писание Твое, Законоположника, но в покаянии мя приими и в разум призови, да не буду стяжание ни брашно чуждему, Спасе, Сам мя ущедри. 

\bfseries 

Преподобная Мати Марие, моли Бога о нас.

\normalfont{}

Великих безместий во глубину низведшися, неодержима была еси; но востекла еси помыслом лучшим к крайней деяньми яве добродетели преславно, ангельское естество, Марие, удививши. 

\itshape Слава\normalfont{}: Нераздельное Существом, Неслитное Лицы богословлю Тя, Троическое Едино Божество, яко Единоцарственное и Сопрестольное, вопию Ти песнь великую, в вышних трегубо песнословимую. 

\itshape И ныне\normalfont{}: И раждаеши, и девствуеши, и пребываеши обоюду естеством Дева, Рождейся обновляет законы естества, утроба же раждает нераждающая. Бог идеже хощет, побеждается естества чин: творит бо, елика хощет.


\medskip\bfseries Песнь 5\normalfont{}\nopagebreak

\itshape Ирмос\normalfont{}: От нощи утренююща, Человеколюбче, просвети, молюся, и настави и мене на повеления Твоя, и научи мя, Спасе, творити волю Твою. 

Низу сничащую подражай, о душе, прииди, припади к ногама Иисусовыма, да тя исправит, и да ходиши право стези Господни. 

Аще и кладязь еси глубокий, Владыко, источи ми воду из пречистых Твоих жил, да, яко самаряныня, не ктому, пияй, жажду: жизни бо струи источаеши. 

Силоам да будут ми слезы моя, Владыко Господи, да умыю и аз зеницы сердца, и вижду Тя, умна Света превечна. 

\bfseries 

Преподобная Мати Марие, моли Бога о нас.

\normalfont{}

Несравненным желанием, всебогатая, древу возжелевши поклонитися животному, сподобилася еси желания, сподоби убо и мене улучити вышния славы. 

\itshape Слава\normalfont{}: Тя, Троице, славим Единаго Бога: Свят, Свят, Свят еси, Отче, Сыне и Душе, Простое Существо, Единице присно покланяемая. 

\itshape И ныне\normalfont{}: Из Тебе облечеся в мое смешение, нетленная, безмужная Мати Дево, Бог, создавый веки, и соедини Себе человеческое естество.


\medskip\bfseries Песнь 6\normalfont{}\nopagebreak

\itshape Ирмос\normalfont{}: Возопих всем сердцем моим к щедрому Богу, и услыша мя от ада преисподняго, и возведе от тли живот мой. 

Аз есмь, Спасе, юже погубил еси древле царскую драхму; но вжег светильник, Предтечу Твоего, Слове, взыщи и обрящи Твой образ. 

Востани и побори, яко Иисус Амалика, плотския страсти, и гаваониты, лестныя помыслы, присно побеждающи. 

\bfseries 

Преподобная Мати Марие, моли Бога о нас.

\normalfont{}

Да страстей пламень угасиши, слез капли источала еси присно, Марие, душею распаляема, ихже благодать подаждь и мне, твоему рабу. 

\bfseries 

Преподобная Мати Марие, моли Бога о нас.

\normalfont{}

Безстрастие Небесное стяжала еси крайним на земли житием, мати. Темже тебе поющим страстей избавитися молитвами твоими молися. 

\itshape Слава\normalfont{}: Троица есмь Проста, Нераздельна, раздельна Личне и Единица есмь естеством соединена, Отец глаголет, и Сын, и Божественный Дух. 

\itshape И ныне\normalfont{}: Утроба Твоя Бога нам роди, воображена по нам: Егоже, яко Создателя всех, моли, Богородице, да Твоими молитвами оправдимся. 

Господи, помилуй. (Трижды.) 

\itshape Слава, и ныне\normalfont{}:


\medskip\bfseries Кондак, глас 6:\normalfont{}\nopagebreak

Душе моя, душе моя, востани, что спиши? конец приближается, и имаши смутитися: воспряни убо, да пощадит тя Христос Бог, везде сый и вся исполняяй.


\medskip\bfseries Песнь 7\normalfont{}\nopagebreak

\itshape Ирмос\normalfont{}: Согрешихом, беззаконновахом, неправдовахом пред Тобою, ниже соблюдохом, ниже сотворихом, якоже заповедал еси нам; но не предаждь нас до конца, отцев Боже. 

Исчезоша дние мои, яко соние востающаго; темже, яко Езекия, слезю на ложи моем, приложитися мне летом живота. Но кий Исаия предстанет тебе, душе, аще не всех Бог? 

Припадаю Ти и приношу Тебе, якоже слезы, глаголы моя: согреших, яко не согреши блудница, и беззаконновах, яко иный никтоже на земли. Но ущедри, Владыко, творение Твое и воззови мя. 

Погребох образ Твой и растлих заповедь Твою, вся помрачися доброта, и страстьми угасися, Спасе, свеща. Но ущедрив, воздаждь ми, якоже поет Давид, радование. 

Обратися, покайся, открый сокровенная, глаголи Богу, вся ведущему: Ты веси моя тайная, Едине Спасе. Но Сам мя помилуй, якоже поет Давид, по милости Твоей. 

\bfseries 

Преподобная Мати Марие, моли Бога о нас.

\normalfont{}

Возопивши к Пречистей Богоматери, первее отринула еси неистовство страстей, нужно стужающих, и посрамила еси врага запеншаго. Но даждь ныне помощь от скорби и мне, рабу твоему. 

\bfseries 

Преподобная Мати Марие, моли Бога о нас.

\normalfont{}

Егоже возлюбила еси, Егоже возжелела еси, Егоже ради плоть изнурила еси, преподобная, моли ныне Христа о рабех: яко да милостив быв всем нам, мирное состояние дарует почитающим Его. 

\itshape Слава\normalfont{}: Троице Простая, Нераздельная, Единосущная и Естество Едино, Светове и Свет, и Свята Три, и Едино Свято поется Бог Троица; но воспой, прослави Живот и Животы, душе, всех Бога. 

\itshape И ныне\normalfont{}: Поем Тя, благословим Тя, покланяемся Ти, Богородительнице, яко Неразлучныя Троицы породила еси Единаго Христа Бога и Сама отверзла еси нам, сущим на земли, Небесная.


\medskip\bfseries Песнь 8\normalfont{}\nopagebreak

\itshape Ирмос\normalfont{}: Егоже воинства Небесная славят, и трепещут херувими и серафими, всяко дыхание и тварь, пойте, благословите и превозносите во вся веки. 

Слезную, Спасе, сткляницу яко миро истощавая на главу, зову Ти, якоже блудница, милости ищущая, мольбу приношу и оставление прошу прияти. 

Аще и никтоже, якоже аз, согреши Тебе, но обаче приими и мене, благоутробне Спасе, страхом кающася и любовию зовуща: согреших Тебе Единому, помилуй мя, Милостиве. 

Пощади, Спасе, Твое создание и взыщи, яко Пастырь, погибшее, предвари заблуждшаго, восхити от волка, сотвори мя овча на пастве Твоих овец. 

Егда, Судие, сядеши, яко благоутробен, и покажеши страшную славу Твою, Спасе, о каковый страх тогда, пещи горящей, всем боящимся нестерпимаго судища Твоего. 

\bfseries 

Преподобная Мати Марие, моли Бога о нас.

\normalfont{}

Света незаходимаго Мати тя просветивши, от омрачения страстей разреши. Темже вшедши в духовную благодать, просвети, Марие, тя верно восхваляющия. 

\bfseries 

Преподобная Мати Марие, моли Бога о нас.

\normalfont{}

Чудо ново видев, ужасашеся божественный в тебе воистинну, мати, Зосима: ангела бо зряше во плоти и ужасом весь исполняшеся, Христа поя во веки. 

\itshape Слава\normalfont{}: Безначальне Отче, Сыне Собезначальне, Утешителю Благий, Душе Правый, Слова Божия Родителю, Отца Безначальна Слове, Душе Живый и Зиждяй, Троице Единице, помилуй мя. 

\itshape И ныне\normalfont{}: Яко от оброщения червленицы, Пречистая, умная багряница Еммануилева внутрь во чреве Твоем плоть исткася. Темже Богородицу воистинну Тя почитаем.


\medskip\bfseries Песнь 9\normalfont{}\nopagebreak

\itshape Ирмос\normalfont{}: Безсеменнаго зачатия Рождество несказанное, Матере безмужныя нетленен Плод, Божие бо Рождение обновляет естества. Темже Тя вси роди, яко Богоневестную Матерь, православно величаем. 

Умилосердися, спаси мя, Сыне Давидов, помилуй, беснующияся словом исцеливый, глас же благоутробный, яко разбойнику, мне рцы: аминь, глаголю тебе, со Мною будеши в раи, егда прииду во славе Моей. 

Разбойник оглаголоваше Тя, разбойник богословяше Тя: оба бо на кресте свисяста. Но, о Благоутробне, яко верному разбойнику Твоему, познавшему Тя Бога, и мне отверзи дверь славнаго Царствия Твоего. 

Тварь содрогашеся, распинаема Тя видящи, горы и камения страхом распадахуся, и земля сотрясашеся, и ад обнажашеся, и соомрачашеся свет во дни, зря Тебе, Иисусе, пригвождена ко Кресту. 

Достойных покаяния плодов не истяжи от мене, ибо крепость моя во мне оскуде; сердце мне даруй присно сокрушенное, нищету же духовную: да сия Тебе принесу яко приятную жертву, Едине Спасе. 

Судие мой и Ведче мой, хотяй паки приити со ангелы, судити миру всему, милостивным Твоим оком тогда видев мя, пощади и ущедри мя, Иисусе, паче всякаго естества человеча согрешивша. 

\bfseries 

Преподобная Мати Марие, моли Бога о нас.

\normalfont{}

Удивила еси всех странным житием твоим, ангелов чины и человеков соборы, невещественно поживши и естество прешедши: имже, яко невещественныма ногама вшедши, Марие, Иордан прешла еси. 

\bfseries 

Преподобная Мати Марие, моли Бога о нас.

\normalfont{}

Умилостиви Создателя о хвалящих тя, преподобная мати, избавитися озлоблений и скорбей, окрест нападающих: да избавившеся от напастей, возвеличим непрестанно прославльшаго тя Господа. 

\bfseries 

Преподобне отче Андрее, моли Бога о нас.

\normalfont{}

Андрее честный и отче треблаженнейший, пастырю Критский, не престай моляся о воспевающих тя: да избавимся вси гнева, и скорби, и тления, и прегрешений безмерных, чтущии твою память верно. 

\itshape Слава\normalfont{}: Отца прославим, Сына превознесем, Божественному Духу верно поклонимся, Троице Нераздельней, Единице по Существу, яко Свету и Светом, и Животу и Животом, Животворящему и Просвещающему концы. 

\itshape И ныне\normalfont{}: Град Твой сохраняй, Богородительнице Пречистая, в Тебе бо сей верно царствуяй, в Тебе и утверждается, и Тобою побеждаяй, побеждает всякое искушение, и пленяет ратники, и проходит послушание. 

Таже оба лика вкупе поют \itshape Ирмос\normalfont{}: 

Безсеменнаго зачатия Рождество несказанное, Матере безмужныя нетленен Плод, Божие бо Рождение обновляет естества. Темже Тя вси роди, яко Богоневестную Матерь, православно величаем.



\mychapterending

\mychapter{В четверг пятой седмицы Великого Поста}
%http://www.molitvoslov.com/text575.htm 
 
\medskip\bfseries Песнь 1\normalfont{}\nopagebreak


\itshape Ирмос\normalfont{}: Помощник и Покровитель бысть мне во спасение, Сей мой Бог, и прославлю Его, Бог Отца моего, и вознесу Его: славно бо прославися.


Откуду начну плакати окаяннаго моего жития деяний? кое ли положу начало, Христе, нынешнему рыданию? но яко благоутробен, даждь ми прегрешений оставление.


Гряди, окаянная душе, с плотию твоею, Зиждителю всех исповеждься, и останися прочее преждняго безсловесия, и принеси Богу в покаянии слезы.


Первозданнаго Адама преступлению поревновав, познах себе обнажена от Бога и присносущнаго Царствия и сладости, грех ради моих.


Увы мне, окаянная душе, что уподобилася еси первей Еве? видела бо еси зле, и уязвилася еси горце, и коснулася еси древа, и вкусила еси дерзостно безсловесныя снеди.


Вместо Евы чувственныя мысленная ми бысть Ева, во плоти страстный помысл, показуяй сладкая и вкушаяй присно горькаго напоения.


Достойно из Едема изгнан бысть, яко не сохранив едину Твою, Спасе, заповедь Адам: аз же что постражду, отметая всегда животная Твоя словеса?


Каиново прешед убийство, произволением бых убийца совести душевней, оживив плоть и воевав на ню лукавыми моими деяньми.


Авелеве, Иисусе, не уподобихся правде, дара Тебе приятна не принесох когда, ни деяния божественна, ни жертвы чистыя, ни жития непорочнаго.


Яко Каин и мы, душе окаянная, всех Содетелю деяния скверная, и жертву порочную, и непотребное житие принесохом вкупе: темже и осудихомся.


Брение Здатель живосоздав, вложил еси мне плоть, и кости, и дыхание, и жизнь; но, о Творче мой, Избавителю мой и Судие, кающася приими мя.


Извещаю Ти, Спасе, грехи, яже содеях, и души и тела моего язвы, яже внутрь убийственнии помыслы разбойнически на мя возложиша.


Аще и согреших, Спасе, но вем, яко Человеколюбец еси, наказуеши милостивно и милосердствуеши тепле: слезяща зриши и притекаеши, яко отец, призывая блуднаго.


Повержена мя, Спасе, пред враты Твоими, поне на старость не отрини мене во ад тща, но прежде конца, яко Человеколюбец, даждь ми прегрешений оставление.


В разбойники впадый аз есмь помышленьми моими, весь от них уязвихся ныне и исполнихся ран, но, Сам ми представ, Христе Спасе, исцели.


Священник, мя предвидев, мимо иде, и левит, видя в лютых нага, презре, но, из Марии возсиявый, Иисусе, Ты, представ, ущедри мя.


Агнче Божий, вземляй грехи всех, возьми бремя от мене тяжкое греховное и, яко благоутробен, даждь ми слезы умиления.


Покаяния время, прихожду Ти, Создателю моему: возьми бремя от мене тяжкое греховное и, яко благоутробен, даждь ми слезы умиления.


Не возгнушайся мене, Спасе, не отрини от Твоего лица, возьми бремя от мене тяжкое греховное и, яко благоутробен, даждь мне грехопадений оставление.


Вольная, Спасе, и невольная прегрешения моя, явленная и сокровенная и ведомая и неведомая, вся простив, яко Бог, очисти и спаси мя.


От юности, Христе, заповеди Твоя преступих, всестрастно небрегий, унынием преидох житие. Темже зову Ти, Спасе: поне на конец спаси мя.


Богатство мое, Спасе, изнурив в блуде, пуст есмь плодов благочестивых, алчен же зову: Отче щедрот, предварив, Ты мя ущедри.


Тебе припадаю, Иисусе, согреших Ти, очисти мя, возьми бремя от мене тяжкое греховное и, яко благоутробен, даждь ми слезы умиления.


Не вниди со мною в суд, нося моя деяния, словеса изыскуя и исправляя стремления. Но в щедротах Твоих презирая моя лютая, спаси мя, Всесильне.


\medskip\bfseries Иный канон преподобныя матере нашея Марии Египетския, глас 6:\normalfont{}\nopagebreak


Преподобная мати Марие, моли Бога о нас.


Ты ми даждь светозарную благодать от Божественнаго свыше промышления избежати страстей омрачения и пети усердно твоего, Марие, жития красная исправления.


\medskip\bfseries Преподобная мати Марие, моли Бога о нас.\normalfont{}\nopagebreak


Приклоньшися Христовым Божественным законом, к сему приступила еси, сладостей неудержимая стремления оставивши, и всякую добродетель всеблагоговейно, яко едину, исправила еси.


Преподобне отче Андрее, моли Бога о нас.


Молитвами твоими нас, Андрее, избави страстей безчестных и Царствия ныне Христова общники верою и любовию воспевающия тя, славне, покажи, молимся.


\itshape Слава\normalfont{}: Пресущественная Троице, во Единице покланяемая, возьми бремя от мене тяжкое греховное и, яко благоутробна, даждь ми слезы умиления.


\itshape И ныне\normalfont{}: Богородице, Надежде и Предстательство Тебе поющих, возьми бремя от мене тяжкое греховное и, яко Владычица Чистая, кающася приими мя. 


\medskip\bfseries Песнь 2\normalfont{}\nopagebreak


\itshape Ирмос\normalfont{}: Вонми, Небо, и возглаголю, и воспою Христа, от Девы плотию пришедшаго.


Вонми, Небо, и возглаголю, земле, внушай глас, кающийся к Богу и воспевающий Его.


Вонми ми, Боже, Спасе мой, милостивным Твоим оком и приими мое теплое исповедание.


Согреших паче всех человек, един согреших Тебе; но ущедри, яко Бог, Спасе, творение Твое.


Буря мя злых обдержит, благоутробне Господи; но яко Петру и мне руку простри.


Слезы блудницы, Щедре, и аз предлагаю, очисти мя, Спасе, благоутробием Твоим.


Омрачих душевную красоту страстей сластьми и всячески весь ум персть сотворих.


Раздрах ныне одежду мою первую, юже ми изтка Зиждитель из начала, и оттуду лежу наг.


Облекохся в раздранную ризу, юже изтка ми змий советом, и стыждуся.


Воззрех на садовную красоту и прельстихся умом: и оттуду лежу наг и срамляюся.


Делаша на хребте моем вси начальницы страстей, продолжающе на мя беззаконие их.


Погубих первозданную доброту и благолепие мое и ныне лежу наг и стыждуся.


Сшиваше кожныя ризы грех мне, обнаживый мя первыя боготканныя одежды.


Обложен есмь одеянием студа, якоже листвием смоковным, во обличение моих самовластных страстей.


Одеяхся в срамную ризу и окровавленную студно течением страстнаго и любосластнаго живота.


Оскверних плоти моея ризу и окалях еже по образу, Спасе, и по подобию.


Впадох в страстную пагубу и в вещественную тлю, и оттоле до ныне враг мне досаждает.


Любовещное и любоименное житие, невоздержанием, Спасе, предпочет ныне, тяжким бременем обложен есмь.


Украсих плотский образ скверных помышлений различным обложением и осуждаюся.


Внешним прилежно благоукрашением единем попекохся, внутреннюю презрев Богообразную скинию.


Вообразив моих страстей безобразие, любосластными стремленьми погубих ума красоту.


Погребох перваго образа доброту, Спасе, страстьми, юже, яко иногда драхму, взыскав, обрящи.


Согреших, якоже блудница, вопию Ти: един согреших Тебе; яко миро, приими, Спасе, и моя слезы.


Поползохся, яко Давид, блудно и осквернихся, но омый и мене, Спасе, слезами.


Очисти, якоже мытарь, вопию Ти, Спасе, очисти мя: никтоже бо сущих из Адама, якоже аз, согреших Тебе.


Ни слез, ниже покаяния имам, ниже умиления. Сам ми сия, Спасе, яко Бог, даруй.


Дверь Твою не затвори мне тогда, Господи, Господи, но отверзи ми сию, кающемуся Тебе.


Человеколюбче, хотяй всем спастися, Ты воззови мя и приими, яко благ, кающагося.


Внуши воздыхания души моея и очию моею приими капли, Спасе, и спаси мя.


\bfseries Пресвятая Богородице, спаси нас.


\normalfont{}


Пречистая Богородице Дево, Едина Всепетая, моли прилежно, во еже спастися нам.


\itshape Иный. Ирмос\normalfont{}: Видите, видите, яко Аз есмь Бог, манну одождивый и воду из камене источивый древле в пустыни людем Моим, десницею единою и крепостию Моею.


Видите, видите, яко Аз есмь Бог, внушай, душе моя, Господа вопиюща, и удалися прежняго греха, и бойся, яко неумытнаго и яко Судии и Бога.


Кому уподобилася еси, многогрешная душе? токмо первому Каину и Ламеху оному, каменовавшая тело злодействы и убившая ум безсловесными стремленьми.


Вся прежде закона претекши, о душе, Сифу не уподобилася еси, ни Еноса подражала еси, ни Еноха преложением, ни Ноя, но явилася еси убога праведных жизни.


Едина отверзла еси хляби гнева Бога Твоего, душе моя, и потопила еси всю, якоже землю, плоть, и деяния, и житие, и пребыла еси вне спасительнаго ковчега.


Мужа убих, глаголет, в язву мне и юношу в струп, Ламех рыдая вопияше; ты же не трепещеши, о душе моя, окалявши плоть и ум осквернивши.


О како поревновах Ламеху, первому убийце, душу, яко мужа, ум, яко юношу, яко брата же моего, тело убив, яко Каин убийца, любосластными стремленьми.


Столп умудрила еси создати, о душе, и утверждение водрузити твоими похотьми, аще не бы Зиждитель удержал советы твоя и низвергл на землю ухищрения твоя.


Уязвихся, уранихся, се стрелы вражия, уязвившия мою душу и тело; се струпи, гноения, омрачения вопиют, раны самовольных моих страстей.


Одожди Господь от Господа огнь иногда на беззаконие гневающее, сожег содомляны; ты же огнь вжегла еси геенский, в немже имаши, о душе, сожещися.


Разумейте и видите, яко Аз есмь Бог, испытаяй сердца и умучаяй мысли, обличаяй деяния, и попаляяй грехи, и судяй сиру, и смирену, и нищу.


\bfseries Преподобная мати Марие, моли Бога о нас.


\normalfont{}Простерла еси руце твои к щедрому Богу, Марие, в бездне зол погружаемая, и якоже Петру человеколюбно руку Божественную простре твое обращение всячески Иский.


\medskip\bfseries Преподобная мати Марие, моли Бога о нас.\normalfont{}\nopagebreak


Всем усердием и любовию притекла еси Христу, первый греха путь отвращши, и в пустынях непроходимых питающися, и Того чисте совершающи Божественныя заповеди.


\medskip\bfseries Преподобне отче Андрее, моли Бога о нас.\normalfont{}\nopagebreak


Видим, видим человеколюбие, о душе, Бога и Владыки; сего ради прежде конца тому со слезами припадем вопиюще: Андрея молитвами, Спасе, помилуй нас.


\itshape Слава\normalfont{}: Безначальная, Несозданная Троице, Нераздельная Единице, кающася мя приими, согрешивша спаси, Твое есмь создание, не презри, но пощади и избави мя огненнаго осуждения.


\itshape И ныне\normalfont{}: Пречистая Владычице, Богородительнице, Надеждо к Тебе притекающих и пристанище сущих в бури, Милостиваго и Создателя и Сына Твоего умилостиви и мне молитвами Твоими. 

\medskip\bfseries Песнь 3\normalfont{}\nopagebreak


\itshape Ирмос\normalfont{}: На недвижимом, Христе, камени заповедей Твоих утверди мое помышление.


Огнь от Господа иногда Господь одождив, землю содомскую прежде попали.


На горе спасайся, душе, якоже Лот оный, и в Сигор угонзай.


Бегай запаления, о душе, бегай содомскаго горения, бегай тления Божественнаго пламене.


Исповедаюся Тебе, Спасе, согреших, согреших Ти, но ослаби, остави ми, яко благоутробен.


Согреших Тебе един аз, согреших паче всех, Христе Спасе, да не презриши мене.


Ты еси Пастырь добрый, взыщи мене, агнца, и заблуждшаго да не презриши мене.

Ты еси сладкий Иисусе, Ты еси Создателю мой, в Тебе, Спасе, оправдаюся.


\medskip\bfseries Пресвятая Троице, Боже наш, слава Тебе.\normalfont{}\nopagebreak


О Троице Единице Боже, спаси нас от прелести, и искушений, и обстояний.

\medskip\bfseries Пресвятая Богородице, спаси нас.\normalfont{}\nopagebreak


Радуйся, Богоприятная утробо, радуйся, престоле Господень, радуйся, Мати Жизни нашея.


\itshape Иный. Ирмос\normalfont{}: Утверди, Господи, на камени заповедей Твоих подвигшееся сердце мое, яко Един Свят еси и Господь.


Источник живота стяжах Тебе, смерти Низложителя, и вопию Ти от сердца моего прежде конца: согреших, очисти, спаси мя.


При Нои, Спасе, блудствовавшия подражах, онех наследствовав осуждение в потопе погружения.


Согреших, Господи, согреших Тебе, очисти мя: несть бо иже кто согреши в человецех, егоже не превзыдох прегрешеньми.


Хама онаго, душе, отцеубийца подражавши, срама не покрыла еси искренняго, вспять зря возвратившися.


Благословения Симова не наследовала еси, душе окаянная, ни пространное одержание, якоже Иафеф, имела еси на земли оставления.


От земли Харран изыди от греха, душе моя, гряди в землю, точащую присноживотное нетление, еже Авраам наследствова.


Авраама слышала еси, душе моя, древле оставльша землю отечества и бывша пришельца, сего произволению подражай.


У дуба Мамврийскаго учредив патриарх ангелы, наследствова по старости обетования ловитву.


Исаака, окаянная душе моя, разумевши новую жертву, тайно всесожженную Господеви, подражай его произволению.


Исмаила слышала еси, трезвися, душе моя, изгнана, яко рабынино отрождение, виждь, да не како подобно что постраждеши, ласкосердствующи.


Агаре древле, душе, египтяныне уподобилася еси, поработившися произволением и рождши новаго Исмаила, презорство.


Иаковлю лествицу разумела еси, душе моя, являемую от земли к Небесем: почто не имела еси восхода тверда, благочестия.


Священника Божия и царя уединена, Христово подобие в мире жития, в человецех подражай.


Не буди столп сланый, душе, возвратившися вспять, образ да устрашит тя содомский, горе в Сигор спасайся.


Запаления, якоже Лот, бегай, душе моя, греха, бегай Содомы и Гоморры, бегай пламене всякаго безсловеснаго желания.


Помилуй, Господи, помилуй мя, вопию Ти, егда приидеши со ангелы Твоими воздати всем по достоянию деяний.


Моления, Владыко, Тебе поющих не отвержи, но ущедри, Человеколюбче, и подаждь верою просящим оставление.


\medskip\bfseries Преподобная мати Марие, моли Бога о нас.\normalfont{}\nopagebreak


Содержим есмь бурею и треволнением согрешений, но сама мя, мати, ныне спаси и к пристанищу Божественнаго покаяния возведи.


\medskip\bfseries Преподобная мати Марие, моли Бога о нас.\normalfont{}\nopagebreak


Рабское моление и ныне, преподобная, принесши ко благоутробней молитвами твоими Богородице, отверзи ми Божественныя входы.


\medskip\bfseries Преподобне отче Андрее, моли Бога о нас.\normalfont{}\nopagebreak


Твоими молитвами даруй и мне оставление долгов, о Андрее, Критский председателю, покаяния бо ты таинник преизрядный.


\itshape Слава\normalfont{}: Троице Простая, Несозданная, Безначальное Естество, в Троице певаемая Ипостасей, спаси ны, верою покланяющияся державе Твоей.


\itshape И ныне\normalfont{}: От Отца безлетна Сына в лето, Богородительнице, неискусомужно родила еси, странное чудо, пребывши Дева доящи.


\itshape Ирмос\normalfont{}: Утверди, Господи, на камени заповедей Твоих подвигшееся сердце мое, яко Един Свят еси и Господь. 

\medskip\bfseries Седален, глас 8:\normalfont{}\nopagebreak


Светила богозрачная, Спасовы апостоли, просветите нас во тьме жития, яко да во дни ныне благообразно ходим, светом воздержания нощных страстей отбегающе, и светлыя страсти Христовы узрим, радующеся. 

\medskip\bfseries Слава, другий седален, глас тойже:\normalfont{}\nopagebreak


Апостольская двоенадесятице Богоизбранная, мольбу Христу ныне принеси, постное поприще всем прейти, совершающим во умилении молитвы, творящим усердно добродетели, яко да сице предварим видети Христа Бога славное Воскресение, славу и хвалу приносяще. 


\medskip\bfseries И ныне, Богородичен:\normalfont{}\nopagebreak



Непостижимаго Бога, Сына и Слово, несказанно паче ума из Тебе рождшееся, моли, Богородице, со апостолы, мир вселенней чистый подати, и согрешений дати нам прежде конца прощение, и Царствия Небеснаго крайния ради благости сподобити рабы Твоя. 


\medskip\bfseries Таже трипеснец, глас 8:\normalfont{}\nopagebreak


\medskip\bfseries Песнь 4\normalfont{}\nopagebreak


\itshape Ирмос\normalfont{}: Услышах, Господи, смотрения Твоего таинство, разумех дела Твоя и прославих Твое Божество.


\medskip\bfseries Святии апостоли, молите Бога о нас.\normalfont{}\nopagebreak


Воздержанием поживше, просвещеннии Христовы апостоли, воздержания время нам ходатайствы Божественными утишают.


\medskip\bfseries Святии апостоли, молите Бога о нас.\normalfont{}\nopagebreak


Двоенадесятострунный орган песнь воспе спасительную, учеников лик Божественный, лукавая возмущая гласования.


\medskip\bfseries Святии апостоли, молите Бога о нас.\normalfont{}\nopagebreak


Одождением духовным всю подсолнечную напоисте, сушу отгнавше многобожия, всеблаженнии.


\medskip\bfseries Пресвятая Богородице, спаси нас.\normalfont{}\nopagebreak


Смирившася спаси мя, высокомудренно пожившаго, рождшая Вознесшаго смиренное естество, Дево Всечистая.


Иный. \itshape Ирмос\normalfont{}, глас тойже: Услышах, Господи, смотрения Твоего таинство, разумех дела Твоя и прославих Твое Божество.


\medskip\bfseries Святии апостоли, молите Бога о нас.\normalfont{}\nopagebreak


Апостольское всечестное ликостояние, Зиждителя всех молящее, проси помиловати ны, восхваляющия тя.


\medskip\bfseries Святии апостоли, молите Бога о нас.\normalfont{}\nopagebreak


Яко делателе суще, Христовы апостоли, во всем мире Божественным словом возделавшии, приносите плоды Ему всегда.


\medskip\bfseries Святии апостоли, молите Бога о нас.\normalfont{}\nopagebreak


Виноград бысте Христов воистинну возлюбленный, вино бо духовное источисте миру, апостоли.


\medskip\bfseries Пресвятая Троице, Боже наш, слава Тебе.\normalfont{}\nopagebreak


Преначальная, Сообразная, Всесильнейшая Троице Святая, Отче, Слове и Душе Святый, Боже, Свете и Животе, сохрани стадо Твое.


\medskip\bfseries Пресвятая Богородице, спаси нас.\normalfont{}\nopagebreak


Радуйся, престоле огнезрачный, радуйся, светильниче свещеносный, радуйся, горо освящения, ковчеже Жизни, святых святая сене.


\itshape Великаго канона Ирмос\normalfont{}: Услыша пророк пришествие Твое, Господи, и убояся, яко хощеши от Девы родитися и человеком явитися, и глаголаше: услышах слух Твой и убояхся, слава силе Твоей, Господи.


Дел Твоих да не презриши, создания Твоего да не оставиши, Правосуде. Аще и един согреших, яко человек, паче всякаго человека, Человеколюбче; но имаши, яко Господь всех, власть оставляти грехи.


Приближается, душе, конец, приближается, и нерадиши, ни готовишися, время сокращается: востани, близ при дверех Судия есть. Яко соние, яко цвет, время жития течет: что всуе мятемся?


Воспряни, о душе моя, деяния твоя, яже соделала еси, помышляй, и сия пред лице твое принеси, и капли испусти слез твоих; рцы со дерзновением деяния и помышления Христу и оправдайся.


Не бысть в житии греха, ни деяния, ни злобы, еяже аз, Спасе, не согреших, умом, и словом, и произволением, и предложением, и мыслию, и деянием согрешив, яко ин никтоже когда.


Отсюду осужден бых, отсюду и препрен бых аз, окаянный, от своея совести, еяже ничтоже в мире нужнейше: Судие, Избавителю мой и Ведче, пощади, и избави, и спаси мя, раба Твоего.


Лествица, юже виде древле великий в патриарсех, указание есть, душе моя, деятельнаго восхождения, разумнаго возшествия: аще хощеши убо деянием, и разумом, и зрением пожити, обновися.


Зной дневный претерпе лишения ради патриарх и мраз нощный понесе, на всяк день снабдения творя, пасый, труждаяйся, работаяй, да две жене сочетает.


Жены ми две разумей, деяние же и разум в зрении, Лию убо деяние, яко многочадную, Рахиль же разум, яко многотрудную; ибо кроме трудов ни деяние, ни зрение, душе, исправится.


Бди, о душе моя, изрядствуй, якоже древле великий в патриарсех, да стяжеши деяние с разумом, да будеши ум, зряй Бога, и достигнеши незаходящий мрак в видении, и будеши великий купец.


Дванадесяте патриархов великий в патриарсех детотворив, тайно утверди тебе лествицу деятельнаго, душе моя, восхождения: дети, яко основания, степени, яко восхождения, премудренно подложив.


Исава возненавиденнаго подражала еси, душе, отдала еси прелестнику твоему первыя доброты первенство, и отеческия молитвы отпала еси, и дважды поползнулася еси, окаянная, деянием и разумом: темже ныне покайся.


Едом Исав наречеся, крайняго ради женонеистовнаго смешения: невоздержанием бо присно разжигаем и сластьми оскверняем, Едом именовася, еже глаголется разжжение души любогреховныя.


Иова на гноищи слышавши, о душе моя, оправдавшагося, того мужеству не поревновала еси, твердаго не имела еси предложения во всех, яже веси, и имиже искусилася еси, но явилася еси нетерпелива.


Иже первее на престоле, наг ныне на гноище гноен, многий в чадех и славный, безчаден и бездомок напрасно: палату убо гноище и бисерие струпы вменяше.


Царским достоинством, венцем и багряницею одеян, многоименный человек и праведный, богатством кипя и стады, внезапу богатства, славы царства, обнищав, лишися.


Аще праведен бяше он и непорочен паче всех, и не убеже ловления льстиваго и сети; ты же грехолюбива сущи, окаянная душе, что сотвориши, аще чесому от недоведомых случится наити тебе?


Тело осквернися, дух окаляся, весь острупихся, но яко врач, Христе, обоя покаянием моим уврачуй, омый, очисти, покажи, Спасе мой, паче снега чистейша.


Тело Твое и кровь, распинаемый о всех, положил еси, Слове: тело убо, да мя обновиши, кровь, да омыеши мя. Дух же предал еси, да мя приведеши, Христе, Твоему Родителю.


Соделал еси спасение посреде земли, Щедре, да спасемся. Волею на древе распялся еси, Едем затворенный отверзеся, горняя и дольняя тварь, языцы вси спасени покланяются Тебе.


Да будет ми купель кровь из ребр Твоих, вкупе и питие, источившее воду оставления, да обоюду очищаюся, помазуяся и пия, яко помазание и питие, Слове, животочная Твоя словеса.


Наг есмь чертога, наг есмь и брака, купно и вечери, светильник угасе, яко безъелейный, чертог заключися мне спящу, вечеря снедеся, аз же по руку и ногу связан, вон низвержен есмь.


Чашу Церковь стяжа, ребра Твоя живоносная, из нихже сугубыя нам источи токи оставления и разума, во образ древняго и новаго, двоих вкупе заветов, Спасе наш.


Время живота моего мало и исполнено болезней и лукавства, но в покаянии мя приими и в разум призови, да не буду стяжание, ни брашно чуждему, Спасе, Сам мя ущедри.


Высокоглаголив ныне есмь, жесток же и сердцем, вотще и всуе, да не с фарисеем осудиши мя. Паче же мытарево смирение подаждь ми, Едине Щедре, Правосуде, и сему мя сочисли.


Согреших, досадив сосуду плоти моея, вем, Щедре, но в покаянии мя приими и в разум призови, да не буду стяжание ни брашно чуждему, Спасе, Сам мя ущедри.


Самоистукан бых страстьми, душу мою вредя, Щедре, но в покаянии мя приими и в разум призови, да не буду стяжание ни брашно чуждему, Спасе, Сам мя ущедри.


Не послушах гласа Твоего, преслушах Писание Твое, Законоположника, но в покаянии мя приими и в разум призови, да не буду стяжание ни брашно чуждему, Спасе, Сам мя ущедри.


\bfseries Пресвятая Троице, Боже наш, слава Тебе.





Преподобная мати Марие, моли Бога о нас.\normalfont{}


Безплотных жительство в плоти преходящи, благодать, преподобная, к Богу велию воистинну прияла еси, верно о чтущих тя предстательствуй. Темже молим тя, от всяких напастей и нас молитвами твоими избави.


\medskip\bfseries Преподобная мати Марие, моли Бога о нас.\normalfont{}\nopagebreak


Великих безместий во глубину низведшися, неодержима была еси, но востекла еси помыслом лучшим к крайней деяньми яве добродетели преславно, ангельское естество, Марие, удививши.


\medskip\bfseries Преподобне отче Андрее, моли Бога о нас.\normalfont{}\nopagebreak


Андрее, отеческая похвало, молитвами твоими не престай, моляся, предстоя Троице Пребожественней, яко да избавимся мучения, любовию предстателя тя Божественнаго, всеблаженне, призывающии, Криту удобрение.


\itshape Слава\normalfont{}: Нераздельное Существом, Неслитное Лицы богословлю Тя, Троическое Едино Божество, яко Единоцарственное и Сопрестольное, вопию Ти песнь великую, в вышних трегубо песнословимую.


\itshape И ныне\normalfont{}: И раждаеши, и девствуеши, и пребываеши обоюду естеством Дева, Рождейся обновляет законы естества, утроба же раждает нераждающая. Бог идеже хощет, побеждается естества чин: творит бо, елика хощет. 


\medskip\bfseries Песнь 5\normalfont{}\nopagebreak


\itshape Ирмос\normalfont{}: От нощи утренююща, Человеколюбче, просвети, молюся, и настави и мене на повеления Твоя, и научи мя, Спасе, творити волю Твою.


В нощи житие мое преидох присно, тьма бо бысть, и глубока мне мгла, нощь греха, но яко дне сына, Спасе, покажи мя.


Рувима подражая, окаянный аз содеях беззаконный и законопреступный совет на Бога Вышняго, осквернив ложе мое, яко отчее он.


Исповедаюся Тебе, Христе Царю: согреших, согреших, яко прежде Иосифа братия продавшии, чистоты плод и целомудрия.


От сродников праведная душа связася, продася в работу сладкий во образ Господень; ты же вся, душе, продалася еси злыми твоими.


Иосифа праведнаго и целомудреннаго ума подражай, окаянная и неискусная душе, и не оскверняйся безсловесными стремленьми, присно беззаконнующи.


Аще и в рове поживе иногда Иосиф, Владыко Господи, но во образ погребения и востания Твоего: аз же что Тебе когда сицевое принесу?


Моисеов слышала еси ковчежец, душе, водами, волнами носим речными, яко в чертозе древле бегающий дела, горькаго совета фараонитска.


Аще бабы слышала еси, убивающия иногда безвозрастное мужеское, душе окаянная, целомудрия деяние, ныне, яко великий Моисей, сси премудрость.


Яко Моисей великий египтянина, ума, уязвивши, окаянная, не убила еси, душе; и како вселишися, глаголи, в пустыню страстей покаянием?


В пустыню вселися великий Моисей; гряди убо, подражай того житие, да и в купине Богоявления, душе, в видении будеши.


Моисеов жезл воображай, душе, ударяющий море и огустевающий глубину во образ Креста Божественнаго: имже можеши и ты великая совершити.


Аарон приношаше огнь Богу непорочный, нелестный; но Офни и Финеес, яко ты, душе, приношаху чуждее Богу, оскверненное житие.


Яко тяжкий нравом, фараону горькому бых, Владыко, Ианни и Иамври, душею и телом, и погружен умом, но помози ми.


Калу примесихся, окаянный, умом, омый мя, Владыко, банею моих слез, молю Тя, плоти моея одежду убелив яко снег.


Аще испытаю моя дела, Спасе, всякаго человека превозшедша грехами себе зрю, яко разумом мудрствуяй, согреших не неведением.


Пощади, пощади, Господи, создание Твое, согреших, ослаби ми, яко естеством чистый Сам сый Един, и ин разве Тебе никтоже есть кроме скверны.


Мене ради Бог сый, вообразился еси в мя, показал еси чудеса, исцелив прокаженныя и разслабленнаго стягнув, кровоточивыя ток уставил еси, Спасе, прикосновением риз.


Кровоточивую подражай, окаянная душе, притецы, удержи ометы Христовы, да избавишися ран и услышиши от Него: вера твоя спасе тя.


Низу сничащую подражай, о душе, прииди, припади к ногама Иисусовыма, да тя исправит, и да ходиши право стези Господни.


Аще и кладязь еси глубокий, Владыко, источи ми воду из пречистых Твоих жил, да, яко самаряныня, не ктому, пияй, жажду: жизни бо струи источаеши.


Силоам да будут ми слезы моя, Владыко Господи, да умыю и аз зеницы сердца и вижду Тя умно, Света превечна.


\medskip\bfseries Преподобная мати Марие, моли Бога о нас.\normalfont{}\nopagebreak


Несравненным желанием, всебогатая, древу возжелевши поклонитися животному, сподобилася еси желания, сподоби убо и мене улучити вышния славы.


\medskip\bfseries Преподобная мати Марие, моли Бога о нас.\normalfont{}\nopagebreak


Струи Иорданския прешедши, обрела еси покой безболезненный, плоти сласти избежавши, еяже и нас изми твоими молитвами, преподобная.


Преподобный отче Андрее, моли Бога о нас.


Яко пастырей изряднейша, Андрее премудре, избранна суща тя, любовию велиею и страхом молю, твоими молитвами спасение улучити и жизнь вечную.


\itshape Слава\normalfont{}: Тя, Троице, славим, Единаго Бога: Свят, Свят, Свят еси, Отче, Сыне и Душе, Простое Существо, Единице присно покланяемая.


\itshape И ныне\normalfont{}: Из Тебе облечеся в мое смешение, нетленная, безмужная Мати Дево, Бог, создавый веки, и соедини Себе человеческое естество. 


\medskip\bfseries Песнь 6\normalfont{}\nopagebreak


\itshape Ирмос\normalfont{}: Возопих всем сердцем моим к щедрому Богу, и услыша мя от ада преисподняго, и возведе от тли живот мой.


Слезы, Спасе, очию моею и из глубины воздыхания чисте приношу, вопиющу сердцу: Боже, согреших Ти, очисти мя.


Уклонилася еси, душе, от Господа твоего, якоже Дафан и Авирон, но пощади, воззови из ада преисподняго, да не пропасть земная тебе покрыет.


Яко юница, душе, разсвирепевшая, уподобилася еси Ефрему, яко серна от тенет сохрани житие, вперивши деянием ум и зрением.


Рука нас Моисеова да уверит, душе, како может Бог прокаженное житие убелити и очистити, и не отчайся сама себе, аще и прокаженна еси.


Волны, Спасе, прегрешений моих, яко в мори Чермнем возвращающеся, покрыша мя внезапу, яко египтяны иногда и тристаты.


Неразумное, душе, произволение имела еси, яко прежде Израиль: Божественныя бо манны предсудила еси безсловесно любосластное страстей объядение.


Кладенцы, душе, предпочла еси хананейских мыслей паче жилы камене, из негоже премудрости река, яко чаша, проливает токи богословия.


Свиная мяса и котлы и египетскую пищу паче Небесныя предсудила еси, душе моя, якоже древле неразумнии людие в пустыни.


Яко удари Моисей, раб Твой, жезлом камень, образно животворивая ребра Твоя прообразоваше, из нихже вси питие жизни, Спасе, почерпаем.


Испытай, душе, и смотряй, якоже Иисус Навин, обетования землю, какова есть, и вселися в ню благозаконием.


Востани и побори, яко Иисус Амалика, плотския страсти, и гаваониты, лестныя помыслы, присно побеждающи.


Прейди времене текущее естество, яко прежде ковчег, и земли оныя буди во одержании обетования, душе, Бог повелевает.


Яко спасл еси Петра, возопивша, спаси, предварив мя, Спасе, от зверя избави, простер Твою руку, и возведи из глубины греховныя.


Пристанище Тя вем утишное, Владыко, Владыко Христе, но от незаходимых глубин греха и отчаяния мя, предварив, избави.


Аз есмь, Спасе, юже погубил еси древле царскую драхму; но вжег светильник, Предтечу Твоего, Слове, взыщи и обрящи Твой образ.


\medskip\bfseries Преподобная мати Марие, моли Бога о нас.\normalfont{}\nopagebreak


Да страстей пламень угасиши, слез капли источала еси присно, Марие, душею распалаема, ихже благодать подаждь и мне, твоему рабу.


\medskip\bfseries Преподобная мати Марие, моли Бога о нас.\normalfont{}\nopagebreak


Безстрастие Небесное стяжала еси крайним на земли житием, мати. Темже тебе поющим, от страстей избавитися молитвами твоими молися.


Преподобный отче Андрее, моли Бога о нас.


Критскаго тя пастыря и председателя и вселенныя молитвенника ведый, притекаю, Андрее, и вопию ти: изми мя, отче, из глубины греха.


\itshape Слава\normalfont{}: Троица есмь Проста, Нераздельна, раздельна Личне и Единица есмь естеством соединена, Отец глаголет, и Сын, и Божественный Дух.


\itshape И ныне\normalfont{}: Утроба Твоя Бога нам роди, воображена по нам: Егоже, яко Создателя всех, моли, Богородице, да молитвами Твоими оправдимся.


\itshape Ирмос\normalfont{}: Возопих всем сердцем моим к щедрому Богу, и услыша мя от ада преисподняго, и возведе от тли живот мой. 

\medskip\bfseries Кондак, глас 6:\normalfont{}\nopagebreak


;Душе моя, душе моя, востани, что спиши? конец приближается, и имаши смутитися: воспряни убо, да пощадит тя Христос Бог, везде сый и вся исполняяй. 


\medskip\bfseries Икос:\normalfont{}\nopagebreak


Христово врачевство видя отверсто и от сего Адаму истекающее здравие, пострада, уязвися диавол и, яко бедствуя, рыдаше и своим другом возопи: что сотворю Сыну Мариину, убивает мя Вифлеемлянин, Иже везде сый и вся исполняяй. 

\medskip\bfseries Таже блаженны, глас 6:\normalfont{}\nopagebreak


Во Царствии Твоем помяни нас, Господи.


Разбойника, Христе, рая жителя сотворил еси, на кресте Тебе возопивша: помяни мя; того покаянию сподоби и мене, недостойнаго.


Блажени нищии духом, яко тех есть Царство Небесное.


Маноя слышавши древле, душе моя, Бога в явлении бывша и из неплодове тогда приемша плод обетования, того благочестие подражай.


Блажени плачущии, яко тии утешатся.


Сампсоновой поревновавши лености, главу остригла еси, душе, дел твоих, предавши иноплеменником любосластием целомудренную жизнь и блаженную.


Блажени кротции, яко тии наследят землю.


Прежде челюстию ослею победивый иноплеменники, ныне пленение ласкосердству страстному обретеся; но избегни, душе моя, подражания, деяния и слабости.


Блажени алчущии и жаждущии правды, яко тии насытятся.


Варак и Иеффай военачальницы, судии Израилевы предпочтени быша, с нимиже Деворра мужеумная; тех доблестьми, душе, вмужившися, укрепися.


Блажени милостивии, яко тии помиловани будут.


Иаилино храбрство познала еси, душе моя, Сисара древле прободшую и спасение соделавшую древом острым, слышиши, имже тебе крест образуется.


Блажени чистии сердцем, яко тии Бога узрят.


Пожри, душе, жертву похвальную, деяние, яко дщерь, принеси от Иеффаевы чистейшую и заколи, яко жертву, страсти плотския Господеви твоему.


Блажени миротворцы, яко тии сынове Божии нарекутся.


Гедеоново руно помышляй, душе моя, с небесе росу подыми и приникни, якоже пес, и пий воду, от закона текущую, изгнетением письменным.


Блажени изгнани правды ради, яко тех есть Царство Небесное.


Илии священника осуждение, душе моя, восприяла еси, лишением ума приобретши страсти себе, якоже он чада, делати беззаконная.


Блажени есте, егда поносят вам и изженут и рекут всяк зол глагол на вы лжуще, Мене ради.


В судиях левит небрежением свою жену дванадесятим коленом раздели, душе моя, да скверну обличит от Вениамина беззаконную.


Радуйтеся и веселитеся, яко мзда ваша многа на Н ебесех.


Любомудренная Анна молящися, устне убо двизаше ко хвалению, глас же ея не слышашеся, но обаче неплодна сущи, сына молитвы раждает достойна.


Помяни нас, Господи, егда приидеши во Царствии Твоем.


В судиях спричтеся Аннино порождение, великий Самуил, егоже воспитала Армафема в дому Господни; тому поревнуй, душе моя, и суди прежде инех дела твоя.


Помяни нас, Владыко, егда приидеши во Царствии Твоем.


Давид на царство избран, царски помазася рогом Божественного мира; ты убо, душе моя, аще хощеши вышняго Царствия, миром помажися слезами.


Помяни нас, Святый, егда приидеши во Царствии Твоем.


Помилуй создание Твое, Милостиве, ущедри руку Твоею творение и пощади вся согрешившия, и мене паче всех, Твоих презревшаго повелений.


\itshape Слава\normalfont{}: Безначальну и рождению же и происхождению Отцу покланяюся рождшему, Сына славлю рожденнаго, пою сопросиявшаго Отцу же и Сыну Духа Святаго.


\itshape И ныне\normalfont{}: Преестественному Рождеству Твоему покланяемся, по естеству славы Младенца Твоего не разделяюще, Богородительнице: Иже бо Един Лицем, сугубыми исповедуется естествы. 


\medskip\bfseries Песнь 7\normalfont{}\nopagebreak


\itshape Ирмос\normalfont{}: Согрешихом, беззаконновахом, неправдовахом пред Тобою, ниже соблюдохом, ниже сотворихом, якоже заповедал еси нам; но не предаждь нас до конца, отцев Боже.


Согреших, беззаконновах и отвергох заповедь Твою, яко во гресех произведохся, и приложих язвам струпы себе; но Сам мя помилуй, яко благоутробен, отцев Боже.


Тайная сердца моего исповедах Тебе, Судии моему, виждь мое смирение, виждь и скорбь мою, и вонми суду моему ныне, и Сам мя помилуй, яко благоутробен, отцев Боже.


Саул иногда, яко погуби отца своего, душе, ослята, внезапу царство обрете к прослутию; но блюди, не забывай себе, скотския похоти твоя произволивши паче Царства Христова.


Давид иногда Богоотец, аще и согреши сугубо, душе моя, стрелою убо устрелен быв прелюбодейства, копием же пленен быв убийства томлением; но ты сама тяжчайшими делы недугуеши, самохотными стремленьми.


Совокупи убо Давид иногда беззаконию беззаконие, убийству же любодейство растворив, покаяние сугубое показа абие; но сама ты, лукавнейшая душе, соделала еси, не покаявшися Богу.


Давид иногда вообрази, списав яко на иконе песнь, еюже деяние обличает, еже содея, зовый: помилуй мя, Тебе бо Единому согреших всех Богу, Сам очисти мя.


Кивот яко ношашеся на колеснице, Зан оный, егда превращшуся тельцу, точию коснуся, Божиим искусися гневом; но того дерзновения убежавши, душе, почитай Божественная честне.


Слышала еси Авессалома, како на естество воста, познала еси того скверная деяния, имиже оскверни ложе Давида отца; но ты подражала еси того страстная и любосластная стремления.


Покорила еси неработное твое достоинство телу твоему, иного бо Ахитофела обретши врага, душе, снизшла еси сего советом; но сия разсыпа Сам Христос, да ты всяко спасешися.


Соломон чудный и благодати премудрости исполненный, сей лукавое иногда пред Богом сотворив, отступи от Него; емуже ты проклятым твоим житием, душе, уподобилася еси.


Сластьми влеком страстей своих, оскверняшеся, увы мне, рачитель премудрости, рачитель блудных жен, и странен от Бога; егоже ты подражала еси умом, о душе, сладострастьми скверными.


Ровоаму поревновала еси, не послушавшему совета отча, купно же и злейшему врагу Иеровоаму, прежнему отступнику, душе, но бегай подражания и зови Богу: согреших, ущедри мя.


Ахаавовым поревновала еси сквернам, душе моя, увы мне, была еси плотских скверн пребывалище и сосуд срамлен страстей, но из глубины твоея воздохни и глаголи Богу грехи твоя.


Попали Илия иногда дващи пятьдесят Иезавелиных, егда студныя пророки погуби во обличение Ахаавово, но бегай подражания двою, душе, и укрепляйся.


Заключися тебе небо, душе, и глад Божий постиже тя, егда Илии Фесвитянина, якоже Ахаав, не покорися словесем иногда, но Сараффии уподобися, напитай пророчу душу.


Манассиева собрала еси согрешения изволением, поставльши, яко мерзости, страсти, и умноживши, душе, негодование, но того покаянию ревнующи тепле, стяжи умиление.


Припадаю Ти и приношу Тебе, якоже слезы, глаголы моя: согреших, яко не согреши блудница, и беззаконновах, яко иный никтоже на земли. Но ущедри, Владыко, творение Твое и воззови мя.


Погребох образ Твой и растлих заповедь Твою, вся помрачися доброта, и страстьми угасися, Спасе, свеща. Но, ущедрив, воздаждь ми, якоже поет Давид, радование.


Обратися, покайся, открый сокровенная, глаголи Богу, вся ведущему: Ты веси моя тайная, едине Спасе. Но Сам мя помилуй, якоже поет Давид, по милости Твоей.


Исчезоша дние мои, яко соние востающаго; темже, яко Езекия, слезю на ложи моем, приложитися мне летом живота. Но кий Исаия предстанет тебе, душе, аще не всех Бог?


\medskip\bfseries Преподобная мати Марие, моли Бога о нас.\normalfont{}\nopagebreak


Возопивши к Пречистей Богоматери, первее отринула еси неистовство страстей, нужно стужающих, и посрамила еси врага запеншаго. Но даждь ныне помощь от скорби и мне, рабу твоему.


\medskip\bfseries Преподобная мати Марие, моли Бога о нас.\normalfont{}\nopagebreak


Егоже возлюбила еси, Егоже возжелела еси, Егоже ради плоть изнурила еси, преподобная, моли ныне Христа о рабех: яко да милостив быв всем нам, мирное состояние дарует почитающим Его.


\medskip\bfseries Преподобне отче Андрее, моли Бога о нас.\normalfont{}\nopagebreak


На камени мя веры молитвами твоими утверди, отче, страхом мя Божественным ограждая, и покаяние, Андрее, подаждь ми, молюся ти, и избави мя от сети врагов, ищущих мя.


\itshape Слава\normalfont{}: Троице Простая, Нераздельная, Единосущная и Естество Едино, Светове и Свет, и Свята Три, и Едино Свято поется Бог Троица; но воспой, прослави Живот и Животы, душе, всех Бога.


\itshape И ныне\normalfont{}: Поем Тя, благословим Тя, покланяемся Ти, Богородительнице, яко Неразлучныя Троицы породила еси Единаго Христа Бога и Сама отверзла еси нам, сущим на земли, Небесная. 


\medskip\bfseries Трипеснец, глас 8:

\normalfont{}


\medskip\bfseries Песнь 8\normalfont{}\nopagebreak


\itshape Ирмос\normalfont{}: Безначальнаго Царя славы, Егоже трепещут Небесныя силы, пойте, священницы, людие, превозносите во вся веки.


\medskip\bfseries Святии апостоли, молите Бога о нас.\normalfont{}\nopagebreak


Яко углие невещественнаго огня, попалите вещественныя страсти моя, возжизающе ныне во мне желание Божественныя любве, апостоли.


\medskip\bfseries Святии апостоли, молите Бога о нас.\normalfont{}\nopagebreak


Трубы благогласныя Слова почтим, имиже падоша стены неутверждены вражия и богоразумия утвердишася забрала.


\medskip\bfseries Святии апостоли, молите Бога о нас.\normalfont{}\nopagebreak


Кумиры страстныя души моея сокрушите, иже храмы и столпы сокрушисте врага, апостоли Господни, храмове освященнии.


Пресвятая Богородице, спаси нас.


Вместила еси Невместимаго естеством, носила еси Носящаго вся, доила еси, Чистая, питающаго тварь Христа Жизнодавца.


Иный трипеснец. \itshape Ирмос\normalfont{}: Безначальнаго Царя:


\medskip\bfseries Святии апостоли, молите Бога о нас.\normalfont{}\nopagebreak


Духа началохитростием создавше всю Церковь, апостоли Христовы, в ней благословите Христа во веки.


\medskip\bfseries Святии апостоли, молите Бога о нас.\normalfont{}\nopagebreak


Вострубивше трубою учений, низвергоша апостоли всю лесть идольскую, Христа превозносяща во вся веки.


\medskip\bfseries Святии апостоли, молите Бога о нас.\normalfont{}\nopagebreak


Апостоли, доброе преселение, назирателие мира и Небеснии жителие, вас присно восхваляющия избавите от бед.


\medskip\bfseries Пресвятая Троице, Боже наш, слава Тебе.\normalfont{}\nopagebreak


Трисолнечное Всесветлое Богоначалие, Единославное и Единопрестольное Естеству, Отче Вседетелю, Сыне и Божественный Душе, пою Тя во веки.


\medskip\bfseries Пресвятая Богородице, спаси нас.\normalfont{}\nopagebreak


Яко честный и превышший престол, воспоим Божию Матерь непрестанно, людие, Едину по рождестве Матерь и Деву.


Великаго канона \itshape Ирмос\normalfont{}: Егоже воинства Небесная славят, и трепещут херувими и серафими, всяко дыхание и тварь, пойте, благословите и превозносите во вся веки.


Согрешивша, Спасе, помилуй, воздвигни мой ум ко обращению, приими мя кающагося, ущедри вопиюща: согреших Ти, спаси, беззаконновах, помилуй мя.


Колесничник Илия колесницею добродетелей вшед, яко на небеса, ношашеся превыше иногда от земных; сего убо, душе моя, восход помышляй.

Иорданова струя первее милотию Илииною Елиссеем ста сюду и сюду; ты же, о душе моя, сея не причастилася еси благодати за невоздержание.


Елиссей иногда прием милоть Илиину, прият сугубую благодать от Бога; ты же, о душе моя, сея не причастилася еси благодати за невоздержание.


Соманитида иногда праведнаго учреди, о душе, нравом благим; ты же не ввела еси в дом ни странна, ни путника. Темже чертога изринешися вон, рыдающи.


Гиезиев подражала еси, окаянная, разум скверный всегда, душе, егоже сребролюбие отложи поне на старость; бегай геенскаго огня, отступивши злых твоих.


Ты Озии, душе, поревновавши, сего прокажение в себе стяжала еси сугубо: безместная бо мыслиши, беззаконная же дееши; остави, яже имаши, и притецы к покаянию.


Ниневитяны, душе, слышала еси кающияся Богу, вретищем и пепелом, сих не подражала еси, но явилася еси злейшая всех, прежде закона и по законе прегрешивших.


В рове блата слышала еси Иеремию, душе, града Сионя рыданьми вопиюща и слез ищуща; подражай сего плачевное житие и спасешися.


Иона в Фарсис побеже, проразумев обращение ниневитянов, разуме бо, яко пророк, Божие благоутробие: темже ревноваше пророчеству не солгатися.


Даниила в рове слышала еси, како загради уста, о душе, зверей; уведела еси, како отроцы, иже о Азарии, погасиша верою пещи пламень горящий.


Ветхаго Завета вся приведох ти, душе, к подобию; подражай праведных боголюбивая деяния, избегни же паки лукавых грехов.


Правосуде Спасе, помилуй и избави мя огня и прещения, еже имам на суде праведно претерпети; ослаби ми прежде конца, добродетелию и покаянием.


Яко разбойник, вопию Ти: помяни мя; яко Петр, плачу горце: ослаби ми, Спасе; зову, яко мытарь, слезю, яко блудница; приими мое рыдание, якоже иногда хананеино.


Гноение, Спасе, исцели смиренныя моея души, Едине Врачу, пластырь мне наложи, и елей, и вино, дела покаяния, умиление со слезами.


Хананею и аз подражая, помилуй мя, вопию, Сыне Давидов; касаюся края ризы, яко кровоточивая, плачу, яко Марфа и Мария над Лазарем.


Слезную, Спасе, сткляницу яко миро истощавая на главу, зову Ти, якоже блудница, милости ищущая, мольбу приношу и оставление прошу прияти.


Аще и никтоже, якоже аз, согреши Тебе, но обаче приими и мене, благоутробне Спасе, страхом кающася и любовию зовуща: согреших Тебе Единому, помилуй мя, Милостиве.


Пощади, Спасе, Твое создание и взыщи, яко Пастырь, погибшее, предвари заблуждшаго, восхити от волка, сотвори мя овча на пастве Твоих овец.

Егда, Судие, сядеши, яко благоутробен, и покажеши страшную славу Твою, Спасе, о каковый страх тогда, пещи горящей, всем боящимся нестерпимаго судища Твоего.


\medskip\bfseries Преподобная мати Марие, моли Бога о нас.\normalfont{}\nopagebreak


Света незаходимаго Мати тя просветивши, от омрачения страстей разреши. Темже вшедши в духовную благодать, просвети, Марие, тя верно восхваляющия.


\medskip\bfseries Преподобная мати Марие, моли Бога о нас.\normalfont{}\nopagebreak


Чудо ново видев, ужасашеся божественный в тебе воистинну, мати, Зосима: ангела бо зряше во плоти и ужасом весь исполняшеся, Христа поя во веки.


\medskip\bfseries Преподобне отче Андрее, моли Бога о нас.\normalfont{}\nopagebreak


Яко дерзновение имый ко Господу, Андрее Критский, честная похвало, молю, молися разрешение от уз беззакония ныне обрести мне молитвами твоими, яко покаяния учитель и преподобных слава.


Благословим Отца и Сына и Святаго Духа Господа.

Безначальне Отче, Сыне Собезначальне, Утешителю Благий, Душе Правый, Слова Божия Родителю, Отца Безначальна Слове, Душе Живый и Зиждяй, Троице Единице, помилуй мя.


\itshape И ныне\normalfont{} и присно и во веки веков.


Яко от оброщения червленицы, Пречистая, умная багряница Еммануилева внутрь во чреве Твоем плоть исткася. Темже Богородицу воистинну Тя почитаем.


Хвалим, благословим, покланяемся Господеви, поюще и превозносяще во вся веки.


\itshape Ирмос\normalfont{}: Егоже воинства Небесная славят, и трепещут херувими и серафими, всяко дыхание и тварь, пойте, благословите и превозносите во вся веки. 

\medskip\bfseries Поем Честнейшую:


Трипеснец, глас 8:\normalfont{}


\medskip\bfseries Песнь 9\normalfont{}\nopagebreak


\itshape Ирмос\normalfont{}: Воистинну Богородицу Тя исповедуем, спасеннии Тобою, Дево чистая, с безплотными лики Тя величающе.


\medskip\bfseries Святии апостоли, молите Бога о нас.\normalfont{}\nopagebreak


Источницы спасительныя воды явльшеся апостоли, истаявшую душу мою греховною жаждою оросите.


\medskip\bfseries Святии апостоли, молите Бога о нас.\normalfont{}\nopagebreak


Плавающаго в пучине погибели и в погружении уже бывша Твоею десницею, якоже Петра, Господи, спаси мя.


\medskip\bfseries Святии апостоли, молите Бога о нас.\normalfont{}\nopagebreak


Яко соли, вкусных суще учений, гнильство ума моего изсушите и неведения тьму отжените.


\medskip\bfseries Пресвятая Богородице, спаси нас.\normalfont{}\nopagebreak


Радость яко родившая, плач мне подаждь, имже Божественное утешение, Владычице, в будущем дни обрести возмогу.


Иный. \itshape Ирмос\normalfont{}: Тя, Небесе и земли Ходатаицу:

\medskip\bfseries Святии апостоли, молите Бога о нас.\normalfont{}\nopagebreak


Тя, благославное апостольское собрание, песньми величаем: вселенней бо светила светлая явистеся, прелесть отгоняще.


\medskip\bfseries Святии апостоли, молите Бога о нас.\normalfont{}\nopagebreak


Благовестною мрежею вашею словесныя рыбы уловивше, сия приносите всегда снедь Христу, апостоли блаженнии.

\medskip\bfseries Святии апостоли, молите Бога о нас.\normalfont{}\nopagebreak


К Богу вашим прошением помяните нас, апостоли, от всякаго избавитися искушения, молимся, любовию воспевающия вас.


\medskip\bfseries Пресвятая Троице, Боже наш, слава Тебе.\normalfont{}\nopagebreak

Тя, Триипостасную Единицу, Отче, Сыне со Духом, Единаго Бога Единосущна пою, Троицу Единосильную Безначальную.


\medskip\bfseries Пресвятая Богородице, спаси нас.\normalfont{}\nopagebreak


Тя, Детородительницу и Деву, вси роди ублажаем, яко Тобою избавльшеся от клятвы: радость бо нам родила еси, Господа.


Великаго канона \itshape Ирмос\normalfont{}: Безсеменнаго зачатия Рождество несказанное, Матере безмужныя нетленен Плод, Божие бо Рождение обновляет естества. Темже Тя вси роди, яко Богоневестную Матерь, православно величаем.


Ум острупися, тело оболезнися, недугует дух, слово изнеможе, житие умертвися, конец при дверех. Темже, моя окаянная душе, что сотвориши, егда приидет Судия испытати твоя?


Моисеово приведох ти, душе, миробытие и от того все заветное Писание, поведающее тебе праведныя и неправедныя; от нихже вторыя, о душе, подражала еси, а не первыя, в Бога согрешивши.


Закон изнеможе, празднует Евангелие, Писание же все в тебе небрежено бысть, пророцы изнемогоша, и все праведное слово; струпи твои, о душе, умножишася, не сущу врачу, исцеляющему тя.


Новаго привожду ти Писания указания, вводящая тя, душе, ко умилению: праведным убо поревнуй, грешных же отвращайся и умилостиви Христа молитвами же, и пощеньми, и чистотою, и говением.


Христос вочеловечися, призвав к покаянию разбойники и блудницы; душе, покайся, дверь отверзеся Царствия уже, и предвосхищают е фарисее, и мытари, и прелюбодеи кающиися.


Христос вочеловечися, плоти приобщився ми, и вся елика суть естества хотением исполни греха кроме, подобие тебе, о душе, и образ предпоказуя Своего снисхождения.


Христос волхвы спасе, пастыри созва, младенец множества показа мученики, старцы прослави и старыя вдовицы, ихже не поревновала еси, душе, ни деянием, ни житию, но горе тебе, внегда будеши судитися.


Постився Господь дний четыредесять в пустыни, последи взалка, показуя человеческое; душе, да не разленишися, аще тебе приложится враг, молитвою же и постом от ног твоих да отразится.


Христос искушашеся, диавол искушаше, показуя камение, да хлеби будут, на гору возведе видети вся царствия мира во мгновении; убойся, о душе, ловления, трезвися, молися на всякий час Богу.


Горлица пустыннолюбная, глас вопиющаго возгласи, Христов светильник, проповедуяй покаяние, Ирод беззаконнова со Иродиадою. Зри, душе моя, да не увязнеши в беззаконныя сети, но облобызай покаяние.


В пустыню вселися благодати Предтеча, и Иудея вся и Самария слышавше течаху и исповедаху грехи своя, крещающеся усердно: ихже ты не подражала еси, душе.


Брак убо честный и ложе нескверно, обоя бо Христос прежде благослови, плотию ядый, и в Кане же на браце воду в вино совершая, и показуя первое чудо, да ты изменишися, о душе.


Разслабленнаго стягну Христос, одр вземша, и юношу умерша воздвиже, вдовиче рождение, и сотнича отрока, и самаряныне явися, в дусе службу тебе, душе, предживописа.


Кровоточивую исцели прикосновением края ризна Господь, прокаженныя очисти, слепыя и хромыя просветив, исправи, глухия же и немыя и ничащия низу исцели словом: да ты спасешися, окаянная душе.


Недуги исцеляя, нищим благовествоваше Христос Слово, вредныя уврачева, с мытари ядяше, со грешники беседоваше, Иаировы дщере душу предумершую возврати осязанием руки.


Мытарь спасашеся, и блудница целомудрствоваше, и фарисей, хваляся, осуждашеся. Ов убо: очисти мя; ова же: помилуй мя; сей же величашеся вопия: Боже, благодарю Тя, и прочия безумныя глаголы.


Закхей мытарь бе, но обаче спасашеся, и фарисей Симон соблажняшеся, и блудница приимаше оставительная разрешения от Имущаго крепость оставляти грехи, юже, душе, потщися подражати.


Блуднице, о окаянная душе моя, не поревновала еси, яже приимши мира алавастр, со слезами мазаше нозе Спасове, отре же власы, древних согрешений рукописание Раздирающаго ея.


Грады, имже даде Христос благовестие, душе моя, уведала еси, како прокляти быша. Убойся указания, да не будеши якоже оны, ихже содомляном Владыка уподобив, даже до ада осуди.


Да не горшая, о душе моя, явишися отчаянием, хананеи веру слышавшая, еяже дщи словом Божиим исцелися; Сыне Давидов, спаси и мене, воззови из глубины сердца, якоже она Христу.


Умилосердися, спаси мя, Сыне Давидов, помилуй, беснующияся словом исцеливый, глас же благоутробный, яко разбойнику, мне рцы: аминь, глаголю тебе, со Мною будеши в раи, егда прииду во славе Моей.


Разбойник оглаголоваше Тя, разбойник богословяше Тя: оба бо на кресте свисяста. Но, о Благоутробне, яко верному разбойнику Твоему, познавшему Тя Бога, и мне отверзи дверь славнаго Царствия Твоего.


Тварь содрогашеся, распинаема Тя видящи, горы и камения страхом распадахуся, и земля сотрясашеся, и ад обнажашеся, и соомрачашеся свет во дне, зря Тебе, Иисусе, пригвождена ко Кресту.


Достойных покаяния плодов не истяжи от мене, ибо крепость моя во мне оскуде; сердце мне даруй присно сокрушенное, нищету же духовную: да сия Тебе принесу яко приятную жертву, едине Спасе.


Судие мой и Ведче мой, хотяй паки приити со ангелы судити миру всему, милостивным Твоим оком тогда видев мя, пощади и ущедри мя, Иисусе, паче всякаго естества человеча согрешивша.


\medskip\bfseries Преподобная мати Марие, моли Бога о нас.\normalfont{}\nopagebreak


Удивила еси всех странным житием твоим, ангелов чины и человеков соборы, невещественно поживши и естество прешедши; имже, яко невещественныма ногама вшедши, Марие, Иордан прешла еси.


\medskip\bfseries Преподобная мати Марие, моли Бога о нас.\normalfont{}\nopagebreak


Умилостиви Создателя о хвалящих тя, преподобная мати, избавитися озлоблений и скорбей, окрест нападающих: да избавившеся от напастей, возвеличим непрестанно прославльшаго тя Господа.


\medskip\bfseries Преподобне отче Андрее, моли Бога о нас.\normalfont{}\nopagebreak


Андрее честный и отче треблаженнейший, пастырю Критский, не престай моляся о воспевающих тя: да избавимся вси гнева, и скорби, и тления, и прегрешений безмерных, чтущии твою память верно.


\itshape Слава\normalfont{}: Троице Единосущная, Единице Триипостасная, Тя воспеваем, Отца славяще, Сына величающе и Духу покланяющеся, Единому Естеству воистинну Богу, Жизни же и живущему Царству безконечному.


\itshape И ныне\normalfont{}: Град Твой сохраняй, Богородительнице Пречистая, в Тебе бо сей верно царствуяй, в Тебе и утверждается, и Тобою побеждаяй, побеждает всякое искушение, и пленяет ратники, и проходит послушание.


Таже оба лика вкупе поют \itshape Ирмос\normalfont{}:


Безсеменнаго зачатия Рождество несказанное, Матере безмужныя нетленен Плод, Божие бо Рождение обновляет естества. Темже Тя вси роди, яко Богоневестную Матерь, православно величаем.\mychapterending


\mypart{ПРИЛОЖЕНИЯ}

\chapter{Словарь и термины}
\markright{Словарь и термины}
%http://www.molitvoslov.com/slovar.php 
 
\begin{mymulticols}{2}\footnotesize

\letter{А}


\noindent\textbf{Абие} "--- немедленно, тотчас. 




\noindent\textbf{Авва} "--- отец. 




\noindent\textbf{Аввадон} "--- евр. «Погубитель»; имя ангела бездны. 




\noindent\textbf{Авраамово недро, лоно} "--- иносказательно: рай, место вечного блаженства. 




\noindent\textbf{Агаряне} "--- потомки Исмаила, сына Агари, наложницы Авраама, иносказательно "--- кочевые восточные племена. 




\noindent\textbf{Агиасма} "--- освященная по церковному чину вода. Освященная на празднике Богоявления вода называется Великой агиасмой. 




\noindent\textbf{Агиос} "--- надписание на древних иконах; греч. «святой». 




\noindent\textbf{Агкира} (читается «анкира») "--- якорь. 




\noindent\textbf{Агнец} "--- ягненок; чистое, кроткое существо; изымаемая на проскомидии часть просфоры для Евхаристии; мн. ч. "--- « \noindent\textbf{агнцы}» "--- иногда значит «христиане». 




\noindent\textbf{Агница} "--- овечка. 




\noindent\textbf{Агня} "--- ягненок. 




\noindent\textbf{Ад} "--- место нахождения душ умерших до освобождения их Господом Иисусом Христом; место вечного мучения грешников; жилище диавола. 




\noindent\textbf{Адамант} "--- алмаз; бриллиант; драгоценный камень. 




\noindent\textbf{Адамантовый} "--- твердый; крепкий; драгоценный. 




\noindent\textbf{Адов} "--- адский. 




\noindent\textbf{Адонаи} "--- евр. «Господь мой». 




\noindent\textbf{Аер} "--- покровец, полагаемый сверху священных сосудов на Литургии. 




\noindent\textbf{Аермонский} "--- связанный с горой Аермон. 




\noindent\textbf{Аз} "--- я. 




\noindent\textbf{Аиромантия} "--- воздуховолхвование, т.~е. суеверное гадание на основании атмосферных явлений. 




\noindent\textbf{Акафист} "--- греч. «неседальное»; церковная служба, во время которой возбраняется сидеть. 




\noindent\textbf{Аки} "--- как будто, как бы. 




\noindent\textbf{Акриды} "--- пища Иоанна Крестителя; по мнению одних "--- род съедобной саранчи, или кузнечиков; по мнению других "--- какое-то растение. 




\noindent\textbf{Аксиос} "--- греч. «достоин». 




\noindent\textbf{Алавастр} "--- каменный сосуд. 




\noindent\textbf{Алектор} "--- петух. 




\noindent\textbf{Алкати} "--- голодать; хотеть есть, сильно желать чего-либо. 




\noindent\textbf{Алкота} "--- голод. 




\noindent\textbf{Аллилуия} "--- евр. «хвалите Бога»; «слава Богу!» 




\noindent\textbf{Аллилуия красная} "--- пение «аллилуйя» на особый умилительный распев. См. Триодь постную. 




\noindent\textbf{Аллилуиарий, аллилуиар} "--- стих, возглашаемый чтецом после чтения Апостола на Литургии. При этом возглашении на клиросах поют «аллилуия». 




\noindent\textbf{Алой, алое} "--- сок благовонного дерева, употреблявшийся для каждения и бальзамирования. 




\noindent\textbf{Алтабас} "--- самая лучшая старинная парча. 




\noindent\textbf{Амалик} "--- народ, живший между Палестиною и Египтом. В церковной поэзии это имя часто прилагается к диаволу. 




\noindent\textbf{Амвон} "--- возвышенная часть храма перед царскими вратами. 




\noindent\textbf{Амвросия} "--- неистлеваемая пища. 




\noindent\textbf{Амигдал} "--- миндаль. 




\noindent\textbf{Аминь} "--- евр. «да будет так»; «истинно»; «подлинно»; «да». 




\noindent\textbf{Амо, аможе} "--- куда. 




\noindent\textbf{Аможе аще} "--- куда бы ни. 




\noindent\textbf{Аналой} (правильнее \noindent\textbf{аналогий}) "--- возвышенный стол, на который полагаются церковные книги при чтении и иконы. 




\noindent\textbf{Анафема} "--- отлучение от общины верных и предание суду Божию; тот, кто подвергся такому отлучению. 




\noindent\textbf{Анафематствовати} "--- предавать анафеме. 




\noindent\textbf{Анахорет} "--- отшельник. 




\noindent\textbf{Ангел} "--- вестник. 




\noindent\textbf{Ангеловидный} "--- внешне напоминающий Ангела. 




\noindent\textbf{Ангелозрачный} "--- внешне напоминающий Ангела. 




\noindent\textbf{Ангелоименитый} "--- знаменитый, почитаемый в лике ангельском; носящий имя какого-либо Ангела. 




\noindent\textbf{Ангелолепный} "--- приличный Ангелу. 




\noindent\textbf{Ангеломудренный} "--- имеющий мудрость Ангела. 




\noindent\textbf{Ангельское житие, ангельский образ} "--- высшая степень монашеского совершенства; греч. «схима». 




\noindent\textbf{Анепсий} "--- племянник, родственник. 




\noindent\textbf{Антидор} "--- благословенный хлеб, т.~е. остатки той просфоры, из которой на проскомидии был изъят Агнец. 




\noindent\textbf{Антиминс} "--- греч. «вместопрестолие», освященный плат с изображением Иисуса Христа во гробе и вшитыми св. мощами. Только на антиминсе может быть совершаема Литургия. 




\noindent\textbf{Антифон} "--- греч. «противугласник»; песнопение, которое должно быть пето попеременно на обоих клиросах. 




\noindent\textbf{Антихрист} "--- греч. «противник Христа». 




\noindent\textbf{Антологион} "--- греч. «Цветослов»; название «Минеи праздничной». 




\noindent\textbf{Анфипат} "--- наместник, проконсул. 




\noindent\textbf{Анфракс} "--- яхонт. 




\noindent\textbf{Апокалипсис} "--- греч. «откровение». 




\noindent\textbf{Аполлион} "--- греч. «Погубитель»; имя ангела бездны. 




\noindent\textbf{Апостол} "--- греч. «посланник». 




\noindent\textbf{Апостасис} "--- отступничество. 




\noindent\textbf{Апостата} "--- отступник. 




\noindent\textbf{Априллий} "--- апрель. 




\noindent\textbf{Ариил} "--- горн у жертвенника всесожжения в храме Иерусалимском. 




\noindent\textbf{Армония} "--- гармония. 




\noindent\textbf{Ароматы} "--- душистая мазь. 




\noindent\textbf{Артос} "--- греч.хлеб квасной; он освящается с особой молитвой в день св. Пасхи. 




\noindent\textbf{Архангел} "--- начальствующий у Ангелов, название одного из чинов ангельских. 




\noindent\textbf{Архиерей} "--- первосвященник, епископ. 




\noindent\textbf{Архимагир} "--- главный повар. 




\noindent\textbf{Архипастырь} "--- первенствующий епископ. 




\noindent\textbf{Архисинагог} "--- начальник синагоги. 




\noindent\textbf{Архистратиг} "--- военачальник, полководец. 




\noindent\textbf{Архитектон} "--- архитектор, художник-строитель; главный строитель. 




\noindent\textbf{Архитриклин} "--- распорядитель пира. 




\noindent\textbf{Асмодей, Азмодеос} "--- «губитель», имя бесовское. 




\noindent\textbf{Аспид} "--- ядовитая змея. 




\noindent\textbf{Аспид парящий} "--- летающий ящер. 




\noindent\textbf{Ассарий} "--- мелкая медная монетка. 




\noindent\textbf{Астерикс} "--- звездица, поставляемая на дискосе при совершении Литургии. 




\noindent\textbf{Афарим} "--- соглядатаи; лазутчики. 




\noindent\textbf{Афедрон} "--- задний проход (Мф. 15, 17). 




\noindent\textbf{Афинеи} "--- афиняне. 




\noindent\textbf{Африкия} "--- Африка. 




\noindent\textbf{Аще} "--- если; хотя; или; ли. 




\noindent\textbf{Аще убо} "--- поскольку; потому что. 




\letter{Б} 





\noindent\textbf{Баальник} "--- волшебник. 




\noindent\textbf{Баба} "--- повивальная женщина. 




\noindent\textbf{Бабити} "--- помогать при родах. 




\noindent\textbf{Багряница} "--- ткань темно-красного цвета; порфира, пурпурная одежда высокопоставленных особ. 




\noindent\textbf{Балия} "--- колдунья; волшебница. 




\noindent\textbf{Баня пакибытия} "--- таинство св. Крещения. 




\noindent\textbf{Баснословити} "--- рассказывать небылицы; лгать. 




\noindent\textbf{Баснь} "--- ложное и бесполезное учение. 




\noindent\textbf{Бдение} "--- бодрствование; продолжительное ночное богослужение. 




\noindent\textbf{Бденно} "--- неусыпно, бодрственно. 




\noindent\textbf{Бденный} "--- неусыпный. 




\noindent\textbf{Бдети} "--- бодрствовать; не спать. 




\noindent\textbf{Бедне} "--- трудно; несносно; тяжело. 




\noindent\textbf{Бесоватися} "--- бесноваться, неистоваться. 




\noindent\textbf{Бедник} "--- калека; увечный. 




\noindent\textbf{Бедный} "--- иногда: увечный; калека. 




\noindent\textbf{Безведрие} "--- ненастье. 




\noindent\textbf{Безвидный} "--- не имеющий вида, образа. 




\noindent\textbf{Безвиновный} "--- не имеющий начала или причины для своего бытия. Одно из Божественных определений. 




\noindent\textbf{Безвозрастное} "--- младенец. 




\noindent\textbf{Безвременне} "--- некстати; неблаговременно. 




\noindent\textbf{Безгласие} "--- немота; молчание. 




\noindent\textbf{Безгодие} "--- бедствие; несчастье; тяжелый период в жизни. 




\noindent\textbf{Безквасный} "--- пресный; не кислый. 




\noindent\textbf{Безкнижный} "--- неученый. 




\noindent\textbf{Безлетно} "--- бесконечно; вечно; прежде всех времен. 




\noindent\textbf{Безматерен} "--- не имеющий матери. 




\noindent\textbf{Безмездник} "--- не принимающий мзды, платы. 




\noindent\textbf{Безмилостивный} "--- не чувствующий или не оказывающий милости, жалости. 




\noindent\textbf{Безмолвник} "--- пустынножитель; отшельник. 




\noindent\textbf{Безмолвный} "--- иногда значит: безопасный; спокойный. 




\noindent\textbf{Безневестный} "--- безбрачный; девственный. 




\noindent\textbf{Безочство} "--- нахальство; бесстыдство; наглость. 




\noindent\textbf{Безпрестани} "--- всегда; непрерывно. 




\noindent\textbf{Безпреткновенный} "--- не имеющий претыкания, соблазна, препятствия. 




\noindent\textbf{Безпутие} "--- совращение с пути; развращение. 




\noindent\textbf{Безсквернен} "--- не имеющий скверны или порока. 




\noindent\textbf{Безсловеснство} "--- скотство; глупость; безумие. 




\noindent\textbf{Безсловесные} "--- животные, скоты. 




\noindent\textbf{Безсребреник} "--- человек, трудящийся даром, бесплатно. 




\noindent\textbf{Безстудие} "--- бесстыдство. 




\noindent\textbf{Безцарный} "--- не имеющий над собой царя. 




\noindent\textbf{Безчадие} "--- неимение, лишение детей. 




\noindent\textbf{Безчаствовати} "--- лишать положенной части; обделять. 




\noindent\textbf{Безчестен} "--- обесчещенный. 




\noindent\textbf{Безчиние} "--- беспорядок; неустройство; смешение. 




\noindent\textbf{Безчинновати} "--- вести беспорядочную жизнь. 




\noindent\textbf{Безчисльство} "--- бесчисленное множество. 




\noindent\textbf{Бервенный} "--- деревянный. 




\noindent\textbf{Бесный} "--- бесноватый. 




\noindent\textbf{Бийца} "--- драчун. 




\noindent\textbf{Било} "--- колотушка, при помощи которой в монастырях созывают на молитву. 




\noindent\textbf{Бисер} "--- жемчуг. 




\noindent\textbf{Благий} "--- хороший; добрый. 




\noindent\textbf{Благовест} "--- удары колокола, созывающие христиан на молитву в храм. От \textbf{«звона»} отличается тем, что благовестят в один колокол, а звонят во многие. 



\noindent\textbf{Благовестити} "--- возвещать доброе; проповедовать. 




\noindent\textbf{Благоверный} "--- исповедующий правую веру; православный. 




\noindent\textbf{Благовещение} "--- добрая весть. 




\noindent\textbf{Благоволити} "--- хорошо относиться к кому-нибудь; принимать участие в ком-либо. 




\noindent\textbf{Благовоние} "--- благоухание, хороший запах. 




\noindent\textbf{Благовременне} "--- в удобное время. 




\noindent\textbf{Благогласник} "--- проповедник слова Божия. 




\noindent\textbf{Благодатный} "--- преисполненный Божественной благодати. 




\noindent\textbf{Благоделие} "--- доброе, богоугодное дело. 




\noindent\textbf{Благодушествовати} "--- радоваться. 




\noindent\textbf{Благоискусный} "--- имеющий знание в полезных вещах. 




\noindent\textbf{Благокласный} "--- приносящий обильную жатву. 




\noindent\textbf{Благоключимый} "--- случившийся вовремя. 




\noindent\textbf{Благокрасный} "--- очень красивый. 




\noindent\textbf{Благолепие} "--- красота; великолепие; богатое убранство. 




\noindent\textbf{Благолепно} "--- красиво; прилично. 




\noindent\textbf{Благолозный} "--- приносящий обильные, хорошие плоды. 




\noindent\textbf{Благолюбец} "--- склонный к добру. 




\noindent\textbf{Благомилостивый} "--- милосердный. 




\noindent\textbf{Благомощие} "--- крепость; сила. 




\noindent\textbf{Благомужство} "--- благоразумная храбрость, доблесть. 




\noindent\textbf{Благонаказательный} "--- направляющий к благонравию. 




\noindent\textbf{Благоодеждный} "--- украшенный изящной одеждой. 




\noindent\textbf{Благоотдатливый} "--- воздающий добром за зло. 




\noindent\textbf{Благопитание} "--- сладкая, вкусная пища. 




\noindent\textbf{Благопослушливый} "--- слушающий со вниманием; послушный. 




\noindent\textbf{Благопослушный} "--- внимательный; легко, хорошо слышимый. 




\noindent\textbf{Благопотребный} "--- хорошо устроенный; угодный; необходимый. 




\noindent\textbf{Благоразтворение} "--- очищение; оздоровление; прояснение. 




\noindent\textbf{Благоразтворити} "--- очищать; оздоровлять. 




\noindent\textbf{Благорасленный} "--- хорошо растущий. 




\noindent\textbf{Благорозгный} "--- ветвистый. 




\noindent\textbf{Благосеннолиственный} "--- тенистый. 




\noindent\textbf{Благосенный} "--- производящий обильную тень. 




\noindent\textbf{Благословенный} "--- прославляемый; восхваляемый; превозносимый. 




\noindent\textbf{Благословити} "--- посвятить Богу; желать добра; хвалить; помолиться о ниспослании Божией благодати на кого-либо; дозволить; пожелать добра. 




\noindent\textbf{Благословная вина} "--- уважительная причина. 




\noindent\textbf{Благостояние} "--- твердость, крепость (в добре, против зла). 




\noindent\textbf{Благостыня} "--- благодеяние; милосердие; добродетель, доброе дело. 




\noindent\textbf{Благость} "--- доброта. 




\noindent\textbf{Благотещи} "--- быстро идти. 




\noindent\textbf{Благоуветие} "--- снисхождение. 




\noindent\textbf{Благоуветливый} "--- снисходительный. 




\noindent\textbf{Благоутишие} "--- тихая, ясная погода. 




\noindent\textbf{Благоутробие} "--- милосердие. 




\noindent\textbf{Благохваление} "--- откровенная похвала. 




\noindent\textbf{Благоцветный} "--- испещренный; изобилующий цветами. 




\noindent\textbf{Благочествовати} "--- благоговеть; благоговейно почитать кого-либо. 




\noindent\textbf{Благочестие} "--- истинное Богопочитание. 




\noindent\textbf{Благочестивый, благочестный} "--- богобоязненный; благоговейный; почитающий Бога. 




\noindent\textbf{Блаженный} "--- счастливый. 




\noindent\textbf{Блажити} "--- ублажать; прославлять. 




\noindent\textbf{Блазнити} "--- соблазнять. 




\noindent\textbf{Блед, бледый} "--- бледный. 




\noindent\textbf{Блещатися} "--- блистать; сиять; светить. 




\noindent\textbf{Близна} "--- рубец; морщина; складка. 




\noindent\textbf{Блистание} "--- сверкание; излияние света, блеска. 




\noindent\textbf{Блудилище} "--- непотребный дом. 




\noindent\textbf{Блудодей} "--- нарушитель супружества. 




\noindent\textbf{Блудопитие} "--- побуждающая к блуду попойка. 




\noindent\textbf{Блужение} "--- неверность Богу истинному, служение идолам (Исх. 34, 15; Суд. 8, 33). Как нарушение брачного союза есть блудодеяние, так в духовном смысле и нарушение союза с Богом есть служение идолам, хождение во след богов иных, то есть блужение, тем более что некоторые виды идолослужения сопровождались блудом в собственном смысле слова. 




\noindent\textbf{Блюдомый} "--- сохраняемый. 




\noindent\textbf{Блюсти} "--- хранить; беречь; соблюдать. 




\noindent\textbf{Блядение} "--- суесловие; ложные слова; вранье. 




\noindent\textbf{Бо} "--- потому что; так как; ибо; поскольку. 




\noindent\textbf{Богатити} "--- обогащать. 




\noindent\textbf{Богоглаголивый} "--- говорящий по внушению от Бога или от Его Имени. 




\noindent\textbf{Богодельне} "--- по действию Бога. 




\noindent\textbf{Боголепно} "--- так, как прилично Богу. 




\noindent\textbf{Боголепный} "--- имеющий Божественную красоту, достоинство. 




\noindent\textbf{Богомужный} "--- Богочеловеческий. 




\noindent\textbf{Богоначальный} "--- имеющий в Боге свое начало. 




\noindent\textbf{Богоотец} "--- это название в церковных книгах усвояется Давиду, от рода которого родился Христос. 




\noindent\textbf{Бодренно} "--- бдительно; неусыпно. 




\noindent\textbf{Болезновати} "--- терпеть боль; страдать. 




\noindent\textbf{Болий} "--- больший. 




\noindent\textbf{Борзе} "--- скоро. 




\noindent\textbf{Борзитися} "--- спешить. 




\noindent\textbf{Боритель} "--- противник. 




\noindent\textbf{Брада} "--- борода. 




\noindent\textbf{Брадатый} "--- бородатый. 




\noindent\textbf{Бразда} "--- борозда. 




\noindent\textbf{Бракоокрадованная} "--- лишенная целомудрия, девственности. 




\noindent\textbf{Бранити} "--- запрещать; оборонять; препятствовать. 




\noindent\textbf{Брань} "--- война; битва. 




\noindent\textbf{Братися} "--- бороться; воевать. 




\noindent\textbf{Брашно} "--- пища; еда. 




\noindent\textbf{Бремя} "--- ноша; тяжесть. 




\noindent\textbf{Брение} "--- глина; грязь. 




\noindent\textbf{Бренный} "--- взятый из земли; слабый; непрочный. 




\noindent\textbf{Брещи} "--- стеречь; хранить. 




\noindent\textbf{Брозда} "--- удила (часть конской сбруи). 




\noindent\textbf{Бряцати} "--- звенеть. 




\noindent\textbf{Будильник} "--- один из монахов в обители, будящий на молитву братию. 




\noindent\textbf{Буесловие} "--- глупые речи; вранье. 




\noindent\textbf{Буесловити} "--- говорить глупые речи. 




\noindent\textbf{Буй, (буий)} "--- безумный; сумасшедший; глупый. 




\noindent\textbf{Буйство} "--- глупость; безумие; сумасшествие. 




\noindent\textbf{Былие} "--- трава. 




\letter{В} 





\noindent\textbf{Вавилонское семя} "--- племя нечестивцев. 




\noindent\textbf{Вага} "--- весы; тяжесть. 




\noindent\textbf{Вадити} "--- делать ложный донос; клеветать; обвинять; приманивать; привлекать. 




\noindent\textbf{Ваия} "--- ветви; листья. 




\noindent\textbf{Вайный} "--- состоящий из ваий. 




\noindent\textbf{Валсамный} "--- благовонный; пахучий; ароматический. 




\noindent\textbf{Вап (а)} "--- краска. 




\noindent\textbf{Вар} "--- зной; жара; кипяток. 




\noindent\textbf{Варити} "--- предварять; упреждать; опереживать, предостерегать. 




\noindent\textbf{Василиск} "--- большая ядовитая змея. 




\noindent\textbf{Вборзе} "--- скоро. 




\noindent\textbf{Ввергати} "--- вбрасывать. 




\noindent\textbf{Вдавати} "--- поручать; передавать; доверять. 




\noindent\textbf{Веглас} "--- знающий; искусный. 




\noindent\textbf{Ведети} "--- знать. 




\noindent\textbf{Ведунство} "--- волхвование; ворожба; чародейство. 




\noindent\textbf{Веельзевул} "--- «повелитель мух»; начальник злых духов; одно из имен сатаны. 




\noindent\textbf{Вежди} "--- веки; ресницы. 




\noindent\textbf{Веие} "--- ветвь; сучок. 




\noindent\textbf{Велелепие} "--- красота; великолепие; украшение. 




\noindent\textbf{Велелепота} "--- красота; великолепие; украшение. 




\noindent\textbf{Велемудренно} "--- высокомудренно. 




\noindent\textbf{Веление} "--- указ; повеление; заповедь; учение. 




\noindent\textbf{Велеречивый} "--- многословный; хвастливый. 




\noindent\textbf{Велеречити} "--- много говорить; хвастать; гордиться. 




\noindent\textbf{Велиар} (или \textbf{Велиал}) "--- одно из имен диавола. 




\noindent\textbf{Велий} "--- великий; сильный. 




\noindent\textbf{Великовыйный} "--- гордый. 




\noindent\textbf{Великое} "--- самый большой, главный колокол. 




\noindent\textbf{Великодарный, великодаровный, великодаровитый} "--- щедро награждающий. 




\noindent\textbf{Величатися} "--- гордиться; хвалиться; кичиться. 




\noindent\textbf{Велми} "--- весьма; очень. 




\noindent\textbf{Вельблуд, велбуд} "--- верблюд; толстый канат. 




\noindent\textbf{Вельблуждь} "--- верблюжий. 




\noindent\textbf{Вено} "--- плата жениха за невесту. 




\noindent\textbf{Венчати} "--- возлагать венок или венец; удостаивать; сподоблять. 




\noindent\textbf{Вепрь} "--- дикий кабан. 




\noindent\textbf{Вербие} "--- ива; лоза. 




\noindent\textbf{Вервица} "--- четки. 




\noindent\textbf{Верея} "--- дверь; перекладина; столб у ворот. 




\noindent\textbf{Вержение} "--- кидание; метание; бросание. 




\noindent\textbf{Вержение камени} "--- расстояние, равное тому, на какое можно бросить камень. 




\noindent\textbf{Верзити} "--- кинуть. 




\noindent\textbf{Верзитися} "--- упасть. 




\noindent\textbf{Вериги} "--- цепи; оковы. 




\noindent\textbf{Верт, вертоград} "--- сад. 




\noindent\textbf{Вертеп} "--- пещера. 




\noindent\textbf{Вертоградарь} "--- садовник. 




\noindent\textbf{Верху} "--- на; над; сверху. 




\noindent\textbf{Весь} "--- селение, деревушка. 




\noindent\textbf{Ветия} "--- оратор; ритор. 




\noindent\textbf{Ветрило} "--- парус. 




\noindent\textbf{Ветхий деньми} "--- имя Божие в Дан. 7, 9. На основании этого пророческого видения в иконографической традиции новозаветной Церкви образ Бога Отца изображается в виде старца. 




\noindent\textbf{Вечеря} "--- ужин; пир. 




\noindent\textbf{Вечеряти} "--- ужинать. 




\noindent\textbf{Вещенеистовный} "--- пристрастившийся до безумия к тленным благам. 




\noindent\textbf{Вещь} "--- дело; событие. 




\noindent\textbf{Взаим} "--- в долг; взаймы. 




\noindent\textbf{Взиматися} "--- подниматься. 




\noindent\textbf{Взыгратися} "--- играть; скакать; веселиться. 




\noindent\textbf{Взыскати} "--- стремиться, искать. 




\noindent\textbf{Взятися} "--- взяться; отвориться; совершиться. 




\noindent\textbf{Вина} "--- причина; обвинение; извинение. 




\noindent\textbf{Винарь} "--- виноградарь. 




\noindent\textbf{Винничина} "--- виноградная лоза. 




\noindent\textbf{Винопийца} "--- пьяница. 




\noindent\textbf{Винопитие} "--- употребление вина. 




\noindent\textbf{Висети} "--- висеть; держаться на чем-либо. 




\noindent\textbf{Виссон} "--- драгоценная тонкая пряжа желтоватого цвета или одежда из этой ткани. 




\noindent\textbf{Виталище} "--- место жительства; жилище. 




\noindent\textbf{Виталница} "--- комната; гостиница; постоялый двор; ночлег. 




\noindent\textbf{Витати} "--- обитать; пребывать; проживать; ночевать. 




\noindent\textbf{Вкупе} "--- вместе. 




\noindent\textbf{Влагалище} "--- мешок; карман; ларец. 




\noindent\textbf{Владычный} "--- господский или Господний. 




\noindent\textbf{Владящий} "--- обладающий; господствующий. 




\noindent\textbf{Власти} "--- имя одного из чинов ангельских. 




\noindent\textbf{Власяница} "--- одежда из жесткого, колючего волоса. 




\noindent\textbf{Влаятися} "--- мыкаться; колебаться; волноваться; носиться по волнам. 




\noindent\textbf{Влещи} "--- тащить. 




\noindent\textbf{Влещися} "--- брести; медленно идти; тащиться. 




\noindent\textbf{Вмале} "--- вскоре; немного спустя; почти; едва. 




\noindent\textbf{Вне(уду)} "--- извне; снаружи. 




\noindent\textbf{Внегда} "--- когда. 




\noindent\textbf{Внезапу} "--- вдруг; неожиданно. 




\noindent\textbf{Внемшийся} "--- загоревшийся. 




\noindent\textbf{Внитие} "--- вхождение; явление; пришествие. 




\noindent\textbf{Внове} "--- недавно. 




\noindent\textbf{Внутрь (уду)} "--- внутри. 




\noindent\textbf{Вняти} "--- обратить внимание; услышать. 




\noindent\textbf{Вняти от} "--- остерегаться. 




\noindent\textbf{Внятися} "--- загореться. 




\noindent\textbf{Воврещи} "--- бросить во что-либо; ввергнуть; внести. 




\noindent\textbf{Водрузити} "--- утвердить; укрепить. 




\noindent\textbf{Во еже} "--- чтобы; ради; для. 




\noindent\textbf{Вожделети} "--- сильно желать. 




\noindent\textbf{Возбеситися} "--- сделаться неистовым. 




\noindent\textbf{Возбранный, взбранный} "--- военный; храбрый в бранях; победоносный. 




\noindent\textbf{Возбраняти} "--- препятствовать; удерживать. 




\noindent\textbf{Возбряцати} "--- воспеть; хвалить в песнях. 




\noindent\textbf{Возвлачити} "--- затащить наверх. 




\noindent\textbf{Возврещи, возвергати} "--- возложить, возлагать. 




\noindent\textbf{Возглавие} "--- подушка; изголовье. 




\noindent\textbf{Возглас} "--- окончательные слова молитвы, тайно творимой священником. 




\noindent\textbf{Возглашение} "--- громкое пение или чтение; Возглас. 




\noindent\textbf{Возглядати} "--- взирать; смотреть. 




\noindent\textbf{Возгнещати} "--- разводить огонь. 




\noindent\textbf{Воздвигнути} "--- поднять; возвысить. 




\noindent\textbf{Воздвижение} "--- поднятие, возвышение. 




\noindent\textbf{Воздвизатися} "--- иногда: отправляться в путь. 




\noindent\textbf{Воздвизати} "--- поднимать; возвышать. 




\noindent\textbf{Воздеяние} "--- поднятие, возвышение. 




\noindent\textbf{Воздух} "--- покровец, полагаемый сверху священных сосудов на Литургии. 




\noindent\textbf{Возлежати} "--- лежать облокотившись; полулежать. 




\noindent\textbf{Возмущение} "--- смятение; бунт. 




\noindent\textbf{Возмятати} "--- возмущать; производить раздор. 




\noindent\textbf{Возниспослати} "--- послать свыше; наградить. 




\noindent\textbf{Возничати} "--- поднять голову. 




\noindent\textbf{Возобразитися} "--- принять образ; олицетвориться; вселиться в видимый образ; вновь быть изображенным. 




\noindent\textbf{Возраст} "--- возраст (число лет); рост. 




\noindent\textbf{Возрастити} "--- вырастить; увеличить. 




\noindent\textbf{Возреяти} "--- поколебать; потрясти. 




\noindent\textbf{Возставити} "--- восстановить; поставить на прежнее место. 




\noindent\textbf{Волити} "--- хотеть; желать; требовать. 




\noindent\textbf{Волна} "--- шерсть; руно; овчина. 




\noindent\textbf{Волхв} "--- мудрец; звездочет; чародей; предсказатель. 




\noindent\textbf{Волчец} "--- колючая трава. 




\noindent\textbf{Вонь} "--- в него. 




\noindent\textbf{Воня} "--- запах; курение. 




\noindent\textbf{Воня злая} "--- смрад. 




\noindent\textbf{Вопити, вопияти} "--- громко кричать; взывать. 




\noindent\textbf{Ворожити} "--- колдовать; предсказывать будущее. 




\noindent\textbf{Ворожея} "--- волшебник; колдун; отравитель. 




\noindent\textbf{Воскликновение} "--- хоровое пение. 




\noindent\textbf{Восклонятися} "--- выпрямляться; подниматься; разгибаться. 




\noindent\textbf{Воскресати} "--- восставать; оживать; возвращаться к жизни. 




\noindent\textbf{Воскресение} "--- восстание из мертвых. 




\noindent\textbf{Воскрешати} "--- оживлять. 




\noindent\textbf{Воскрилие} "--- подол; край одежды; пола верхней одежды. 




\noindent\textbf{Восперяти} "--- оперять; окрылять (надеждой). 




\noindent\textbf{Восписовати} "--- изображать письменно; изъявлять. 




\noindent\textbf{Воспросити} "--- попросить. 




\noindent\textbf{Воспрянути} "--- вскочить; вспрыгнуть; приходить в себя. 




\noindent\textbf{Востерзати} "--- извлекать; выдергивать. 




\noindent\textbf{Восторгати} "--- рвать; щипать; полоть. 




\noindent\textbf{Востягнути} "--- подтянуть; укрепить; подтащить. 




\noindent\textbf{Востязати, востязовати} "--- исследовать; испытывать; интересоваться. 




\noindent\textbf{Восхитити} "--- изловить; поймать; не законно захватить; похитить; увлечь в высоту; привести в восторг. 




\noindent\textbf{Вотще} "--- понапрасну; впустую; даром; тщетно. 




\noindent\textbf{Воутрие, воутрий} "--- на другой день. 




\noindent\textbf{Впасти} "--- упасть; попасть; ввалиться; подвергнуться; подпасть. 




\noindent\textbf{Вперити} "--- возвысить; поднять; устремить вверх наподобие пера. 




\noindent\textbf{Вперитися} "--- воспарить; взлететь. 




\noindent\textbf{Вперсити} "--- внутрь себя принять. 




\noindent\textbf{Впрямо} "--- прямо; напротив. 




\noindent\textbf{Врабий} "--- воробей. 




\noindent\textbf{Вран} "--- ворон. 




\noindent\textbf{Врата красныя} "--- западные церковные двери. 




\noindent\textbf{Вратарь, вратник} "--- сторож у ворот. 




\noindent\textbf{Врачба} "--- лекарство; врачевание. 




\noindent\textbf{Врачебница} "--- больница. 




\noindent\textbf{Вреды} "--- кожное заболевание. 




\noindent\textbf{Вресноту} "--- вправду; по достоинству; пристойно. 




\noindent\textbf{Вретище} "--- плохая, грубая одежда; дерюга; скорбное одеяние. 




\noindent\textbf{Вреяти} "--- кипеть; пениться; разгорячаться; бить ключом; выкипать. 




\noindent\textbf{Вреяти} "--- ввергать; вметать; вталкивать. 




\noindent\textbf{Всеблаголепный} "--- великолепный. 




\noindent\textbf{Всевидно} "--- всенародно. 




\noindent\textbf{Вседетельный} "--- все создавший. 




\noindent\textbf{Всеконечне} "--- совершенно. 




\noindent\textbf{Всекрасный} "--- самый красивый. 




\noindent\textbf{Вселукавый} "--- самый коварный, т.~е. диавол. 




\noindent\textbf{Всеоружие} "--- полное вооружение. 




\noindent\textbf{Всепетый} "--- препрославленный; всеми или всюду восхваляемый. 




\noindent\textbf{Всесожжение} "--- жертвоприношение, при котором жертва сжигалась целиком. 




\noindent\textbf{Всечреждение} "--- богатое угощение. 




\noindent\textbf{Всеядец} "--- тот, кто всех поедает, т.~е. ад, или смерть. 




\noindent\textbf{Всеяти} "--- посеять. 




\noindent\textbf{Вскрай} "--- по краю; близ; около. 




\noindent\textbf{Вскую} "--- почему? из-за чего? за что? 




\noindent\textbf{Вспять} "--- назад. 




\noindent\textbf{Всуе} "--- напрасно. 




\noindent\textbf{Всяко} "--- совершенно; совсем; вовсе. 




\noindent\textbf{Вторицею} "--- вторично; усиленно. 




\noindent\textbf{Вчиняти} "--- учреждать; узаконять. 




\noindent\textbf{Выну} "--- всегда. 




\noindent\textbf{Выспренний} "--- высокий; гордый. 




\noindent\textbf{Выспрь} "--- вверх. 




\noindent\textbf{Высоковыйный} "--- гордый; надменный; кичливый. 




\noindent\textbf{Высокомудрствовати} "--- высокоумствовать; гордиться. 




\noindent\textbf{Вышелетный} "--- предвечный. 




\noindent\textbf{Выя} "--- шея. 




\noindent\textbf{Вящший} "--- больший. 




\noindent\textbf{Вящше} "--- больше. 




\letter{Г} 





\noindent\textbf{Гаггрена} (читается «гангрена») "--- гангрена; антонов огонь; рак. 




\noindent\textbf{Гадание} "--- загадка; неясность. 




\noindent\textbf{Гади} (мн. ч.) "--- пресмыкающиеся. 




\noindent\textbf{Гаждение} "--- поношение; бесчестие; ругань. 




\noindent\textbf{Газофилакия} "--- казнохранилище во храме Иерусалимском. 




\noindent\textbf{Ганание} "--- загадка; притча. 




\noindent\textbf{Гастримаргия} "--- чревобесие; обжорство. 




\noindent\textbf{Гащи} "--- штаны; нижнее мужское белье. 




\noindent\textbf{Геенна} "--- долина Гинном около Иерусалима, где идолопоклонствующие иудеи при царе Ахазе сжигали своих детей в честь идола Молоха. Иносказательно: место будущих мучений, загробных наказаний. 




\noindent\textbf{Генварь} "--- январь. 




\noindent\textbf{Гибель} "--- трата; расход. 




\noindent\textbf{Главизна} "--- глава; начало; причина. 




\noindent\textbf{Главотяж} "--- головная повязка у иудеев. 




\noindent\textbf{Глагол} "--- слово; речь. 




\noindent\textbf{Глаголати} "--- говорить; рассказывать. 




\noindent\textbf{Глаголемый} "--- называемый; так называемый. 




\noindent\textbf{Глас} "--- голос; напев. 




\noindent\textbf{Глезна} "--- голень; ступня. 




\noindent\textbf{Глоба} "--- казнь; наказание. 




\noindent\textbf{Глумец} "--- кощун; пересмешник. 




\noindent\textbf{Глумилище} "--- место для скачек, плясок, маскарадов и т. п. 




\noindent\textbf{Глумитися} "--- забавляться; тешиться; получать удовольствие. 




\noindent\textbf{Глумы} (мн. ч.) "--- шутки; смех; игры. 




\noindent\textbf{Глядати} "--- смотреть; глядеть. 




\noindent\textbf{Гнати} "--- гнать; преследовать; идти; следовать за кем или чем-либо. 




\noindent\textbf{Гной, гноище} "--- навозная куча; раны. 




\noindent\textbf{Гнушатися} "--- считать гнусным; презирать. 




\noindent\textbf{Гобзование} "--- изобилие, довольство. 




\noindent\textbf{Гобзовати} "--- изобиловать; избыточествовать; быть богатым. 




\noindent\textbf{Гобзующий} "--- живущий в довольстве. 




\noindent\textbf{Говение} "--- почитание (например, поста). 




\noindent\textbf{Говети} "--- чтить; почитать (например, пост). 




\noindent\textbf{Говядо} "--- рогатый скот. 




\noindent\textbf{Годе} "--- угодно; приятно; подходяще. 




\noindent\textbf{Година, год} "--- час; время; пора. 




\noindent\textbf{Голоть} "--- гололедица, лед. 




\noindent\textbf{Гомола} "--- ком; комок; катыш; кусок. 




\noindent\textbf{Гонзати} "--- убегать; спасаться бегством. 




\noindent\textbf{Гонзнутие} "--- избежание. 




\noindent\textbf{Гонитель} "--- преследователь. 




\noindent\textbf{Горе} "--- ввысь; вверх. 




\noindent\textbf{Горее} "--- хуже; бедственнее. 




\noindent\textbf{Горлица, горличишь} "--- дикий голубь. 




\noindent\textbf{Горнец} "--- горшок; котелок; жаровня. 




\noindent\textbf{Горнило} "--- кузнечный горн; плавильня; место для плавки или очищения огнем. 




\noindent\textbf{Горница} "--- верхняя комната; столовая. 




\noindent\textbf{Горний} "--- высокий; вышний; небесный. 




\noindent\textbf{Горохищный} "--- пасущийся; блуждающий по горам; украденный диким зверем с горного пастбища. 




\noindent\textbf{Гортанобесие} "--- пристрастие к лакомствам. 




\noindent\textbf{Горушный} "--- горчичный. 




\noindent\textbf{Горший} "--- злейший; худший. 




\noindent\textbf{Господствия} "--- один из чинов ангельских. 




\noindent\textbf{Господыня} "--- госпожа. 




\noindent\textbf{Гостинник} "--- корчмарь; содержатель постоялого двора. 




\noindent\textbf{Градарь} "--- садовник; огородник. 




\noindent\textbf{Градеж} "--- оплот; забор. 




\noindent\textbf{Грезн} "--- гроздь виноградная. 




\noindent\textbf{Гривна} "--- ожерелье; носимая на шее цепь. 




\noindent\textbf{Гроздие} "--- кисть плодов; ветвь (винограда). 




\noindent\textbf{Грясти} "--- идти; шествовать. 




\noindent\textbf{Гугнивый} "--- гнусавый; заика; косноязычный; картавый; говорящий в нос. 




\noindent\textbf{Гудение} "--- игра на гуслях или арфе. 




\noindent\textbf{Гудец} "--- гуслист; музыкант. 




\letter{Д} 





\noindent\textbf{Да} "--- пусть; чтобы. 




\noindent\textbf{Дабы} "--- чтобы. 




\noindent\textbf{Далечайше} "--- гораздо далее. 




\noindent\textbf{Далечен} "--- далекий; трудный. 




\noindent\textbf{Даннословие} "--- обещание; обязательство. 




\noindent\textbf{Двакраты} "--- дважды. 




\noindent\textbf{Дващи} "--- дважды. 




\noindent\textbf{Дверь адова} "--- смерть. 




\noindent\textbf{Двизати} "--- двигать; шевелить. 




\noindent\textbf{Двоедушный} "--- нетвердый в вере. 




\noindent\textbf{Двоица} "--- пара. 




\noindent\textbf{Дворище} "--- небольшой или запустелый дом. 




\noindent\textbf{Дебельство} "--- тучность; полнота; дородность. 




\noindent\textbf{Дебрь} "--- долина; ложбина; овраг; ущелье. 




\noindent\textbf{Девствовати} "--- хранить девство, целомудрие. 




\noindent\textbf{Действо} "--- действие; представление. 




\noindent\textbf{Декемврий} "--- декабрь. 




\noindent\textbf{Делатилище} "--- купеческая лавка; орудие в чьих-либо руках. 




\noindent\textbf{Делва} "--- бочка; кадка. 




\noindent\textbf{Делма} "--- для. 




\noindent\textbf{Деля} "--- для; ради. 




\noindent\textbf{Демественник} "--- певчий. 




\noindent\textbf{Демоноговение} "--- почитание бесов. 




\noindent\textbf{Демоночтец} "--- идолопоклонник. 




\noindent\textbf{Денница} "--- утренняя заря, утренняя звезда; отпадший ангел. 




\noindent\textbf{Денносветлый} "--- подобный дневному свету. 




\noindent\textbf{Держава} "--- сила; крепость; власть; государство. 




\noindent\textbf{Державно} "--- властно; могущественно. 




\noindent\textbf{Дерзать} "--- осмеливаться; полагаться. 




\noindent\textbf{Дерзновение} "--- смелость. 




\noindent\textbf{Дерзословие} "--- наглая речь. 




\noindent\textbf{Дерзостник} "--- наглец; нахал. 




\noindent\textbf{Дерзый} "--- смелый; бесстыдный; дерзкий. 




\noindent\textbf{Десница} "--- правая рука. 




\noindent\textbf{Десный} "--- правый; находящийся с правой стороны. 




\noindent\textbf{Десятина} "--- десятая часть. 




\noindent\textbf{Десятословие} "--- десять заповедей Божиих, данных через Моисея. 




\noindent\textbf{Детищ} "--- младенец; дитя; отроча. 




\noindent\textbf{Детосаждение} "--- зачатие во утробе младенца. 




\noindent\textbf{Диадима} "--- венец; диадема. 




\noindent\textbf{Дивий} "--- дикий; лесной. 




\noindent\textbf{Дивьячитися} "--- зверствовать. 




\noindent\textbf{Дидрахма} "--- греч.двойная драхма", древнегреч. серебряная монета. 




\noindent\textbf{Динарий} "--- монета. 




\noindent\textbf{Длань} "--- ладонь. 




\noindent\textbf{Дмение} "--- гордость. 




\noindent\textbf{Дмитися} "--- гордиться; кичиться. 




\noindent\textbf{Дне} "--- относящийся к числу песнопений из Октоиха, а в дни пения Триоди "--- из этой книги. 




\noindent\textbf{Дненощно} "--- в течение целых суток. 




\noindent\textbf{Днесь} "--- сегодня, ныне; теперь. 




\noindent\textbf{Днешний} "--- нынешний; сегодняшний. 




\noindent\textbf{Доблий, добльственный, доблестный} "--- крепкий в добре; твердый в добродетели. 




\noindent\textbf{Доброзрачие} "--- красота; благообразие. 




\noindent\textbf{Доброкласный} "--- приносящий обильную жатву. 




\noindent\textbf{Добропобедный} "--- прославленный победами. 




\noindent\textbf{Доброта} "--- красота. 




\noindent\textbf{Доброхвальный} "--- заслуживающий похвалы; похвальный. 




\noindent\textbf{Довлесотворити} "--- удовлетворить. 




\noindent\textbf{Довлети} "--- доставать; быть достаточным; хватать. 




\noindent\textbf{Доволний} "--- достаточный. 




\noindent\textbf{Догмат} "--- греч. одно из основных положений веры. 




\noindent\textbf{Дождити} "--- посылать дождь; кропить; орошать. 




\noindent\textbf{Дозде} "--- доселе; до сего дня; досюда. 




\noindent\textbf{Дозела} "--- чрезвычайно. 




\noindent\textbf{Доилица} "--- кормилица; мамка. 




\noindent\textbf{Доити} "--- кормить грудью. 




\noindent\textbf{Доколе} "--- до какого времени? долго ли? 




\noindent\textbf{Долний} "--- нижний; земной (как противоп. «небесный, горний»). 




\noindent\textbf{Долу, доле} "--- внизу; вниз. 




\noindent\textbf{Долувлекущий} "--- тянущий вниз. 




\noindent\textbf{Дондеже} "--- пока. 




\noindent\textbf{Донележе} "--- пока. 




\noindent\textbf{Дориносити} "--- сопровождать кого-либо в качестве стражи, свиты. 




\noindent\textbf{Досаждение} "--- делание неугодного; нечестие; оскорбление. 




\noindent\textbf{Достижно} "--- понятно. 




\noindent\textbf{Достояние} "--- имение; наследство; власть. 




\noindent\textbf{Драхма} "--- древнегреч. серебряная монета. 




\noindent\textbf{Драчие} "--- сорная трава. 




\noindent\textbf{Древле} "--- давно. 




\noindent\textbf{Древодель} "--- плотник; столяр. 




\noindent\textbf{Дреколие} "--- колья. 




\noindent\textbf{Дрождие} "--- дрожжи; отстой. 




\noindent\textbf{Другиня} "--- подруга. 




\noindent\textbf{Дружина} "--- общество (товарищей, сверстников). 




\noindent\textbf{Дручити} "--- удручать; томить; изнурять. 




\noindent\textbf{Дряселовати} "--- быть пасмурным, мрачным, печальным. 




\noindent\textbf{Дряхлование} "--- печаль. 




\noindent\textbf{Дряхлый} "--- печальный. 




\noindent\textbf{Дска, дщица} "--- доска; дощечка. 




\noindent\textbf{Дуга} "--- радуга. 




\noindent\textbf{Дхнути} "--- дохнуть; дунуть. 




\noindent\textbf{Дщи, дщерь} "--- дочь. 




\letter{Е}





\noindent\textbf{Евнух} "--- скопец; сторож при гареме; придворный. 




\noindent\textbf{Егда} "--- когда. 




\noindent\textbf{Егов} "--- его (притяжательный падеж от местоимения «он»). 




\noindent\textbf{Еда} "--- разве? неужели? 




\noindent\textbf{Едем} "--- Эдем; рай земной. 




\noindent\textbf{Единако} "--- согласно; одинаково. 




\noindent\textbf{Единаче} "--- одинаково; равно; еще. 




\noindent\textbf{Единаче ли} "--- неужели еще? 




\noindent\textbf{Единовидный} "--- одновидный; однообразный. 




\noindent\textbf{Единою} "--- однажды. 




\noindent\textbf{Еже} "--- что; кое. 




\noindent\textbf{Езеро} "--- озеро. 




\noindent\textbf{Ей} "--- да; истинно; верно. 




\noindent\textbf{Ексапсалмы} "--- шестопсалмие. 




\noindent\textbf{Ектения} "--- усиленное моление; прошение. 




\noindent\textbf{Елей} "--- оливковое, деревянное масло. 




\noindent\textbf{Елень} "--- олень; лань. 




\noindent\textbf{Елеонский} "--- оливковый. 




\noindent\textbf{Елижды аще} "--- когда бы ни. 




\noindent\textbf{Елижды, еликожды} "--- всегда как; всякий раз, когда. 




\noindent\textbf{Еликий} "--- кто; который. 




\noindent\textbf{Елико} "--- сколько. 




\noindent\textbf{Елико-елико} "--- через короткое время; очень скоро. 




\noindent\textbf{Еликомощно} "--- по возможности; сколько дозволяют силы. 




\noindent\textbf{Еллин} "--- грек; язычник; прозелит иудаизма. 




\noindent\textbf{Елма} "--- поскольку; насколько. 




\noindent\textbf{Епендит} "--- верхнее платье. 




\noindent\textbf{Епистолия} "--- письмо; послание. 




\noindent\textbf{Еродий} "--- цапля. 




\noindent\textbf{Есмирнисменный} "--- смешанный вместе со смирной. 




\noindent\textbf{Ехидна} "--- ядовитая змея. 




\letter{Ж} 





\noindent\textbf{Жаждати} "--- хотеть пить; сильно желать чего-либо. 




\noindent\textbf{Жалость} "--- ревность; рвение. 




\noindent\textbf{Жатель} "--- жнец. 




\noindent\textbf{Жегомый} "--- тот, кого жгут огнем; больной огнем; больной огневицей, горячкой. 




\noindent\textbf{Жезл} "--- посох; трость; палка. 




\noindent\textbf{Женитва} "--- бракосочетание; супружество; брак. 




\noindent\textbf{Женонеистовый} "--- похотливый; блудный; сластолюбивый. 




\noindent\textbf{Жестоковыйный} "--- бесчувственный; упрямый. 




\noindent\textbf{Живити} "--- животворить; давать жизнь; оживлять. 




\noindent\textbf{Живодавец} "--- податель жизни. 




\noindent\textbf{Живоначалие} "--- начало; причина жизни. 




\noindent\textbf{Живот} "--- жизнь. 




\noindent\textbf{Животный} "--- живущий; одушевленный. 




\noindent\textbf{Жребя} "--- жеребенок. 




\noindent\textbf{Жрети} "--- заколать; приносить жертвоприношение. 




\noindent\textbf{Жупел} "--- горячая сера. 




\letter{З} 





\noindent\textbf{Забавати} "--- заговаривать; заколдовывать. 




\noindent\textbf{Забавление} "--- промедление; мешкание; ожидание. 




\noindent\textbf{Забавляти} "--- удерживать; замедлять. 




\noindent\textbf{Забобоны} "--- самовольная служба, бесчиние. 




\noindent\textbf{Забрало} "--- стена; забор. 




\noindent\textbf{Завет} "--- союз; договор; условие. 




\noindent\textbf{Завида} "--- зависть. 




\noindent\textbf{Завистно} "--- мало; недостаточно. 




\noindent\textbf{За еже} "--- для того, чтобы. 




\noindent\textbf{Заздати} "--- загородить. 




\noindent\textbf{Зазрети} "--- заглянуть; заметить; осудить; упрекнуть. 




\noindent\textbf{Заимование} "--- заем; долг. 




\noindent\textbf{Заимовати} "--- занимать; заимствовать. 




\noindent\textbf{Заклание} "--- жертвоприношение. 




\noindent\textbf{Заклеп} "--- запор; замок; задвижка. 




\noindent\textbf{Заколение} "--- жертвоприношение. 




\noindent\textbf{Законописец} "--- составитель законов. 




\noindent\textbf{Законополагати} "--- давать закон. 




\noindent\textbf{Закров} "--- место для укрытия. 




\noindent\textbf{Залещи} "--- быть в засаде; скрываться. 




\noindent\textbf{Заматорети} "--- устареть; зачерстветь; состариться. 




\noindent\textbf{Замреженый} "--- пойманный в сети. 




\noindent\textbf{Зане} "--- так как; потому что. 




\noindent\textbf{Занеже} "--- поскольку. 




\noindent\textbf{Зань} "--- за него. 




\noindent\textbf{Запаление} "--- загорание; пожар. 




\noindent\textbf{Запев} "--- краткий стих, предваряющий стихиры (на «Господи, воззвах», хвалитны, стиховны) или тропари канона. 




\noindent\textbf{Запечатствовати} "--- запечатать; утвердить; связать; скрепить. 




\noindent\textbf{Запинание} "--- враждебное действие. 




\noindent\textbf{Запойство} "--- пьянство. 




\noindent\textbf{Запона} "--- завеса. 




\noindent\textbf{Запрение} "--- отрицание; запирание. 




\noindent\textbf{Запретити} "--- запретить; опечалиться; скорбеть. 




\noindent\textbf{Запустение} "--- опустение; пустыня. 




\noindent\textbf{Запустети} "--- придти в запущение или запустение, запустеть. 




\noindent\textbf{Запяти, запнути} "--- остановить; задержать; обольститься. 




\noindent\textbf{Запятие} "--- препинание; препятствие; преткновение. 




\noindent\textbf{Заревидный} "--- подобный заре. 




\noindent\textbf{Зарелучный} "--- лучезарный. 




\noindent\textbf{Застояти} "--- останавливать на дороге; удерживать; наскучивать; утруждать. 




\noindent\textbf{За ся} "--- за себя. 




\noindent\textbf{Затвор} "--- замок; запор; место молитвенного подвига некоторых иноков, давших обет не исходить из своей келлии. 




\noindent\textbf{Заточаемый} "--- обуреваемый ветром; носимый; гонимый. 




\noindent\textbf{Затулити} "--- закрыть; спрятать; укрыть. 




\noindent\textbf{Затуне} "--- даром; без причины. 




\noindent\textbf{Зауститися} "--- закрыть уста; замолчать. 




\noindent\textbf{Заутра} "--- до восхода солнца; поутру; рано; завтра. 




\noindent\textbf{Заутрие} "--- завтрашний день. 




\noindent\textbf{Заушение} "--- пощечина; удар рукой по лицу. 




\noindent\textbf{Заушати} "--- заграждать уста; запрещать говорить. 




\noindent\textbf{Захленутися} "--- погрузиться. 




\noindent\textbf{Заходный} "--- западный. 




\noindent\textbf{Зачало} "--- начало; название отрезков текста в книгах Священного Писания Нового Завета. 




\noindent\textbf{Заяти} "--- взять взаймы; занять. 




\noindent\textbf{Звездоблюститель} "--- астроном. 




\noindent\textbf{Звездоволхвовати} "--- гадать по звездам; заниматься астрологией. 




\noindent\textbf{Звездозаконие} "--- астрономия. 




\noindent\textbf{Звездослов} "--- астролог. 




\noindent\textbf{Звездословие} "--- астрология. 




\noindent\textbf{Звездословити} "--- заниматься астрологией. 




\noindent\textbf{Звероядина} "--- скот, поврежденный хищным зверем. 




\noindent\textbf{Звиздание} "--- свист; посвист. 




\noindent\textbf{Звиздати} "--- свистеть. 




\noindent\textbf{Звонец} "--- колокольчик. 




\noindent\textbf{Звонница} "--- колокольня. 




\noindent\textbf{Звяцати} "--- звенеть; бренчать. 




\noindent\textbf{Здати} "--- строить. 




\noindent\textbf{Зде} "--- здесь. 




\noindent\textbf{Здо} "--- здание; стена; крыша. 




\noindent\textbf{Зелейник} "--- знахарь, лечащий травами и заговором. 




\noindent\textbf{Зелейничество} "--- напоение отравой. 




\noindent\textbf{Зелейный} "--- состоящий из зелия, т.~е. травы или других растений. 




\noindent\textbf{Зелие} "--- трава; растение. 




\noindent\textbf{Зело, зельне} "--- весьма; очень сильно. 




\noindent\textbf{Зельный} "--- сильный; великий. 




\noindent\textbf{Земен} "--- земной. 




\noindent\textbf{Земстий} "--- земной. 




\noindent\textbf{Зеница} "--- зрачок в глазе. 




\noindent\textbf{Зепь} "--- карман; мешок. 




\noindent\textbf{Зерцало} "--- зеркало. 




\noindent\textbf{Зиждитель} "--- создатель; творец. 




\noindent\textbf{Зиждити} "--- строить. 




\noindent\textbf{Зима} "--- зима; холод; плохая погода. 




\noindent\textbf{Злак} "--- растение; зелень; овощ. 




\noindent\textbf{Златарь} "--- золотых дел мастер. 




\noindent\textbf{Златица, златница} "--- золотая монета. 




\noindent\textbf{Златозарный} "--- яркоблестящий. 




\noindent\textbf{Злато} "--- золото. 




\noindent\textbf{Златокованный} "--- отчеканенный из золота. 




\noindent\textbf{Златокровный} "--- имеющий позлащенную крышу. 




\noindent\textbf{Злачный} "--- травный; богатый растительностью, злаками. 




\noindent\textbf{Зле} "--- зло; жестоко; худо. 




\noindent\textbf{Злоба} "--- забота. 




\noindent\textbf{Злокозненный} "--- исполненный злобы; лукавства. 




\noindent\textbf{Злокоман} "--- злодей; зложелатель; враг. 




\noindent\textbf{Злонравие} "--- развратный или дурной нрав. 




\noindent\textbf{Злообстояние} "--- беда ;несчастье. 




\noindent\textbf{Злопомнение} "--- злопамятство. 




\noindent\textbf{Злоречети} "--- бранить; ругать; злословить; поносить. 




\noindent\textbf{Злосердный} "--- безжалостный. 




\noindent\textbf{Злосмрадие} "--- зловоние. 




\noindent\textbf{Злосоветие} "--- злой умысел. 




\noindent\textbf{Злострастие} "--- сильные и порочные страсти. 




\noindent\textbf{Злостужати} "--- сильно досаждать. 




\noindent\textbf{Злостужный} "--- причиняющий большое беспокойство, мучение. 




\noindent\textbf{Злотечение} "--- развратные или злые поступки. 




\noindent\textbf{Злоумерший} "--- претерпевший тяжелую смерть. 




\noindent\textbf{Злоухищряти} "--- замышлять зло. 




\noindent\textbf{Злохитренный} "--- коварный. 




\noindent\textbf{Злохудожный} "--- лукавый; злобный; беззаконный. 




\noindent\textbf{Злый} "--- злой; плохой; негодный; худой; жестокий. 




\noindent\textbf{Знаемый} "--- знакомый, близкий человек. 




\noindent\textbf{Знаменательне} "--- прообразовательно. 




\noindent\textbf{Знаменательный} "--- прообразовательный; обозначающий нечто. 




\noindent\textbf{Знаменати} "--- обозначать знаком; помечать; изображать; показывать; являть. 




\noindent\textbf{Знамение} "--- знак; признак; явление; чудо. 




\noindent\textbf{Знаменоносец} "--- чудотворец. 




\noindent\textbf{Знаменоносный} "--- чудотворный. 




\noindent\textbf{Зобати} "--- наполнять зоб; клевать; есть; поглощать. 




\noindent\textbf{Зрак} "--- лицо; вид; образ. 




\noindent\textbf{Зрети} "--- смотреть. 




\noindent\textbf{Зрети к смерти} "--- находиться при последнем издыхании. 




\noindent\textbf{Зыбати} "--- шевелить; двигать; качать. 




\letter{И}





\noindent\textbf{И} "--- его. 




\noindent\textbf{Игемон} "--- вождь; начальник; правитель. 




\noindent\textbf{Иго} "--- ярмо; ноша. 




\noindent\textbf{Игралище} "--- место для представления. 




\noindent\textbf{Игрище} "--- смешное или непристойное представление. 




\noindent\textbf{Идеже} "--- где; когда. 




\noindent\textbf{Идолобесие} "--- неистовое идолопоклонство. 




\noindent\textbf{Иерей} "--- священник. 




\noindent\textbf{Иждивати} "--- проживать; тратить; издерживать. 




\noindent\textbf{Иже} "--- который. 




\noindent\textbf{Изблистати} "--- осиять; облистать; излить свет. 




\noindent\textbf{Избодати} "--- пропороть; поразить; пронзить; проколоть; выколоть. 




\noindent\textbf{Изборение} "--- поражение. 




\noindent\textbf{Изборати} "--- побеждать; поражать. 




\noindent\textbf{Избременяти} "--- облегчать; освобождать от бремени; выгружать. 




\noindent\textbf{Избутелый} "--- согнивший; испортившийся. 




\noindent\textbf{Избыти} "--- остаться в избытке, излишестве; изобиловать; освободиться. 




\noindent\textbf{Избыток} "--- довольство; изобилие. 




\noindent\textbf{Изваяние} "--- идол; кумир. 




\noindent\textbf{Извержение} "--- исключение из церковного клира или лишение сана. 




\noindent\textbf{Извесити} "--- свесить; вывесить. 




\noindent\textbf{Извествовати} "--- объявлять; оглашать; удостоверять. 




\noindent\textbf{Известно} "--- точно; тщательно. 




\noindent\textbf{Извет} "--- донос; извещение. 




\noindent\textbf{Извещен} "--- уверен. 




\noindent\textbf{Извещение} "--- удостоверение. 




\noindent\textbf{Извитие словес} "--- красноречие; витийство. 




\noindent\textbf{Извитийствовати} "--- красноречиво рассказать. 




\noindent\textbf{Извлачитися} "--- раздеться; разоблачиться. 




\noindent\textbf{Изволение} "--- воля; желание. 




\noindent\textbf{Изволити} "--- дозволять; захотеть; пожелать. 




\noindent\textbf{Извращати} "--- выворачивать; изменять; превращать. 




\noindent\textbf{Изврещи} "--- выбросить; вымести. 




\noindent\textbf{Извыкати} "--- научиться, познавать. 




\noindent\textbf{Изгвоздити} "--- выдернуть, вынуть гвозди. 




\noindent\textbf{Изгибающий} "--- погибающий, пропадающий. 




\noindent\textbf{Изгибнути} "--- погибнуть; пропасть; потеряться. 




\noindent\textbf{Изглаждати} "--- исключать; уничтожать. 




\noindent\textbf{Издетска} "--- с детства. 




\noindent\textbf{Издревле} "--- издавна; исстари. 




\noindent\textbf{Издручитися} "--- изнурить себя. 




\noindent\textbf{Изженяти} "--- изгонять; выгонять. 




\noindent\textbf{Излазити} "--- выходить; сходить (например, с корабля). 




\noindent\textbf{Излиха} "--- чрезмерно; еще более. 




\noindent\textbf{Излишше} "--- до излишества; паче меры. 




\noindent\textbf{Изляцати} "--- протягивать; простирать. 




\noindent\textbf{Измерети} "--- умереть. 




\noindent\textbf{Изметати} "--- извергать; выкидывать; выбрасывать. 




\noindent\textbf{Измена} "--- замена; перемена; выкуп. 




\noindent\textbf{Изменяти} "--- заменять; переменять. 




\noindent\textbf{Изменяти лице} "--- притворяться. 




\noindent\textbf{Измлада} "--- смолоду. 




\noindent\textbf{Измовение} "--- омытие; очищение. 




\noindent\textbf{Измолкати} "--- перестать говорить; замолкать. 




\noindent\textbf{Изнесение, изношение} "--- вынос. 




\noindent\textbf{Изницати} "--- возникать; появляться. 




\noindent\textbf{Износити} "--- выносить; произносить; производить; произращать; приносить. 




\noindent\textbf{Изнуждати} "--- выводить из нужды. 




\noindent\textbf{Изобнажати} "--- обнаруживать; являть; открывать. 




\noindent\textbf{Изостати} "--- остаться где-либо. 




\noindent\textbf{Изощряти} "--- обострить; наточить. 




\noindent\textbf{Изращение} "--- вырощение; произведение; порождение. 




\noindent\textbf{Изриновение} "--- выбрасывание; извержение; исключение. 




\noindent\textbf{Изриновенный} "--- изверженный; выкинутый; прогнанный. 




\noindent\textbf{Изринути} "--- столкнуть; опрокинуть; повалить; погубить. 




\noindent\textbf{Изрок} "--- изречение; осуждение. 




\noindent\textbf{Изрыти} "--- вырыть; выкопать. 




\noindent\textbf{Изрядно} "--- особенно; преимущественно. 




\noindent\textbf{Изсунути} "--- вынуть; исторгнуть; вырвать; изъять. 




\noindent\textbf{Изступление} "--- изумление; восторг. 




\noindent\textbf{Изуведети} "--- уразуметь; познать. 




\noindent\textbf{Изуздитися} "--- освободиться; получить волю. 




\noindent\textbf{Изумевати} "--- недоумевать; не понимать. 




\noindent\textbf{Изумителен} "--- буйствующий; беснующийся. 




\noindent\textbf{Изумитися} "--- сойти с ума; обезуметь. 




\noindent\textbf{Изути} "--- разуть; снять обувь. 




\noindent\textbf{Изчленити} "--- лишить членов; сокрушить члены; изуродовать. 




\noindent\textbf{Изъядати} "--- проматывать; растрачивать. 




\noindent\textbf{Иконом} "--- домоправитель. 




\noindent\textbf{Иконоратный} "--- иконоборственный. 




\noindent\textbf{Икос} "--- пространная песнь, написанная в похвалу святого или праздника. 




\noindent\textbf{Имати} "--- брать. 




\noindent\textbf{Иматисма} "--- верхнее платье, плащ. 




\noindent\textbf{Именный} "--- сокровищный; касающийся имения. 




\noindent\textbf{Имуществительно} "--- преимущественно. 




\noindent\textbf{Ин} "--- иной; другой. 




\noindent\textbf{Инамо} "--- в ином месте. 




\noindent\textbf{Иноковати} "--- жить по-иночески. 




\noindent\textbf{Инуде, инде} "--- в ином месте, в иное место. 




\noindent\textbf{Ипакои} "--- песнопение, положенное по малой ектении после полиелея на воскресной утрени. 




\noindent\textbf{Ипарх} "--- начальник области; градоначальник; наместник. 




\noindent\textbf{Ипостась} "--- лицо. 




\noindent\textbf{Ирмос} "--- песнопение, стоящее в начале каждой из песен канона. 




\noindent\textbf{Иродианы} "--- сторонники Ирода. 




\noindent\textbf{Ирой} "--- греч. миф. герой. 




\noindent\textbf{Исказити} "--- испортить; оскопить. 




\noindent\textbf{Искапати} "--- источать; испускать каплями; истечь. 




\noindent\textbf{Исковати} "--- выковать. 




\noindent\textbf{Искони} "--- изначала; вначале; всегда. 




\noindent\textbf{Исконный} "--- бывший искони; всегдашний. 




\noindent\textbf{Искренний} "--- ближний. 




\noindent\textbf{Искус} "--- испытание; искушение; проверка. 




\noindent\textbf{Исперва} "--- сначала; искони. 




\noindent\textbf{Исплевити} "--- выполоть; вырвать; исторгнуть; выдернуть; собрать. 




\noindent\textbf{Исплести} "--- сплести; сложить; составить. 




\noindent\textbf{Исповедатися} "--- признаваться; открыто выражать свою веру. 




\noindent\textbf{Исповедник} "--- человек, подвергавшийся страданиям или гонению за веру Христову. 




\noindent\textbf{Исполнение} "--- полнота; наполнение; совершение. 




\noindent\textbf{Исполнь} "--- наполненный; исполненный. 




\noindent\textbf{Исполняти} "--- наполнять; совершать. 




\noindent\textbf{Исполу} "--- вполовину; пополам; частию. 




\noindent\textbf{Исправити} "--- выпрямить; исправить; направить; укрепить. 




\noindent\textbf{Исправление} "--- восстановление; правый образ жизни. 




\noindent\textbf{Испразднити} "--- ниспровергнуть; уничтожить; умалить. 




\noindent\textbf{Испрати, исперити} "--- вытоптать; вымыть. 




\noindent\textbf{Испытно} "--- тщательно. 




\noindent\textbf{Испытовати} "--- выведывать. 




\noindent\textbf{Иссоп} "--- растение, употребляемое в пучках для кропления. 




\noindent\textbf{Истее} "--- точнее; яснее. 




\noindent\textbf{Истесы} "--- чресла, лядвеи. 




\noindent\textbf{Истицание} "--- истечение; истечение семени; поллюция. 




\noindent\textbf{Истаевати} "--- растаять; исчезать. 




\noindent\textbf{Исторгнути} "--- вырвать; вывести. 




\noindent\textbf{Истощание} "--- изнурение; унижение; снисхождение. 




\noindent\textbf{Истрезвлятися} "--- протрезвляться. 




\noindent\textbf{Истукан} "--- статуя; болван; идол. 




\noindent\textbf{Истый, истовый} "--- точный; подлинный; истинный. 




\noindent\textbf{Истязати} "--- вытягивать; получать; допрашивать. 




\noindent\textbf{Исходище} "--- место выхода; исток; начало. 




\noindent\textbf{Исходище вод(ное)} "--- ручей; поток; река. 




\noindent\textbf{Исходище путей} "--- распутье; перекресток. 




\noindent\textbf{Исчадие} "--- детище; плод; род; потомки. 




\noindent\textbf{Иулий} "--- июль. 




\noindent\textbf{Иуний} "--- июнь. 




\letter{К}





\noindent\textbf{Кадило} "--- возносимое во славу Божию благовонное курение. 




\noindent\textbf{Кадильница} "--- сосуд, в котором на горящие угли возлагается фимиам для совершения каждения. 




\noindent\textbf{Кадь} "--- кадка, ушат. 




\noindent\textbf{Каженик} "--- скопец; сторож при гареме; придворный. 




\noindent\textbf{Казатель} "--- учитель, наставник. 




\noindent\textbf{Казати} "--- наставлять; поучать. 




\noindent\textbf{Казити} "--- искажать; повреждать. 




\noindent\textbf{Како} "--- как. 




\noindent\textbf{Камара} "--- шатер; скиния; горница; покои. 




\noindent\textbf{Камо} "--- куда? 




\noindent\textbf{Кампан} "--- колокол. 




\noindent\textbf{Камы, камык} "--- камень. 




\noindent\textbf{Камык горящ} "--- сера. 




\noindent\textbf{Кандило} "--- лампада. 




\noindent\textbf{Кандиловжигатель} "--- пономарь. 




\noindent\textbf{Кандия} "--- небольшая чаша. 




\noindent\textbf{Капище} "--- идольский храм. 




\noindent\textbf{Катапетасма} "--- завеса. 




\noindent\textbf{Кафисма} "--- один из 20 разделов, на которые разделена Псалтирь. 




\noindent\textbf{Кацея} "--- кадильница не на цепочках, а на ручке. 




\noindent\textbf{Кацы} "--- каковые; которые; какие. 




\noindent\textbf{Квас} "--- закваска; дрожжи. 




\noindent\textbf{Квасный} "--- приготовленный на дрожжах. 




\noindent\textbf{Келарня, келарница} "--- помещение в монастыре для сохранения вещей, необходимых келарю. 




\noindent\textbf{Келарь} "--- старшая хозяйственная должность в монастыре. 




\noindent\textbf{Кивот} "--- ящик для икон. 




\noindent\textbf{Кидар} "--- головной убор ветхозаветного первосвященника. 




\noindent\textbf{Кимвал} "--- музыкальный инструмент. 




\noindent\textbf{Кимин} "--- тмин. 




\noindent\textbf{Киновия} "--- общежительный монастырь. 




\noindent\textbf{Кинсон} "--- дань; подать; ценз. 




\noindent\textbf{Кириопасха} "--- название праздника Пасхи, пришедшегося на день Благовещения Пресвятой Богородицы 25 марта. 




\noindent\textbf{Кичение} "--- гордость. 




\noindent\textbf{Клада} "--- колода (орудие пытки). 




\noindent\textbf{Кладенец} "--- яма; клад. 




\noindent\textbf{Кладязь} "--- колодец. 




\noindent\textbf{Клас} "--- колос. 




\noindent\textbf{Клеврет} "--- товарищ; собрат. 




\noindent\textbf{Клепало} "--- колотушка, при помощи которой в монастырях созывают на молитву. 




\noindent\textbf{Клепати} "--- звонить; стучать или бить в клепало. 




\noindent\textbf{Клеть} "--- изба; покои; кладовая; комната. 




\noindent\textbf{Клирос} "--- возвышение в храме, на котором располагаются певчие. 




\noindent\textbf{Клич} "--- крик, гам. 




\noindent\textbf{Клобук} "--- покрывало, носимое монашествующими поверх камилавки. 




\noindent\textbf{Ключимый} "--- годный; хороший; случившийся кстати; полезный. 




\noindent\textbf{Ключитися} "--- приключиться; случиться. 




\noindent\textbf{Книгочий} "--- судья; приставник. 




\noindent\textbf{Книжник} "--- ученый. 




\noindent\textbf{Ков} "--- умысел; заговор. 




\noindent\textbf{Ковчег} "--- кованый ящик: сундук; ларец. 




\noindent\textbf{Кодрант} "--- мелкая римская монета. 




\noindent\textbf{Козлогласование} "--- бесчинные крики на пиршестве. 




\noindent\textbf{Козни} "--- лукавство; хитрость. 




\noindent\textbf{Кокош} "--- наседка. 




\noindent\textbf{Колено} "--- род; поколение. 




\noindent\textbf{Колесницегонитель} "--- возница; преследователь на колеснице. 




\noindent\textbf{Коливо} "--- вареная пшеница с медом, приносимая для благословения в церковь на праздники. 




\noindent\textbf{Колиждо} "--- когда; как. 




\noindent\textbf{Колико} "--- сколько. 




\noindent\textbf{Колия} "--- яма; ров. 




\noindent\textbf{Колми} "--- сколько. 




\noindent\textbf{Колми паче} "--- тем более; особенно. 




\noindent\textbf{Коло} "--- колесо. 




\noindent\textbf{Колобродити} "--- ходить вокруг; уклоняться. 




\noindent\textbf{Коль} "--- сколько; насколько; как. 




\noindent\textbf{Колькраты} "--- сколько раз; как часто. 




\noindent\textbf{Комбоста} "--- сырая капуста. 




\noindent\textbf{Кондак} "--- короткая песнь в честь святого или праздника. 




\noindent\textbf{Коноб} "--- котел; горшок; умывальница. 




\noindent\textbf{Конура} "--- небольшой мешочек, носимый суеверными людьми вместе с кореньями или другими амулетами. 




\noindent\textbf{Копр} "--- укроп; анис. 




\noindent\textbf{Кораблец} "--- небольшой корабль. 




\noindent\textbf{Корван} "--- дар; жертва Богу. 




\noindent\textbf{Корвана} "--- казнохранилище при храме Иерусалимском. 




\noindent\textbf{Кормило} "--- руль. 




\noindent\textbf{Кормильствовати} "--- править; руководить. 




\noindent\textbf{Кормление} "--- правление; управление. 




\noindent\textbf{Корчаг} "--- лохань. 




\noindent\textbf{Корчемница} "--- корчма;кабак. 




\noindent\textbf{Коснити} "--- медлить. 




\noindent\textbf{Косноязычный} "--- медленноязычный; заика. 




\noindent\textbf{Косный} "--- медленный; нерешительный; упорно остающийся в одном и том же состоянии. 




\noindent\textbf{Котва} "--- якорь. 




\noindent\textbf{Кош} "--- кошель; корзина. 




\noindent\textbf{Кошница} "--- кошель, корзина. 




\noindent\textbf{Кощунник} "--- шут, балагур. 




\noindent\textbf{Кощунница} "--- актриса; танцовщица. 




\noindent\textbf{Кощуны} "--- смехотворство. 




\noindent\textbf{Крабица} "--- коробочка; ящичек; ковчежец; ларчик. 




\noindent\textbf{Крава} "--- корова. 




\noindent\textbf{Краегранесие, краестрочие} "--- акростих, т.~е. поэтическое произведение, в котором начальные буквы каждой строчки составляют слово, фразу или следуют порядку алфавита. 




\noindent\textbf{Крамола} "--- смута; заговор; бунт. 




\noindent\textbf{Красный} "--- красивый; прекрасный; непорочный. 




\noindent\textbf{Красовул} "--- мерная чаша в монастырях, вмещающая более 200 г . 




\noindent\textbf{Крастель} "--- перепел. 




\noindent\textbf{Крата} "--- раз. 




\noindent\textbf{Крепкий} "--- сильный; крепкий. 




\noindent\textbf{Креплий} "--- крепчайший, сильнейший. 




\noindent\textbf{Кресати} "--- извлекать; высекать огонь; оживлять. 




\noindent\textbf{Крин} "--- лилия. 




\noindent\textbf{Кроме} "--- вне; извне; отдельно; кроме. 




\noindent\textbf{Кромешный} "--- внешний; запредельный; отдаленный; лишенный. 




\noindent\textbf{Кропило} "--- кисть для окропления освященной водой. 




\noindent\textbf{Ктитор} "--- создатель; строитель или снабдитель храма или монастыря; церковный староста. 




\noindent\textbf{Ктому} "--- впредь; затем; еще; уже; более. 




\noindent\textbf{Купа} "--- кипа; груда; куча; ворох. 




\noindent\textbf{Купель} "--- озеро; пруд; садок; сосуд для совершения Таинства Крещения. 




\noindent\textbf{Купина} "--- соединение нескольких однородных предметов: куст, сноп; терновый куст. 




\noindent\textbf{Купно} "--- вместе. 




\noindent\textbf{Купный} "--- совместный. 




\noindent\textbf{Кустодия} "--- стража; караул; охрана; 




\noindent\textbf{Кутия} "--- вареная пшеница с медом, приносимая в церковь на поминовение усопших христиан. 




\noindent\textbf{Куща} "--- шатер; палатка; шалаш. 




\noindent\textbf{Кущник} "--- человек, делающий палатки или живущий в шалаше. 




\letter{Л}





\noindent\textbf{Ладия} "--- небольшое судно; кораблик; ладья. 




\noindent\textbf{Ладан} "--- благоуханная смола, влагаемая в кадильницу на горящие угли для благовонного курения. 




\noindent\textbf{Лазарома} "--- гробная одежда; повой; плащаница, в которую повивали усопших у иудеев. 




\noindent\textbf{Лазня} "--- баня. 




\noindent\textbf{Лай} "--- хула; поношение. 




\noindent\textbf{Лакать} "--- евр. мера длины. 




\noindent\textbf{Ланита} "--- щека. 




\noindent\textbf{Лаятель} "--- ругатель; хулитель; седящий в засаде. 




\noindent\textbf{Лвичищ} "--- львенок. 




\noindent\textbf{Левиафан} "--- крокодил. 




\noindent\textbf{Легеон} "--- полк; толпа; множество. 




\noindent\textbf{Лежание} "--- лежание; опочивание. 




\noindent\textbf{Лемаргия} "--- гортанобесие, т.~е. гурманство. 




\noindent\textbf{Лентион, лентий} "--- полотенце. 




\noindent\textbf{Лепо} "--- красиво. 




\noindent\textbf{Лепоподобно} "--- благопристойно; по достоинству. 




\noindent\textbf{Лепота} "--- красота; изящество. 




\noindent\textbf{Лепта} "--- мелкая монетка. 




\noindent\textbf{Лествица} "--- лестница. 




\noindent\textbf{Лестчий} "--- льстивый, ложный. 




\noindent\textbf{Лесть} "--- обман; хитрость; коварность. 




\noindent\textbf{Лето} "--- год; время. 




\noindent\textbf{Леторасль} "--- выросшее за год, годовой побег дерева. 




\noindent\textbf{Леть} "--- льзя; можно. 




\noindent\textbf{Леха} "--- гряда, ряд. 




\noindent\textbf{Лечба} "--- лекарство; врачевство. 




\noindent\textbf{Лечец} "--- лекарь; врач. 




\noindent\textbf{Лжа} "--- ложь. 




\noindent\textbf{Лжесловесие} "--- лживые речи. 




\noindent\textbf{Лив} "--- полдень; юг; юго-западный ветер. 




\noindent\textbf{Ливан} "--- иногда значит то же, что и Ладан. 




\noindent\textbf{Лик} "--- собрание; хор. 




\noindent\textbf{Ликование} "--- многолюдное пение; пляска; танцы. 




\noindent\textbf{Ликоватися} "--- приветствовать чрез соприкосновение правой щекой. 




\noindent\textbf{Ликовне} "--- с ликованием. 




\noindent\textbf{Ликостояние} "--- бдение на молитве церковной. 




\noindent\textbf{Лития} "--- исхождение из церкви на молитву. 




\noindent\textbf{Литра} "--- мера веса. 




\noindent\textbf{Литургисати} "--- совершать Литургию. 




\noindent\textbf{Лихва} "--- прибыль; проценты. 




\noindent\textbf{Лихоимец} "--- ростовщик; сребролюбец. 




\noindent\textbf{Лице} "--- лицо; вид; человек. 




\noindent\textbf{Личина} "--- маскарадная или шутовская маска. 




\noindent\textbf{Лишатися} "--- нуждаться. 




\noindent\textbf{Лишше} "--- больше, сверх того. 




\noindent\textbf{Лобзание} "--- устное целование. 




\noindent\textbf{Ловитва} "--- ловля; охота; сети; добыча; грабеж. 




\noindent\textbf{Ловительство} "--- засада, ловушка. 




\noindent\textbf{Ложе} "--- постель, одр. 




\noindent\textbf{Ложесна} "--- утроба женщины. 




\noindent\textbf{Лоза} "--- виноград. 




\noindent\textbf{Ломимый} "--- преломляемый. 




\noindent\textbf{Лоно} "--- пазуха; грудь; колени. 




\noindent\textbf{Луновение} "--- месячный цикл у женщин. 




\noindent\textbf{Лысто} "--- голень; икры; лытка. 




\noindent\textbf{Льстивый} "--- обманчивый. 




\noindent\textbf{Льщение} "--- обман; коварство; лесть. 




\noindent\textbf{Любо} "--- либо, или. 




\noindent\textbf{Любомудрие} "--- философия. 




\noindent\textbf{Любомятежный} "--- склонный к мятежу. 




\noindent\textbf{Любоначалие} "--- властолюбие. 




\noindent\textbf{Любопразднственный} "--- любящий празднствовать. 




\noindent\textbf{Любопрение} "--- любовь состязаться, спорить. 




\noindent\textbf{Любосластие} "--- сластолюбие; любовь к плотским утехам. 




\noindent\textbf{Любочестие} "--- почитание; чествование. 




\noindent\textbf{Любочестный} "--- достойный похвалы, чести. 




\noindent\textbf{Любы} "--- любовь. 




\noindent\textbf{Люте} "--- жестоко; тяжко. 




\noindent\textbf{Лютый} "--- свирепый; жестокий; злой; мучительный. 




\noindent\textbf{Лядвея} "--- ляжка; верхняя половина ноги; промежность. 




\noindent\textbf{Лярва} "--- маска; личина. 




\letter{М}





\noindent\textbf{Маание} "--- знак рукой, головою, глазами или иного рода, содержащий приказание; повеление; воля. 




\noindent\textbf{Маий} "--- май. 




\noindent\textbf{Малакия} "--- грех рукоблудия. 




\noindent\textbf{Малимый} "--- уменьшаемый. 




\noindent\textbf{Малобрещи} "--- нерадеть о чем-либо. 




\noindent\textbf{Мамона} "--- богатство; имение. 




\noindent\textbf{Мандра} "--- ограда. 




\noindent\textbf{Мание} "--- знак рукой, головою, глазами или иного рода, содержащий приказание; повеление; воля. 




\noindent\textbf{Манна} "--- небесный хлеб, данный израильтянам в пустыне. 




\noindent\textbf{Манноприемный} "--- содержащий манну. 




\noindent\textbf{Маслина} "--- олива; оливковое дерево. 




\noindent\textbf{Масличный} "--- оливковый. 




\noindent\textbf{Мастити} "--- намазывать. 




\noindent\textbf{Маститый} "--- обильный; тучный; заслуженный. 




\noindent\textbf{Масть} "--- мазь; масло. 




\noindent\textbf{Матеродевственный} "--- одновременно относящийся и к матери, и к деве. 




\noindent\textbf{Матеролепне} "--- по-матерински. 




\noindent\textbf{Матерский} "--- материнский. 




\noindent\textbf{Матерь градовом} "--- столица; первопрестольный град. 




\noindent\textbf{Мгляный} "--- окруженный или покрытый мглой. 




\noindent\textbf{Медвен} "--- медовый. 




\noindent\textbf{Медленоязычный} "--- косноязычный; заика. 




\noindent\textbf{Медница} "--- медная монетка. 




\noindent\textbf{Медовина} "--- вареный мед с хмелем. 




\noindent\textbf{Медоточный} "--- источающий, изливающий мед. 




\noindent\textbf{Медоязычный} "--- сладкословесный. 




\noindent\textbf{Междорамие} "--- пространство между плечами. 




\noindent\textbf{Мездник} "--- наемник. 




\noindent\textbf{Мерзость} "--- скверна; гнусность; беззаконие;нечестие; иногда "--- идол. 




\noindent\textbf{Мерило} "--- мера; весы. 




\noindent\textbf{Меск} "--- полуосел; мул; лошак. 




\noindent\textbf{Мессия} "--- евр. помазанник. 




\noindent\textbf{Метание} "--- поясной поклон. 




\noindent\textbf{Мех} "--- кожаный мешок для сохранения и перевоза жидкостей. 




\noindent\textbf{Мжа} "--- мигание; прищур. 




\noindent\textbf{Мжати} "--- жмурить глаза; щуриться; плохо видеть. 




\noindent\textbf{Мзда} "--- награда; плата. 




\noindent\textbf{Мздовоздаятель} "--- оплачивающий работу, дающий награду. 




\noindent\textbf{Мздоимание} "--- взяточничество. 




\noindent\textbf{Мила ся деяти} "--- низко припадать к земле; просить сжалиться над собой. 




\noindent\textbf{Милоть} "--- овчина; грубый шерстяной плащ из овечьей шерсти. 




\noindent\textbf{Милый} "--- жалкий; заслуживающий сожаления. 




\noindent\textbf{Мимотещи} "--- идти, проходить мимо, не останавливаясь. 




\noindent\textbf{Мирная} "--- название великой ектении. 




\noindent\textbf{Миро} "--- благовонная жидкость или мазь. 




\noindent\textbf{Мироподательне} "--- подавая мир. 




\noindent\textbf{Мироточец} "--- источающий чудотворное миро. 




\noindent\textbf{Мироявленный} "--- явленный, открытый миру. 




\noindent\textbf{Мирсина} "--- название красивого дерева. 




\noindent\textbf{Младодеяти, младодействовати} "--- принимать образ младенца; облекаться в плоть. 




\noindent\textbf{Младоумие} "--- незрелость ума. 




\noindent\textbf{Млат} "--- молот. 




\noindent\textbf{Млеко} "--- молоко. 




\noindent\textbf{Мнас} "--- мина, древнегреч. серебряная монета. 




\noindent\textbf{Мнее} "--- менее. 




\noindent\textbf{Мнети, мнити} "--- думать; предполагать; казаться. 




\noindent\textbf{Мний} "--- меньший. 




\noindent\textbf{Мних} "--- монах. 




\noindent\textbf{Многажды, множицею} "--- часто; много раз. 




\noindent\textbf{Многобезсловесие} "--- невежество. 




\noindent\textbf{Многобогатый} "--- изобилующий во всем. 




\noindent\textbf{Многоболезненный} "--- подъявший многие труды, подвиги, беды, страдания. 




\noindent\textbf{Многоборимый} "--- подвергаемый сильным искушениям, нападениям. 




\noindent\textbf{Многобурный} "--- тревожный. 




\noindent\textbf{Многогобзенный} "--- весьма обильный. 




\noindent\textbf{Многогубо} "--- многократно. 




\noindent\textbf{Многокласный} "--- колосистый. 




\noindent\textbf{Многомятущий} "--- преисполненный суетою. 




\noindent\textbf{Многонарочитый} "--- весьма знаменитый. 




\noindent\textbf{Многообразне} "--- во многих видах; различно. 




\noindent\textbf{Многооранный} "--- многократно возделанный. 




\noindent\textbf{Многоочитый} "--- имеющий множество глаз. 




\noindent\textbf{Многоплодие} "--- плодоносие; многочадие. 




\noindent\textbf{Многоплотие} "--- тучность. 




\noindent\textbf{Многопрелестный} "--- исполненный прелестей и соблазнов. 




\noindent\textbf{Многосветлый} "--- радостный; торжественный. 




\noindent\textbf{Многослезный} "--- исполненный печали и горя. 




\noindent\textbf{Многоснедный} "--- изобилующий многообразием пищи. 




\noindent\textbf{Многосугубый} "--- усугубленный; умноженный; усиленный. 




\noindent\textbf{Многосуетный} "--- совершенно пустой, бесполезный. 




\noindent\textbf{Многоуветливый} "--- очень снисходительный. 




\noindent\textbf{Многоцелебный} "--- подающий многие исцеления. 




\noindent\textbf{Многочастне} "--- много раз. 




\noindent\textbf{Многочудесный} "--- источающий многие чудеса; прославленный чудотворениями. 




\noindent\textbf{Многоязычный} "--- состоящий из множества племен. 




\noindent\textbf{Молва} "--- говор; ропот; слух; забота; волнение. 




\noindent\textbf{Молвити} "--- заботиться; суетиться; волноваться; роптать. 




\noindent\textbf{Молие} "--- моль. 




\noindent\textbf{Молниезрачный} "--- напоминающий молнию. 




\noindent\textbf{Мочащийся к стене} "--- пес. 




\noindent\textbf{Мощи} "--- нетленное тело угодника Божия. 




\noindent\textbf{Мравий} "--- муравей. 




\noindent\textbf{Мраз} "--- мороз. 




\noindent\textbf{Мрежа} "--- рыболовная сеть. 




\noindent\textbf{Мужатая} "--- замужняя. 




\noindent\textbf{Мужатица} "--- замужняя женщина. 




\noindent\textbf{Муженеискусная} "--- не познавшая мужа; не причастная браку. 




\noindent\textbf{Мурин} "--- эфиоп; арап; негр; чернокожий; дух тьмы; бес. 




\noindent\textbf{Мусийский, мусикийский} "--- музыкальный. 




\noindent\textbf{Мусикия} "--- музыка. 




\noindent\textbf{Мшела} "--- взятка. 




\noindent\textbf{Мшелоимство} "--- корыстолюбие. 




\noindent\textbf{Мшица} "--- мошка; мошкара. 




\noindent\textbf{Мытарь} "--- сборщик подати. 




\noindent\textbf{Мытница} "--- таможня; дом или двор для сбора пошлин. 




\noindent\textbf{Мыто} "--- пошлина; сбор; налог. 




\noindent\textbf{Мышца} "--- рука; плечо; сила. 




\noindent\textbf{Мясопуст} "--- последний день вкушения мясной пищи. 




\noindent\textbf{Мясоястие, мясоед} "--- время, когда Устав разрешает вкушение мяса. 




\noindent\textbf{Мятва} "--- мята. 




\letter{Н}





\noindent\textbf{Набдевати} "--- снабжать; наделять; хранить. 




\noindent\textbf{Наваждати} "--- научать; подстрекать. 




\noindent\textbf{Навет} "--- наговор; клевета; козни. 




\noindent\textbf{Навклир} "--- хозяин корабля. 




\noindent\textbf{Навыкнути} "--- приучиться; привыкнуть. 




\noindent\textbf{Наготовати} "--- ходить без одежды. 




\noindent\textbf{Нагствовати} "--- см. Наготовати. 




\noindent\textbf{Надходити} "--- внезапно постигнуть, случиться. 




\noindent\textbf{Наздати} "--- надстроить; укрепить; утвердить. 




\noindent\textbf{Назирати} "--- примечать; наблюдать. 




\noindent\textbf{Назнаменовати} "--- назначать; обозначать; осенять Крестом. 




\noindent\textbf{Наипаче} "--- особенно; преимущественно. 




\noindent\textbf{Наитие} "--- нисшествие; нашествие; сошествие. 




\noindent\textbf{Наказание} "--- иногда: учение. 




\noindent\textbf{Наляцати} "--- натянуть. 




\noindent\textbf{На мале} "--- малое время; дешево. 




\noindent\textbf{Намащати} "--- намазывать; втирать. 




\noindent\textbf{На мнозе} "--- на долгое время; дорого. 




\noindent\textbf{Наопак} "--- наоборот; вопреки. 




\noindent\textbf{Напаствуемый} "--- находящийся в напасти. 




\noindent\textbf{Наперсник} "--- друг, доверенное лицо. 




\noindent\textbf{Напоследок} "--- недавно. 




\noindent\textbf{Нард} "--- колосистое ароматическое растение. 




\noindent\textbf{Нарекованный} "--- предопределенный; предуставленный; назначенный. 




\noindent\textbf{Нарицати} "--- называть. 




\noindent\textbf{Нарок} "--- определенное или назначенное время. 




\noindent\textbf{Нарочитый} "--- особый; славный. 




\noindent\textbf{Наругатися} "--- насмеяться; пренебречь; опозорить. 




\noindent\textbf{Насмертник} "--- осужденный на смерть 




\noindent\textbf{Насущный} "--- настоящий; нынешний; существенный; необходимый. 




\noindent\textbf{На толице} "--- в такое время; за такую цену, за столько. 




\noindent\textbf{Началозлобный} "--- виновник зла. 




\noindent\textbf{Начаток} "--- начало; первый плод. 




\noindent\textbf{Начертавати} "--- изобразить. 




\noindent\textbf{Наясне} "--- наружу; открыто. 




\noindent\textbf{Наяти} "--- нанять. 




\noindent\textbf{Неблазненный} "--- безопасный; непогрешимый. 




\noindent\textbf{Неблазный} "--- непрельщаемый. 




\noindent\textbf{Небрещи} "--- нерадеть; пренебрегать. 




\noindent\textbf{Невеглас} "--- невежда; простак; неученый. 




\noindent\textbf{Невеститель} "--- снабжающий бедных невест приданым. 




\noindent\textbf{Невестоукрасити} "--- украсить как невесту. 




\noindent\textbf{Невечерний} "--- непомрачаемый; светлый. 




\noindent\textbf{Невиновный} "--- беспричинный; самобытный. 




\noindent\textbf{Невозбранно} "--- беспрепятственно. 




\noindent\textbf{Невозносительно} "--- смиренно. 




\noindent\textbf{Негли} "--- неужели; может быть; авось. 




\noindent\textbf{Неделя} "--- церковное название воскресного дня. 




\noindent\textbf{Недремлющий} "--- неусыпный. 




\noindent\textbf{Недристый} "--- имеющий широкую грудь. 




\noindent\textbf{Недро} "--- нутро; утроба; грудь; внутренность; залив. 




\noindent\textbf{Недуг} "--- болезнь. 




\noindent\textbf{Неже} "--- нежели; чем. 




\noindent\textbf{Независтный} "--- неиспорченный; невредимый; довольный; обильный. 




\noindent\textbf{Неиждиваемый} "--- не могущий быть истрачен или использован до конца. 




\noindent\textbf{Неизводимый} "--- непрекращаемый. 




\noindent\textbf{Неизгиблемый} "--- не подлежащий тлению или времени. 




\noindent\textbf{Неизреченный} "--- невыразимый. 




\noindent\textbf{Неискусобрачный} "--- не испытавший брака. 




\noindent\textbf{Неискусомужная} "--- не познавшая мужа. 




\noindent\textbf{Неиспытный} "--- сокровенный; тайный. 




\noindent\textbf{Неистовно} "--- с ожесточением; с яростию. 




\noindent\textbf{Неистовый} "--- вышедший из себя; находящийся не в должном состоянии. 




\noindent\textbf{Неисследимый} "--- непостижимый. 




\noindent\textbf{Неключимый} "--- бесполезный; негодный. 




\noindent\textbf{Некосненно} "--- немедленно. 




\noindent\textbf{Не ктому} "--- более не; еще не; уже не. 




\noindent\textbf{Нелестный} "--- необманчивый; нелукавый. 




\noindent\textbf{Нелеть} "--- нельзя. 




\noindent\textbf{Неможение} "--- болезнь; немощь; бессилие. 




\noindent\textbf{Немокренно} "--- по суху. 




\noindent\textbf{Немощствующий} "--- больной. 




\noindent\textbf{Необименный} "--- необъятный. 




\noindent\textbf{Не обинутися} "--- поступать смело. 




\noindent\textbf{Необинуяся} "--- смело; дерзновенно. 




\noindent\textbf{Неопальный} "--- несгораемый. 




\noindent\textbf{Неописанне} "--- изобразимо. 




\noindent\textbf{Неопределенный} "--- беспредельный. 




\noindent\textbf{Неоранный} "--- непаханный; невозделанный; нетронутый. 




\noindent\textbf{Неотметный} "--- неотчужденный. 




\noindent\textbf{Неплоды, неплодовь} "--- бесплодная женщина. 




\noindent\textbf{Неподобный} "--- непристойный. 




\noindent\textbf{Непорочны} "--- название 17-й кафизмы псалма 118. 




\noindent\textbf{Непорочный} "--- беспорочный; святой; чистый. 




\noindent\textbf{Неправдовати} "--- поступать нечестиво. 




\noindent\textbf{Непраздная} "--- беременная. 




\noindent\textbf{Непревратный} "--- непременный; неизменяемый. 




\noindent\textbf{Непреложно} "--- неизменно; без изменения. 




\noindent\textbf{Непременный} "--- неизменяемый. 




\noindent\textbf{Непреоборимый} "--- неодолимый; непобедимый. 




\noindent\textbf{Непщевание} "--- мнение; подлог; выдумка. 




\noindent\textbf{Непщевати} "--- думать; придумывать; считать. 




\noindent\textbf{Неразседный} "--- неразрушаемый. 




\noindent\textbf{Нерешимый} "--- несокрушимый; неразвязываемый. 




\noindent\textbf{Неседальное} "--- церковная служба, во время которой возбраняется сидеть. 




\noindent\textbf{Несланый} "--- несоленый. 




\noindent\textbf{Неслиянне} "--- неслитно. 




\noindent\textbf{Несмесне} "--- не смешиваясь. 




\noindent\textbf{Нестареемый} "--- вечный; неизменный. 




\noindent\textbf{Несть} "--- нет. 




\noindent\textbf{Нестояние} "--- непостоянство; смущение. 




\noindent\textbf{Несумненный} "--- несомненный; надежный; беспристрастный. 




\noindent\textbf{Несущий} "--- не имеющий бытия. 




\noindent\textbf{Нетление} "--- неуничтожимость; вечность; несокрушимость. 




\noindent\textbf{Нетребе} "--- не нужно. 




\noindent\textbf{Нетреный} "--- непротертый; непроходимый. 




\noindent\textbf{Нетесноместно} "--- удобовместительно. 




\noindent\textbf{Нетяжестне} "--- без труда. 




\noindent\textbf{Не у} "--- еще не. 




\noindent\textbf{Неудобоприятный} "--- невместимый; непонятный; непостижимый. 




\noindent\textbf{Неудобь} "--- неудобно; трудно. 




\noindent\textbf{Неумытный} "--- неподкупный. 




\noindent\textbf{Нечаяние} "--- неожиданность; беспечность. 




\noindent\textbf{Неясыть} "--- пеликан. 




\noindent\textbf{Ниже} "--- тем более не…; ни даже…; и не… 




\noindent\textbf{Николиже} "--- никогда. 




\noindent\textbf{Ни ли} "--- разве не? неужели? или не? 




\noindent\textbf{Ниц} "--- вниз; лицем на землю. 




\noindent\textbf{Нищетный} "--- нищенский; униженный; бедный. 




\noindent\textbf{Новемврий} "--- ноябрь. 




\noindent\textbf{Новина} "--- новость. 




\noindent\textbf{Новозданный} "--- вновь построенный. 




\noindent\textbf{Новопросвещенный} "--- недавно крещеный. 




\noindent\textbf{Новосаждение} "--- почки; отпрыски 




\noindent\textbf{Ноемврий} "--- ноябрь. 




\noindent\textbf{Ножница} "--- ножны. 




\noindent\textbf{Нощный вран} "--- филин; сова. 




\noindent\textbf{Нудитися} "--- неволиться; принуждаться; достигаться с усилием. 




\noindent\textbf{Нудить} "--- пытать. 




\noindent\textbf{Нудма} "--- насильно. 




\noindent\textbf{Нуждник} "--- употребляющий усилие. 




\noindent\textbf{Нырище} "--- развалины; руины; нежилое место. 




\noindent\textbf{Ню} "--- ее. 




\letter{О}





\noindent\textbf{Обавание} "--- ворожба; нашептывание; волхвование; колдовство. 




\noindent\textbf{Обаватель} "--- обаятель; чародей; ворожея. 




\noindent\textbf{Обавати} "--- обаять; очаровывать; ворожить; колдовать; заговаривать. 




\noindent\textbf{Обада} "--- оболгание; оклеветание. 




\noindent\textbf{Обажаемый} "--- оклеветаемый. 




\noindent\textbf{Обанадесять} "--- двенадцать. 




\noindent\textbf{Обапо} "--- с обеих сторон; по обеим сторонам. 




\noindent\textbf{Обаче} "--- однако; впрочем; но. 




\noindent\textbf{Обвеселити} "--- обрадовать. 




\noindent\textbf{Обвечеряти} "--- ночевать; переночевать. 




\noindent\textbf{Обглядати} "--- смотреть; оглядывать. 




\noindent\textbf{Обдержание} "--- сдерживание; управление; стеснение; грусть; впадение. 




\noindent\textbf{Обдесноручный} "--- человек, свободно владеющий как правой, так и левой рукой. 




\noindent\textbf{Обезвинити} "--- остаться без наказания; не знать за собой вины. 




\noindent\textbf{Обезжилити} "--- лишить сил, крепости. 




\noindent\textbf{Обезплодствити} "--- лишить плода, успеха. 




\noindent\textbf{Обезтлити} "--- сделать нетленным. 




\noindent\textbf{Обесити} "--- повесить на чем-либо. 




\noindent\textbf{Обет, обетование} "--- обещание. 




\noindent\textbf{Обетшати} "--- придти в ветхость; состариться; сделаться негодным; ослабеть; сокрушиться. 




\noindent\textbf{Обещник} "--- сообщник; товарищ. 




\noindent\textbf{Обжадати} "--- доносить; клеветать. 




\noindent\textbf{Обзорище} "--- высокая башня для наблюдения за местностью. 




\noindent\textbf{Обидитель} "--- обидчик. 




\noindent\textbf{Обиматель} "--- собиратель винограда. 




\noindent\textbf{Обиноватися} "--- колебаться; сомневаться; робеть; говорить непрямо, намеками. 




\noindent\textbf{Обиновение} "--- отступление. 




\noindent\textbf{Обиталище} "--- жилище. 




\noindent\textbf{Обитель} "--- гостиница. 




\noindent\textbf{Облагати} "--- ублажать; говорить ласково. 




\noindent\textbf{Облагодатити} "--- ниспослать благодать. 




\noindent\textbf{Облагоухати} "--- исполнить благовонием. 




\noindent\textbf{Облазнити} "--- направить по ложному следу; ввести в заблуждение. 




\noindent\textbf{Облазнитися} "--- впасть в заблуждение. 




\noindent\textbf{Область} "--- власть; сила; господство. 




\noindent\textbf{Облачити} "--- одеть. 




\noindent\textbf{Облещи} "--- облечь; одеть; лечь вокруг; окружить; сделать привал; остановиться; остаться. 




\noindent\textbf{Облистание} "--- озарение; яркий свет. 




\noindent\textbf{Облистати} "--- осветить; озарить. 




\noindent\textbf{Обличати} "--- показывать чье-либо подлинное лицо; выказывать; обнаруживать. 




\noindent\textbf{Обложити} "--- окружить. 




\noindent\textbf{Обноществовати} "--- ночевать; препроводить ночь. 




\noindent\textbf{Обнощь} "--- всю ночь. 




\noindent\textbf{Обожати} "--- обоготворять; чествовать как Бога; делать причастным Божественной благодати. 




\noindent\textbf{Оболгати} "--- обмануть. 




\noindent\textbf{Обон пол} "--- по ту сторону; за. 




\noindent\textbf{Обочие} "--- висок. 




\noindent\textbf{Обоюду} "--- по обе стороны; с обеих сторон. 




\noindent\textbf{Обрадованный} "--- приветствованный. 




\noindent\textbf{Образовати} "--- изображать; приобретать образ. 




\noindent\textbf{Обращати} "--- поворачивать; перевертывать; перемещать; вращать. 




\noindent\textbf{Обрести} "--- найти. 




\noindent\textbf{Обретаемый} "--- находимый. 




\noindent\textbf{Обретение} "--- находка; открытие. 




\noindent\textbf{Оброк} "--- плата за службу. 




\noindent\textbf{Обручник} "--- жених, помолвленный с невестой, но еще не вступивший с ней в брак. 




\noindent\textbf{Обсолонь} "--- против солнца. 




\noindent\textbf{Обстояние} "--- осада; беда; напасть. 




\noindent\textbf{Обушие} "--- мочка у уха. 




\noindent\textbf{Обуяти} "--- обезуметь; испортиться; обессилить. 




\noindent\textbf{Объюродити} "--- обезуметь; поглупеть. 




\noindent\textbf{Ов} "--- иной; один. 




\noindent\textbf{Овамо} "--- там; туда. 




\noindent\textbf{Овен} "--- баран. 




\noindent\textbf{Ово} "--- или; либо. 




\noindent\textbf{Овогда} "--- иногда. 




\noindent\textbf{Овоуду} "--- с другой стороны; оттуда. 




\noindent\textbf{Огласити} "--- объявить всенародно; научить; просветить. 




\noindent\textbf{Оглохновение} "--- глухота. 




\noindent\textbf{Огневица} "--- горячка. 




\noindent\textbf{Огненосный} "--- носимый в вихрях огня. 




\noindent\textbf{Огнепальный} "--- пылающий; горящий; палящий. 




\noindent\textbf{Огребатися} "--- удаляться; остерегаться. 




\noindent\textbf{Огустети} "--- сгустить; сделать густым; свернуться (о молоке). 




\noindent\textbf{Одебелети} "--- растолстеть; огрубеть. 




\noindent\textbf{Одесную} "--- справа; по правую руку. 




\noindent\textbf{Одесятствовати} "--- выделять десятую часть. 




\noindent\textbf{Одигитрия} "--- путеводительница. 




\noindent\textbf{Одождити} "--- окропить; оросить; послать в виде дождя; в большом количестве. 




\noindent\textbf{Одр} "--- постель; кровать. 




\noindent\textbf{Ожестети} "--- сделаться жестким; засохнуть. 




\noindent\textbf{Озимение} "--- зимовка. 




\noindent\textbf{Озлобление} "--- несчастье; гнев. 




\noindent\textbf{Озлобляти} "--- причинять несчастье; гневить; распалять гневом. 




\noindent\textbf{Озобати} "--- пожирать. 




\noindent\textbf{Окаивати} "--- признавать отверженным. 




\noindent\textbf{Окаляти} "--- пачкать; осквернять; марать. 




\noindent\textbf{Окаменяти} "--- делать каменным. 




\noindent\textbf{Окаянный} "--- достойный проклятия; нечестивый; грешник. 




\noindent\textbf{Окаянство} "--- преступность; богоборчество; грех. 




\noindent\textbf{Око} "--- глаз. 




\noindent\textbf{Окованный} "--- обложенный оковами. 




\noindent\textbf{Окормитель} "--- кормчий; правитель. 




\noindent\textbf{Окормляти} "--- направлять; руководить; править. 




\noindent\textbf{Окоявленне} "--- очевидно; откровенно. 




\noindent\textbf{Окрастовети} "--- покрыться коростою. 




\noindent\textbf{Окрест} "--- кругом; около. 




\noindent\textbf{Окриляемый} "--- ограждаемый крыльями. 




\noindent\textbf{Оле} "--- О! 




\noindent\textbf{Оловина} "--- любое хмельное питие, отличное от виноградного вина. 




\noindent\textbf{Олтарь} "--- алтарь, жертвенник. 




\noindent\textbf{Оляденети} "--- зарасти тернием, сорняками. 




\noindent\textbf{Омакати} "--- обливать. 




\noindent\textbf{Ометы} "--- полы; края одежды. 




\noindent\textbf{Она} "--- они (двое). 




\noindent\textbf{Онагр} "--- дикий осел. 




\noindent\textbf{Онамо, онуду} "--- там; туда. 




\noindent\textbf{Онде} "--- в ином месте; там. 




\noindent\textbf{Онема} "--- им (двоим). 




\noindent\textbf{Он пол} "--- противоположный берег. 




\noindent\textbf{Онсица} "--- такой-то. 




\noindent\textbf{Опасно} "--- осмотрительно; тщательно; осторожно; опасно. 




\noindent\textbf{Оплазивый} "--- любопытный; пустословный; лазутчик. 




\noindent\textbf{Оплазнство} "--- ухищрение; пустословие. 




\noindent\textbf{Оплот} "--- ограда; забор; тын. 




\noindent\textbf{Ополчатися} "--- готовиться к сражению. 




\noindent\textbf{Оправдание} "--- заповедь; устав; закон. 




\noindent\textbf{Опреснок} "--- пресный хлеб, испеченный без использования дрожжей. 




\noindent\textbf{Орало} "--- плуг; соха. 




\noindent\textbf{Оранный} "--- распаханный. 




\noindent\textbf{Оратай} "--- пахарь. 




\noindent\textbf{Орати} "--- пахать. 




\noindent\textbf{Орган} "--- орган, музыкальный инструмент. 




\noindent\textbf{Осанна} "--- молитвенное восклицание у евреев-«спасение (от Бога)». 




\noindent\textbf{Оселский} "--- ослиный. 




\noindent\textbf{Жернов оселский} "--- верхний большой жернов в мельнице, приводимый в движение ослом. 




\noindent\textbf{Осенити} "--- покрыть тенью. 




\noindent\textbf{Осклабитися} "--- усмехнуться; улыбнуться. 




\noindent\textbf{Оскорбети} "--- опечалиться; соскучиться. 




\noindent\textbf{Оскорд} "--- топор. 




\noindent\textbf{Ослаба} "--- облегчение; льгота. 




\noindent\textbf{Осля} "--- молодой осел. 




\noindent\textbf{Осмица} "--- восемь. 




\noindent\textbf{Осмоктати} "--- обсосать; облизать. 




\noindent\textbf{Оставити} "--- оставить; простить; позволить. 




\noindent\textbf{Остенити} "--- огородить стеной, защитить. 




\noindent\textbf{Острастший} "--- обидящий. 




\noindent\textbf{Острог} "--- земляной вал. 




\noindent\textbf{Острупити} "--- поразить проказой. 




\noindent\textbf{Осуществовати} "--- осуществлять; давать бытие. 




\noindent\textbf{Осьмерицею} "--- восемь раз. 




\noindent\textbf{Отай} "--- тайно; скрытно. 




\noindent\textbf{Отверзати} "--- открывать; отворять. 




\noindent\textbf{Отвнеуду} "--- снаружи. 




\noindent\textbf{Отдати} "--- иногда: простить. 




\noindent\textbf{Отдоенное} "--- грудной младенец. 




\noindent\textbf{Отдоитися} "--- воскормить грудью. 




\noindent\textbf{Отерпати} "--- делаться твердым (терпким); деревенеть; отвердевать; неметь. 




\noindent\textbf{Отити в путь всея земли} "--- умереть. 




\noindent\textbf{Откосненно} "--- наискось. 




\noindent\textbf{Откровение} "--- открытие; просветление; просвещение. 




\noindent\textbf{Отлог} "--- ущерб; урон. 




\noindent\textbf{Отложение} "--- отвержение; отступление. 




\noindent\textbf{Отметатися} "--- отрекаться; не признавать; отвергаться; отпадать. 




\noindent\textbf{Отметный} "--- отвергнутый; запрещенный. 




\noindent\textbf{Отнелиже} "--- с тех пор как; с того времени как. 




\noindent\textbf{Отнюд} "--- совершенно; отнюдь. 




\noindent\textbf{Отнюдуже, отонюдуже} "--- откуда; почему. 




\noindent\textbf{Отобоюду} "--- с той и с другой стороны. 




\noindent\textbf{Отонуду} "--- с другой стороны. 




\noindent\textbf{Отполу} "--- от половины; с середины. 




\noindent\textbf{Отре} "--- сор; мякина; кожура. 




\noindent\textbf{Отребить} "--- очистить; ощипать. 




\noindent\textbf{Отрешати} "--- отвязывать; освобождать. 




\noindent\textbf{Отрешатися} "--- разлучаться. 




\noindent\textbf{Отреяти} "--- отбрасывать; отвергать. 




\noindent\textbf{Отрицатися} "--- отвергать; отметать. 




\noindent\textbf{Отрождение} "--- возрождение. 




\noindent\textbf{Отрок} "--- раб; служитель; мальчик до двенадцати лет; ученик; воин. 




\noindent\textbf{Отроковица} "--- девица до двенадцати лет. 




\noindent\textbf{Отроча} "--- дитя; младенец. 




\noindent\textbf{Отрыгнути} "--- извергнуть. 




\noindent\textbf{Отрыгнуть слово} "--- произнести. 




\noindent\textbf{Оттоле} "--- с того времени. 




\noindent\textbf{Отторгати} "--- открывать; отталкивать. 




\noindent\textbf{Оцеждати} "--- процеживать. 




\noindent\textbf{Оцет} "--- уксус. 




\noindent\textbf{Отщетевати} "--- отнимать; удалять. 




\noindent\textbf{Отщетити} "--- потерять; погубить. 




\noindent\textbf{Очепие} "--- ошейник. 




\noindent\textbf{Очеса} "--- очи, глаза. 




\noindent\textbf{Ошаяватися} "--- устраняться, удаляться. 




\noindent\textbf{Ошиб} "--- хвост. 




\noindent\textbf{Ошуюю} "--- слева; по левую руку. 




\letter{П}





\noindent\textbf{Павечерня, павечерница} "--- малая вечерня. 




\noindent\textbf{Паволока} "--- покрывало; чехол; пелена; покров. 




\noindent\textbf{Пагуба} "--- гибель; моровая язва. 




\noindent\textbf{Пажить} "--- луг; нива; пастбище; поле; корм для скота. 




\noindent\textbf{Пазнокти} (мн. ч.) "--- копыта; когти; ногти. 




\noindent\textbf{Паки} "--- опять; еще; снова. 




\noindent\textbf{Пакибытие} "--- духовное обновление. 




\noindent\textbf{Пакости деяти} "--- бить руками; ударять по щеке; оскорблять; вредить. 




\noindent\textbf{Пакостник} "--- причинитель зла, вреда; болезнь; боль; жало. 




\noindent\textbf{Пакость} "--- гадость; нечистота; мерзость. 




\noindent\textbf{Палата} "--- дворец. 




\noindent\textbf{Иже в палате суть} "--- правительство. 




\noindent\textbf{Палестра} "--- место для соревнований. 




\noindent\textbf{Палителище} "--- сильный огонь. 




\noindent\textbf{Палительный} "--- сожигающий. 




\noindent\textbf{Палица} "--- трость; дубина; палка. 




\noindent\textbf{Паличник} "--- ликтор; телохранитель; полицейский пристав. 




\noindent\textbf{Памятозлобие} "--- злопамятство. 




\noindent\textbf{Панфирь} "--- пантера или лев. 




\noindent\textbf{Пара} "--- пар; мгла; дым. 




\noindent\textbf{Параекклесиарх} "--- кандиловжигатель; пономарь. 




\noindent\textbf{Параклис} "--- усердная молитва. 




\noindent\textbf{Параклит} "--- утешитель. 




\noindent\textbf{Паримия} "--- притча; чтения из Священного Писания на вечерне или царских часах. 




\noindent\textbf{Парити} "--- лететь; висеть в воздухе (подобно пару). 




\noindent\textbf{Парусия} "--- торжественное шествие; второе славное пришествие Господа нашего Иисуса Христа; торжественное архиерейское богослужение. 




\noindent\textbf{Пасомый} "--- пасущийся; находящийся в ведении пастыря. 




\noindent\textbf{Паствити} "--- пасти. 




\noindent\textbf{Паствуемый} "--- имеющий пастыря. 




\noindent\textbf{Пастися} "--- согрешить (особенно против седьмой заповеди). 




\noindent\textbf{Пастыреначальник} "--- начальник над пастырями. 




\noindent\textbf{Пастырь} "--- пастух. 




\noindent\textbf{Паучина} "--- паутина. 




\noindent\textbf{Паче} "--- лучше; больше. 




\noindent\textbf{Паче естества} "--- сверхъестественно. 




\noindent\textbf{Паче слова} "--- невыразимо. 




\noindent\textbf{Паче ума} "--- непостижимо. 




\noindent\textbf{Певк, певг} "--- хвойное дерево. 




\noindent\textbf{Педагогон} "--- детородный член. 




\noindent\textbf{Пекло} "--- горючая сера, смола; неперестающий огонь. 




\noindent\textbf{Пентикостарий} "--- название «Триоди цветной». 




\noindent\textbf{Пентикостия} "--- Пятидесятница. 




\noindent\textbf{Пеняжник} "--- меняла. 




\noindent\textbf{Пенязь} "--- мелкая монета. 




\noindent\textbf{Первее} "--- прежде; сперва; вначале; наперед. 




\noindent\textbf{Первоверховный} "--- первый из верховных. 




\noindent\textbf{Первовозлежание} "--- возлежание, восседание на первых, почетных местах в собраниях. 




\noindent\textbf{Первоначаток} "--- первородное животное или первоснятый плод. 




\noindent\textbf{Первостоятель} "--- первенствующий священнослужитель. 




\noindent\textbf{Пернатый} "--- имеющий перья. 




\noindent\textbf{Перси} (мн. ч.) "--- грудь; передняя часть тела. 




\noindent\textbf{Перст} "--- палец. 




\noindent\textbf{Перст возложити на уста} "--- замолчать. 




\noindent\textbf{Перстный} "--- земляной; сделанный из земли. 




\noindent\textbf{Перстосозданный} "--- сотворенный из персти. 




\noindent\textbf{Персть} "--- прах; земля; пыль. 




\noindent\textbf{Песнопети} "--- прославлять в песнях. 




\noindent\textbf{Песнословити} "--- см. Песнопети. 




\noindent\textbf{Пестовати} "--- нянчить; воспитывать. 




\noindent\textbf{Пестротный} "--- разноцветный; нарядный. 




\noindent\textbf{Пестун} "--- воспитатель; педагог; дядька. 




\noindent\textbf{Петель} "--- петух. 




\noindent\textbf{Петлоглашение} "--- пение петуха; раннее утро; время от 12 до 3 часов ночи, по народному счету времени у иудеев. 




\noindent\textbf{Печаловати} (ся) "--- сетовать, тужить; печалиться. 




\noindent\textbf{Печаловник} "--- опекун. 




\noindent\textbf{Печатствовати} "--- запечатывать; утверждать; сокрывать. 




\noindent\textbf{Печать} "--- перстень. 




\noindent\textbf{Пешешествовати} "--- идти пешком. 




\noindent\textbf{Пещися} "--- заботиться; иметь попечение. 




\noindent\textbf{Пивный} "--- то, что можно выпить. 




\noindent\textbf{Пиво} "--- питие; напиток. 




\noindent\textbf{Пиган} "--- рута, трава. 




\noindent\textbf{Пира} "--- сума; котомка. 




\noindent\textbf{Пирга} "--- башня; столп. 




\noindent\textbf{Писало} "--- остроконечная трость для писания на вощаной дощечке. 




\noindent\textbf{Писание ставильное} "--- ставленная грамота, даваемая архиереем новопосвященному пресвитеру или диакону. 




\noindent\textbf{Пискати} "--- играть на свирели. 




\noindent\textbf{Писмя} "--- буква; графический знак; буквальный смысл. 




\noindent\textbf{Пистикия} "--- чистый; беспримесный. 




\noindent\textbf{Питенный} "--- возлелеянный; выращенный в неге. 




\noindent\textbf{Питомый} "--- откормленный; дебелый. 




\noindent\textbf{Пищный} "--- содержащий обильную пищу; питательный. 




\noindent\textbf{Пиянство} "--- пьянство. 




\noindent\textbf{Плавы} "--- нивы. 




\noindent\textbf{Плавый} "--- зрелый; спелый, соломенного цвета. 




\noindent\textbf{Плат} "--- лоскут; заплатка. 




\noindent\textbf{Плащаница} "--- погребальные пелены; покрывало; полотно; плащ. 




\noindent\textbf{Плевел} "--- сорняк; негодная трава. 




\noindent\textbf{Плежити} "--- ползать на чреве; пресмыкаться. 




\noindent\textbf{Плежущий} "--- пресмыкающийся. 




\noindent\textbf{Пленица} "--- косичка; цепочка; ожерелье; корзина; цепь; оковы; узы. 




\noindent\textbf{Плескати} "--- бить в ладоши; аплодировать. 




\noindent\textbf{Плесна} "--- стопа; ступня. 




\noindent\textbf{Плесница} "--- обувь типа сандалий. 




\noindent\textbf{Плещи} "--- плечи. 




\noindent\textbf{Плещущий} "--- ударяющий в ладоши. 




\noindent\textbf{Плинфа} "--- кирпич. 




\noindent\textbf{Плинфоделание} "--- обжигание кирпичей. 




\noindent\textbf{Плищ} "--- крик; шум. 




\noindent\textbf{Плодствовати} "--- приносить плоды. 




\noindent\textbf{Плод устен} "--- слово; голос. 




\noindent\textbf{Плод чрева} "--- ребенок; дети. 




\noindent\textbf{Плотолюбие} "--- забота о теле. 




\noindent\textbf{Плотски} "--- плотью; телесно. 




\noindent\textbf{Плотский} "--- плотский; чувственный; телесный. 




\noindent\textbf{Плоть} "--- тело; человек; немощь или слабость человека; страсть. 




\noindent\textbf{Плюновение} "--- слюна. 




\noindent\textbf{Плясалище} "--- балаган. 




\noindent\textbf{Плясица, плясавица} "--- танцовщица; актриса. 




\noindent\textbf{Победительно} "--- торжественно; победоносно. 




\noindent\textbf{Поболети} "--- тужить; сожалеть. 




\noindent\textbf{Поборник} "--- защитник. 




\noindent\textbf{Повапленный} "--- покрашенный; побеленный. 




\noindent\textbf{Повергнути} "--- бросить; опрокинуть. 




\noindent\textbf{Повесть} "--- рассказ. 




\noindent\textbf{Повити} "--- принять роды или обвить пеленами. 




\noindent\textbf{Повои} "--- повязка; пелена. 




\noindent\textbf{Поглумитися} "--- рассуждать; размышлять; подумать. 




\noindent\textbf{Подвигнуться} "--- трепетать; двинуться. 




\noindent\textbf{Подвизати} "--- побуждать; поощрять. 




\noindent\textbf{Подвизатися} "--- совершать подвиги; трудиться. 




\noindent\textbf{Подвои} "--- косяки дверей. 




\noindent\textbf{Подникати} "--- наклоняться; нагибаться. 




\noindent\textbf{Подобитися} "--- напоминать что-либо. 




\noindent\textbf{Подобник} "--- подражатель. 




\noindent\textbf{Подобозрачен} "--- внешне похожий. 




\noindent\textbf{Подточилие} "--- сосуд для собирания выжатого сока. 




\noindent\textbf{Подъяремник} "--- находящийся под ярмом (например, осел). 




\noindent\textbf{Подъяремничий} "--- принадлежащий подъяремнику. 




\noindent\textbf{Поелику} "--- поскольку; потому что; так как; насколько. 




\noindent\textbf{Поелику аще} "--- сколько бы ни. 




\noindent\textbf{Пожрети} "--- принести в жертву. 




\noindent\textbf{Позде} "--- поздно; не рано. 




\noindent\textbf{Позобати} "--- склевать. 




\noindent\textbf{Позорище, позор} "--- многолюдное зрелище. 




\noindent\textbf{Поимати} "--- брать. 




\noindent\textbf{Полма} "--- пополам; надвое. 




\noindent\textbf{Польский} "--- полевой. 




\noindent\textbf{Помавати, поманути} "--- делать знаки; изъясняться без слов. 




\noindent\textbf{Помале} "--- вскоре; немного погодя. 




\noindent\textbf{Поматы} "--- скрижали на мантиях архиерейских. 




\noindent\textbf{Пометати} "--- мести; выметати; бросать. 




\noindent\textbf{Помизати} "--- мигать. 




\noindent\textbf{Поне} "--- хотя; по крайней мере; так как. 




\noindent\textbf{Понеже} "--- потому что; так как. 




\noindent\textbf{Понос, поношение} "--- позор; бесславие. 




\noindent\textbf{Понт} "--- море; большое озеро. 




\noindent\textbf{Понява, понявица} "--- полотенце. 




\noindent\textbf{Пооблещися} "--- надеть сверху другую одежду. 




\noindent\textbf{Поострити} "--- наточить. 




\noindent\textbf{Поползнутися} "--- поскользнуться; совратиться; соблазниться. 




\noindent\textbf{Поприще} "--- мера длины, равная тысяче шагов или суточному переходу. 




\noindent\textbf{Пореватися} "--- порываться; стремиться; двигаться. 




\noindent\textbf{Поречение} "--- обвинение; жалоба; упрек; попрек. 




\noindent\textbf{Порещи} "--- обвинить; укорить; осудить. 




\noindent\textbf{Поругание} "--- бесчестие; поношение; воспаление; язва. 




\noindent\textbf{Поругати} "--- обесчестить. 




\noindent\textbf{Порфира} "--- ткань темно-красного цвета; порфира, пурпурная одежда высокопоставленных особ. 




\noindent\textbf{Порча} "--- яд; отрава. 




\noindent\textbf{Поряду} "--- по порядку. 




\noindent\textbf{Поскору} "--- скоро; бегло; без пения (о службе). 




\noindent\textbf{Последи} "--- затем; в конце концов. 




\noindent\textbf{Последний} "--- остальной; конечный; окончательный. 




\noindent\textbf{Последование} "--- изложение молитвословий только одного рода, т.~е. или изменяемых, или неизменяемых. 




\noindent\textbf{Последовати} "--- исследовать; следовать. 




\noindent\textbf{Послушествовати} "--- свидетельствовать; давать показания. 




\noindent\textbf{Посолонь} "--- по-солнечному; как солнце; от востока на запад. 




\noindent\textbf{Поспешествовати} "--- помогать; пособлять. 




\noindent\textbf{Поспешник} "--- пособник; помощник. 




\noindent\textbf{Посреде} "--- посередине. 




\noindent\textbf{Поставление} "--- посвящение в сан. 




\noindent\textbf{Постриг} "--- пострижение в монашество. 




\noindent\textbf{Посупление} "--- наклонение головы в печали; печаль; грусть; сетование. 




\noindent\textbf{Посягати} "--- вступать в брак. 




\noindent\textbf{Потворник} "--- угодник; льстец; чародей; колдун. 




\noindent\textbf{Потворы} "--- чародейство; колдовство. 




\noindent\textbf{Поткнутися} "--- споткнуться. 




\noindent\textbf{Потреба} "--- потребность; необходимость; случай. 




\noindent\textbf{Потребник} "--- блин; лепешка. 




\noindent\textbf{Потщитися} "--- поспешить; постараться. 




\noindent\textbf{Поуститель} "--- подстрекатель. 




\noindent\textbf{Поущати} "--- поощрять; побуждать; наставлять; поучать. 




\noindent\textbf{Похотствовати} "--- иметь вожделение, похоть. 




\noindent\textbf{Почерпало} "--- бадья; кошель; ведро. 




\noindent\textbf{Починатися} "--- начинаться. 




\noindent\textbf{Почити} "--- успокоиться. 




\noindent\textbf{Пояти} "--- взять. 




\noindent\textbf{Правый} "--- прямой; истинный; правильный; праведный. 




\noindent\textbf{Праг} "--- порог. 

\noindent\textbf{Празднословие} "--- пустой; вздорный разговор. 



\noindent\textbf{Праздный, празден} "--- беспредельный; ленивый; пустой; незанятый. 




\noindent\textbf{Прати} "--- попирать; давить. 




\noindent\textbf{Превзятися} "--- превознестись; возгордиться. 




\noindent\textbf{Превитати} "--- странствовать. 




\noindent\textbf{Превозвышенное око} "--- высокоумие; гордость. 




\noindent\textbf{Превратити} "--- изменить; поворотить; разрушить. 




\noindent\textbf{Превременный} "--- предвечный, существовавший до начала времени. 




\noindent\textbf{Предвзыграти} "--- предвозвестить радостью. 




\noindent\textbf{Предвозгласити} "--- начать пение; предвозвестить. 




\noindent\textbf{Предградие} "--- пригород; оплот; защита; ограда. 




\noindent\textbf{Преддворие} "--- передний, внешний двор в восточном доме. 




\noindent\textbf{Предзаклатися} "--- прежде других вкусить смерть, принести себя в жертву. 




\noindent\textbf{Преди} "--- впереди. 




\noindent\textbf{Предитещи} "--- бежать впереди. 




\noindent\textbf{Предложение} "--- жертвенник; то место в алтаре, где стоит жертвенник и хранятся священные сосуды. 




\noindent\textbf{Предначинательный псалом} "--- название псалма 103, поскольку им начинается вечерня. 




\noindent\textbf{Предний} "--- первый; вящий; изящный; старший. 




\noindent\textbf{Предпряда} "--- ткань темно-красного цвета; порфира, пурпурная одежда высокопоставленных особ. 




\noindent\textbf{Предстательство} "--- ходатайство; заступничество; усердная молитва. 




\noindent\textbf{Предстолпие} "--- укрепление. 




\noindent\textbf{Предстоятель} "--- настоятель. 




\noindent\textbf{Предтеча, предитеча} "--- идущий или бегущий впереди. 




\noindent\textbf{Предуставити} "--- предназначить. 




\noindent\textbf{Предусрести} "--- встретить заранее. 




\noindent\textbf{Предуведети} "--- предвидеть; знать заранее. 




\noindent\textbf{Предъявленне} "--- предображая. 




\noindent\textbf{Презорливый} "--- гордый; надменный. 




\noindent\textbf{Преизбыточествовати} "--- быть довольну; жить в изобилии. 




\noindent\textbf{Преизлиха} "--- сильно; очень; жестоко. 




\noindent\textbf{Преимение} "--- преимущество; превосходство. 




\noindent\textbf{Преиспещренный} "--- разукрашенный. 




\noindent\textbf{Преисподний} "--- самый низкий. 




\noindent\textbf{Преисподняя} "--- место нахождения душ умерших до освобождения их Господом Иисусом Христом; место вечного мучения грешников; жилище диавола. 




\noindent\textbf{Преитие} "--- превосхождение. 




\noindent\textbf{Прелагати пределы} "--- портить межи; нарушать границы. 




\noindent\textbf{Прелесть} "--- обман. 




\noindent\textbf{Прелюбы} "--- прелюбодейство. 




\noindent\textbf{Премудрость} "--- высшее знание; мудрость. 




\noindent\textbf{Преначальный} "--- доначальный; превышающий всякое начало. 




\noindent\textbf{Преогорчити} "--- противиться; быть непокорным; упрямиться. 




\noindent\textbf{Преодеян} "--- обильно украшен. 




\noindent\textbf{Преоруженный} "--- слишком вооруженный; гордый. 




\noindent\textbf{Препирати} "--- опровергать; отражать; отбивать; одолевать; увещевать. 




\noindent\textbf{Преподобие} "--- святость. 




\noindent\textbf{Преполовение} "--- половина; середина. 




\noindent\textbf{Преполовити} "--- переполовинить; разделить пополам; пройти половину пути. 




\noindent\textbf{Препона} "--- препятствие. 




\noindent\textbf{Препоясатися} "--- подпоясаться; приготовиться к чему-либо. 




\noindent\textbf{Препретельный} "--- спорный; убедительный. 




\noindent\textbf{Препростый} "--- неученый; невежда. 




\noindent\textbf{Препяти} "--- остановить. 




\noindent\textbf{Пререкаемый} "--- спорный. 




\noindent\textbf{Пререкати} "--- прекословить; говорить наперекор; перечить. 




\noindent\textbf{Пресецающий} "--- пересекающий; перерубающий. 




\noindent\textbf{Преслушание} "--- неповиновение. 




\noindent\textbf{Пресмыкаться} "--- ползти по земле. 




\noindent\textbf{Преспевати} "--- иметь успех. 




\noindent\textbf{Преставити} "--- переставить; переместить; переселить в вечность. 




\noindent\textbf{Престоли} "--- один из чинов ангельских. 




\noindent\textbf{Пресущественный} "--- предвечный; исконный. 




\noindent\textbf{Пресущный} "--- сверхъестественный. 




\noindent\textbf{Претися} "--- спорить; тягаться. 




\noindent\textbf{Претити} "--- запрещать; скорбеть; смущаться. 




\noindent\textbf{Преткновение} "--- помеха; соблазн; задержка; остановка. 




\noindent\textbf{Претор} "--- претория, резиденция представителя римской власти в Иерусалиме. 




\noindent\textbf{Претыкание} "--- помеха; соблазн; задержка; остановка. 




\noindent\textbf{Прещати} "--- грозить; устрашать. 




\noindent\textbf{Прещение} "--- угроза; страх; запрет. 




\noindent\textbf{Прибежище} "--- убежище; приют; покров; спасение. 




\noindent\textbf{Приведение} "--- доступ. 




\noindent\textbf{Привещевати} "--- приветствовать. 




\noindent\textbf{Привлещи} "--- притащить; позвать; призвать. 




\noindent\textbf{Привременный} "--- временный; непостоянный. 




\noindent\textbf{Придевати} "--- прицеплять; приближаться; подносить. 




\noindent\textbf{Придел} "--- небольшая церковь, пристроенная к главному храму. 




\noindent\textbf{Приделати} "--- прирастить; увеличить; принести. 




\noindent\textbf{Придеяти} "--- подносить; приносить. 




\noindent\textbf{Призрети} "--- милостиво посмотреть; принять; приютить. 




\noindent\textbf{Прииждивати} "--- расходовать; издерживать. 




\noindent\textbf{Приискренне} "--- точно так же; равно; точь-в-точь. 




\noindent\textbf{Прикровение} "--- прикрытие; предлог; выдуманная причина для сокрытия чего-либо. 




\noindent\textbf{Прикуп} "--- барыш; прибыль. 




\noindent\textbf{Прикупование} "--- купечество; торговля. 




\noindent\textbf{Прилог} "--- приложение; желание сделать зло; злоба; клевета. 




\noindent\textbf{Приложение} "--- заплатка; лоскут. 




\noindent\textbf{Приметати} "--- прибрасывать; отдавать; уступать. 




\noindent\textbf{Приметатися} "--- припадать; отдаваться; лежать у порога. 




\noindent\textbf{Примешатися} "--- присоединяться. 




\noindent\textbf{Приникнути} "--- пригнуться; наклониться; припасть; проникнуть. 




\noindent\textbf{Приобряща} "--- польза; плод; корысть. 




\noindent\textbf{Приразитися} "--- напасть; удариться. 




\noindent\textbf{Приревание} "--- устремление. 




\noindent\textbf{Приристати} "--- подбегать. 




\noindent\textbf{Пририщущий} "--- подбегающий. 




\noindent\textbf{Присвянути} "--- завянуть; засохнуть. 




\noindent\textbf{Приседение} "--- угнетение; окружение. 




\noindent\textbf{Приседети} "--- находиться около чего-либо; замышлять зло; нападать. 




\noindent\textbf{Присно} "--- непрестанно; всегда. 




\noindent\textbf{Присноживотный} "--- всегда живущий. 




\noindent\textbf{Присносущий} "--- вечный; всегдашний. 




\noindent\textbf{Присносущный} "--- всегда существующий. 




\noindent\textbf{Приснотекущий} "--- неиссякаемый. 




\noindent\textbf{Присный} "--- родной; близкий. 




\noindent\textbf{Приставление} "--- заплатка; назначение; управление; присмотр. 




\noindent\textbf{Пристанище} "--- приют; убежище; пристань. 




\noindent\textbf{Пристати} "--- прибегнуть; подбежать. 




\noindent\textbf{Пристрашен} "--- испуган. 




\noindent\textbf{Притвор} "--- вход в храм. 




\noindent\textbf{Прителный} "--- спорный. 




\noindent\textbf{Притча} "--- иносказание; загадка. 




\noindent\textbf{Причаститися} "--- стать участником. 




\noindent\textbf{Причастник} "--- участник. 




\noindent\textbf{Пришлец} "--- приезжий; пришелец. 




\noindent\textbf{Приятилище} "--- вместилище; поместилище; хранилище. 




\noindent\textbf{Пробавити} "--- продолжить; протянуть. 




\noindent\textbf{Продерзивый} "--- дерзкий. 




\noindent\textbf{Прозябение} "--- произрастание; росток. 




\noindent\textbf{Прозябнути} "--- расцвести; вырасти; произрастить. 




\noindent\textbf{Произникнути} "--- произойти; вырасти. 




\noindent\textbf{Пролитися стопам} "--- поскользнуться; иносказательно "--- согрешить. 




\noindent\textbf{Пронарещи} "--- предсказать; предназначить. 




\noindent\textbf{Проникнути} "--- вырасти; процвесть. 




\noindent\textbf{Проничение} "--- племя; род; стебель; росток. 




\noindent\textbf{Проповедати} "--- учить; провозглашать; проповедовать. 




\noindent\textbf{Прорещи} "--- предсказать. 




\noindent\textbf{Пророкованный} "--- предсказанный; предвозвещенный. 




\noindent\textbf{Пророковещательный} "--- говоримый пророком. 




\noindent\textbf{Проручествовати} "--- посвящать; рукополагать. 




\noindent\textbf{Просаждатися} "--- разрываться. 




\noindent\textbf{Просветительный} "--- светлый; просвещающий. 




\noindent\textbf{Просветити лице} "--- весело или милостивно взглянуть. 




\noindent\textbf{Проскомисати} "--- совершать проскомидию. 




\noindent\textbf{Прослутие} "--- притча; пословица; осмеяние. 




\noindent\textbf{Простый} "--- стоящий прямо; прямой; чистый; несмешанный. 




\noindent\textbf{Простыня} "--- сострадание. 




\noindent\textbf{Просядати} "--- разрываться; разваливаться; трескаться. 




\noindent\textbf{Протерзатися} "--- прорываться. 




\noindent\textbf{Противозрети} "--- смотреть прямо. 




\noindent\textbf{Противу, прямо} "--- против; напротив. 




\noindent\textbf{Проуведети} "--- узнать заранее; предвидеть. 




\noindent\textbf{Проявленне} "--- явно. 




\noindent\textbf{Пругло} "--- силок; петля; сеть. 




\noindent\textbf{Прудный} "--- неровный; каменистый. 




\noindent\textbf{Пружатися} "--- сопротивляться (отсюда "--- пружина); биться в припадке. 




\noindent\textbf{Пружие} "--- пища Иоанна Крестителя; по мнению одних "--- род съедобной саранчи, или кузнечиков; по мнению других "--- какое-то растение. 




\noindent\textbf{Пря} "--- спор; тяжба; беспорядок. 




\noindent\textbf{Пряжмо} "--- жареная пища. 




\noindent\textbf{Прямный} "--- находящийся напротив. 




\noindent\textbf{Пустити} "--- отпустить; развестись. 




\noindent\textbf{Пустыня} "--- уединенное, малообитаемое место. 




\noindent\textbf{Пустыня} "--- монастырь, расположенный вдалеке от населенных мест. 




\noindent\textbf{Путесотворити} "--- сохранять в пути; проложить дорогу. 




\noindent\textbf{Путы} "--- узы; кандалы; цепи; оковы. 




\noindent\textbf{Пучина} "--- водоворот; море. 




\noindent\textbf{Пучинородный} "--- морской; родившийся в море. 




\noindent\textbf{Пущеница} "--- разведенная с мужем женщина. 




\noindent\textbf{Птицеволхвование} "--- суеверие, состоящее в гадании по полету птиц или по их внутренностям. 




\noindent\textbf{Пядь, пядень} "--- мера длины, равная трем дланям, а каждая длань равна четырем перстам, а перст равен четырем граням или зернам. 




\noindent\textbf{Пясть} "--- кулак. 




\noindent\textbf{Пяток} "--- пятница. 




\letter{Р}





\noindent\textbf{Рабий} "--- рабский. 




\noindent\textbf{Работа} "--- рабство. 




\noindent\textbf{Работен} "--- покорен; порабощен. 




\noindent\textbf{Равви, раввуни} "--- учитель. 




\noindent\textbf{Равноангельно} "--- подобно Ангелам. 




\noindent\textbf{Равноапостольный} "--- сравниваемый с апостолами. 




\noindent\textbf{Равнодушный, равнодушевный} "--- единодушный; имеющий одинаковое усердие. 




\noindent\textbf{Равночестный} "--- достойный равного почитания. 




\noindent\textbf{Радованный} "--- радостный. 




\noindent\textbf{Радоватися} "--- радоваться; наслаждаться. 




\noindent\textbf{Радощи} "--- радости (мн. ч.); веселье. 




\noindent\textbf{Радуйся} "--- здравствуй; прощай. 




\noindent\textbf{Раждежение} "--- горение; воспламенение. 




\noindent\textbf{Разботети} "--- растолстеть; разбухнуть. 




\noindent\textbf{Разве} "--- кроме. 




\noindent\textbf{Развет} "--- мятеж; заговор. 




\noindent\textbf{Разврат} "--- волнение; возмущение. 




\noindent\textbf{Разгбенный} "--- разогнутый. 




\noindent\textbf{Разгнутие} "--- разгибание; раскрытие книги. 




\noindent\textbf{Раздолие} "--- долина. 




\noindent\textbf{Разжизати} "--- разжигать; раскалять; расплавлять. 




\noindent\textbf{Размыслити} "--- усомниться; задуматься; остановиться. 




\noindent\textbf{Разнство} "--- различие. 




\noindent\textbf{Разрешити} "--- развязать; освободить. 




\noindent\textbf{Разслабленный} "--- паралитик. 




\noindent\textbf{Разум} "--- ум; познание; разумение. 




\noindent\textbf{Разумети телом} "--- почувствовать. 




\noindent\textbf{Рака} "--- евр. дурак; пустой человек. 




\noindent\textbf{Рака} "--- гробница; ковчег с мощами святого угодника Божия. 




\noindent\textbf{Рало} "--- соха; плуг. 




\noindent\textbf{Рамо} "--- плечо. 




\noindent\textbf{Рамена} "--- плечи. 




\noindent\textbf{Расплощатися} "--- развертываться. 




\noindent\textbf{Распростирати} "--- расстеливать; разворачивать. 




\noindent\textbf{Распудити} "--- распугать; разогнать; рассеять. 




\noindent\textbf{Распутие} "--- перекресток. 




\noindent\textbf{Раст} "--- росток. 




\noindent\textbf{Растерзати} "--- разорвать. 




\noindent\textbf{Растнити} "--- рассечь. 




\noindent\textbf{Расточати} "--- рассеивать; рассыпать; проматывать; беспутно проживать. 




\noindent\textbf{Растренный} "--- перепиленный. 




\noindent\textbf{Расчинити} "--- расположить по порядку. 




\noindent\textbf{Ратай} "--- воин. 




\noindent\textbf{Ратовати} "--- воевать; сражаться; отстаивать. 




\noindent\textbf{Ратовище} "--- древко копья. 




\noindent\textbf{Рать} "--- война; воинство. 




\noindent\textbf{Рачитель} "--- попечитель; любитель. 




\noindent\textbf{Рачительный} "--- заботливый; достойный заботы. 




\noindent\textbf{Рвение, ревность} "--- ярость; страстное желание; страсть. 




\noindent\textbf{Ребра северова} "--- северный склон горы Сион. 




\noindent\textbf{Ревновати} "--- завидовать. 




\noindent\textbf{Рек} "--- ты, он сказал. 




\noindent\textbf{Рекла} "--- сказала. 




\noindent\textbf{Рекомый} "--- прозываемый. 




\noindent\textbf{Рекох} "--- я сказал. 




\noindent\textbf{Репие} "--- репейник; колючее растение. 




\noindent\textbf{Ресно} "--- ресницы; глаз. 




\noindent\textbf{Реснота} "--- действительность; истина. 




\noindent\textbf{Реть} "--- ссора; спор. 




\noindent\textbf{Рещи} "--- сказать; говорить. 




\noindent\textbf{Реяти} "--- отталкивать; отбрасывать. 




\noindent\textbf{Риза} "--- одежда; священное облачение. 




\noindent\textbf{Ризница} "--- помещение для сохранения риз. 




\noindent\textbf{Ризничий} "--- начальник над ризницей; хранитель церковной утвари. 




\noindent\textbf{Ристалище} "--- стадион; цирк. 




\noindent\textbf{Ристати} "--- рыскать; бегать. 




\noindent\textbf{Рог} "--- рок; иносказательно: сила; власть; защита. 




\noindent\textbf{Род} "--- происхождение; племя; поколение. 




\noindent\textbf{Родостама} "--- розовая вода, которой по обычаю в праздник Воздвижения производят омовение Честнаго и Животворящего Креста Господня при его воздвизании. 




\noindent\textbf{Рожаный} "--- роговой; напоминающий рог. 




\noindent\textbf{Рожец} "--- сладкий стручок. 




\noindent\textbf{Розга} "--- молодая ветвь; побег; отпрыск. 




\noindent\textbf{Росодательный} "--- росоносный; дающий росу. 




\noindent\textbf{Рота} "--- божба; клятва. 




\noindent\textbf{Ротитель} "--- клятвопреступник. 




\noindent\textbf{Ротитися} "--- клясться; божиться. 




\noindent\textbf{Ругатися} "--- насмехаться. 




\noindent\textbf{Рукописание} "--- список; письмо; письменный договор; свиток; расписка; обязательство. 




\noindent\textbf{Рукоять} "--- горсть; охапка. 




\noindent\textbf{Руно} "--- шерсть; овчина. 




\noindent\textbf{Ручка} "--- сосуд. 




\noindent\textbf{Рцем} "--- скажем (повел.наклонение). 




\noindent\textbf{Рцы} "--- скажи. 




\noindent\textbf{Рыбарь} "--- рыбак. 




\noindent\textbf{Рясно} "--- ожерелье; подвески. 




\letter{С}





\noindent\textbf{Самвик} "--- музыкальный инструмент. 




\noindent\textbf{Самовидец} "--- очевидец. 




\noindent\textbf{Самогласная стихира} "--- имеющая свой особый распев. 




\noindent\textbf{Самоохотие} "--- по собственному желанию. 




\noindent\textbf{Самоподобен} "--- стихира, имеющая свой особый распев. 




\noindent\textbf{Самочиние} "--- бесчиние; беспорядок. 




\noindent\textbf{Сата} "--- мера сыпучих тел. 




\noindent\textbf{Сбодати} "--- пронзить; заколоть. 




\noindent\textbf{Свара} "--- ссора; брань. 




\noindent\textbf{Сваритися} "--- ссориться. 




\noindent\textbf{Сведети} "--- ведать; знать. 




\noindent\textbf{Светильничное} "--- начало вечерни. 




\noindent\textbf{Светлопозлащен} "--- великолепно украшен. 




\noindent\textbf{Светлость} "--- светящаяся красота. 




\noindent\textbf{Светозарный} "--- озаряющий светом. 




\noindent\textbf{Светолитие} "--- сияние. 




\noindent\textbf{Светоначальник} "--- создатель светил. 




\noindent\textbf{Светоносец} "--- несущий свет. 




\noindent\textbf{Свечеряти} "--- совместно с кем-либо участвовать в пиру. 




\noindent\textbf{Свидение} "--- наставление; приказание. 




\noindent\textbf{Свирепоустие} "--- необузданность языка. 




\noindent\textbf{Свиток} "--- сверток; рукопись, намотанная на палочку. 




\noindent\textbf{Связание злата} "--- впрядение золотых нитей. 




\noindent\textbf{Связень} "--- узник; невольник. 




\noindent\textbf{Святилище} "--- алтарь; храм. 




\noindent\textbf{Святитель} "--- архиерей; епископ. 




\noindent\textbf{Святотатство} "--- похищение священных вещей. 




\noindent\textbf{Святцы} "--- месяцеслов (книга, содержащая имена святых, расположенных по дням года); икона «Всех святых». 




\noindent\textbf{Священнотаинник} "--- посвященный в Божественные тайны. 




\noindent\textbf{Се} "--- вот. 




\noindent\textbf{Седмерицею} "--- семикратно. 




\noindent\textbf{Седмица} "--- семь дней, которые в современном языке принято называть «неделя». 




\noindent\textbf{Седмичный} "--- относящийся к любому из дней седмицы, кроме Недели (воскресного дня); будничный. 




\noindent\textbf{Секира} "--- топор. 




\noindent\textbf{Секраты} "--- недавно; только что. 




\noindent\textbf{Секратый} "--- свежий; новый. 




\noindent\textbf{Селный} "--- полевой; дикий. 




\noindent\textbf{Село} "--- поле. 




\noindent\textbf{Семидал} "--- мелкая пшеничная мука; крупчатка. 




\noindent\textbf{Семо} "--- сюда. 




\noindent\textbf{Семя} "--- семя; потомки; род. 




\noindent\textbf{Сеннописание} "--- неясное изображение. 




\noindent\textbf{Сень} "--- тень; покров над престолом. 




\noindent\textbf{Септемврий} "--- сентябрь. 




\noindent\textbf{Серповидец} "--- наименование святого пророка Захарии. 




\noindent\textbf{Серядь} "--- монашеское рукоделие; пряжа. 




\noindent\textbf{Сеть} "--- западня. 




\noindent\textbf{Сечиво} "--- топор. 




\noindent\textbf{Сигклит} (читается «синклит») "--- собрание, сенат. 




\noindent\textbf{Сиесть} "--- то есть. 




\noindent\textbf{Сикарий} "--- убийца; разбойник. 




\noindent\textbf{Сикелия} "--- о. Сицилия. 




\noindent\textbf{Сикер} "--- хмельной напиток, изготовленный не из винограда. 




\noindent\textbf{Силы} "--- название одного из чинов ангельских; иногда значит чудеса. 




\noindent\textbf{Синаксарий} "--- сокращенное изложение житий святых или праздников. 




\noindent\textbf{Синедрион} "--- верховное судилище у иудеев. 




\noindent\textbf{Синфрог} "--- сопрестолие, т.~е. скамьи по обе стороны горнего места для сидения сослужащих архиерею священников. 




\noindent\textbf{Сиречь} "--- то есть; именно. 




\noindent\textbf{Сирини} "--- (в Ис. 13, 21) "--- страусы; сирены. 




\noindent\textbf{Сирт} "--- отмель; мель. 




\noindent\textbf{Сирый} "--- сиротливый; одинокий; беспомощный; бедный. 




\noindent\textbf{Сице} "--- так; таким образом. 




\noindent\textbf{Сицевый} "--- такой; таковой. 




\noindent\textbf{Скверна} "--- нечистота; грязь; порок. 




\noindent\textbf{Сквозе} "--- сквозь; через. 




\noindent\textbf{Скимен} "--- молодой лев; львенок. 




\noindent\textbf{Скиния} "--- палатка; шатер. 




\noindent\textbf{Скинотворец} "--- делатель палаток. 




\noindent\textbf{Скит} "--- маленький монастырь. 




\noindent\textbf{Склабитися} "--- улыбаться, усмехаться. 




\noindent\textbf{Скнипа} "--- вошь. 




\noindent\textbf{Сковник} "--- соучастник; сообщник. 




\noindent\textbf{Скоктание} "--- щекотание; подстрекательство. 




\noindent\textbf{Скопчий} "--- скопческий. 




\noindent\textbf{Скоротеча} "--- скороход; гонец. 




\noindent\textbf{Скорпия} "--- скорпион. 




\noindent\textbf{Скрания} "--- висок. 




\noindent\textbf{Скрижаль} "--- доска; таблица. 




\noindent\textbf{Скудель} "--- глина; то, что сделано из глины; кувшин; черепица. 




\noindent\textbf{Скудельник} "--- горшечник. 




\noindent\textbf{Скудный} "--- бедный; тощий. 




\noindent\textbf{Скураты} "--- маски; личины. 




\noindent\textbf{Славник} "--- молитвословие, положенное по уставу после «Славы». 




\noindent\textbf{Славословие} "--- прославление. 




\noindent\textbf{Сладковонный} "--- благоуханный. 




\noindent\textbf{Сладкогласие} "--- стройное пение. 




\noindent\textbf{Сладкопение} "--- тихое, умильное пение. 




\noindent\textbf{Слана} "--- гололедица; мороз; ледник; замерзший иней. 




\noindent\textbf{Сланость} "--- соленая морская вода; солончак, т.~е. сухая, пропитанная солью земля; гололед. 




\noindent\textbf{Сластотворный} "--- обольщающий плотскими удовольствиями. 




\noindent\textbf{Сликовствовати} "--- совместно играть; веселиться. 




\noindent\textbf{Словесный} "--- разумный. 




\noindent\textbf{Словоположение} "--- договор; условия. 




\noindent\textbf{Сложитися} "--- уговориться; определить. 




\noindent\textbf{Слота} "--- ненастье; дурная погода. 




\noindent\textbf{Слух} "--- слава; народная молва. 




\noindent\textbf{Слякий, слукий} "--- согнутый; скорченный; горбатый. 




\noindent\textbf{Сляцати} "--- сгибать; горбить. 




\noindent\textbf{Смарагд} "--- изумруд. 




\noindent\textbf{Смежити} "--- сблизить; соединять края, межи; закрывать. 




\noindent\textbf{Смерчие} "--- кедр. 




\noindent\textbf{Смеситися} "--- перемещаться; совокупиться плотски. 




\noindent\textbf{Смиряти} "--- унижать. 




\noindent\textbf{Смоква} "--- плод фигового дерева. 




\noindent\textbf{Смотрение} "--- промысел; попечение; забота. 




\noindent\textbf{Смясти} "--- привести в смятение; встревожить. 




\noindent\textbf{Снабдевати} "--- сберегать; сохранять. 




\noindent\textbf{Снедати} "--- съедать; разорять; сокрушать. 




\noindent\textbf{Снедь} "--- пища. 




\noindent\textbf{Сниматися} "--- сходиться; собираться. 




\noindent\textbf{Снисхождение} "--- снисшествие. 




\noindent\textbf{Снитися} "--- вступить в брак; сойтись. 




\noindent\textbf{Снуждею} "--- поневоле; насильно; по принуждению. 




\noindent\textbf{Собесити} "--- вместе повесить. 




\noindent\textbf{Соблюдение} "--- точное исполнение; темница. 




\noindent\textbf{С соблюдением} "--- видимым образом; явно. 




\noindent\textbf{Совет} "--- совет; решение; определение. 




\noindent\textbf{Советный} "--- рассудительный. 




\noindent\textbf{Совлачити} "--- разоблачить; снять. 




\noindent\textbf{Совлечение} "--- раздевание. 




\noindent\textbf{Совлещися} "--- раздеться. 




\noindent\textbf{Совозвести} "--- возвести вместе с собой. 




\noindent\textbf{Совоздыхати} "--- печалиться вместе. 




\noindent\textbf{Совопрошатися} "--- беседовать; состязаться в споре. 




\noindent\textbf{Совоспитанный} "--- воспитанный совместно с кем-либо. 




\noindent\textbf{Согласно} "--- единодушно. 




\noindent\textbf{Соглядатай} "--- разведчик; шпион. 




\noindent\textbf{Соглядати} "--- рассматривать; наблюдать; разведывать. 




\noindent\textbf{Сограждати} "--- сооружать; строить. 




\noindent\textbf{Содеватися} "--- сделаться. 




\noindent\textbf{Соделование} "--- дело; превращение. 




\noindent\textbf{Содетель} "--- творец. 




\noindent\textbf{Содетельный} "--- творческий. 




\noindent\textbf{Сокровище} "--- потаенное место; задняя комната; хранилище; клад; драгоценность; погреб. 




\noindent\textbf{Сокровиществовати} "--- собирать сокровища. 




\noindent\textbf{Сокрушение} "--- уничтожение. 




\noindent\textbf{Сокрушение сердца} "--- раскаяние. 




\noindent\textbf{Солило} "--- солонка; чаша; блюдо. 




\noindent\textbf{Соние} "--- сон; сновидение. 




\noindent\textbf{Сонм} "--- собрание; множество. 




\noindent\textbf{Сонмище} "--- синагога. 




\noindent\textbf{Сопель} "--- свирель, дудка. 




\noindent\textbf{Сопети} "--- играть на дудке. 




\noindent\textbf{Сопец} "--- сопельщик-музыкант, играющий на сопели, флейте (при похоронах у иудеев). 




\noindent\textbf{Сопретися} "--- ссориться; тягаться. 




\noindent\textbf{Соприбывати} "--- увеличиваться. 




\noindent\textbf{Соприсносущный} "--- совместно существующий в вечности. 




\noindent\textbf{Сопрягати} "--- соединять браком. 




\noindent\textbf{Сорокоустие} "--- пшеница, вино, фимиам, свечи и пр., приносимые в церковь на 40 дней поминовения усопших христиан. 




\noindent\textbf{Соскание} "--- шнурок; веревочка. 




\noindent\textbf{Сосканый} "--- витой; крученый 




\noindent\textbf{Соскутовати} "--- спеленать; окутать. 




\noindent\textbf{Состреляти} "--- поразить стрелой. 




\noindent\textbf{Сосудохранительница} "--- помещение для сохранения церковной утвари. 




\noindent\textbf{Сосуды смертные} "--- орудия смерти. 




\noindent\textbf{Сосцы} "--- иногда иносказательно так называются водные источники. 




\noindent\textbf{Сотница} "--- сотня; пение «Господи, помилуй» сто раз при воздвизании Честнаго Креста Господня. 




\noindent\textbf{Сотово тело} "--- мед. 




\noindent\textbf{Соуз} "--- союз; связь. 




\noindent\textbf{Сочетаватися} "--- вступать в союз, в брак. 




\noindent\textbf{Сочиво} "--- чечевица; вареная пшеница с медом. 




\noindent\textbf{Сочинение} "--- составление; собрание. 




\noindent\textbf{Спекулатор} "--- телохранитель. 




\noindent\textbf{Спира} "--- отряд; рота; полк. 




\noindent\textbf{Сплавати} "--- сопутствовать в плавании. 




\noindent\textbf{Споболети} "--- вместе печалиться; тужить. 




\noindent\textbf{Споборать} "--- вместе воевать. 




\noindent\textbf{Спод} "--- ряд; куча; отдел. 




\noindent\textbf{Спона} "--- препятствие. 




\noindent\textbf{Спослушествовати} "--- свидетельствовать; подтверждать. 




\noindent\textbf{Споспешник} "--- помощник. 




\noindent\textbf{Спостник} "--- вместе постящийся. 




\noindent\textbf{Спострадати} "--- вместе страдать. 




\noindent\textbf{Споющий} "--- вместе или одновременно поющий. 




\noindent\textbf{Спротяженный} "--- продолжительный. 




\noindent\textbf{Спуд} "--- сосуд; ведерко; мера сыпучих тел; покрышка; плита. 




\noindent\textbf{Спяти} "--- низвергнуть; опрокинуть 




\noindent\textbf{Срамословие} "--- сквернословие. 




\noindent\textbf{Срасленный} "--- сросшийся. 




\noindent\textbf{Срачица} "--- сорочка; рубаха. 




\noindent\textbf{Сребреник} "--- серебряная монета. 




\noindent\textbf{Сребропозлащенный} "--- позолоченный по серебру. 




\noindent\textbf{Средоградие} "--- перегородка; простенок; преграда. 




\noindent\textbf{Средостение} "--- перегородка; средняя стена. 




\noindent\textbf{Сретение} "--- встреча. 




\noindent\textbf{Сристатися} "--- стекаться; сбегаться. 




\noindent\textbf{Срящь} "--- неприятная встреча; нападение; зараза; мор; гаданье; приметы. 




\noindent\textbf{Ставленник} "--- человек, подготовляемый к посвящению в духовный сан. 




\noindent\textbf{Стадия} "--- мера длины, равная 100--125 шагам. 




\noindent\textbf{Стаинник} "--- сопричастный с кем-либо одной тайне. 




\noindent\textbf{Стакти} "--- благовонный сок. 




\noindent\textbf{Стамна} "--- сосуд; ведерко; кувшин. 




\noindent\textbf{Старей} "--- начальник; старший. 




\noindent\textbf{Статир} "--- серебряная или золотая монета. 




\noindent\textbf{Статия} "--- глава; подраздел. 




\noindent\textbf{Стегно} "--- верхняя половина ноги; бедро; ляжка. 




\noindent\textbf{Стезя} "--- тропинка; дорожка. 




\noindent\textbf{Стень, сень} "--- тень; отражение; образ. 




\noindent\textbf{Степени} "--- ступени. 




\noindent\textbf{Стерти} "--- стереть; разрушить. 




\noindent\textbf{Стихира} "--- песнопение. 




\noindent\textbf{Стих началу} "--- первый возглас священника при богослужении общественном или частном. 




\noindent\textbf{Стихологисати} "--- петь избранные стихи из Псалтири при богослужении. 




\noindent\textbf{Стихология} "--- чтение или пение Псалтири. 




\noindent\textbf{Стихословити} "--- петь избранные стихи из Псалтири при богослужении. 




\noindent\textbf{Стицатися} "--- стекаться; сходиться. 




\noindent\textbf{Сткляница} "--- стакан. 




\noindent\textbf{Сткляный} "--- стеклянный. 




\noindent\textbf{Столп} "--- башня; крепость. 




\noindent\textbf{Столпостена} "--- башня; крепость. 




\noindent\textbf{Стомах} "--- желудок. 




\noindent\textbf{Стопа} "--- ступня. 




\noindent\textbf{Сторицею} "--- во сто раз. 




\noindent\textbf{Стогна} "--- улица, дорога. 




\noindent\textbf{Страдальчество} "--- мученичество. 




\noindent\textbf{Стража} "--- караул; охрана; мера времени для ночи. 




\noindent\textbf{Страннолепный} "--- необычный. 




\noindent\textbf{Странный} "--- сторонний; чужой; прохожий; необычайный. 




\noindent\textbf{Странь} "--- напротив; против. 




\noindent\textbf{Страсть} "--- страдание; страсть; душевный порыв. 




\noindent\textbf{Стратиг} "--- военачальник. 




\noindent\textbf{Стратилат} "--- военачальник; воевода. 




\noindent\textbf{Страхование} "--- угроза; страх; ужас. 




\noindent\textbf{Стрекало} "--- спица; палочка с колючкой для управления скотом. 




\noindent\textbf{Стрещи} "--- стеречь. 




\noindent\textbf{Стрищи} "--- стричь; подстригать. 




\noindent\textbf{Стропотный} "--- кривой; извилистый; строптивый; упрямый; злой. 




\noindent\textbf{Стрыти} "--- стереть; сокрушить. 




\noindent\textbf{Студ} "--- стыд; срам. 




\noindent\textbf{Студенец} "--- колодец; родник; источник. 




\noindent\textbf{Студень} "--- холод; стужа; мороз. 




\noindent\textbf{Стужаемый} "--- беспокоимый. 




\noindent\textbf{Стужание} "--- стеснение; гонение; досаждение. 




\noindent\textbf{Стужати} "--- докучать; надоедать; теснить. 




\noindent\textbf{Стужатися} "--- скорбеть; печалиться. 




\noindent\textbf{Стужный} "--- тревожный. 




\noindent\textbf{Стягнути} "--- обвязать; собрать; исцелить. 




\noindent\textbf{Стязатися} "--- спорить; препираться. 




\noindent\textbf{Сугубый} "--- двойной; удвоенный; увеличенный; усиленный. 




\noindent\textbf{Сударь} "--- плат; пелена. 




\noindent\textbf{Судище} "--- суд; приговор. 




\noindent\textbf{Суемудренный} "--- софистический; пустословный. 




\noindent\textbf{Суеслов} "--- пустослов. 




\noindent\textbf{Суета} "--- пустота; ничтожность; мелочность; бессмысленность. 




\noindent\textbf{Суетие} "--- суетность; суета. 




\noindent\textbf{Сулица} "--- копье; кинжал; кортик. 




\noindent\textbf{Супостат} "--- противник; враг. 




\noindent\textbf{Супротивный} "--- противник. 




\noindent\textbf{Супруг} "--- пара; чета. 




\noindent\textbf{Супря} "--- спор; тяжба. 




\noindent\textbf{Суровый} "--- зеленый; свежий; сырой. 




\noindent\textbf{Сходник} "--- лазутчик; разведчик; шпион. 




\noindent\textbf{Счиневати} "--- соединять. 




\noindent\textbf{Сыноположение} "--- усыновление. 




\letter{Т}





\noindent\textbf{Таже} "--- потом; затем. 




\noindent\textbf{Тай} "--- тайно; скрытно. 




\noindent\textbf{Таинник} "--- посвященный в чьи-либо тайны. 




\noindent\textbf{Тайноядение} "--- тайное невоздержание от пищи в пост. 




\noindent\textbf{Таити} "--- скрывать. 




\noindent\textbf{Тако} "--- так. 




\noindent\textbf{Такожде} "--- равно; также. 




\noindent\textbf{Талант} "--- древнегреч. мера веса и монета. 




\noindent\textbf{Тамо} "--- там; туда. 




\noindent\textbf{Тартар} "--- место нахождения душ умерших до освобождения их Господом Иисусом Христом; место вечного мучения грешников; жилище диавола. 




\noindent\textbf{Татаур} "--- ремень для привешивания языка к колоколу; кожаный ремень, носимый монашествующими . 




\noindent\textbf{Тать} "--- вор. 




\noindent\textbf{Татьба} "--- кража; воровство. 




\noindent\textbf{Тафта} "--- тонкая шелковая материя. 




\noindent\textbf{Тацы} "--- таковы. 




\noindent\textbf{Таче} "--- для того; также; тогда. 




\noindent\textbf{Тварь} "--- творение; создание; произведение. 




\noindent\textbf{Твердыня} "--- крепость; цитадель; тюрьма. 




\noindent\textbf{Твердь} "--- основание; видимый небосклон, принимаемый глазом за твердую сферу, купол небес. 




\noindent\textbf{Твержа} "--- крепость; цитадель; тюрьма. 




\noindent\textbf{Тезоименитство} "--- одноименность; именины; день Ангела. 




\noindent\textbf{Тектон} "--- плотник; столяр. 




\noindent\textbf{Телец} "--- теленок; бычок. 




\noindent\textbf{Темже, темже убо} "--- поэтому; следовательно; итак. 




\noindent\textbf{Темник} "--- начальник над десятью тысячами человек. 




\noindent\textbf{Темнозрачный} "--- черный. 




\noindent\textbf{Темнонеистовство} "--- мракобесие; непросвещенность. 




\noindent\textbf{Теревинф} "--- дубрава; чаща; лес; большое ветвистое дерево с густой листвой. 




\noindent\textbf{Терние} "--- терновник; колючее растение. 




\noindent\textbf{Терноносный} "--- плодоносящий терние; иносказательно: не имеющий добрых дел. 




\noindent\textbf{Терпкий} "--- кислый; вяжущий; суровый. 




\noindent\textbf{Теснина} "--- узкий проход. 




\noindent\textbf{Теснота} "--- беда; напасть. 




\noindent\textbf{Тетива} "--- туго натянутая веревка. 




\noindent\textbf{Тещи} "--- быстро идти. 




\noindent\textbf{Тещити} "--- источать; испускать. 




\noindent\textbf{Тимение} "--- болото; топь; тина. 




\noindent\textbf{Тимпан} "--- литавра; бубен. 




\noindent\textbf{Тимпанница} "--- девушка, играющая на тимпане. 




\noindent\textbf{Тирон} "--- молодой воин, солдат. 




\noindent\textbf{Титло} "--- надпись; ярлык; знак для сокращения слова. 




\noindent\textbf{Тихообразно} "--- спокойно; кротко. 




\noindent\textbf{Тление, тля} "--- гниение; уничтожение; разрушение. 




\noindent\textbf{Тлети} "--- растлевать; гнить; разрушаться. 




\noindent\textbf{Тлити} "--- повреждать; губить. 




\noindent\textbf{Тма} "--- темнота; мрак; десять тысяч. 




\noindent\textbf{Тоболец} "--- мешок; котомка; сумка. 




\noindent\textbf{Ток} "--- течение. 




\noindent\textbf{Токмо} "--- только. 




\noindent\textbf{Толико} "--- столько. 




\noindent\textbf{Толк} "--- толкование; учение; особое мнение. 




\noindent\textbf{Толковник} "--- переводчик; истолкователь. 




\noindent\textbf{Толковый} "--- объясняющий; содержащий объяснения. 




\noindent\textbf{Толмач} "--- переводчик. 




\noindent\textbf{Толь} "--- столько. 




\noindent\textbf{Томитель} "--- мучитель. 




\noindent\textbf{Томити} "--- мучить; пытать. 




\noindent\textbf{Томление} "--- мучение; пытка. 




\noindent\textbf{Топазий} "--- топаз. 




\noindent\textbf{Торжище} "--- площадь; рынок. 




\noindent\textbf{Торжник} "--- меняла; торговец. 




\noindent\textbf{Торчаный} "--- растерзанный. 




\noindent\textbf{Точило} "--- пресс для выжимания виноградного сока. 




\noindent\textbf{Точию} "--- только. 




\noindent\textbf{Тощно} "--- усердно; точно. 




\noindent\textbf{Трапеза} "--- стол; кушание; столовая, трапезная; святой престол. 




\noindent\textbf{Требе} "--- потребно; надобно. 




\noindent\textbf{Требище} "--- жертвенник; языческий храм. 




\noindent\textbf{Треблаженный} "--- весьма блаженный. 




\noindent\textbf{Требование} "--- нужда; потребность. 




\noindent\textbf{Требовати} "--- нуждаться; иметь потребность. 




\noindent\textbf{Трегубо} "--- трояко; трижды. 




\noindent\textbf{Трекровник} "--- третий этаж. 




\noindent\textbf{Тресна} "--- украшение на одежде. 




\noindent\textbf{Третицею} "--- трижды; в третий раз. 




\noindent\textbf{Тридевять} "--- двадцать семь. 




\noindent\textbf{Трисиянный} "--- светящий от трех Светил. 




\noindent\textbf{Тристат} "--- военачальник. 




\noindent\textbf{Трищи} "--- трижды. 




\noindent\textbf{Тропарь} "--- краткое песнопение, выражающее характеристику праздника или события в жизни святого. 




\noindent\textbf{Трость} "--- тростник (использовавшийся в качестве пишущего инструмента). 




\noindent\textbf{Труд} "--- болезнь; недуг. 




\noindent\textbf{Труждатися} "--- трудиться; затрудняться. 




\noindent\textbf{Трус} "--- землетрясение. 




\noindent\textbf{Трыти} "--- тереть; омывать. 




\noindent\textbf{Ту} "--- тут; там; здесь. 




\noindent\textbf{Туга} "--- скорбь. 




\noindent\textbf{Тук} "--- жир; сало; богатство; пресыщение. 




\noindent\textbf{Тул} "--- колчан для стрел. 




\noindent\textbf{Туне} "--- напрасно; даром; впустую. 




\noindent\textbf{Тунегиблемый} "--- истрачиваемый noнапрасну. 




\noindent\textbf{Тщание} "--- усердие; старание. 




\noindent\textbf{Тщатися} "--- стараться; спешить. 




\noindent\textbf{Тщета} "--- урон; вред; убыток. 




\noindent\textbf{Тщий} "--- пустой; бесполезный; неудовлетворенный. 




\noindent\textbf{Тягота} "--- тяжесть; обременение. 




\noindent\textbf{Тяжание} "--- работа; дело; пашня; поле. 




\noindent\textbf{Тяжатель} "--- работник. 




\noindent\textbf{Тяжати} "--- работать. 




\noindent\textbf{Тяжкосердый} "--- бесчувственный. 




\letter{У}





\noindent\textbf{У} "--- еще; 




\noindent\textbf{не у} "--- еще не. 




\noindent\textbf{Убо} "--- а; же; вот; хотя; почему; поистине; подлинно. 




\noindent\textbf{Убрус} "--- плат; полотенце. 




\noindent\textbf{Убудитися} "--- пробудиться; очнуться. 




\noindent\textbf{Уведети} "--- узнать. 




\noindent\textbf{Увет} "--- увещание. 




\noindent\textbf{Увязение} "--- возложение на голову венца. 




\noindent\textbf{Увясло} "--- головная повязка. 




\noindent\textbf{Углебнути} "--- тонуть; погружаться; погрязать. 




\noindent\textbf{Угобзити} "--- обогатить, одарить. 




\noindent\textbf{Угонзнути} "--- убежать; ускользнуть; уйти. 




\noindent\textbf{Уготоватися} "--- приготовиться. 




\noindent\textbf{Угрызнути} "--- укусить зубами. 




\noindent\textbf{Уд} "--- телесный член. 




\noindent\textbf{Удава} "--- веревка. 




\noindent\textbf{Удица} "--- удочка. 




\noindent\textbf{Удобие} "--- удобнее. 




\noindent\textbf{Удобострастие} "--- склонность к угождению страстям. 




\noindent\textbf{Удобрение} "--- украшение. 




\noindent\textbf{Удобрити} "--- наполнить; украсить. 




\noindent\textbf{Удобь} "--- легко; удобно. 




\noindent\textbf{Удовлитися} "--- удовлетвориться. 




\noindent\textbf{Удолие, удоль, юдоль} "--- долина. 




\noindent\textbf{Удручати} "--- утомлять; оскорблять. 




\noindent\textbf{Уже} "--- веревка; цепь; узы. 




\noindent\textbf{Ужик} "--- родственник; родственница. 




\noindent\textbf{Узилище} "--- тюрьма. 




\noindent\textbf{Укорение} "--- бесславие; бесчестие. 




\noindent\textbf{Укрой} "--- повязка; пелена. 




\noindent\textbf{Укроп} "--- теплота, т.~е. горячая вода, вливаемая во святой потир на Литургии. 




\noindent\textbf{Укрух} "--- ломоть; кусок. 




\noindent\textbf{Улучити} "--- застать; найти; получить. 




\noindent\textbf{Умащати} "--- намазывать; натирать. 




\noindent\textbf{Умет} "--- помет; кал; сор. 




\noindent\textbf{Умовредие} "--- безумие. 




\noindent\textbf{Умучити} "--- укротить. 




\noindent\textbf{Уне} "--- лучше. 




\noindent\textbf{Унзнути} "--- воткнуть; вонзить. 




\noindent\textbf{Унше} "--- лучше; полезнее. 




\noindent\textbf{Упитанная} "--- откормленные животные. 




\noindent\textbf{Упование} "--- твердая надежда. 




\noindent\textbf{Упразднити} "--- уничтожить; отменить; исчезнуть. 




\noindent\textbf{Уранити} "--- встать рано утром. 




\noindent\textbf{Урок} "--- урок; подать; оброк. 




\noindent\textbf{Усекновение} "--- отсечение. 




\noindent\textbf{Усерязь} "--- серьга. 




\noindent\textbf{Усма} "--- выделанная кожа. 




\noindent\textbf{Усмарь} "--- кожевенник; скорняк. 




\noindent\textbf{Усмен} "--- кожаный. 




\noindent\textbf{Уста} "--- рот; губы; речь. 




\noindent\textbf{Устранити} "--- лишить; избежать; устранить. 




\noindent\textbf{Устудити} "--- охладить; остудить. 




\noindent\textbf{Усырити} "--- сделать сырым, твердым, мокрым. 




\noindent\textbf{Утварне} "--- по порядку; нарядно. 




\noindent\textbf{Утварь} "--- одежда; украшение; убранство. 




\noindent\textbf{Утверждение} "--- основание; опора. 




\noindent\textbf{Утешение} "--- угощение. 




\noindent\textbf{Утолити} "--- успокоить; утешить; умерить. 




\noindent\textbf{Утреневати} "--- рано вставать; совершать утреннюю молитву. 




\noindent\textbf{Утроба} "--- чрево; живот; сердце; душа. 




\noindent\textbf{Ухание} "--- обоняние; запах. 




\noindent\textbf{Ухлебити} "--- накормить. 




\noindent\textbf{Учреждати} "--- угостить. 




\noindent\textbf{Учреждение} "--- пир; обед; угощение. 




\noindent\textbf{Ушеса} "--- уши. 




\noindent\textbf{Ущедрити} "--- обогатить; помиловать; пожалеть. 




\letter{Ф}





\noindent\textbf{Факуд} "--- евр. начальник. 




\noindent\textbf{Фарос} "--- маяк. 




\noindent\textbf{Фаска} "--- Пасха. 




\noindent\textbf{Февруарий} "--- февраль. 




\noindent\textbf{Фелонь} "--- плащ; верхняя одежда; одно из священных облачений пресвитера. 




\noindent\textbf{Фиала} "--- чаша; бокал с широким дном. 




\noindent\textbf{Фимиам} "--- благовонная смола для воскурения при каждении. 




\noindent\textbf{Финикс} "--- пальма. 




\letter{Х}





\noindent\textbf{Халван} "--- благовонная смола. 




\noindent\textbf{Халколиван} "--- ливанская медь; янтарь. 




\noindent\textbf{Халуга} "--- плетень; забор; закоулок. 




\noindent\textbf{Харатейный} "--- написанный на прегаменте или папирусной бумаге. 




\noindent\textbf{Хартия} "--- пергамент или папирусная бумага; рукописный список. 




\noindent\textbf{Харя} "--- маска; личина. 




\noindent\textbf{Хврастие} "--- хворост. 




\noindent\textbf{Херет} "--- училище, темница. 




\noindent\textbf{Хитон} "--- нижняя одежда; рубашка. 




\noindent\textbf{Хитрец} "--- художник; ремесленник. 




\noindent\textbf{Хитрогласница} "--- риторика. 




\noindent\textbf{Хитрословесие} "--- риторика. 




\noindent\textbf{Хитрость} "--- художество; ремесло. 




\noindent\textbf{Хитрость, ухищренная вымыслом} "--- стенобитные и метательные военные машины. 




\noindent\textbf{Хищноблудие} "--- насильственное привлечение к блуду. 




\noindent\textbf{Хламида} "--- верхнее мужское платье; плащ; мантия. 




\noindent\textbf{Хлептати} "--- лакать. 




\noindent\textbf{Хлябь} "--- водопад; пропасть; бездна; простор; подъемная дверь. 




\noindent\textbf{Ходатай} "--- посредник, примиритель. 




\noindent\textbf{Хоругвь} "--- военное знамя. 




\noindent\textbf{Хотение} "--- воля. 




\noindent\textbf{Храмляти} "--- хромать. 




\noindent\textbf{Храм набдящий} "--- казнохранилище. 




\noindent\textbf{Храм, храмина} "--- дом; помещение; место богослужения. 




\noindent\textbf{Хранилище} "--- повязка на лбу или на руках со словами Закона. 




\noindent\textbf{Худогий} "--- искусный; умелый. 




\noindent\textbf{Худогласие} "--- косноязычие; заикание. 




\noindent\textbf{Художество} "--- наука; причуда; выкрутаса. 




\noindent\textbf{Худость} "--- скудость; недостоинство. 




\noindent\textbf{Хула} "--- злословие; нарекание. 




\letter{Ц}





\noindent\textbf{Цветник} "--- луг. 




\noindent\textbf{Цевница} "--- флейта; свирель. 




\noindent\textbf{Целование} "--- приветствие. 




\noindent\textbf{Целовати} "--- приветствовать. 




\noindent\textbf{Целомудрие} "--- благоразумие, непорочность и чистота телесная. 




\noindent\textbf{Целый} "--- здоровый, невредимый. 




\noindent\textbf{Цельбоносный} "--- врачебный; целительный. 




\letter{Ч}





\noindent\textbf{Чадо} "--- дитя. 




\noindent\textbf{Чадородие} "--- рождение детей. 




\noindent\textbf{Чадце} "--- деточка. 




\noindent\textbf{Чарование} "--- яд; отрава; волхвование; заговаривание. 




\noindent\textbf{Чаровник} "--- отравитель; волхв; колдун. 




\noindent\textbf{Чары} "--- волшебство; колдовство; заговаривание. 




\noindent\textbf{Часть} "--- часть; жребий; участь. 




\noindent\textbf{Чаяти} "--- надеяться; ждать. 




\noindent\textbf{Чван} "--- сосуд; штоф; кружка; фляжка. 




\noindent\textbf{Чванец} "--- сосуд; штоф; кружка; фляжка. 




\noindent\textbf{Чело} "--- лоб. 




\noindent\textbf{Челядь} "--- слуги; домочадцы. 




\noindent\textbf{Чепь} "--- цепь. 




\noindent\textbf{Червленица} "--- ткань темно-красного цвета; порфира, пурпурная одежда высокопоставленных особ. 




\noindent\textbf{Червленый} "--- багряный. 




\noindent\textbf{Чермноватися} "--- краснеть. 




\noindent\textbf{Чермный} "--- красный. 




\noindent\textbf{Чернец} "--- монах. 




\noindent\textbf{Черничие} "--- лесная смоква. 




\noindent\textbf{Чертог} "--- палата; покои. 




\noindent\textbf{Чесати} "--- собирать плоды. 




\noindent\textbf{Чесо} "--- чего; что. 




\noindent\textbf{Чести} "--- читать. 




\noindent\textbf{Честный} "--- уважаемый; прославляемый. 




\noindent\textbf{Четверовластник} "--- управляющий четвертою частью страны. 




\noindent\textbf{Чин} "--- порядок; полное изложение или указание всех молитвословий. 




\noindent\textbf{Чинити} "--- составлять; делать. 




\noindent\textbf{Чревобесие} "--- объядение; обжорство. 




\noindent\textbf{Чревоношение} "--- зачатие и ношение в утробе младенца. 




\noindent\textbf{Чреда} "--- порядок; очередь; черед. 




\noindent\textbf{Чреждение} "--- угощение. 




\noindent\textbf{Чресла} "--- поясница; бедра; пах. 




\noindent\textbf{Чтилище} "--- идол; кумир. 




\noindent\textbf{Чудитися} "--- удивляться. 




\noindent\textbf{Чудотворити} "--- творить чудеса. 




\noindent\textbf{Чути, чуяти} "--- чувствовать; слышать; ощущать. 




\letter{Ш}





\noindent\textbf{Шелом} "--- шлем; каска. 




\noindent\textbf{Шепотник} "--- наушник; клеветник. 




\noindent\textbf{Шептание} "--- клевета. 




\noindent\textbf{Шипок} "--- цветок шиповника. 




\noindent\textbf{Шуий} "--- левый. 




\noindent\textbf{Шуйца} "--- левая рука. 




\letter{Щ}





\noindent\textbf{Щедрота} "--- милость; великодушие; снисхождение. 




\noindent\textbf{Щедрый} "--- милостивый. 




\noindent\textbf{Щогла} "--- мачта; веха; жердь. 




\letter{Ю}





\noindent\textbf{Ю} "--- ее. 




\noindent\textbf{Юг} "--- зной; название южного ветра; иносказательно; несчастие. 




\noindent\textbf{Юдоль} "--- долина. 




\noindent\textbf{Юдоль плачевная} "--- мир сей. 




\noindent\textbf{Юдуже} "--- там, где. 




\noindent\textbf{Южик (а)} "--- родственник; родственница. 




\noindent\textbf{Юзник} "--- узник; заключенный. 




\noindent\textbf{Юзы} "--- кандалы; узы. 




\noindent\textbf{Юнец} "--- телок; молодой бычок. 




\noindent\textbf{Юница} "--- телка; молодая корова. 




\noindent\textbf{Юнота} "--- молодой человек. 




\noindent\textbf{Юродивый} "--- глупый; принявший духовный подвиг юродства. 




\letter{Я}





\noindent\textbf{Ягодичина} "--- фиговое дерево. 




\noindent\textbf{Ядрило} "--- мачта. 




\noindent\textbf{Ядца} "--- лакомка; гурман; обжора. 




\noindent\textbf{Ядь} "--- пища; еда. 




\noindent\textbf{Язвина} "--- нора. 




\noindent\textbf{Язвити} "--- жечь; ранить. 




\noindent\textbf{Язык} "--- народ; племя; орган речи. 




\noindent\textbf{Язя, язва} "--- рана; ожог. 




\noindent\textbf{Яко} "--- ибо; как; так как; потому что; когда. 




\noindent\textbf{Яковый} "--- каковой. 




\noindent\textbf{Якоже} "--- так чтобы; как; так как. 




\noindent\textbf{Яннуарий} "--- январь. 




\noindent\textbf{Ярем} "--- ярмо; груз; тяжесть; служение. 




\noindent\textbf{Ярина} "--- волна; шерсть. 




\noindent\textbf{Ясти} "--- есть; кушать. 




\noindent\textbf{Яти} "--- брать. 
\normalfont\end{mymulticols}
\longpage\mychapterending


 



\renewcommand{\ornament}{uzor_begin_9}
\mychapter{Список изменений}
\markright{Список изменений}

Этот раздел содержит информацию об изменениях и исправлениях в текстах молитв сайта \url{http://www.molitvoslov.com}, вошедших в эту и предыдущие публикации. Информация о более новых версиях располагается в начале, а о более старых "--- в конце списка.

%\begin{center}
%\begin{tabular}{lp{0.7\textwidth}l}
%\toprule
%Версия & Изменения\\
%\midrule
%20110606 &
%Обновление текстов с сайта от 06 мая 2011 г.
%\\
%20110503.02 & Содержимое сайта от 03 мая 2011 г. \\
%\bottomrule
%\end{tabular}
%\end{center}

\small

\section*{Версия 20120531}

\begin{itemize}

\item Обновление текстов с сайта от 31-го мая 2012 г.
\item Исправлен русскоязычный поиск по тексту, а также копирование и вставка из PDF в программе Adobe Acrobat Reader.
\item Исправлены опечатки в следующих текстах:

\begin{itemize}
\item Молитва пред иконою «Неупиваемая Чаша».
\item Молитвы задержания.
\item Молитва против антихриста.
\item Молитва апостолу Симону Зилоту.
\item Молитва о спасении державы Российской и утолении в ней раздоров и нестроений.
\item Тропарь Кресту и молитва за отечество (Молитвы утренние).
\item Молитвы на сон грядущим.
\item Святому праведному Иоанну, Кронштадтскому чудотворцу.
\end{itemize}

\item В следующих молитвах и канонах добавлены недостающие припевы, а вместо сокращений (например, «Слава:», «И ныне:» и др.) используются полные тексты молитв:

\begin{itemize}
\item Последование ко Святому Причащению.
\item Благодарственные молитвы по Святом Причащении.
\end{itemize}

\end{itemize}

\section*{Версия 20111108}

Обновление текстов с сайта от 8-го ноября 2011 г.

\begin{itemize}

\item Добавлен раздел «Утешение сродникам самоубийц».
\item Исправлены опечатки в следующих текстах:

\begin{itemize}

\item «Богородичное правило».
\item «Канон за болящего, глас 3-й».
\item Молитва «Пророку, Предтече и Крестителю Господню Иоанну».
\item Молитва «Святителю Димитрию, Ростовскому чудотворцу».
\item Молитва «Святой блаженной Ксении Петербургской».
\item «Молитва покаянная, читаемая в церквах России во дни смуты».
\item «КАНОН ПРЕПОДОБНОМУ ПАИСИЮ ВЕЛИКОМУ ОБ ИЗБАВЛЕНИИ ОТ МУК УМЕРШИХ БЕЗ ПОКАЯНИЯ».
\item Псалом 50-й, в разделе «ЧАСЫ ТРЕТИЙ И ШЕСТОЙ».
\item Псалом 50-й, в «Молитвы утренние».
\item Псалом 50-й и 26-й, в разделе «ЕДИНОВЕРЧЕСКИЕ МОЛИТВЫ».
\item Псалом 26-й, в разделе «МОЛИТВЕННОЕ ПРАВИЛО ПРЕПОДОБНОГО АМВРОСИЯ OПТИНСКОГО».
\item Молитва «Святому преподобному Иосифу Волоцкому».
\item «Молитва от осквернения».
\item Молитва «Великомученику Георгию Победоносцу».
\end{itemize}

\end{itemize}

\section*{Версия 20110709}

Обновление текстов с сайта от 9 июля 2011 г.

\begin{itemize}

\item Добавлен новый раздел «О Молитве», а также «Словарь и термины» (около 2500 слов и выражений).

\item Небольшое изменение в разделе «Божественная литургия».

\item Добавлены заглавия тропарей, Богородичных, канонов и песен в

\begin{itemize}

\item «Канон Ангелу Хранителю» («Молитвы»).

\item «Канон молебный ко Господу Иисусу Христу и Пречистой Богородице Матери Господни при разлучении души от тела всякого правоверного» («НАПУТСТВИЕ ХРИСТИАНИНА ПЕРЕД СМЕРТЬЮ И ЗАУПОКОЙНЫЕ МОЛИТВЫ»).

\item в разделе «КАНОН ПРЕПОДОБНОМУ ПАИСИЮ ВЕЛИКОМУ ОБ ИЗБАВЛЕНИИ ОТ МУК УМЕРШИХ БЕЗ ПОКАЯНИЯ».

\end{itemize}
\end{itemize}

\section*{Версия 20110616}

Обновление текстов с сайта от 16 июня 2011 г.
\begin{itemize}

\item В разделе «РАЗЛИЧНЫЕ МОЛИТВЫ»
\begin{itemize}
\item Исправлены опечатки в «Молитвах ко Господу Св. Иоанна Кронштадского».
\end{itemize}

\item В разделе «МОЛИТВЫ СВЯТЫМ»
\begin{itemize}
\item Убраны повторы молитв прп. Александру Свирскому и прп. Серафиму Саровскому.
\item Добавлены молитвы св. первомученику архидиакону Стефану; св. царице Грузинской Тамаре; святителю Иоанну Шанхайскому и Сан-Францисскому.
\item Исправлены опечатки в молитвах св. Царю Страстотерпцу Николаю; св. прп. Силуану Афонскому; cвятителю Тихону, Патриарху Московскому и Всея Руси.
\end{itemize}

\item В разделе «МОЛИТВЫ В СКОРБЯХ И ИСКУШЕНИЯХ ТВОРИМЫЕ»
\begin{itemize}
\item Добавлена «Молитва преследуемого человеками (свт. Игнатия Брянчанинова)».
\end{itemize}

\item В разделе «МОЛИТВЫ ЗА БОЛЯЩИХ»
\begin{itemize}
\item Исправлена опечатка в «Каноне за болящего, глас 3-й», Песнь 9.
\end{itemize}

\item В разделе «НАПУТСТВИЕ ХРИСТИАНИНА ПЕРЕД СМЕРТЬЮ И ЗАУПОКОЙНЫЕ МОЛИТВЫ»
\begin{itemize}
\item Исправлен «Чин литии, совершаемой мирянином дома и на кладбище».
\item Исправлен тропарь прп. Паисию Великому в молитве «Об ослаблении вечных мук умерших без покаяния», а также в разделе «КАНОН ПРЕПОДОБНОМУ ПАИСИЮ ВЕЛИКОМУ ОБ ИЗБАВЛЕНИИ ОТ МУК УМЕРШИХ БЕЗ ПОКАЯНИЯ».
\item Исправлен «Канон молебный ко Господу Иисусу Христу и Пречистой Богородице Матери Господни при разлучении души от тела всякого правоверного».
\end{itemize}

\item В разделе «В ПОМОЩЬ КАЮЩЕМУСЯ»
\begin{itemize}
\item Исправлены опечатки в «Десяти заповедях».
\end{itemize}


\end{itemize}


\section*{Версия 20110503.02}
Содержимое сайта от 03 мая 2011 г.
\mychapterending
%\normalfont\nopagebreak\bigskip\bigskip\begin{center}\includegraphics[width=0.15\textwidth]{uzor_end_4}\end{center}



\end{document}
